% !TEX program = pdflatex
\documentclass[12pt,a4paper,openany]{report}

\usepackage[utf8]{inputenc}
\usepackage{CJKutf8}
\usepackage{amsmath,amssymb,mathtools,amsthm}
\usepackage{geometry}
\geometry{top=25mm,bottom=25mm,left=20mm,right=20mm}
\usepackage{xcolor}
\usepackage{url}
\usepackage{iftex}
\ifPDFTeX
  \usepackage[scaled=0.84]{beramono}
\else
  \usepackage{fontspec}
  \setmainfont{Consolas}
  \setsansfont{Consolas}
  \setmonofont{Consolas}
\fi
\usepackage[unicode,colorlinks=true,linkcolor=blue,citecolor=blue,urlcolor=blue]{hyperref}
\usepackage{xurl}
\usepackage{enumitem}
\usepackage{booktabs}
\usepackage{longtable}
\usepackage{tocloft}

\Urlmuskip=0mu plus 1mu
\def\UrlFont{\ttfamily\small}
\DeclareRobustCommand{\urlunderscore}{\raisebox{-0.35ex}{\rule{0.55em}{0.35pt}}}
\makeatletter
\g@addto@macro\UrlSpecials{\do\_{\urlunderscore}}
\makeatother
\DeclareUrlCommand{\codeurl}{\urlstyle{tt}}
\DeclareRobustCommand{\code}[1]{\ifmmode\text{\codeurl{#1}}\else\codeurl{#1}\fi}
\sloppy

\renewcommand{\cftchapleader}{\cftdotfill{\cftdotsep}}
\renewcommand{\bibname}{参考文献}
\newtheorem{definition}{定义}[chapter]
\newtheorem{assumption}{假设}[chapter]
\newtheorem{axiom}{公理}[chapter]
\newtheorem{theorem}{定理}[chapter]
\newtheorem{lemma}{引理}[chapter]
\newtheorem{proposition}{命题}[chapter]
\newtheorem{corollary}{推论}[chapter]
\newtheorem{remark}{备注}[chapter]
\newtheorem{Certificate}{核验记录}[chapter]

\newcommand{\NN}{\mathbb N}
\newcommand{\RR}{\mathbb R}
\newcommand{\EE}{\mathbb E}
\newcommand{\DeltaN}{\Delta^{n-1}}
\newcommand{\Gov}{\mathcal G}
\newcommand{\Law}{\mathcal L}
\newcommand{\Mech}{\mathcal M}
\newcommand{\Feas}{\mathcal F}
\newcommand{\Part}{\mathcal P}
\newcommand{\Act}{\mathcal A}
\newcommand{\Goods}{\mathcal J}
\newcommand{\Rules}{\mathcal R}
\newcommand{\States}{\mathcal X}
\newcommand{\Budget}{\mathcal B}
\newcommand{\Info}{\mathcal I}
\newcommand{\DefEq}{\coloneqq}
\newcommand{\Abs}[1]{\left\lvert #1\right\rvert}
\newcommand{\Norm}[1]{\left\lVert #1\right\rVert}
\newcommand{\Tuple}[1]{\left\langle #1\right\rangle}
\DeclareMathOperator*{\argmax}{arg\,max}
\DeclareMathOperator*{\argmin}{arg\,min}
\DeclareMathOperator{\supp}{supp}
\DeclareMathOperator{\CapacitySlack}{CapacitySlack}
\DeclareMathOperator{\AsymptoticSlack}{AsymptoticSlack}
\DeclareMathOperator{\PostAI}{PostAI}
\DeclareMathOperator{\Binding}{Binding}
\DeclareMathOperator{\Eligible}{Eligible}
\DeclareMathOperator{\Measured}{Measured}
\DeclareMathOperator{\Verified}{Verified}
\DeclareMathOperator{\Valued}{Valued}
\DeclareMathOperator{\Entitled}{Entitled}
\DeclareMathOperator{\Paid}{Paid}
\DeclareMathOperator{\Auditable}{Auditable}
\DeclareMathOperator{\Appealable}{Appealable}
\DeclareMathOperator{\BudgetBalanced}{BudgetBalanced}
\DeclareMathOperator{\IC}{IC}
\DeclareMathOperator{\IR}{IR}
\DeclareMathOperator{\DuplicateClaimResistant}{DuplicateClaimResistant}
\DeclareMathOperator{\CollusionResistant}{CollusionResistant}
\DeclareMathOperator{\OpportunityFair}{OpportunityFair}
\DeclareMathOperator{\RuleTransparent}{RuleTransparent}
\DeclareMathOperator{\LedgerConsistent}{LedgerConsistent}
\DeclareMathOperator{\ExecutionConsistent}{ExecutionConsistent}
\DeclareMathOperator{\GovernanceFair}{GovernanceFair}
\DeclareMathOperator{\Sustainable}{Sustainable}
\DeclareMathOperator{\PrimaryDominant}{PrimaryDominant}
\DeclareMathOperator{\Financialized}{Financialized}
\DeclareMathOperator{\RuleSensitive}{RuleSensitive}
\DeclareMathOperator{\Falsifiable}{Falsifiable}

\setlength{\emergencystretch}{8em}
\setlength{\parindent}{0pt}
\setlength{\parskip}{0.6em}
\setcounter{tocdepth}{2}
\setcounter{secnumdepth}{3}

\begin{document}
\begin{CJK*}{UTF8}{gkai}

\title{市场参与的形式系统\\
\large 参与式分配、能者多得与最小救济的演化机制}
\author{Gao Zheng（高政）}
\date{2026-09-29}
\pagenumbering{roman}
\maketitle
\setcounter{page}{2}

\chapter*{前言}
\addcontentsline{toc}{chapter}{前言}

本文研究后 AI 条件下的生产--消费内核、市场参与、价值形成与分配机制。市场经济的配置路径由分散主体在一般法律约束内
通过供求、价格、盈亏、声誉和进入退出共同生成；计划配置则由中心主体在参与者选择之前预先指定关键产量、
价格或个体配置。某类可复制产品的生产能力约束松弛以后，有效需求和互补稀缺要素仍使最优产出有限，
生产--消费循环对就业、收入与新增需求的边际传导可能衰减。具有真实服务、但不以新增物质产出为核心的参与--分配
行业以竞技为主导形成双边注意力市场：赛事组织连接对局、观赛、教学、解说、创作与服务，多个生产者通过
广告注意力竞争投入赞助资金。组织以有来源的资金池覆盖分配，完成约定的有效参与取得正基本份额，已验证
贡献形成差异份额，使重在参与与能者多得同时成立。参与收入回到商品购买，销售回款覆盖产品成本并形成销售
利润；竞技经历产生的产品偏好带来增量销售利润，用以覆盖赞助及配套营销成本，并形成新增净利润。已实现
净利润再支持下一期生产与竞争性赞助，从而连接竞技、注意力、组织分配、消费和再投入。
本文的核心机制不是注意力变现本身，也不是核验、确权和申诉工具本身，而是
\[
\boxed{\text{参与式分配}\;+\;\text{能者多得}\;+\;\text{坐标最小救济}}.
\]
三项功能分别扩大取得正份额的有效参与范围，保留能力、贡献与稀缺性的差异激励，并使市场收入不足者仍能维持
生存和真实市场接入。基于注意力的偏好回流为它提供市场资金与演化选择压力；资格、计量、核验、合同确权、
实际支付、复算和申诉则提供可实施、可纠错的稳定结构。后两者分别是经济与演化基础、实施基础，不能取代三项
核心功能本身。
现有自媒体、直播和平台创作表明注意力可以转化为广告、订阅、打赏或平台收入，却同时提供一个不充分性反例：
若收入集中于少数头部主体，多数有效参与者没有正份额且生存底线没有保障，则
\[
  \text{注意力可以变现}
  \nRightarrow
  \text{广泛吸收就业与参与并修复收入--需求循环}.
\]
因此，注意力偏好回流只说明资金来源和选择动力，不自动完成分配。该反例排除“单纯注意力经济已经充分”的
推论；它不单独证明本文所列组织形式具有历史唯一性。本文后续证明的是：一旦目标同时固定为广泛正份额、
贡献差异激励和全员底线接入，任何达成该目标的载体都须具有与上述三项相等价的功能坐标。
在最低物质需要、正超额预算和可识别需求系统下，兴趣消费还可能在声明收入域内取得大于一的收入弹性；当
付费覆盖、行业净收入、供给进入、真实性稀缺和净需求桥接同时越过事前阈值时，它可以形成分层的广义奢侈品
产业，并在局部稳定域内填补可识别需求缺口。经验证的长期参与还可以在稳定规则下沉积为共同历史、共享术语、
身份结构、制度记忆与共同体认同；文化意义映射再把参与体验转化为身份、归属、声誉等心理效用，并进入持续
需求。围棋为这一文化沉积提供已经存在的参照模型，新形成的参与组织须从共同经历中积累自身的社会文化背景。
该结论不把兴趣、注意力、高价或平台单方面宣布的文化标签本身当作宏观规模证明。
持续正必要品需要与同一需要的多个可行载体还给出一个更基础的结论：只要满足最低需要的具体配置由消费者
分散选择，已实现购买必然形成载体间的非零选择份额投射。本文只把这一可观测份额简称为弱偏好投射，不把它
等同于心理偏好、无约束偏好或福利排序；市场负责显露投向、集中度和变化。该结构必然性不自动推出注意力传导
为正；注意力能否改变原有选择分布仍须由反事实证据识别。
当自动化同时压低必要品成本和传统劳动收入时，最低物质束、市场接入束及其资金缺口还会成为循环自身的状态
变量。本文证明：若没有救济的市场收入不能维持生存与真实进入，且由此使大众需求、参与密度或注意力桥接跌破
非退化阈值，那么在全员底线承诺及生产剩余可持续融资条件下，坐标最小救济可与稳定市场循环相容。逐人救济下界来自全员底线承诺；总量稳定本身不决定逐人额度。它只抬升
生存与接入下限，不规定兴趣、价格、投入、赢家或下限以上收入，因而不等于按需分配或收入平均化。
市场回应先形成交易与可分配剩余，组织合同再对其内部资格、计量、验证、初始确权、支付、流动、复核和修订
作出安排。规范治理的内部关键，是不以新增物质产出或简单工时为核心的参与--分配行业，能否对参与事实、
贡献比较与收入份额作出可重放、横向一致、允许纠错的相对公正判定；外部边界才是在可识别风险外溢不超过
阈值的条件下尽量释放合法进入、竞争与试验空间。治理不替市场主体选择具体价格、交易对象、偏好或最终赢家，
也不以消灭利润、租金、合同剩余索取或贫富差距为目标。资本、组织者或高能力参与者依公开合同取得的正余额，按剩余索取权分析；仅在本文声明的分散投资--组织域、自愿交易、真实竞争、真实退出和正边际激励条件
同时成立时，后文才把它识别为风险承担、投资和能者多得的必要激励结构。该术语在本文定义内不承载负面道德
判断。当新型参与尚无稳定的资格、证据、收益权和争议处理结构时，
治理还可以先行建立市场得以运行的制度条件；这种构成性制度供给只生成非单点的合法选择空间，不预先指定
其中的价格、投入、偏好或配置结果。
本文另把生产自动化状态与治理自动化状态分开：产能约束松弛不推出审计、计量、反作弊、申诉或法律执行成本
自动下降。只有固定制度成本可以随规模摊薄、标准化记录与复算成本随治理自动化下降、不可自动化残余成本保持
受控时，才能确认单位治理成本与收入占比的上极限界；残余单位成本及人均收入另有极限时，才确认相应收敛。总治理成本可以随参与规模上升。
参与组织依合同把真实参与转化为可验证事实、可执行收益权与实际货币支付。收益权登记保存主体、金额、
来源、付款义务和履行状态；直接合同支付构成收入进入市场的基本路径。证券化、收益权转让及其他金融安排
仅在当事人选择且适用条件成立时构成附加分支。制度有效性取决于事实核验、资金覆盖、权利可执行、独立审计
与争议救济；线下赛事、协会、企业或普通服务平台均可在这些条件下组织参与。组织的登记与执行工具不决定
其分配规则是否公正，也不能替代真实收入、付款义务和结算证据。
这些治理环节是三项核心功能可持续实施的基础设施；它们提高可执行性，却不把制度工具本身置于分配目标之上。
本文区分狭义物质生产、广义生产、经济价值、市场价格、组织内支付与价值归属，并在附录中给出生产活动与围棋行业循环的条件性有向同伦表示。该扩展不承担预算、激励或收入闭合的证明。分析工具
来自生产理论、信息品经济学、机制设计、制度经济学、双边平台、公共决策、一般均衡与有向代数拓扑
\cite{AcemogluRestrepo2019,Varian1998InformationGoods,Myerson1981OptimalAuction,
Coase1937Firm,North1990Institutions,RochetTirole2003TwoSided,ArrowDebreu1954Equilibrium,Grandis2003}。

据此，本文研究的制度环节范围超出注意力经济、文化消费、平台分配、公共救济或组织治理中的任一单一部门投影，
并考察人工智能使生产能力与人类劳动参与分离以后，生产、收入、需求、参与、分配、救济、激励、文化、市场
构成与风险边界如何联合闭合。本文把“后 AI 基础经济范式候选”保留为待外部比较和同行评议的理论位置假设，
而不由系统内部自行确认。第~\ref{sec:theoretical-position}节给出声明制度环节覆盖与不可约性的条件性判据，并记录经验参数、
宏观规模、外部理论比较和学术共同体确认仍属独立证明责任。

本文不以日历年代、技术名称或生产成本绝对归零定义后 AI。固定有限非空产品域以后，只有每项产品通过精确产能
松弛或带预注册容限的渐近松弛核验记录，才能确认全域后 AI 适用状态；只有部分产品通过时，结论只投影到对应子域；
没有产品通过时，后 AI 总体结论不适用。正的能源、算力、维护、环境和责任成本与后 AI 状态相容，真正排除适用
的是这些成本或其他生产约束仍在边际上持续截断有限最优产量。局部治理和分配命题可以在域外单独检验，但不能
替代后 AI 适用域前件。

适用域通过也不授权整体替换现行制度。本文进一步把现行市场、法币、产权、合同、税收、公共供给、救济与监管
中仍然有效的坐标固定为声明基线；只有在传统生产--就业--劳动收入--消费传导的具体失灵点另经核验记录确认以后，
才允许从预注册候选族中选择能够修复该失灵点、保持其余兼容坐标且制度变更最小的路径。该路径称为“后 AI
条件触发后的最小兼容适用路径”。它可以与按劳分配、按需分配、基本收入、公共服务、合作社、数据分红及其他
经济机制组合；本文不由内部一致性推出自身唯一、全局最优或优于这些替代制度。

全文证明结构分为四个相互衔接而不可互相代替的层级。第一层由有限经济对象、可行集、最优值与影子价格说明
产能约束何时松弛、稀缺性向何处迁移，并保持生产--消费物质内核。第二层以生产者竞争性投入为起点，分别经
商品偏好与参与收入形成需求双通道，再加入兴趣品收入弹性、供给进入、文化沉积和动态稳定，判定局部产业规模
是否足以覆盖需求缺口。第三层把任一参与者的行为放入资格--计量--验证--估值--确权--支付链，并把
相对公正、初始权利、可选流动性、可选证券化和实际法币结算严格分层。第四层把逐活动留痕、跨活动来源保持、
内部公正、外部安全域、救济底盘、剩余激励、经验验证与停止条件合成为总核验记录，最后判定其条件性循环联合效应
和基础经济环节不可约性。因而本文不是由一个总公式覆盖全部问题，而是由局部定理、失败反例和制度环节核验记录组成
有向依赖图。会计恒等式、合同可行性、行为均衡和经验成立分别验收；第~\ref{sec:constructive-participation-economy}节给出连续投入基准，第~\ref{sec:contract-scale-extension}节求解含误差核验、组织合同和有限规模，第~\ref{sec:aggregate-income-closure}节计算资金机会成本、自动化收入变化和共同可行的总量实例。

在各定理已经列明的前提核验记录下，宏观与文化主路径写成
\[
\boxed{
\begin{aligned}
  \mathfrak C_{\mathrm{slack},j}=1
  &\Longrightarrow
  \CapacitySlack(j,A)
  \Longrightarrow Y_j(A,r)=Y_j^*,\\
  \mathfrak C_{\mathrm{bridge}}=1
  &\Longrightarrow
  \Pi_t^P\to\mathbf I_t^{B*}\to
  (\boldsymbol\vartheta_{t+1}^P,\mathbf T_t^B)
  \to D_{t+1}^P,
  \quad
  \frac{\partial D_{t+1}^P}{\partial\Pi_t^P}>0,\\
  \mathfrak C_{\mathrm{culture}}=1
  &\Longrightarrow
  Z_t^C\to C_t\to\mu_{C_t}\to W_h\to P_h\to D^I.
\end{aligned}}
\]
这些路径仍不混同经济价值、价格和支付：
\[
  \operatorname{PositiveMarginalContribution}_i
  \not\Longrightarrow \Paid_i,
  \qquad
  q_M>0\not\Longrightarrow\nu>0,
  \qquad
  \nu>0\not\Longrightarrow q_M>0.
\]

参与者层不使用
\(\Eligible\Rightarrow\Measured\Rightarrow\cdots\) 的裸蕴含，而使用同一对象、同一规则版本和连续来源引用的
逐项履约核验记录：
\[
\boxed{
\begin{aligned}
  \mathfrak C_i^{\mathrm{atom}}=1
  \Longrightarrow{}&
  \Eligible_i\wedge\Measured_i\wedge\Verified_i\\
  &{}\wedge\Valued_i\wedge\Entitled_i\wedge\Paid_i\\
  &{}\wedge\Auditable_i\wedge\Appealable_i.
\end{aligned}}
\]
其中每个谓词只由其对应的域条件、条件值和输出等式确认；任一前件单独都不推出下一后件。权利与法币路径
进一步分支为
\[
\boxed{
\begin{gathered}
  \mathfrak C_i^{\mathrm{part}}=1
  \wedge\mathfrak C_i^{\mathrm{claim}}=1
  \Longrightarrow b_i^0\Longrightarrow\beta_i,\\
  \beta_i
  \begin{cases}
    \xrightarrow{\text{直接合同、奖金、服务费或分成}}\Psi_i,\\
    \xrightarrow{\mathfrak C^{\mathrm{finance}}=1}\sigma(\beta_i)
    \xrightarrow{\text{受监管金融分支}}\Psi_i,
  \end{cases}\\
  \Psi_i\wedge\mathfrak C^{\mathrm{fiat}}=1
  \wedge y_i^F>0
  \Longrightarrow
  \text{实际法币到账}
  \Longrightarrow
  \text{项目外消费、储蓄或投资}.
\end{gathered}}
\]
所以证券化不是法币支付的必要前件，治理认证也不是正权利、正价格、流动性或偿付能力的充分前件。全文任何
蕴含都只在其量词域与核验记录前提内成立；逻辑核验记录、参数核验记录、因果核验记录、规模核验记录和外部有效性核验记录分别确认，
不得以技术相关性、法律许可、合同履行、文献稀缺或公式闭合替代现实成立。

章节安排与上述依赖图一一对应：第~\ref{chap:language-objects}章给出类型、谓词、制度算子、核验记录及证明义务；
第~\ref{chap:capacity-transfer}章刻画产能约束松弛、最优产出与稀缺性转移；
第~\ref{chap:market-go-homotopy}章建立生产--消费内核、兴趣消费、注意力双通道、文化沉积、救济底盘及围棋行业参照；
第~\ref{chap:governance-participation}章处理市场内制度选择、有效参与、相对公正与初始收益权；
第~\ref{chap:primary-finance}章区分初始确权、直接支付、可选流动性、条件证券化与金融规则的适用边界；
第~\ref{chap:platform-contribution}章形式化可验证参与计价、贡献分配、合同支付与可复算核验记录；
第~\ref{chap:activity-governance}章把活动规则、事实记录、跨活动收益权衔接、审计与申诉组合为组织治理条件；
第~\ref{chap:stability-synthesis}章最后给出动态稳定、福利边界、可证伪预测、总核验记录、条件性循环联合效应与停止条件。
前一章的结论只有在后一章明确引用其核验记录时才构成证明前件；章节相邻不自动产生逻辑蕴含。

\begin{remark}[适用条件]
自动化会改变任务分配与劳动需求，而非机械地产生“劳动不再稀缺”的普遍结论
\cite{AcemogluRestrepo2019}。信息产品具有高固定成本与低复制边际成本，但竞争结构、差异化、产权和
注意力仍影响价格与租金~\cite{Varian1998InformationGoods}。因此，本文所有“后 AI”结论均限定在
明确的产品域、技术域、制度域和核验记录前提内。
\end{remark}

\noindent\textbf{本地配套文献与引用约定。}
本仓库集中保存的配套资料包括 G-Framework 与 G-Algebra 的纯粹数学文献
\cite{Gao2025PureMathLocal}、应用数学三卷
\cite{Gao2025AppMathVol1Local,Gao2025AppMathVol2Local,Gao2025AppMathVol3Local}
及中文工作笔记 v5.2\cite{Gao2026NotebookLocal}。
这些本地资料统一位于 \code{docs/project_docs/}；参考文献所列本地路径均相对于仓库根目录，
并同时给出对应的 Zenodo DOI 与引用版本。本地文件版本与所列发布版本不同时，另在条目中注明。
本地条目提供配套理论文献与工作笔记的检索入口；各项经济结论的适用前提、证明与经验核验按正文各节逐项给出。

\clearpage
\tableofcontents
\clearpage
\pagenumbering{arabic}

\chapter{形式语言、经济对象与市场内制度算子}
\label{chap:language-objects}

\section{规范化核心陈述}

\begin{remark}[竞技主导的组织参与经济]
生产--消费是满足正物质需要并维持再生产的必要内核；这一宏观必要性不等于每项正价值活动都必须新增物质产出。
生产技术给出可行集合，分散市场回应形成交易路径、价格与剩余，组织合同在其内部收入池上选择资格、计量、
产权、转移和执行结构。生产--消费内核对参与、收入分配和新增需求的边际贡献衰减时，经济出现停滞风险而非
逻辑必然停滞。本文把核心功能规范化为
\(
\mathbf g^{\mathrm{core}}
=(g_{\mathrm{part}},g_{\mathrm{diff}},g_{\mathrm{relief}})
\)：参与式分配使每名被接纳且完成约定的有效参与者取得正基本份额，能者多得使同一比较域内更高的已验证
贡献取得更高分配，坐标最小救济使市场收入不足者达到生存--市场接入底线。核心组织以竞技活动汇聚注意力，
以教学、解说、创作、观赛和组织服务支持赛事及持续参与，多个生产者通过广告注意力竞争把既有盈利投入赞助。
资金池先覆盖正基本份额，再按可比较的已验证贡献分配差异份额；赛果不取消已经完成约定的基本份额。
商品购买形成覆盖产品成本的销售回款与销售利润，偏好引导产生的增量销售利润覆盖广告营销投入并扩大利润，
已实现净利润再进入下一期生产与赞助。规范治理首先检验参与--分配机制能否在声明
目标、规则版本和市场收入池内重放、由独立复算者一致复核并达到相对公正。兴趣消费的“广义奢侈品”属性另由
最低需要、超额预算、收入
弹性和付费需求判定；“广阔产业”及“足以支持循环”还分别要求进入规模、净余额、需求缺口覆盖和动态稳定，
不能从兴趣标签、注意力数量或单笔高价推出。经验证参与还须通过可重放的文化沉积、共同体承认与心理效用
传导，才能把短期刺激变成持续需求；稳定、可审计的底层规则与开放、内生的文化演化必须同时成立。公共治理
再给出非单点合法路径空间，并以风险外溢阈值约束该
空间。若市场收入跌破最低物质与市场接入束，并使需求和参与跌破循环阈值，则经济载体还须在跨期预算可持续且不
消灭正剩余激励的范围内提供坐标最小的充分救济。它保障生存和真实进入，却不配置下限以上收入。二者都不
替代需求、真实成本、价格发现或消费者选择，也不要求把全部合作剩余按某一唯一伦理基准分尽；正剩余索取
在正边际回报条件下保留投资、组织与能力激励。注意力偏好回流为三项核心功能提供竞争性资金与选择压力；
平台、企业、合作组织、公共救济或其组合都可以成为载体，故该功能结构不以某个中心一次设计为必要条件。
在新型参与缺少可执行市场身份时，资格、计量、验证、确权、支付、申诉和安全边界需要经制度构造进入
经济状态；构造的是分散选择能够发生的规则与路径集合，而不是由中心主体事先给出唯一市场结果。制度研究把规则视为塑造行动集合、
激励与集体结果的对象；机制设计则把“互动规则”作为实现给定结果的可设计映射
\cite{Ostrom1990Commons,Myerson1981OptimalAuction}。
\end{remark}

\section{对象域与基本经济条件}

\begin{definition}[经济状态]
令参与者集合、产品集合与行动空间分别为
\[
  \Part=\{1,\ldots,n\},\qquad
  \Goods=\{1,\ldots,m\},\qquad
  \Act=\prod_{i\in\Part}\Act_i.
\]
令技术状态为 \(A\in\RR_+\)，真实资源状态为 \(r\in\RR_+^d\)，一般法律状态为
\(\ell\in\Law\)，市场内组织制度状态为 \(G\in\Gov\)。
定义经济状态
\[
  x=(A,r,\ell,G,a,q,y)\in\States,
\]
其中 \(a\in\Act\) 为参与行为，\(q\in\RR_+^m\) 为产出，\(y\in\RR^n\) 为净资源请求权。
\end{definition}

\begin{definition}[生产可行集与市场内制度机制]
技术--资源可行集定义为
\[
  \Feas(A,r)
  =\{q\in\RR_+^m:F(q;A,r)\le0\}.
\]
市场内组织、平台、赛事、企业或合作项目的制度机制写成八元组
\[
  G=(E,M,V,W,\phi,T,L,Q),
\]
对每个 \(i\in\Part\)，令测量误差、计量信号、真实性证据、组织权重、初始收益权、支付结果、
可转移权利和历史记录空间依次为
\[
  \mathcal E_i,\quad \mathcal Z_i,\quad \Omega_i,
  \quad \mathcal W_i\subseteq\RR_+,
  \quad \mathcal B_i\subseteq\RR_+,
  \quad \mathcal Y_i^{\mathrm{pay}},
  \quad \mathcal B_i^{\mathrm{tr}},
  \quad \mathcal H_G.
\]
令合格行为域与带证据的已验证信号域为
\[
  \Act_i^{E,G}=\{a_i\in\Act_i:E_G(i,a_i)=1\},
  \qquad
  \mathcal Z_i^{V,G}
  =\{(z_i,\omega_i)\in\mathcal Z_i\times\Omega_i:
  V_G(z_i,\omega_i)=1\}.
\]
七个算子在各自声明域上具有类型
\[
\begin{aligned}
  E_G&:\Part\times\Act_i\longrightarrow\{0,1\},&
  M_G&:\Act_i\times\mathcal E_i\longrightarrow\mathcal Z_i,\\
  V_G&:\mathcal Z_i\times\Omega_i\longrightarrow\{0,1\},&
  W_G&:\mathcal Z_i\longrightarrow\mathcal W_i,\\
  \phi_G&:\mathcal W_i\longrightarrow\mathcal B_i,&
  T_G&:\mathcal B_i\longrightarrow\mathcal Y_i^{\mathrm{pay}},\\
  L_G&:\mathcal B_i\rightharpoonup\mathcal B_i^{\mathrm{tr}},&
  Q_G&:\mathcal H_G\longrightarrow\{0,1\}^{2}.
\end{aligned}
\]
其中 \(E\) 为资格规则，\(M\) 为计量规则，\(V\) 为真实性验证规则，\(W\) 是由已验证信号生成组织权重的
估值算子。
\(\phi\) 为初始收益权生成规则，\(T\) 为支付规则，\(L\) 为可选流动性规则，\(Q\) 输出审计与申诉条件。
符号 \(\rightharpoonup\) 表示只有满足适用法律和合同前提时流动性分支才定义；它不属于初始支付的必经路径。
组织内分配映射为
\[
  \Gamma_G:\Feas(A,r)\times\Act\longrightarrow\RR_+^n.
\]
这里的 \(G\) 只支配声明组织的资格、权利和内部支付，不等同于公共监管算子，也不预先指定整个市场的
产量、价格或个体配置；公共监管及其合法路径空间由第~\ref{chap:market-go-homotopy}章另行定义。
\end{definition}

\begin{definition}[基本经济条件及其因果语义]
对 \(i\in\Part\) 与 \(a_i\in\Act_i\)，定义
\[
\begin{aligned}
  \Eligible_i(a_i;G)&\Longleftrightarrow E_G(i,a_i)=1,\\
  \Measured_i(a_i,\varepsilon_i,z_i;G)&\Longleftrightarrow
  a_i\in\Act_i^{E,G}\ \wedge\ z_i=M_G(a_i,\varepsilon_i)\in\mathcal Z_i,\\
  \Verified_i(z_i;G)&\Longleftrightarrow V_G(z_i,\omega_i)=1,\\
  \Valued_i(z_i;G)&\Longleftrightarrow w_i=W_G(z_i)\ge0,\\
  \Entitled_i(w_i;G)&\Longleftrightarrow b_i=\phi_G(w_i)\ge0,\\
  \Paid_i(b_i;G)&\Longleftrightarrow y_i=T_G(b_i)\text{ 已执行},\\
  \Auditable_i(h_i;G)&\Longleftrightarrow \pi_1Q_G(h_i)=1,\\
  \Appealable_i(h_i;G)&\Longleftrightarrow \pi_2Q_G(h_i)=1.
\end{aligned}
\]
式中 \(\varepsilon_i\) 为测量误差，\(\omega_i\) 为身份、来源或行为真实性证据。
这里的 \(\Valued_i\) 专指规则 \(G\) 已赋予组织识别权重，不等同于第~\ref{chap:market-go-homotopy}章
相对于社会剩余基线定义的边际经济价值 \(\nu_i\)。组织估值可以漏认、错认或拒绝支付已经存在的经济价值。
本文的谓词因果箭头
\[
  P\Longrightarrow_G Q
\]
表示“在规则 \(G\) 与明确列出的附加前提下，谓词 \(P\) 蕴含谓词 \(Q\)”；它不表示经验相关性，
也不省略中间条件。
\end{definition}

\begin{theorem}[资格至支付的逐项履约链的充分核验记录与不可裸推定理]
\label{thm:typed-atomic-execution-chain}
对固定 \(i\) 和规则版本 \(G\)，令
\[
  \mathfrak C_{i}^{\mathrm{atom}}
  =\left\langle
  \begin{gathered}
  i,a_i,\varepsilon_i,\omega_i,z_i,w_i,b_i,y_i,h_i;\\
  E_G,M_G,V_G,W_G,\phi_G,T_G,Q_G
  \end{gathered}
  \right\rangle
\]
并要求同一见证依次满足
\[
\begin{gathered}
  E_G(i,a_i)=1,
  \qquad z_i=M_G(a_i,\varepsilon_i)\in\mathcal Z_i,
  \qquad V_G(z_i,\omega_i)=1,\\
  w_i=W_G(z_i)\in\mathcal W_i,
  \qquad b_i=\phi_G(w_i)\in\mathcal B_i,
  \qquad y_i=T_G(b_i)\in\mathcal Y_i^{\mathrm{pay}},\\
  Q_G(h_i)=(1,1),
  \qquad
  h_i\text{ 连续引用 }(a_i,z_i,\omega_i,w_i,b_i,y_i)\text{ 及规则版本}.
\end{gathered}
\]
则该核验记录充分推出
\[
  \Eligible_i\wedge\Measured_i\wedge\Verified_i
  \wedge\Valued_i\wedge\Entitled_i\wedge\Paid_i
  \wedge\Auditable_i\wedge\Appealable_i.
\]
反之，任一单独基本经济条件只确认本坐标，不推出下一坐标；下一步只有在其算子定义域、条件、输出等式和同一见证
引用均成立时才可确认。流动性算子 \(L_G\) 若适用，是从 \(b_i\) 出发的条件分支，不是生成 \(b_i\) 或
执行初始支付的前提。
\end{theorem}

\begin{Certificate}[履约链域--值--来源三重闭合]
\label{cert:typed-atomic-execution-chain}
验收程序对每一步同时检查：上一输出属于下一算子的定义域；实际值等于声明算子的输出；记录中的前序哈希、
签名或等价引用指向同一行为和同一规则版本。任何一项失败，核验记录停在最后一个已通过坐标并记录
\(\operatorname{Rejected}\) 或 \(\operatorname{AppealPending}\)，不得把后续谓词补写为真。
\end{Certificate}

\begin{proof}[证明（有限函数复合与反例）]
核验记录第一式给出 \(\Eligible_i\)。由于 \(a_i\in\Act_i^{E,G}\) 且
\(M_G:\Act_i\times\mathcal E_i\to\mathcal Z_i\)，第二式给出同一见证的 \(\Measured_i\)。其后
\(V_G(z_i,\omega_i)=1\) 给出验证条件，而 \(W_G:\mathcal Z_i\to\mathcal W_i\)，故依次得到 \(\Verified_i\)、\(\Valued_i\)、
\(\Entitled_i\) 与 \(\Paid_i\)。连续历史引用和 \(Q_G(h_i)=(1,1)\) 再给出可审计与可申诉，合取结论成立。

不可裸推由逐层反例得到：合格行为可以因测量失败而没有 \(z_i\)；已有测量可以因伪造证据而未验证；已验证
信号可以被 \(W_G\) 赋零权重；正权重可以被 \(\phi_G\) 赋零请求权；正请求权可以因支付池或履约失败而未支付；
已支付记录也可能因历史缺失而不可审计或申诉。因此任何前一谓词单独都不是后一谓词的充分条件。证毕。
\end{proof}

\section{经济对象的有限构造}

\begin{theorem}[市场内参与分配对象存在]
\label{thm:governance-economy-object-exists}
若 \(\Part,\Goods,\Act,\Feas(A,r)\) 非空，八个组织制度算子
\(E,M,V,W,\phi,T,L,Q\) 均为其声明域上的可定义映射，则
\[
  \Mech_{\mathrm{MP}}
  \DefEq
  \Tuple{\Part,\Goods,\Act,\Feas,E,M,V,W,\phi,T,L,Q}
\]
是一个可定义的市场内参与分配对象。
\end{theorem}

\begin{Certificate}[有限构造：对象层级递归闭合]
\label{cert:governance-economy-object-exists}
令 \(P(k)\) 表示前 \(k\) 层对象已经构造：
\[
\begin{aligned}
P(0)&:\Part,\Goods\text{ 已给定},&
P(1)&:\Act\text{ 已给定},\\
P(2)&:\Feas(A,r)\text{ 已给定},&
P(3)&:E,M,V\text{ 已给定},\\
P(4)&:W,\phi,T,L\text{ 已给定},&
P(5)&:Q\text{ 已给定}.
\end{aligned}
\]
核验记录转移为
\[
  P(k)\wedge\operatorname{Defined}(O_{k+1})
  \Longrightarrow P(k+1),\qquad 0\le k<5.
\]
终端核验记录为 \(P(5)\Longrightarrow\operatorname{Defined}(\Mech_{\mathrm{MP}})\)。
\end{Certificate}

\begin{proof}[证明（有限递推）]
基例由
\[
  \operatorname{Finite}(\Part)\wedge\operatorname{Finite}(\Goods)
  \Longrightarrow P(0)
\]
成立。假设 \(P(k)\) 成立。因下一层算子 \(O_{k+1}\) 在已构造域上可定义，故
\[
  P(k)\wedge\operatorname{Defined}(O_{k+1})
  \Longrightarrow P(k+1).
\]
依次得到
\[
  P(0)\Longrightarrow P(1)\Longrightarrow P(2)
  \Longrightarrow P(3)\Longrightarrow P(4)\Longrightarrow P(5).
\]
因此十二元组中的每一坐标均有已声明的域、陪域和作用顺序，
\(\Mech_{\mathrm{MP}}\) 可定义。证毕。
\end{proof}

\section{技术边界与制度选择的逻辑分离}

\begin{theorem}[可行集合--选择点分离定理]
\label{thm:feasible-selection-separation}
若技术状态 \((A,r)\) 固定，且存在两个市场内制度机制 \(G_1,G_2\) 使
\[
  \Gamma_{G_1}(q,a)\ne\Gamma_{G_2}(q,a)
\]
对某个 \((q,a)\in\Feas(A,r)\times\Act\) 成立，则生产可行集不能唯一决定分配结果。
\end{theorem}

\begin{Certificate}[同一技术截面上的双制度见证]
核验记录为
\[
  \mathfrak C_{\mathrm{sep}}
  =\Tuple{A,r,q,a,G_1,G_2,
  \Gamma_{G_1}(q,a),\Gamma_{G_2}(q,a)}.
\]
验收谓词为
\[
  q\in\Feas(A,r)
  \wedge
  \Gamma_{G_1}(q,a)\ne\Gamma_{G_2}(q,a).
\]
\end{Certificate}

\begin{proof}[证明]
固定技术与资源得到
\[
  (A,r)\text{ 固定}
  \Longrightarrow
  \Feas(A,r)\text{ 固定}.
\]
再取核验记录中的同一个 \((q,a)\)，则
\[
  G_1\ne G_2
  \Longrightarrow
  \Gamma_{G_1}(q,a)\ne\Gamma_{G_2}(q,a).
\]
若 \(\Feas(A,r)\) 能唯一决定分配，则同一 \(q\) 必有唯一 \(y\)，与双制度见证矛盾。
故技术给出可行边界，但不能唯一给出组织内分配点。证毕。
\end{proof}

\begin{lemma}[法律许可的必要而不充分引理]
\label{lem:permission-not-sufficient}
设新型参与行为 \(a_i\) 的持续供给需要
\[
  \Eligible_i(a_i;G)=1,\quad
  \Verified_i(z_i;G)=1,\quad
  \EE[T_G(\phi_G(W_G(z_i)))]\ge c_i(a_i),\quad
  D_G>0,\quad I_G=1,
\]
其中 \(D_G\) 表示有效需求，\(I_G\) 表示履约基础设施可用。则法律或规则许可
\(\Eligible_i=1\) 只是持续行业形成的必要条件之一，不是充分条件。
\end{lemma}

\begin{Certificate}[缺失后件反例]
取
\[
  \Eligible_i=1,\qquad D_G=0,qquad F_G^{\mathrm{external}}=0,
\]
并令声明支付规则只由当期经营需求收入与外部资金
\(F_G^{\mathrm{external}}\) 供资，即
\[
  D_G=0\wedge F_G^{\mathrm{external}}=0
  \Longrightarrow
  T_G(\phi_G(W_G(z_i)))=0.
\]
则在 \(c_i(a_i)>0\) 时，
\[
  \EE[T_G(\phi_G(W_G(z_i)))]-c_i(a_i)<0.
\]
该实例是“许可成立而持续供给失败”的见证。
\end{Certificate}

\begin{proof}[证明]
持续参与要求净收益非负：
\[
  \Sustainable_i(a_i)
  \Longrightarrow
  \EE[T_G(\phi_G(W_G(z_i)))]-c_i(a_i)\ge0.
\]
但
\[
  \Eligible_i=1\wedge D_G=0\wedge F_G^{\mathrm{external}}=0
  \wedge c_i(a_i)>0
  \Longrightarrow
  \EE[T_G(\phi_G(W_G(z_i)))]=0<c_i(a_i).
\]
因此 \(\Eligible_i=1\not\Longrightarrow\Sustainable_i(a_i)\)。证毕。
\end{proof}

\begin{proposition}[经济约束的乘积结构命题]
\label{prop:economic-product-structure}
在一般状态下，发展或福利结果应写成
\[
  W=\mathcal W\bigl(\Feas(A,r),\Omega_\ell,G,\operatorname{Exec}(G),D,\xi\bigr),
\]
其中 \(\Omega_\ell\) 是一般法律状态 \(\ell\) 定义的合法路径空间，\(G\) 是市场内组织制度，
\(\operatorname{Exec}(G)\) 是执行能力，\(D\) 是需求，\(\xi\) 是外部冲击。
除非另有退化条件，任一坐标均不能从该函数中删除。
\end{proposition}

\begin{Certificate}[坐标敏感性]
对每个坐标 \(x_k\in\{\Feas,\Omega_\ell,G,\operatorname{Exec},D,\xi\}\)，给定见证
\[
  \exists x_k',\quad
  \mathcal W(x_k',x_{-k})\ne\mathcal W(x_k,x_{-k}).
\]
若六类见证均存在，则 \(\mathcal W\) 对所有坐标都非退化。
\end{Certificate}

\begin{proof}[证明]
由坐标敏感性定义，
\[
  \exists x_k':\mathcal W(x_k',x_{-k})\ne\mathcal W(x_k,x_{-k})
  \Longrightarrow
  x_k\text{ 不能被常数替代}.
\]
逐坐标应用该蕴含，得到技术、法律可行域、组织制度、执行、需求与冲击都可能改变结果。证毕。
\end{proof}

\begin{corollary}[组织内分配规则敏感性推论]
\label{cor:distribution-governance-layer}
若在某产品域 \(j\) 上 \(\Feas(A,r)=\Feas^*\) 对 \(A\) 局部不变，并且存在可微规则路径
\(g\mapsto G(g)\) 及点 \(g_0\)，使
\(\left.\partial_g\Gamma_{G(g)}\right|_{g=g_0}\ne0\)，则
\[
  \frac{\partial \Gamma_G}{\partial A}=0,\qquad
  \left.\frac{\partial \Gamma_{G(g)}}{\partial g}\right|_{g=g_0}\ne0
\]
分别在该局部域和指定规则点成立。该结论只刻画声明组织内部的分配规则敏感性；市场路径、总剩余与
公共监管仍分别由分散选择、需求和成本以及合法路径空间决定。
\end{corollary}

\begin{Certificate}[偏导数域与非零见证]
核验记录由局部常值邻域 \(U_A\)、可微规则路径 \(g\mapsto G(g)\) 和非零导数见证点
\(g_0\) 构成，其中
\[
  \Feas(A,r)=\Feas^*\quad(A\in U_A),\qquad
  \left.\partial_g\Gamma_{G(g)}\right|_{g=g_0}\ne0.
\]
\end{Certificate}

\begin{proof}[证明]
局部常值直接给出
\[
  \Feas(A,r)=\Feas^*\Longrightarrow\partial_A\Feas=0.
\]
规则路径的非零导数见证给出
\(\left.\partial_g\Gamma_{G(g)}\right|_{g=g_0}\ne0\)。所以该域中被核验记录覆盖的个体分配边际变化来自组织制度选择
而非产能边界。证毕。
\end{proof}

\begin{remark}[本章核验记录不确认的结论]
本章不证明现实社会已经进入后 AI 状态，不证明任何具体法律必然产生行业，也不证明任一组织制度公平。
它只完成对象存在、技术边界与分配选择的逻辑分层，并为后续定理提供基本经济条件语言。
\end{remark}
\clearpage
\chapter{生产能力约束松弛与稀缺性转移}
\label{chap:capacity-transfer}

\section{单产品社会剩余模型}

\begin{definition}[产能、效用、真实成本与剩余]
固定一个可复制产品域 \(j\)。令 \(A\ge0\) 为 AI 相关生产能力状态，\(K_j(A)\) 为最大产能，
\(V_j(q)\) 为总效用，\(C_j(q;r)\) 为能源、算力、环境、维护和其他互补资源成本。定义
\[
  S_j(q;r)=V_j(q)-C_j(q;r),
\qquad
  Y_j(A,r)=\max_{0\le q\le K_j(A)}S_j(q;r).
\]
假设 \(S_j(\cdot;r)\) 连续、严格凹，并存在唯一有限无约束最优点
\[
  q_j^*(r)=\argmax_{q\ge0}S_j(q;r).
\]
人的注意力、使用时间与边际效用饱和可以使 \(q_j^*\) 有限；信息产品的低复制边际成本并不消除
固定成本、注意力约束或差异化市场结构~\cite{Varian1998InformationGoods}。
\end{definition}

\begin{definition}[约束状态谓词]
定义
\[
\begin{aligned}
  \Binding(j,A,r)
  &\Longleftrightarrow K_j(A)<q_j^*(r),\\
  \CapacitySlack(j,A,r)
  &\Longleftrightarrow K_j(A)\ge q_j^*(r).
\end{aligned}
\]
二者分别表示产能上界截断最优点与产能上界不再截断最优点。
\end{definition}

\begin{definition}[后 AI 产品域、精确适用与渐近适用]
\label{def:post-ai-applicability}
固定有限非空产品域 \(\Goods_0\subseteq\Goods\) 与容限 \(\epsilon>0\)。对 \(j\in\Goods_0\)，定义
\[
\begin{aligned}
  \AsymptoticSlack_\epsilon(j,A,r)
  \Longleftrightarrow{}&
  \Binding(j,A,r)
  \wedge
  \lim_{\widetilde A\to\infty}K_j(\widetilde A)=q_j^*(r)\\
  &{}\wedge S_j(\cdot;r),K_j\in C^1
  \wedge
  \sup_{\widetilde A\ge A}\Abs{K_j'(\widetilde A)}<\infty\\
  &{}\wedge
  \Abs{S_j'(K_j(A);r)K_j'(A)}\le\epsilon,
\end{aligned}
\]
并令
\[
  \Goods_{\mathrm{post}}^\epsilon(A,r)
  =
  \left\{
  j\in\Goods_0:
  \CapacitySlack(j,A,r)
  \vee
  \AsymptoticSlack_\epsilon(j,A,r)
  \right\}.
\]
其覆盖率为
\[
  \kappa_{\mathrm{post}}^\epsilon
  =
  \frac{\Abs{\Goods_{\mathrm{post}}^\epsilon(A,r)}}{\Abs{\Goods_0}}.
\]
精确后 AI 适用谓词与容限后 AI 适用谓词分别定义为
\[
\begin{aligned}
  \PostAI^{\mathrm{exact}}(\Goods_0,A,r)
  &\Longleftrightarrow
  \forall j\in\Goods_0,\ \CapacitySlack(j,A,r),\\
  \PostAI_\epsilon(\Goods_0,A,r)
  &\Longleftrightarrow
  \kappa_{\mathrm{post}}^\epsilon=1.
\end{aligned}
\]
当 \(\kappa_{\mathrm{post}}^\epsilon=1\) 时，后 AI 结论在整个声明产品域内适用；当
\(0<\kappa_{\mathrm{post}}^\epsilon<1\) 时，只允许在
\(\Goods_{\mathrm{post}}^\epsilon\) 上确认部分适用结论；当
\(\kappa_{\mathrm{post}}^\epsilon=0\) 时，不得在该产品域确认后 AI 总体结论。
\end{definition}

\begin{Certificate}[后 AI 适用域、覆盖率与成本边界]
\label{cert:post-ai-applicability}
定义
\[
  \mathfrak C_{\mathrm{post}}^{\mathrm{exact}}
  =
  \Tuple{
  \Goods_0,A,r,
  \{\mathfrak C_{\mathrm{slack},j}\}_{j\in\Goods_0}
  },
\]
以及
\[
  \mathfrak C_{\mathrm{post}}^\epsilon
  =
  \Tuple{
  \Goods_0,A,r,\epsilon,
  \{q_j^*,K_j,S_j,C_j\}_{j\in\Goods_0},
  \Goods_{\mathrm{post}}^\epsilon,
  \kappa_{\mathrm{post}}^\epsilon
  }.
\]
核验记录逐产品保存有限最优点、产能函数、真实成本、精确松弛见证或渐近松弛的链式法则见证。允许
\(C_j(q;r)>0\)：能源、算力、维护、环境、责任及其他互补成本可以始终存在。后 AI 前提要求的是复制产能不再
构成主要边际约束，而不是全部生产成本绝对归零。若成本或其他生产约束持续使
\(\Binding(j,A,r)\) 成立，且渐近松弛条件也失败，则该产品不属于
\(\Goods_{\mathrm{post}}^\epsilon\)。
\end{Certificate}

\begin{remark}[整体不适用不撤销局部制度命题]
资格、计量、验证、初始确权、支付、审计、申诉、二级转移守恒和风险约束等局部命题，可以在生产约束尚未松弛
时单独检验；但它们不能据此确认“后 AI 市场参与经济已经整体适用”。适用域核验记录回答前提是否成立，局部经验
核验记录回答机制在该域内是否成立，二者不得互相替代。
\end{remark}

\section{后 AI 条件触发后的最小兼容适用路径}

\begin{definition}[声明制度基线、失灵点与变更支集]
\label{def:minimal-compatible-baseline}
固定声明产品域 \(D\subseteq\Goods_0\)、制度坐标有限集 \(\mathcal U_0\) 与现行制度基线
\[
  \Sigma_0=(\sigma_u^0)_{u\in\mathcal U_0},
  \qquad
  \sigma_u^0=\pi_u(\Sigma_0)\in\mathcal X_u\cup\{\bot\}.
\]
这些坐标可以包括市场价格、法币结算、产权、合同、工资与按劳分配、税收、公共供给、救济、监管以及其他
已经存在的制度环节；原来不存在而可能由候选路径新增的制度环节以 \(\bot\) 记录。把某一坐标列入基线不等于预先
判定其有效或不可修改。令
\[
  \mathcal F_D
  =\Tuple{f_1,\ldots,f_m}
\]
为在 \(D\) 上经预注册证据确认的失灵点，例如传统劳动收入传导衰减、合格参与无法确权、参与收入不能支付、
注意力不能回接商品偏好或风险向制度外溢。对候选制度 \(\Sigma\)，定义变更支集与加权制度距离
\[
\begin{aligned}
  \Delta(\Sigma;\Sigma_0)
  &=\{u\in\mathcal U_0:\pi_u(\Sigma)\ne\pi_u(\Sigma_0)\},\\
  d_w(\Sigma,\Sigma_0)
  &=\sum_{u\in\mathcal U_0}w_u
  \mathbf 1\{\pi_u(\Sigma)\ne\pi_u(\Sigma_0)\},
  \qquad w_u>0.
\end{aligned}
\]
令 \(\mathcal R(\mathcal F_D)\subseteq\mathcal U_0\) 为修复失灵点直接需要改变的坐标，
\(\operatorname{Dep}(\mathcal R)\) 为经逻辑、预算、法律或技术依赖核验记录证明必须联动的坐标。候选制度保持基线
兼容，当且仅当
\[
  \Delta(\Sigma;\Sigma_0)
  \subseteq
  \mathcal R(\mathcal F_D)\cup
  \operatorname{Dep}(\mathcal R(\mathcal F_D)).
\]
因此，“兼容”不是永不修改现行制度，而是不修改与已证失灵及其必要依赖无关的坐标。
\end{definition}

\begin{definition}[后 AI 最小兼容候选族]
\label{def:minimal-compatible-family}
固定容限 \(\epsilon>0\)。令 \(\mathfrak A_D(\Sigma_0,\mathcal F_D)\) 为预注册的有限候选制度族，其中
每个 \(\Sigma\) 均满足：
\begin{enumerate}[label=\textnormal{A\arabic*:},leftmargin=3.5em]
  \item \(D\subseteq\Goods_{\mathrm{post}}^\epsilon(A,r)\)，且 \(\mathcal F_D\) 通过独立失灵核验记录；
  \item \(\Sigma\) 对 \(\mathcal F_D\) 的每个坐标给出可验证修复，并满足预算、资源、法律与技术可行性；
  \item \(\Sigma\) 保持定义~\ref{def:minimal-compatible-baseline} 的基线兼容，并保留非单点市场路径；
  \item \(\Sigma\) 不以排除其他分配理论或制度为可行条件；只有共同预算、资源、权利或法律制度环节发生实际冲突时，
  才把特定组合标记为不兼容。
\end{enumerate}
若 \(\kappa_{\mathrm{post}}^\epsilon=1\)，可取 \(D=\Goods_0\)；若
\(0<\kappa_{\mathrm{post}}^\epsilon<1\)，只能取
\(D\subseteq\Goods_{\mathrm{post}}^\epsilon\)；若 \(\kappa_{\mathrm{post}}^\epsilon=0\)，不得确认本定义的
后 AI 路径核验记录。
\end{definition}

\begin{definition}[最小兼容适用路径]
\label{def:minimal-compatible-path}
在 \(\mathfrak A_D(\Sigma_0,\mathcal F_D)\ne\varnothing\) 时，定义
\[
  \Sigma_D^*
  \in
  \argmin_{\Sigma\in\mathfrak A_D(\Sigma_0,\mathcal F_D)}
  \Bigl(
  \Abs{\Delta(\Sigma;\Sigma_0)},
  d_w(\Sigma,\Sigma_0)
  \Bigr),
\]
其中最小化按字典序进行：先使变更坐标数量最小，再在同数量候选中使预注册加权距离最小。若存在多个并列解，
则全部保留为最小兼容路径族，不由中心主体预先指定唯一赢家。
\end{definition}

\begin{theorem}[后 AI 条件触发后的最小兼容扩展定理]
\label{thm:minimal-compatible-extension}
若定义~\ref{def:minimal-compatible-family} 的 A1--A4 成立，且
\(\mathfrak A_D(\Sigma_0,\mathcal F_D)\) 有限非空，则：
\begin{enumerate}[label=\textup{(M\arabic*)},leftmargin=4em]
  \item 最小兼容路径族非空；
  \item 任一 \(\Sigma_D^*\) 都修复已声明失灵点，并保持与修复无关的基线坐标；
  \item 不存在 \(\widehat\Sigma\in\mathfrak A_D\) 使
  \(\Delta(\widehat\Sigma;\Sigma_0)\subsetneq\Delta(\Sigma_D^*;\Sigma_0)\)，也不存在变更坐标数相同而
  \(d_w(\widehat\Sigma,\Sigma_0)<d_w(\Sigma_D^*,\Sigma_0)\) 的候选；
  \item 对任何按需分配、基本收入、公共供给、按劳分配、合作社、数据分红或其他制度 \(H\)，若组合
  \(\Sigma_D^*\oplus H\) 通过共同预算、资源、权利、法律与制度环节一致性核验记录，则该组合仍属于允许制度域，因而
  \[
    \Sigma_D^*\not\Longrightarrow\neg H.
  \]
\end{enumerate}
因此本系统给出的是后 AI 条件和具体失灵共同触发的最小兼容扩展，不是对现行制度的整体替换，也不是对其他
经济理论的逻辑否定。
\end{theorem}

\begin{proof}[证明（有限候选集上的字典序极小与组合封闭）]
有限非空集合上的自然数坐标 \(\Abs{\Delta}\) 必有最小值；限制到取得该值的非空子集以后，有限个实数
\(d_w\) 仍有最小值，故 M1 成立。A2--A3 直接给出修复与基线兼容，故 M2 成立。若 M3 失败，则严格更小的
变更支集必有更少坐标，或在坐标数相同时存在更小加权距离，均与字典序最小定义矛盾。A4 没有把任何制度名称
作为排除条件；对给定 \(H\)，共同核验记录通过即表示组合没有实际资源、预算、权利、法律或制度环节冲突，所以组合仍
可行，M4 成立。该证明只给出声明候选族内、相对于 \(\Sigma_0\)、\(\mathcal F_D\) 与权重 \(w\) 的极小性，
不推出唯一性、全局福利最优或跨理论优越性。证毕。
\end{proof}

\begin{Certificate}[最小兼容适用路径]
\label{cert:minimal-compatible-path}
核验记录 \(\mathfrak C_{\mathrm{MCP}}\) 必须保存：后 AI 适用域与覆盖率、声明制度基线及版本、传统传导或其他
具体失灵点、修复坐标与必要依赖、有限候选族、逐候选预算与法律可行性、变更支集、预注册权重、字典序比较、
全部并列极小解、保留的基线制度环节，以及与按需分配、基本收入、公共供给、按劳分配、合作社、数据分红和其他
声明制度的组合测试。未提交较小候选的排除见证，不得声称“最小”；只证明内部一致而未比较制度基线，不得声称
“兼容”。
\end{Certificate}

\begin{remark}[相对最小性、适用停止与非排他边界]
“最小兼容”只相对于已声明的 \(D,\Sigma_0,\mathcal F_D,\mathfrak A_D,w\) 成立，不等于国家作用最小、规则
数量最少、道德要求最低或全局福利最优。若后 AI 适用域核验记录失败，或具体失灵点没有通过，不能以本节为由启动
整体制度扩展；若发现变更支集严格更小的可行候选，原核验记录立即失效。按需分配及其他制度可以独立成立、与本路径
互补或在共同指标下参与反事实比较；只有可证明的实际冲突限制组合，而不是理论名称本身限制组合。
\end{remark}

\section{有限最优产出的构造}

\begin{theorem}[有限最优点存在]
\label{thm:finite-unconstrained-optimum}
若 \(S_j(\cdot;r)\) 连续、严格凹，并满足
\[
  \lim_{q\to\infty}S_j(q;r)=-\infty,
\]
则存在唯一有限 \(q_j^*(r)\) 使
\[
  S_j(q_j^*;r)=\max_{q\ge0}S_j(q;r).
\]
\end{theorem}

\begin{Certificate}[有限构造：紧区间扩张]
\label{cert:finite-unconstrained-optimum}
令
\[
  q_N\in\argmax_{0\le q\le N}S_j(q;r),\qquad N\in\NN.
\]
由强制下降条件，存在 \(N_0\) 使
\[
  q>N_0\Longrightarrow S_j(q;r)<S_j(0;r).
\]
归纳谓词为
\[
  P(N):\quad N\ge N_0\Longrightarrow
  \argmax_{[0,N]}S_j=\argmax_{[0,N+1]}S_j.
\]
\end{Certificate}

\begin{proof}[证明（有限递推）]
连续性与紧致性给出
\[
  [0,N]\text{ 紧}\wedge S_j\text{ 连续}
  \Longrightarrow
  \argmax_{[0,N]}S_j\ne\varnothing.
\]
由 \(S_j(q;r)\to-\infty\)，可取 \(N_0\) 使所有 \(q>N_0\) 的值低于 \(S_j(0;r)\)。
故 \(N=N_0\) 时，新增区间 \((N_0,N_0+1]\) 不可能包含最大点，即 \(P(N_0)\) 成立。
若 \(P(N)\) 成立，则新增区间 \((N+1,N+2]\subset(N_0,\infty)\) 同样不包含最大点，故
\[
  P(N)\Longrightarrow P(N+1).
\]
数学归纳法推出所有 \(N\ge N_0\) 的紧区间最大点一致。严格凹性推出该最大点唯一，记为
\(q_j^*(r)\)。证毕。
\end{proof}

\section{精确阈值与渐近阈值}

\begin{theorem}[生产能力边际约束下降定理]
\label{thm:capacity-marginal-relaxation}
设 \(K_j(A)\) 单调不减，并存在 \(\bar A_j\) 使
\[
  K_j(\bar A_j)\ge q_j^*(r).
\]
则对一切 \(A\ge\bar A_j\)，有
\[
  q_j^{\mathrm{opt}}(A,r)=q_j^*(r),\qquad
  Y_j(A,r)=S_j(q_j^*(r);r).
\]
若 \(K_j\) 与 \(Y_j\) 在该域可微，则
\[
  \frac{\partial Y_j}{\partial A}(A,r)=0
  \qquad(A>\bar A_j).
\]
\end{theorem}

\begin{Certificate}[阈值、可行性与唯一最优点]
核验记录为
\[
  \mathfrak C_{\mathrm{slack},j}
  =\Tuple{\bar A_j,K_j(\bar A_j),q_j^*,S_j(q_j^*)}.
\]
验收谓词为
\[
  K_j\text{ 单调不减}
  \wedge K_j(\bar A_j)\ge q_j^*
  \wedge q_j^*=\argmax_{q\ge0}S_j(q;r).
\]
\end{Certificate}

\begin{proof}[证明]
对 \(A\ge\bar A_j\)，
\[
  K_j\text{ 单调不减}
  \Longrightarrow
  K_j(A)\ge K_j(\bar A_j)\ge q_j^*.
\]
因此
\[
  q_j^*\in[0,K_j(A)]
  \Longrightarrow
  \max_{0\le q\le K_j(A)}S_j(q;r)
  =S_j(q_j^*;r).
\]
右端不含 \(A\)，所以在可微点
\[
  \partial_A Y_j(A,r)=0.
\]
这证明的是指定产品域内的产能约束松弛，不是所有真实资源约束消失。证毕。
\end{proof}

\begin{lemma}[渐近松弛的充分条件引理]
\label{lem:asymptotic-capacity-relaxation}
若 \(K_j(A)<q_j^*\)、\(K_j(A)\uparrow q_j^*\)，且
\[
  S_j\in C^1,\qquad K_j\in C^1,\qquad
  \sup_A\Abs{K_j'(A)}<\infty,\qquad
  S_j'(K_j(A);r)\to0,
\]
则
\[
  \frac{\partial Y_j}{\partial A}
  =S_j'(K_j(A);r)K_j'(A)\to0.
\]
\end{lemma}

\begin{Certificate}[链式法则双因子界]
取 \(B_K=\sup_A\Abs{K_j'(A)}<\infty\)。对任意 \(\epsilon>0\)，存在 \(A_\epsilon\) 使
\[
  A\ge A_\epsilon
  \Longrightarrow
  \Abs{S_j'(K_j(A);r)}<\frac{\epsilon}{\max\{1,B_K\}}.
\]
\end{Certificate}

\begin{proof}[证明]
在 \(\Binding(j,A,r)\) 域，严格凹性给出受约束最优点 \(q=K_j(A)\)，故
\[
  Y_j(A,r)=S_j(K_j(A);r).
\]
链式法则推出
\[
  \Abs{\partial_A Y_j}
  =\Abs{S_j'(K_j(A);r)}\Abs{K_j'(A)}
  \le\Abs{S_j'(K_j(A);r)}B_K.
\]
代入核验记录界即得 \(\Abs{\partial_A Y_j}<\epsilon\)，故极限为零。证毕。
\end{proof}

\begin{proposition}[稀缺性转移命题]
\label{prop:scarcity-transfer}
设总生产系统包含可复制品 \(q\) 与互补稀缺要素
\[
  h=(h_{\mathrm{energy}},h_{\mathrm{compute}},h_{\mathrm{land}},
  h_{\mathrm{attention}},h_{\mathrm{identity}},h_{\mathrm{trust}}).
\]
即使 \(\CapacitySlack(j,A,r)\) 成立，只要某个约束
\[
  g_\ell(q,h)\le\bar h_\ell
\]
在最优点具有正拉格朗日乘子 \(\mu_\ell^*>0\)，该要素仍是边际稀缺要素。
\end{proposition}

\begin{Certificate}[KKT 正影子价格]
在适用的约束资格条件下，核验记录为 KKT 元组
\[
  \mathfrak C_{\mathrm{KKT}}
  =\Tuple{q^*,h^*,\mu^*,
  \nabla S-\sum_\ell\mu_\ell^*\nabla g_\ell,
  \mu_\ell^*(g_\ell-\bar h_\ell)}.
\]
若 \(\mu_\ell^*>0\)，则第 \(\ell\) 个约束在边际上具有正影子价格。
\end{Certificate}

\begin{proof}[证明]
KKT 互补松弛给出
\[
  \mu_\ell^*>0
  \Longrightarrow
  g_\ell(q^*,h^*)=\bar h_\ell.
\]
包络定理给出最优值对资源上界的边际响应
\[
  \frac{\partial Y^*}{\partial\bar h_\ell}=\mu_\ell^*>0.
\]
因此，复制产能上界的乘子可以降为零，而能源、注意力、身份或信任约束的乘子仍可为正。
稀缺性发生位置迁移，而非整体消失。证毕。
\end{proof}

\begin{corollary}[互补稀缺约束迁移推论]
\label{cor:constraint-governance-transfer}
当产能约束乘子 \(\mu_K\to0\)，而验证、身份、注意力或执行约束的乘子保持为正时，
新增验证、身份、合约执行等制度能力对声明最优社会剩余的边际作用可超过新增复制产能的边际作用。
\end{corollary}

\begin{Certificate}[边际作用比较]
若制度投入 \(g\) 放松约束 \(\bar h_\ell(g)\)，并满足
\[
  \mu_\ell^*\frac{\partial\bar h_\ell}{\partial g}
  >
  \mu_K^*\frac{\partial K_j}{\partial A},
\]
则制度投入对最优值的一阶效应更大。
\end{Certificate}

\begin{proof}[证明]
包络分解为
\[
  dY^*
  =
  \mu_K^*\,dK_j
  +\sum_\ell\mu_\ell^*\,d\bar h_\ell.
\]
当 \(\mu_K^*\to0\) 且核验记录中的制度投入项严格较大时，
\[
  \partial_gY^*>\partial_AY^*.
\]
因此边际瓶颈从复制产能转向市场主体与组织制度需要识别、验证或配置的互补稀缺要素。证毕。
\end{proof}

\begin{remark}[自动化任务边界]
自动化同时可能产生任务替代效应与新任务恢复效应~\cite{AcemogluRestrepo2019}。
因此，本章不能被解释为“AI 必然使一切劳动失去经济作用”。它只证明：在已满足阈值核验记录的指定产品域，
进一步扩产的边际作用下降，而稀缺与分配问题迁移到需求、互补资源及市场内制度条件。
\end{remark}
\clearpage
\chapter{生产--消费内核、兴趣消费、注意力桥接与参与价值}
\label{chap:market-go-homotopy}

\section{经济包含生产而不被生产穷尽}

\begin{definition}[经济活动、广义生产与狭义物质生产]
令 \(\Act_E\) 为交换、分配、消费、投资、风险承担以及生产在内的全部经济行动集合，
\(\Act_B\) 为产生商品、服务、体验或信息的广义生产行动，\(\Act_M\) 为产生物质产出的狭义生产行动。
本文采用层级
\[
  \Act_M\subseteq\Act_B\subsetneq\Act_E.
\]
严格包含来自至少一种不直接生成新产品的经济行动，例如既有请求权的交换、既有剩余的分配或消费选择。
围棋对局、教学、赛事和解说不属于 \(\Act_M\)，但只要产生可消费的服务或体验，就属于 \(\Act_B\)。
\end{definition}

\begin{definition}[经济价值、市场价格与支付]
固定状态 \(a=(a_i,a_{-i})\)，令
\[
\mathcal W(a)
  =\sum_{k\in\Part}u_k(a)-C(a)-H(a)
\]
为已把真实资源成本 \(C\) 与可识别外部损害 \(H\) 计入后的社会剩余。这里假设各项效用已经转换为同一
货币等价或其他可加基数尺度；若缺少共同尺度或人际可比假设，\(\mathcal W\) 与下述边际量不能被经验识别。
以不参与行动 \(a_i^0\) 为基线，
并以全部相关参与均不发生的声明基线 \(a^0\) 为活动基线，定义活动净经济价值与参与者边际经济价值
\[
  \nu(a)=\mathcal W(a)-\mathcal W(a^0),
  \qquad
  \nu_i(a)
  =\mathcal W(a_i,a_{-i})-\mathcal W(a_i^0,a_{-i}).
\]
令 \(q_M(a)\ge0\) 为该活动新增的物质产出。
另以 \(p_i\) 表示市场报价或成交价格，以 \(T_i\) 表示参与者实际取得的支付。本文把
\(\nu_i\)、\(p_i\) 与 \(T_i\) 分开：前者是相对于声明基线的净剩余增量，后二者分别是交换信号与
分配结果。该定义是一项可审计的操作口径，不声称货币价格穷尽文化价值、人格价值或全部社会价值。
\end{definition}

\begin{theorem}[物质生产与正经济价值的非充分非必要定理]
\label{thm:material-production-neither}
在上述定义下，存在经济状态使
\[
  q_M(a)>0
  \quad\text{而}\quad
  \nu(a)\le0,
\]
也存在经济状态使
\[
  q_M(a)=0
  \quad\text{而}\quad
  \nu(a)>0.
\]
因此，对单项活动的正经济价值而言，狭义物质生产既非充分条件，也非必要条件。
\end{theorem}

\begin{Certificate}[双反例核验记录]
充分性反例取一单位物质产出、总使用效用为零、真实生产与处置成本为 \(c>0\)，则
\[
  q_M=1,\qquad \Delta\mathcal W=0-c=-c<0.
\]
必要性反例取一场不增加物质存量的围棋服务，两名参与者各得效用 \(u>0\)，组织与时间成本为
\(0<c<2u\)，则
\[
  q_M=0,\qquad \Delta\mathcal W=2u-c>0.
\]
\end{Certificate}

\begin{proof}[证明（反例构造）]
第一份核验记录满足物质产出为正而净剩余为负，故
\[
  q_M>0\not\Longrightarrow \nu>0.
\]
第二份核验记录满足物质产出为零而净剩余为正，故
\[
  \nu>0\not\Longrightarrow q_M>0.
\]
两个反例共同否定充分性与必要性。结论只针对狭义物质生产；若采用包含服务与体验的广义生产定义，
第二个反例本身就是服务生产。它也不否定社会总体经济需要物质基础。证毕。
\end{proof}

\begin{theorem}[价值、价格与支付分离定理]
\label{thm:value-price-payment-separation}
若参与者对社会剩余有正边际贡献，则其参与具有本文定义下的正经济价值，但不必已经获得支付：
\[
  \nu_i>0\not\Longrightarrow T_i>0.
\]
反之，补贴、纯转移、操纵或错误计量可以产生正支付而不产生正净剩余：
\[
  T_i>0\not\Longrightarrow \nu_i>0.
\]
\end{theorem}

\begin{Certificate}[未支付贡献与无贡献支付核验记录]
未支付贡献核验记录为
\[
  \Tuple{a_i,a_i^0,\mathcal W(a_i,a_{-i}),\mathcal W(a_i^0,a_{-i}),T_i=0},
  \qquad \nu_i>0.
\]
无贡献支付核验记录为一项预算内纯转移 \(T_i>0\)，同时保持总使用效用与真实成本不变，使
\(\nu_i=0\)。
\end{Certificate}

\begin{proof}[证明（逻辑独立性）]
第一份核验记录直接见证 \(\nu_i>0\wedge T_i=0\)，第二份核验记录直接见证
\(T_i>0\wedge\nu_i=0\)。因此价值存在、价格发现、初始确权和实际支付是四个不同谓词。
后续资格--计量--验证--估值--确权--支付链决定价值如何被识别和分配，不反向决定价值是否已经由参与产生。
证毕。
\end{proof}

\section{生产--消费内核及其边际传导衰减}

\begin{definition}[生产--消费内核与参与--分配桥接]
令 \(q_t^M\) 与 \(c_t^M\) 分别表示第 \(t\) 期物质生产与最终物质消费，\(s_t^M\) 为可用存量，
\(\delta\in(0,1]\) 为折旧率，物质平衡为
\[
  s_{t+1}^M=(1-\delta)s_t^M+q_t^M-c_t^M,
  \qquad s_t^M\ge0.
\]
令 \(\Pi_t^P\) 为生产部门已经由市场交易实现的净盈利，\(Z_t^B\) 为不以新增物质产出为核心的
参与--分配行业的有效规模；\(Z_t^B\) 同时记录通过验证的参与、服务、注意力与实际收入分配，不把虚假流量
计入。定义生产盈利对桥接行业的边际传导系数
\[
  \mu_t^{PB}
  =\frac{\partial Z_{t+1}^B}{\partial\Pi_t^P},
\]
以及桥接行业对下一期生产品有效需求的边际回流系数
\[
  \chi_t^{BP}
  =\frac{\partial D_{t+1}^P}{\partial Z_{t+1}^B}.
\]
这里的参与--分配行业可以生产赛事、内容、体验、教学或社群等服务，只要求 \(q_{B,t}^M=0\)；
“不以生产为核心”只排除新增物质产出，不否认服务生产，也不否认组织内部存在收入分配。
\end{definition}

\begin{theorem}[宏观物质内核与微观价值非必要性的相容定理]
\label{thm:material-core-compatibility}
假设存在持续的正物质需要 \(c_t^M\ge\underline c>0\)，初始存量有限，且 \(\delta>0\)。则任一无限期
满足该需要的可行经济路径都必须在无限多个时期具有 \(q_t^M>0\)，并在每期具有 \(c_t^M>0\)。因此，
生产--消费是维持宏观物质再生产的必要内核；该结论与定理~\ref{thm:material-production-neither} 所证明的
“单项活动取得正经济价值不以新增物质产出为必要条件”同时成立。
\end{theorem}

\begin{Certificate}[物质存量耗尽核验记录]
给定 \(s_0^M<\infty\)、\(\delta>0\) 与 \(c_t^M\ge\underline c>0\)。反设存在 \(T\)，使
\(q_t^M=0\) 对所有 \(t\ge T\) 成立。核验记录记录从 \(T\) 开始的存量递推及首次出现
\(s_t^M<\underline c\) 的时期。
\end{Certificate}

\begin{proof}[证明（物质平衡与反证）]
若 \(q_t^M=0\) 对所有 \(t\ge T\) 成立，则
\[
  s_{t+1}^M\le(1-\delta)s_t^M-\underline c.
\]
有限初始存量在有限期后不足以支持 \(c_t^M\ge\underline c\)，并使非负存量约束失效，与无限期可行性
矛盾。因此物质生产不能最终永久消失；正物质需要又直接给出正消费。该证明是宏观持续路径结论，不把每项
文化、竞技或信息活动重新定义为物质生产。证毕。
\end{proof}

\begin{proposition}[边际传导衰减与停滞风险命题]
\label{prop:core-bridge-stagnation-risk}
设下一期实际活动变化可局部写成
\[
  dY_{t+1}
  =\beta_C\,dC_{t+1}^P
  +\beta_I\,dI_{t+1}^P
  +\beta_B\,dZ_{t+1}^B
  +dX_{t+1},
\]
其中 \(X\) 汇总净出口、公共支出及其他已声明通道。若产能已经松弛，生产盈利增加却同时满足
\[
  \frac{\partial C_{t+1}^P}{\partial\Pi_t^P}\to0,
  \qquad
  \frac{\partial I_{t+1}^P}{\partial\Pi_t^P}\to0,
  \qquad
  \mu_t^{PB}\to0,
  \qquad
  \frac{\partial X_{t+1}}{\partial\Pi_t^P}\to0,
\]
则 \(\partial Y_{t+1}/\partial\Pi_t^P\to0\)。在存在闲置产能、弱有效需求或就业收入收缩的联合证据时，
该状态确认“停滞风险”而非“必然停滞”结论。
\end{proposition}

\begin{Certificate}[边际传导与停滞风险核验记录]
核验记录必须同时记录产能松弛、生产盈利、居民收入与消费、生产投资、参与--分配行业规模、其他需求通道以及
相应局部导数。只有各传导项均接近零且存在需求或利用率弱化证据时，才允许报告风险；盈利集中、单期低增长或
参与不足中的任一现象单独都不充分。
\end{Certificate}

\begin{proof}[证明（全微分投影）]
对局部活动恒等式关于 \(\Pi_t^P\) 求导，得到
\[
  \frac{\partial Y_{t+1}}{\partial\Pi_t^P}
  =\beta_C\frac{\partial C_{t+1}^P}{\partial\Pi_t^P}
  +\beta_I\frac{\partial I_{t+1}^P}{\partial\Pi_t^P}
  +\beta_B\mu_t^{PB}
  +\frac{\partial X_{t+1}}{\partial\Pi_t^P}.
\]
在系数有限且四项边际传导都趋近于零时，右侧趋近于零。该代数结论只证明盈利增量失去活动传导；停滞还需
闲置能力、需求、就业、投资或产出证据，因此命题只确认风险。证毕。
\end{proof}

\section{计划配置、市场调节与监管的法律可行域}

\begin{definition}[计划配置与市场配置]
令经济路径为
\[
  \gamma=(a_0,p_0,y_0;a_1,p_1,y_1;\ldots),
\]
其中 \(a_t\) 是行动，\(p_t\) 是价格，\(y_t\) 是配置。若中心算子 \(\Psi_C\) 在参与者选择之前
直接给出关键产量、价格或配置，称为计划配置：
\[
  (q_t,p_t,y_t)=\Psi_C(x_t).
\]
若每个主体依据局部信息、价格和自身约束选择
\[
  a_{i,t}\in BR_i(p_t,a_{-i,t};\Info_{i,t}),
\]
价格与配置再由供求、成交、盈亏、声誉以及进入退出的联合反馈更新，称为市场配置。市场路径的结果不由
任何单一主体在路径开始前完整指定。分散知识与价格协调的经典讨论见~\cite{Hayek1945Knowledge}。
\end{definition}

\begin{definition}[监管型法律约束]
令 \(\Feas_\ell(x)\subseteq\Feas(A,r)\times\Act\) 为产权、合同、信息披露、竞争秩序、责任承担与
外部风险规则共同定义的逐期合法可行域。令 \(\Omega\) 为候选经济路径全集，并定义合法路径空间
\[
  \Omega_\ell
  =\{\gamma:(q_t,a_t)\in\Feas_\ell(x_t)\text{ 对每个 }t\}.
\]
监管算子
\[
  R_\ell:\Omega\longmapsto
  \{\operatorname{admit},\operatorname{block},\operatorname{remedy},
  \operatorname{internalize}\}
\]
只判断路径是否合法、是否需要救济或是否应内部化风险，不为市场主体选择具体产量、价格、交易对象或赢家。
\end{definition}

\begin{proposition}[市场调节与监管不替代命题]
\label{prop:market-regulation-nonsubstitution}
若 \(\Omega_\ell\) 至少包含两条不同的合法路径 \(\gamma_1,\gamma_2\)，且 \(R_\ell\) 对二者都返回
\(\operatorname{admit}\)，则监管规则不能唯一决定市场配置结果。若监管者进一步指定唯一的产量、价格或
个体配置，则它在该坐标上已从一般约束者转为直接配置者。
\end{proposition}

\begin{Certificate}[双合法路径核验记录]
给出
\[
  \gamma_1\ne\gamma_2,\qquad
  \gamma_1,\gamma_2\in\Omega_\ell,\qquad
  R_\ell(\gamma_1)=R_\ell(\gamma_2)=\operatorname{admit}.
\]
\end{Certificate}

\begin{proof}[证明（集合非单点性）]
双路径核验记录说明法律只把结果限制在含有至少两个元素的集合内。集合约束不能在没有附加选择算子的情况下
推出唯一元素。因此合法域与实际市场路径逻辑分离。反之，若规则把可行结果压缩为某一指定坐标的单点，
该规则就在该坐标上执行直接配置。证毕。
\end{proof}

\begin{remark}[“无形之手”的条件边界]
本文所谓非计划调节，是价格、供求、声誉、盈亏和进入退出共同生成未被中心预设的结果；它不是“没有法律”，
也不推出市场自动达到福利最优。外部性、市场势力、信息不对称、公共品和系统性风险仍可使市场结果偏离
社会最优。监管应守住一般规则并约束风险外溢，但不能据此冒充价格发现或替参与者作全部选择。
\end{remark}

\begin{definition}[参与市场的构成性制度算子]
\label{def:constitutive-institution-operator}
对尚未形成稳定参与--分配市场的状态 \(x\)，定义构成性制度算子
\[
  \mathcal K^{\mathrm{inst}}(x)
  =\bigl(G_x,\Omega_{\ell,x},R_{\ell,x}\bigr),
\]
其中 \(G_x=(E_x,M_x,V_x,W_x,\phi_x,T_x,L_x,Q_x)\) 分别给出资格、计量、验证、组织估值、确权、支付、
账本和申诉规则，\(\Omega_{\ell,x}\) 为产权、合同、竞争与责任规则限定的合法市场路径集合，
\(R_{\ell,x}\) 为监管型合法性与救济算子。称 \(\mathcal K^{\mathrm{inst}}\) 构造了一个非计划的参与市场，
当且仅当：
\begin{enumerate}[label=\textnormal{K\arabic*:},leftmargin=3.5em]
  \item 每条被允许的支付路径都能从参与证据依次重放到资格、计量、验证、确权、支付和申诉状态；
  \item \(\Omega_{\ell,x}\) 至少含两条不同的合法路径，且制度算子不在主体行动之前指定唯一的产量、价格、
  投入、偏好或个体配置；
  \item 实际路径由主体在局部信息、价格、盈亏、声誉及进入退出约束下分散选择；
  \item 内部分配与外部风险仍分别接受 \(\mathfrak C_{\mathrm{RPF}}\) 和 \(\mathfrak C_{\mathrm{safe}}\) 验收。
\end{enumerate}
这里的“构成”是人为建立可执行的市场基础设施，不是把市场结果预先写入制度算子。
\end{definition}

\begin{proposition}[构成性治理与市场选择相容命题]
\label{prop:constitutive-governance-market-compatible}
若定义~\ref{def:constitutive-institution-operator} 的 K1--K4 成立，则
\(\mathcal K^{\mathrm{inst}}\) 可以把原先无法稳定确权和结算的参与行为变为可进入、可比较、可支付和可救济的
候选市场路径，同时不唯一决定实际配置。若制度算子在主体选择前把相关产量、价格或个体配置压缩为单点，
则它在相应坐标上属于计划配置；若只建立规则而不存在真实需求、可分配收入或分散选择，则也不能据此确认
市场循环成立结论。
\end{proposition}

\begin{Certificate}[构成性制度供给核验记录]
\label{cert:constitutive-institution}
核验记录 \(\mathfrak C_{\mathrm{inst}}\) 保存规则版本、资格与证据字段、确权和支付链、申诉与救济路径、候选合法
路径全集、至少两条实际可行的不同路径，以及主体分散选择的记录。规则文本存在、单一行政试点运行或一笔
支付完成，均不足以替代该核验记录。
\end{Certificate}

\begin{proof}[证明（集合构造与选择分离）]
K1 使参与行为取得可执行的制度表示，K2 使合法结果集合保持非单点，K3 把集合内的实际路径交由分散主体
选择，K4 分别约束内部判定与外部损害。因此制度算子决定的是市场能够在哪些规则和路径内运行，而不是市场
最终选择哪一条路径。单点反例直接落入定义中的计划配置；需求或收入池为空的反例则说明制度可行性不蕴含
经济可持续性。证毕。
\end{proof}

\section{竞争性盈利注入、注意力传导与生产偏好指引}

\begin{definition}[注意力型参与--分配桥接行业]
令 \(\mathcal K=\{1,\ldots,K\}\) 为生产者集合，\(K\ge2\)，\(\mathcal H\) 为消费者集合，
\(\mathcal B\) 为赛事、体育、文化、游戏、内容或社群等参与--分配行业。生产者 \(k\) 从已经实现的盈利
\(\Pi_{k,t}^P\) 中选择注意力投入
\[
  I_{k,t}^B\in[0,\bar I_{k,t}],
  \qquad
  \bar I_{k,t}\le\Pi_{k,t}^P.
\]
桥接行业的收入池与可分配余额为
\[
  R_t^B=F_t^B+\sum_{k\in\mathcal K}I_{k,t}^B,
  \qquad
  B_t^B=R_t^B-C_t^B-H_t^B,
\]
其中 \(F_t^B\) 为门票、会员、教学和其他自愿交易收入，\(C_t^B\) 为真实成本，\(H_t^B\) 为已货币化的
外部损害，且 \(C_t^B\) 不重复计入参与者转移支付。组织合同只能在 \(B_t^B\ge0\) 时向有效参与者支付，
并满足
\[
  T_{i,t}^B\ge0,
  \qquad
  \sum_{i\in\Part_B}T_{i,t}^B\le B_t^B.
\]

令 \(a_{hk,t}^{\mathrm{att}}=\operatorname{Attn}_{hk}(I_{k,t}^B,I_{-k,t}^B,z_t)\) 为消费者 \(h\)
分配给生产者 \(k\) 的相对注意力，其中 \(z_t\) 汇总位置、媒介、披露与匹配环境，并假设在内点
\[
  a_{hk,t}^{\mathrm{att}}\ge0,
  \qquad
  \sum_{k\in\mathcal K}a_{hk,t}^{\mathrm{att}}\le1,
  \qquad
  \frac{\partial a_{hk,t}^{\mathrm{att}}}{\partial I_{k,t}^B}>0,
  \qquad
  \frac{\partial a_{hk,t}^{\mathrm{att}}}{\partial I_{\ell,t}^B}<0
  \quad(\ell\ne k).
\]
消费者对生产者及其产品的偏好状态按
\[
  \vartheta_{hk,t+1}^P
  =\Phi_{hk}(\vartheta_{hk,t}^P,a_{hk,t}^{\mathrm{att}},p_{k,t},q_{k,t},\zeta_{hk,t})
\]
更新，下一期生产品需求为
\[
  D_{k,t+1}^P
  =D_k^P(p_{k,t+1},\vartheta_{\cdot k,t+1}^P,Y_{\mathcal H,t+1}).
\]
其中 \(Y_{\mathcal H,t+1}\) 是消费者集合的可支配收入，\(\zeta_{hk,t}\) 是未由投入决定的商品质量、
既有声誉和外部信息状态。
因此，“偏好指引”专指消费者对投入方生产者及其商品的选择方向，不是消费者对桥接行业本身的偏好。
注意力是生产盈利、参与活动与生产商品需求之间的传导载体。
\end{definition}

\begin{definition}[竞技主导的组织域]
\label{def:competition-led-domain}
固定事前公开的赛事规则、结算期与组织合同。核心桥接行业由竞技事件集合及其支持网络构成：对局和赛事
提供共同的规则、比较结果与注意力中心，教学、解说、创作、观赛、审核和社区服务通过可追踪关系连接到
这些事件及其持续参与。竞技主导由事件关系、真实参与与注意力来源共同识别；付费推广以赞助标识、赛事
传播和产品联系进入注意力竞争。支持角色依其约定任务取得参与资格，竞技排名只决定相应贡献维度。
更一般的文化或社群行业沿用后文的抽象经济对象；本系统的核心循环在本竞技组织域内实例化。
\end{definition}

\begin{definition}[有效参与与资金覆盖的基本、差异分配]
\label{def:base-differential-distribution}
组织在接受参与以前声明任务、完成标准、身份规则、复核入口和最大接纳规模 \(N_t^{\max}\)。令
\(\mathcal U_t\) 为按此规则接纳的参与者集合，\(|\mathcal U_t|\le N_t^{\max}\)，定义
\[
  \mathcal E_t=
  \{i\in\mathcal U_t:E_{it}=1,\ V_{it}=1,\ \operatorname{Fulfilled}_{it}=1\},
  \qquad 1\le N_t:=|\mathcal E_t|\le N_t^{\max}.
\]
这里分别验证资格、真实身份与行为、约定完成情况；差异贡献得分不充当基本参与资格的替代门槛。
固定每名有效参与者的基本份额 \(\beta_t^{\mathrm{base}}>0\) 与差异池 \(D_t^{\mathrm{diff}}>0\)，并在
接受合同前以已到账、可用于本期支付的资金覆盖
\[
  N_t^{\max}\beta_t^{\mathrm{base}}+D_t^{\mathrm{diff}}\le B_t^B.
\]
组织成本、平台收费与已声明的共同负担已计入 \(C_t^B,H_t^B\)，不得在结算时再次冲减所承诺的基本份额。
资金不足时，组织调整尚未订立的下一期合同或接纳规模；已经完成约定者的份额仍是到期义务。

令 \(c_{it}\ge0\) 为同一规则域内可比较的已验证贡献，包含竞技表现、协作、专业服务或组织贡献。
固定严格递增函数 \(g_t\)，满足 \(g_t(0)=0\)。定义
\[
  \lambda_{it}=
  \begin{cases}
    \displaystyle\frac{g_t(c_{it})}{\sum_{j\in\mathcal E_t}g_t(c_{jt})},
       &\sum_{j\in\mathcal E_t}g_t(c_{jt})>0,\\[1.2ex]
    1/N_t,&\sum_{j\in\mathcal E_t}g_t(c_{jt})=0,
  \end{cases}
  \qquad
  T_{it}^B=\beta_t^{\mathrm{base}}+D_t^{\mathrm{diff}}\lambda_{it}.
\]
\(T_{it}^B\) 为组织合同的净应付分配；差异得分全为零时，各人仍取得基本份额，差异池按已公开的等份规则
支付。资格、贡献或支付的误判通过独立复核修复。无有效参与者的结算期不确认参与循环核验记录，预留资金按合同
留存或返还；未验证行为与复制身份不产生基本份额。
\end{definition}

\begin{theorem}[参与即分配与能者多得的资金覆盖定理]
\label{thm:base-differential-coverage}
在定义~\ref{def:base-differential-distribution} 的规则域内，对每个 \(i\in\mathcal E_t\) 都有
\[
  T_{it}^B\ge\beta_t^{\mathrm{base}}>0,
  \qquad
  \sum_{i\in\mathcal E_t}T_{it}^B
  =N_t\beta_t^{\mathrm{base}}+D_t^{\mathrm{diff}}\le B_t^B.
\]
对同一期、同一可比规则域内的 \(i,j\in\mathcal E_t\)，
\[
  c_{it}=c_{jt}\Longrightarrow T_{it}^B=T_{jt}^B,
  \qquad
  c_{it}>c_{jt}\Longrightarrow T_{it}^B>T_{jt}^B.
\]
因此完成约定的有效参与取得正分配，能力经已验证贡献表现为更多分配；基本份额覆盖普通参与，差异份额
维持能力、协作与专业投入的激励。应付额依约定结算期限兑现，实际到账另由支付最终性回执证明。
\end{theorem}

\begin{Certificate}[基本份额、差异份额与资金覆盖核验记录]
\label{cert:base-differential-coverage}
记核验记录为 \(\mathfrak C_{\mathrm{part}}^+\)。它保存赛事及支持关系、事前接纳名单与人数上限、完成标准、
真实性验证、基本份额、差异池、已到账资金及其用途限制、成本归属、贡献证据与比较规则、零差异处置、
逐人应付额、结算期限、实际回执与申诉结果。应付与到账分别核对，全体有效参与者均在覆盖名单内。
\end{Certificate}

\begin{proof}[证明（正基本份额与归一分配）]
两种得分状态均有 \(\lambda_{it}\ge0\) 且 \(\sum_i\lambda_{it}=1\)，故逐人正下界与总额恒等式成立。
由 \(N_t\le N_t^{\max}\) 得预算覆盖。若 \(c_{it}>c_{jt}\ge0\)，则总得分为正，并且
\[
  T_{it}^B-T_{jt}^B
  =D_t^{\mathrm{diff}}
    \frac{g_t(c_{it})-g_t(c_{jt})}{\sum_{\ell\in\mathcal E_t}g_t(c_{\ell t})}>0.
\]
等贡献给出相同权重；全零状态由等份规则得到同额。赛果或差异贡献为零均不改变已经满足的有效参与资格，
故不影响基本份额。若未预留足额资金，或允许事后任意扣减基本份额，则前述预算或正下界条件失败，不能确认
本定理核验记录。证毕。
\end{proof}

\begin{definition}[产品销售利润与偏好营销净收益]
\label{def:preference-marketing-return}
固定生产者 \(k\) 与投入至回款的归属窗口，所有跨期金额折算到同一计价时点。令
\[
  S_k^P=p_kq_k,\qquad
  \Pi_k^{P,\mathrm{sale}}=S_k^P-C_k^{P,\mathrm{prod}},\qquad
  M_k(I_k)=I_k+J_k(I_k),
\]
其中 \(C_k^{P,\mathrm{prod}}\) 为产品生产、交付及售后的完整归属成本，\(M_k\) 为该次注意力竞争的
全部增量营销成本，包含赞助 \(I_k\) 和创意、投放、执行等配套成本 \(J_k\)。取
\(J_k(0)=0\)、\(J_k\ge0\)，且在可微域内 \(J_k'\ge0\)。各成本只归属一次；赞助进入组织收入池，
组织从池内支出的赛事成本不再作为生产者的另一笔营销支出重复扣除。

产品销售回款的分解为
\[
  S_k^P=C_k^{P,\mathrm{prod}}+\Pi_k^{P,\mathrm{sale}}.
\]
在持续经营域内 \(\Pi_k^{P,\mathrm{sale}}\ge0\)，消费回流因而覆盖产品成本并形成销售利润。
固定竞争者投入、消费者收入、产品价格、质量及其他可比状态，以未实施本次营销的偏好状态为反事实，定义
\[
  \Delta\Pi_k^{P,\mathrm{pref}}(I_k;I_{-k})
  =\Pi_k^{P,\mathrm{sale}}(\vartheta_k(I_k,I_{-k}))
   -\Pi_k^{P,\mathrm{sale}}(\vartheta_k(0,I_{-k})),
  \qquad
  \Gamma_k^{\mathrm{mkt}}=\Delta\Pi_k^{P,\mathrm{pref}}-M_k(I_k).
\]
此处省略的需求和成本映射均保持同一归属口径。\(\Gamma_k^{\mathrm{mkt}}>0\) 表示偏好带来的增量销售
利润覆盖全部增量营销成本，并留下新增净利润。总销售额、既有销售利润与未扣产品成本的增量收入分别记录，
不能替代这一增量营销回报指标。下一期可用于赞助的 \(\Pi_{k,t+1}^P\) 是实际回款并扣除产品与营销成本后的
可支配已实现利润，投入上限服从其现金可用额。
\end{definition}

\begin{proposition}[消费回款与营销净增利润的分账恒等式]
\label{prop:sales-marketing-accounting}
记 \(P_k(\vartheta,Y)\) 为同一价格、质量、竞争者投入和归属成本下的产品销售利润函数。
令 \(\vartheta^0,\vartheta^1\) 分别为本次营销未实施与已实施的偏好状态，\(Y^0,Y^1\) 为本轮参与分配
前后的可支配收入状态，定义
\[
\begin{aligned}
  \Pi_k^{P,0}&=P_k(\vartheta^0,Y^0),\\
  \Delta\Pi_k^{P,\mathrm{inc}}
    &=P_k(\vartheta^0,Y^1)-P_k(\vartheta^0,Y^0),\\
  \Delta\Pi_k^{P,\mathrm{pref}}
    &=P_k(\vartheta^1,Y^1)-P_k(\vartheta^0,Y^1).
\end{aligned}
\]
按这一事前固定的分解顺序，偏好与收入的交互效应归入后一项，实际产品销售与营销净利润满足
\[
\begin{aligned}
  \Pi_k^{P,\mathrm{net}}
  &=S_k^P-C_k^{P,\mathrm{prod}}-M_k\\
  &=\Pi_k^{P,0}+\Delta\Pi_k^{P,\mathrm{inc}}
    +\Gamma_k^{\mathrm{mkt}},\\
  S_k^P&=C_k^{P,\mathrm{prod}}+M_k+\Pi_k^{P,\mathrm{net}}.
\end{aligned}
\]
若 \(P_k(\vartheta^0,Y^1)\ge0\) 且 \(\Gamma_k^{\mathrm{mkt}}>0\)，则实际产品回款同时覆盖产品成本和
营销成本，并形成正净利润；相对于同一收入状态下不实施本次营销的反事实，净利润增加
\(\Gamma_k^{\mathrm{mkt}}\)。同一笔销售只计一次，营销引起的品牌份额转移也不直接计作全市场新增消费。
\end{proposition}

\begin{proof}[证明（同一销售利润的有序分解）]
从实际销售利润扣除营销支出后，依次插入同一收入状态的偏好反事实与原始基线，得到
\[
  P_k(\vartheta^1,Y^1)-M_k
  =P_k(\vartheta^0,Y^0)+\Delta\Pi_k^{P,\mathrm{inc}}
    +\Delta\Pi_k^{P,\mathrm{pref}}-M_k.
\]
即得净利润恒等式。非负基线利润与正营销净增利润之和为正，移项得到成本
覆盖。各项均以已经扣除产品成本的销售利润计量，故无需再次扣产品成本；两项增益相加恰好等于同一销售
利润变化，故没有重复计算。恒等式的因果解释另依赖已声明反事实的识别证据。证毕。
\end{proof}

\begin{definition}[必要品载体、多选集与已实现选择份额投射]
令 \(\mathcal J_M\) 为声明期内不可消除的必要需要集合。对消费者 \(h\) 和需要
\(g\in\mathcal J_M\)，令 \(\bar q_{hg}>0\) 为必须满足的最低标准化数量，
\(\mathcal C_{hg}^{F}\) 为同时满足最低质量、合法可得、预算可行且能够承载同一需要的有限商品集合。具体
商品、品牌、技术路线和文化属性可以相互替代，必要需要本身不因此消失。

令 \(x_{hgc}\ge0\) 为消费者经分散选择实际配置到载体
\(c\in\mathcal C_{hg}^{F}\) 的标准化数量，并要求
\[
  \sum_{c\in\mathcal C_{hg}^{F}}x_{hgc}=\bar q_{hg}.
\]
定义已实现选择份额投射
\[
  \vartheta_{hgc}^{M}
  =\frac{x_{hgc}}{\bar q_{hg}},
  \qquad
  \boldsymbol\vartheta_{hg}^{M}
  \in\Delta\!\left(\mathcal C_{hg}^{F}\right).
\]
该投射只记录实际选择份额、选择概率或轮换频率。本文后文可将其简称为“弱偏好投射”，但该简称不要求消费者
在全部载体之间具有严格排序，也不把观察份额解释为心理偏好、脱离预算和可得性约束的标准显示偏好、基数效用
或福利大小。
\end{definition}

\begin{theorem}[正刚需下多选选择份额投射必然性]
\label{thm:essential-multichoice-projection}
固定通过后 AI 精确或渐近适用域核验记录的必要品子域。若对某一 \((h,g)\) 有
\[
  \bar q_{hg}>0,
  \qquad
  \left|\mathcal C_{hg}^{F}\right|\ge2,
\]
最低束在预算内可达，且制度没有在消费者行动以前指定唯一载体，则任一满足最低需要的已实现市场配置都唯一
诱导一个
\[
  \boldsymbol\vartheta_{hg}^{M}
  \in\Delta\!\left(\mathcal C_{hg}^{F}\right),
  \qquad
  \sum_{c}\vartheta_{hgc}^{M}=1,
  \qquad
  \max_c\vartheta_{hgc}^{M}>0.
\]
因此，刚需保证相关消费不归零，多选与分散选择保证非零选择份额必然投射到至少一个现实载体；本文只在上述
定义意义下将其简称为弱偏好投射。市场的任务是显露投射方向、集中度和跨期变化。该结论允许正生产成本；它只
要求对应载体位于已经通过
精确或容限适用核验记录的子域，并且最低配置真实可行。
\end{theorem}

\begin{proof}[证明（标准化与单纯形）]
由 \(\bar q_{hg}>0\) 和最低需要约束，已实现向量 \(x_{hg}\) 非负且总量严格为正。逐坐标除以
\(\bar q_{hg}\) 后得到非负向量 \(\boldsymbol\vartheta_{hg}^{M}\)，其坐标和为一，故属于对应概率单纯形；
坐标和为一又排除全部坐标同时为零。标准化映射对给定配置唯一。载体不少于两个使消费者面对真实多选集，
非单点制度条件使实际份额由分散选择而非中心预配形成。即使消费者近似无差异，随机选择、混合购买或跨期
轮换仍给出一个选择份额投射；没有附加行为识别时，不再推出更强的心理或福利解释。证毕。
\end{proof}

\begin{corollary}[选择份额投射存在与注意力诱导偏好迁移分离]
\label{cor:projection-versus-attention-shift}
令 \(\mathbf a^{\mathrm{att}}\) 为注意力状态，定义注意力诱导的偏好迁移为
\[
  \Delta_{\mathrm{att}}\boldsymbol\vartheta_{hg}^{M}
  =
  \boldsymbol\vartheta_{hg}^{M}(\mathbf a')
  -\boldsymbol\vartheta_{hg}^{M}(\mathbf a).
\]
定理~\ref{thm:essential-multichoice-projection} 保证基线投射存在，但不推出
\[
  \Delta_{\mathrm{att}}\boldsymbol\vartheta_{hg}^{M}\ne0
  \quad\text{或}\quad
  \frac{\partial\boldsymbol\vartheta_{hg}^{M}}
       {\partial\mathbf a^{\mathrm{att}}}\ne0.
\]
后两者才是注意力改变生产商品偏好方向或强度的传导命题，必须以可识别的市场变化、购买和复购数据检验。
因此，“选择份额投射在声明条件下必然存在”与“注意力传导系数为正”是两个不同层次的结论；前者只在正刚需、
多载体可行、预算可达和分散选择等适用条件持续成立期间提供偏好传导的现实载体，后者决定竞争性盈利注入能否
利用该载体形成内生市场回报。
\end{corollary}

\begin{Certificate}[刚需、多选与注意力迁移分层核验记录]
核验记录保存必要需要及其最低数量、载体的功能等价标准、质量门槛、价格、可得性、预算可行性、制度是否保留
非单点选择、实际购买份额或选择频率，以及注意力变化前后的考虑集、份额、购买和复购。只见
\(\boldsymbol\vartheta_{hg}^{M}\in\Delta(\mathcal C_{hg}^{F})\) 只能确认选择份额投射存在；只有排除价格、质量、
收入和同期促销等替代解释以后，才能确认注意力诱导偏好迁移。载体集合退化为单点、最低束不可达或由中心
预先指定唯一配置时，不得把已实现份额解释为多选市场选择份额投射。
\end{Certificate}

\begin{assumption}[竞争性注意力投入的正回报条件]
生产者 \(k\) 在固定竞争者投入及可比收入状态下的跨期净收益写为
\[
  V_k(I_k,I_{-k})
  =\Pi_{k,t+1}^{P,\mathrm{sale}}\!\left(D_k^P\!\left(
  \Phi_k\!\left(\operatorname{Attn}_k(I_k,I_{-k})\right)\right)\right)-M_k(I_k).
\]
策略集 \([0,\bar I_k]\) 满足 \(0<\bar I_k<\infty\)，因而非空、紧且凸；\(V_k\) 连续、关于 \(I_k\) 严格凹；唯一最优回应
\(BR_k(I_{-k})\) 连续。并且对相关 \(I_{-k}\) 有
\[
  \left.\frac{\partial V_k}{\partial I_k}\right|_{I_k=0}
  =
  \frac{\partial\Pi_k^{P,\mathrm{sale}}}{\partial D_k^P}
  \frac{\partial D_k^P}{\partial\vartheta_k^P}
  \frac{\partial\vartheta_k^P}{\partial a_k^{\mathrm{att}}}
  \frac{\partial a_k^{\mathrm{att}}}{\partial I_k}-M_k'(0)
  >0.
\]
令 \(L_{k,t}^{\mathrm{cash}}>0\) 为扣除已承诺用途后可用于本轮营销的已实现利润现金，则投入上限满足
\(M_k(\bar I_k)\le L_{k,t}^{\mathrm{cash}}\le\Pi_{k,t}^P\)。本节均衡使用声明需求模型中的收益；决策时的期望收益和
结算后的实际收益分别保存，后者须以真实回款、完整成本及可识别的偏好反事实独立验收。
\end{assumption}

\begin{theorem}[竞争性盈利注入的内生市场均衡定理]
\label{thm:competitive-attention-injection}
在竞争性注意力投入的正回报条件下，存在注意力投入纳什均衡
\(I^*=(I_1^*,\ldots,I_K^*)\)，且对每个 \(k\) 都有
\[
  I_k^*>0,\qquad
  \Gamma_k^{\mathrm{mkt}}(I_k^*;I_{-k}^*)
  =V_k(I_k^*,I_{-k}^*)-V_k(0,I_{-k}^*)>0.
\]
因此模型内的偏好增量销售利润覆盖赞助及配套营销成本，并形成净增利润。由于竞争者投入满足
\(\partial a_k^{\mathrm{att}}/\partial I_\ell<0\)，生产者不仅追求绝对曝光，也为避免相对注意力份额落后而持续回应
竞争者。该均衡由预期产品需求和盈利反馈诱导，不需要中心主体指定投入额、赞助对象或消费者偏好。
\end{theorem}

\begin{Certificate}[竞争性注意力投入核验记录]
核验记录保存紧凸策略域、收益函数、严格凹性、连续最优回应、边界正导数，以及至少两家生产者之间的交叉注意力
导数；分别登记销售回款、产品成本、销售利润、赞助与配套营销成本、偏好反事实和营销净增利润。
若只观察到单一生产者转账，或无法把注意力变化连接到其产品需求与预期盈利，则不能确认“竞争性市场
注入”结论。
\end{Certificate}

\begin{proof}[证明（最优回应不动点）]
各生产者严格凹的目标在紧区间上有唯一最优回应；连续性使联合最优回应映射
\[
  BR:\prod_k[0,\bar I_k]\longrightarrow\prod_k[0,\bar I_k]
\]
连续。由 Brouwer 不动点定理存在 \(I^*=BR(I^*)\)。边界导数严格为正排除 \(I_k^*=0\)，故每个均衡
投入为正。边界正导数还保证存在可行小投入使 \(V_k(I_k,I_{-k}^*)>V_k(0,I_{-k}^*)\)，而最优回应的
收益不低于该小投入；由 \(M_k(0)=0\) 及定义~\ref{def:preference-marketing-return}，即得
\(\Gamma_k^{\mathrm{mkt}}>0\)。交叉导数为负说明竞争者投入降低本生产者的相对注意力份额；但自身最优投入的反应符号取决于收益交叉偏导，不能由份额下降推出追加投入增加；均衡投入也不自动达到社会最优。证毕。
\end{proof}

\begin{proposition}[注意力挤出与最优投入反应的区分]
\label{prop:sponsor-response-sign}
设自身收益连续二次可微且 \(\pi_{k,kk}<0\)，最优回应在内点可微，则
\[
 \frac{dBR_k}{dI_j}=-\frac{\pi_{k,kj}}{\pi_{k,kk}}.
\]
因此回应方向与收益交叉偏导同号，不由注意力份额的交叉偏导决定。
在定义~\ref{def:constructive-market} 的连续投入模型内，
\[
 \pi_{k,kj}=\frac{B(I_k-a-I_j)}{(a+I_k+I_j)^3}.
\]
定理~\ref{thm:constructive-equilibrium-cycle} 的参数点给出 \(\pi_{k,kj}=-7/48\)、
\(\pi_{k,kk}=-7/12\)，故 \(BR_k'(1)=-1/4\)。边界解使用相应的单侧比较或最优回应对应，
不直接套用内点公式。
\end{proposition}
\begin{proof}[证明]
对一阶条件 \(\partial_{I_k}\pi_k(BR_k(I_j),I_j)=0\) 求导，移项即得回应式。
对连续模型的边际收益再求交叉偏导得到所列分式，代入参数即得数值。
\end{proof}

\begin{remark}[收入内生时的赞助收益]
前一定理固定消费者收入，因而最优回应的净收益差等于偏好营销净收益。
若赞助同时改变参与者收入，生产者的目标应包括由此产生的销售回报；最优回应的总收益差
不再等于固定收入下的偏好回报，须按收入与偏好分账分别检查。
定理~\ref{thm:constructive-equilibrium-cycle} 对收入内生情形同时求解均衡并验证正偏好回报。
\end{remark}

\begin{proposition}[偏好传导缺失时无内生注资命题]
\label{prop:no-guidance-no-market-injection}
若注意力投入对生产者没有其他可货币化收益，且在声明域内
\[
  \frac{\partial D_k^P}{\partial\vartheta_k^P}
  \frac{\partial\vartheta_k^P}{\partial a_k^{\mathrm{att}}}
  \frac{\partial a_k^{\mathrm{att}}}{\partial I_k}=0,
\]
则 \(\partial V_k/\partial I_k=-M_k'(I_k)\le-1\)，生产者唯一最优选择为 \(I_k^*=0\)。因此，不能形成消费者对
生产者商品偏好指引的盈利转账，不具有由市场“无形之手”持续推动的注意力投入动力；行政拨付、慈善、
一次性补贴或依靠后来者资金的转账必须另行分类。
\end{proposition}

\begin{proof}[证明（边际收益为零）]
由链式法则，投入的生产侧边际毛收益等于三个偏导数与利润--需求偏导数之积。所列乘积为零且不存在其他
可货币化收益时，\(M_k'=1+J_k'\ge1\)，故 \(\partial V_k/\partial I_k=-M_k'<0\)。严格递减目标在非负区间的
唯一最优点为零。证毕。
\end{proof}

\begin{theorem}[注意力与收入双通道桥接定理]
\label{thm:dual-channel-bridge}
令偏好通道与收入通道的局部增益分别为
\[
  \kappa_{k}^{\mathrm{pref}}
  =
  \left.
  \frac{\partial\sum_{\ell\in\mathcal K_0}D_\ell^P}{\partial I_k}
  \right|_{Y,p,q},
  \qquad
  \kappa^{\mathrm{inc}}
  =\sum_{i\in\Part_B}
  \frac{\partial C_i^P}{\partial T_i^B}
  \frac{\partial T_i^B}{\partial I^B},
\]
其中 \(\Part_B\subseteq\Part\) 为取得桥接行业支付的参与者，\(C_i^P\) 为参与者 \(i\) 对生产品的消费，
且 \(I^B=\sum_k I_k^B\)。\(\mathcal K_0\) 是事前声明的产品集合，偏好系数包括该域内所有品牌的直接效应与
交叉挤出效应；若研究全市场需求，还须纳入域外商品的替代变化。单个赞助商品的销量增加不直接等于这一总效应。
令 \(\Pi_t^P=\sum_k\Pi_{k,t}^P\)，并在固定价格、质量、外部收入与其他
需求状态的局部比较域内，假设约化需求的全微分可分解为
\[
  dD_{t+1}^P
  =\sum_k\kappa_k^{\mathrm{pref}}\,dI_k^{B*}
  +\kappa^{\mathrm{inc}}\,dI^{B*}.
\]
若所有 \(\kappa_k^{\mathrm{pref}}\ge0\)、\(\kappa^{\mathrm{inc}}\ge0\)，所有均衡投入响应
\(\partial I_k^{B*}/\partial\Pi_t^P\ge0\)，且至少一个正通道系数与严格正投入响应相乘为正，则
\[
  \frac{\partial D_{t+1}^P}{\partial\Pi_t^P}
  =
  \sum_k\kappa_k^{\mathrm{pref}}
  \frac{\partial I_k^{B*}}{\partial\Pi_t^P}
  +\kappa^{\mathrm{inc}}
  \frac{\partial I^{B*}}{\partial\Pi_t^P}
  >0.
\]
故竞争性注资一方面通过注意力形成消费者对生产商品的偏好方向，另一方面通过向参与者分配收入增加购买力，
使生产盈利重新取得对参与--分配与下一期消费的正边际贡献。
\end{theorem}

\begin{Certificate}[双通道桥接核验记录]
核验记录分别保存生产者盈利来源、每家投入、竞争者投入、真实注意力、产品偏好或考虑集、后续购买、参与者实际
支付与消费倾向，并用来源可追踪的时间顺序估计两条链。只见赞助和曝光、只见曝光与购买相关、或只见参与者
取得收入，均不能单独识别完整桥接因果。
资金链逐人核对基本与差异份额；产品账逐项核对销售回款、产品成本与销售利润；营销账按命题
~\ref{prop:sales-marketing-accounting} 核对偏好增量利润、赞助及配套成本和净增利润。
\end{Certificate}

\begin{proof}[证明（链式法则与通道相加）]
在声明的局部比较域内，把约化需求全微分除以 \(d\Pi_t^P\)，偏好路径给出
\(\kappa_k^{\mathrm{pref}}\partial I_k^{B*}/\partial\Pi_t^P\)，收入路径给出
\(\kappa^{\mathrm{inc}}\partial I^{B*}/\partial\Pi_t^P\)。各乘积非负且至少一个严格为正，故总和严格为正。
该结论恢复边际传导，不保证乘数足以消除停滞，也不保证注意力投入达到福利最优。证毕。
\end{proof}

\begin{corollary}[制度桥接恢复正边际传导的条件推论]
\label{cor:institutional-bridge-restores-transmission}
令 \(D_{t+1}^{P,0}\) 表示没有合格参与桥接时的生产品需求，\(D_{t+1}^{P,\mathcal K}\) 表示
\(\mathcal K^{\mathrm{inst}}\) 已构造合格候选路径、且市场在其中形成定理~\ref{thm:competitive-attention-injection}
之竞争均衡后的需求。若除桥接外的状态保持可比，没有未计入的挤出效应，并且定理~\ref{thm:dual-channel-bridge}
的正导数条件成立，则
\[
  \frac{\partial D_{t+1}^{P,\mathcal K}}{\partial\Pi_t^P}
  -
  \frac{\partial D_{t+1}^{P,0}}{\partial\Pi_t^P}
  =
  \sum_k\kappa_k^{\mathrm{pref}}
  \frac{\partial I_k^{B*}}{\partial\Pi_t^P}
  +\kappa^{\mathrm{inc}}
  \frac{\partial I^{B*}}{\partial\Pi_t^P}
  >0.
\]
因此，构成性治理与市场选择的组合可以恢复生产盈利通向需求的一个正边际通道；制度构造本身、行政注资本身
或行业存在本身均不充分。在传统生产--就业--收入传导衰减的声明域内，这一增量为维持有效需求、人类参与和
物质再生产提供一条可识别的制度机制；它是否足以支持总量复苏与长期发展，仍取决于规模、分配、净福利和
动态稳定条件。
\end{corollary}

\begin{Certificate}[边际传导恢复核验记录]
核验记录同时保存无桥接基线、制度规则与合法路径、市场竞争形成的均衡投入、两条通道的局部导数以及挤出项。
它只确认声明域内的边际增量。投入规模过小、参与收入被储蓄或泄漏、注意力只转移既有销量、分配失真、
外部损害抵消收益或动态系统不稳定，均可使宏观活动仍未恢复。
\end{Certificate}

\begin{proof}[证明（相对于无桥接基线的全导数差）]
由无桥接基线的定义，新增偏好通道和参与收入通道在
\(D_{t+1}^{P,\mathcal K}-D_{t+1}^{P,0}\) 中恰为定理~\ref{thm:dual-channel-bridge} 声明的两项，其他状态可比且
没有未计入挤出使基线共同项相消。对 \(\Pi_t^P\) 求全导数，得到推论所列等式；由各项非负且至少一项严格
为正，导数差严格为正。若制度只建立候选路径而市场没有形成竞争投入，或两条通道的系数均为零，则严格正号
消失，说明构成性制度本身不是充分条件。证毕。
\end{proof}

\begin{remark}[竞争性注意力投入不是自动福利改进]
多个生产者的注意力竞赛可以内生维持桥接资金，也可能产生广告拥挤、成瘾设计、隐私侵害、虚假流量、
排他赞助、市场封锁或纯粹争夺既有份额的耗散。真实购买检验只能证明私人回报链存在，不能自动证明净社会
剩余为正；仍须从收益中扣除资源成本、被操纵选择和第三方损害。
\end{remark}

\section{兴趣消费的广义奢侈品化与产业规模条件}
\label{sec:interest-luxury-sector}

\begin{definition}[兴趣消费与超额预算]
令消费者集合为 \(\mathcal H\)，消费者 \(h\) 的可支配收入为 \(Y_h>0\)。令 \(m_h\) 为满足食物、住房、
能源、基本医疗与其他声明期内物质需要的复合消费品，价格为 \(p_M(A)>0\)，最低需要量为
\(\bar m_h>0\)。定义最低物质支出与超额预算
\[
  E_h^{\min}(A)=p_M(A)\bar m_h,
  \qquad
  S_h(A)=Y_h(A)-E_h^{\min}(A).
\]
令 \(\mathcal J_I\) 为竞技、文化、艺术、游戏、探索、学习、社群与其他兴趣服务的有限集合，消费者对
服务 \(j\in\mathcal J_I\) 的消费量与价格分别为 \(x_{hj}\ge0\) 与 \(p_j>0\)。若 \(j\) 不进入最低物质
需要束、但通过体验、技能、身份、关系、声誉或共同历史等可消费特征进入消费者效用，则称 \(j\) 为兴趣
消费品。该定义把商品视为可消费特征的载体，而不要求同一兴趣品对所有消费者具有相同效用权重
\cite{Lancaster1966Characteristics}。
\end{definition}

\begin{definition}[经济学意义上的广义奢侈品属性]
在预先声明的价格、收入与制度域 \(\Omega_I\) 内，若兴趣消费品 \(j\) 的 Marshall 需求满足
\[
  \eta_{hj}^{Y}
  =\frac{\partial\ln x_{hj}^{*}}{\partial\ln Y_h}>1,
\]
则称 \(j\) 对消费者 \(h\) 在 \(\Omega_I\) 内具有收入弹性意义上的广义奢侈品属性。该名称只表示其需求
随可支配收入增长得快于收入，不推出炫耀动机、高绝对价格、社会优越性或福利优越性。

给定事前阈值 \(\bar R_I>0\)、\(\bar H_I>0\)，若兴趣消费部门同时满足
\[
  R_I\ge\bar R_I,
  \qquad
  \#\left\{h:\sum_{j\in\mathcal J_I}x_{hj}^{*}>0\right\}\ge\bar H_I,
  \qquad
  R_I-C_I-H_I\ge0,
\]
且存在两个以上独立供给者或可替代品类，则称其在该声明域内形成广义奢侈品产业。这里 \(R_I\) 是来自
消费者付费、竞争性赞助、授权和服务合同且不重复计算的行业收入，\(C_I\) 是真实资源成本，\(H_I\) 是
已货币化外部损害。“广阔”因而同时要求收入规模、付费覆盖、净余额与供给多样性，不以播放量、注册量或
单个高价交易代替。
\end{definition}

\begin{assumption}[最低需要之上的 Stone--Geary 兴趣需求]
对每个 \(h\)，令 \(\mathcal J_h\subseteq\mathcal J_I\) 为其具有正偏好的兴趣集合，并假设
\[
  U_h(m_h,x_h)
  =(m_h-\bar m_h)^{\alpha_h}
  \prod_{j\in\mathcal J_h}x_{hj}^{\beta_{hj}},
  \qquad
  \alpha_h>0,
  \quad
  \beta_{hj}>0,
  \quad
  \alpha_h+\sum_{j\in\mathcal J_h}\beta_{hj}=1,
\]
预算约束为
\[
  p_M(A)m_h+
  \sum_{j\in\mathcal J_h}p_jx_{hj}\le Y_h(A),
  \qquad
  S_h(A)>0.
\]
这是一类线性支出系统的显式消费模型；它不主张所有家庭或全部兴趣品在现实中都服从同一参数
\cite{Stone1954LES}。
\end{assumption}

\begin{theorem}[兴趣消费的收入弹性与预算份额定理]
\label{thm:interest-luxury-engel}
在最低需要之上的 Stone--Geary 兴趣需求假设下，唯一内点需求为
\[
  m_h^{*}
  =\bar m_h+\frac{\alpha_h S_h}{p_M},
  \qquad
  x_{hj}^{*}
  =\frac{\beta_{hj}S_h}{p_j}.
\]
因此，对任意 \(j\in\mathcal J_h\)，
\[
  \eta_{hj}^{Y}
  =\frac{Y_h}{S_h}>1,
\]
而消费者 \(h\) 的兴趣消费总支出及其收入份额为
\[
  E_h^I
  =\sum_{j\in\mathcal J_h}p_jx_{hj}^{*}
  =\beta_h^I S_h,
  \qquad
  s_h^I
  =\frac{E_h^I}{Y_h}
  =\beta_h^I\left(1-\frac{p_M\bar m_h}{Y_h}\right),
\]
其中 \(\beta_h^I=\sum_j\beta_{hj}=1-\alpha_h\)。保持价格和偏好参数不变时，
\[
  \frac{\partial s_h^I}{\partial Y_h}
  =\frac{\beta_h^I p_M\bar m_h}{Y_h^2}>0.
\]
故在该声明域内，兴趣消费不是因名称而被定义成奢侈品，而是由收入弹性和预算份额导数严格取得广义奢侈品
属性。
\end{theorem}

\begin{Certificate}[兴趣消费 Engel 核验记录]
核验记录保存产品域、最低需要束、价格、可支配收入、兴趣品数量、偏好或估计参数、需求系统与适用收入区间，
并分别估计 \(\eta_{hj}^{Y}\) 和 \(\partial s_h^I/\partial Y_h\)。若最低需要束未经声明、收入变化与价格变化
混同，或只观察到高价而没有收入弹性，则不能确认广义奢侈品属性。
\end{Certificate}

\begin{proof}[证明（预算约束、一级条件与 Engel 导数）]
对数化效用并令 \(\lambda_h\) 为预算约束乘子，内点一级条件为
\[
  \frac{\alpha_h}{m_h-\bar m_h}=\lambda_hp_M,
  \qquad
  \frac{\beta_{hj}}{x_{hj}}=\lambda_hp_j.
\]
将
\(
  p_M(m_h-\bar m_h)=\alpha_h/\lambda_h
\)
和
\(
  p_jx_{hj}=\beta_{hj}/\lambda_h
\)
相加，再用参数和为一以及预算约束取等，得到 \(1/\lambda_h=S_h\)，故需求式成立。由于
\[
  \frac{\partial x_{hj}^{*}}{\partial Y_h}=\frac{\beta_{hj}}{p_j},
\]
有
\[
  \eta_{hj}^{Y}
  =\frac{\beta_{hj}}{p_j}\frac{Y_h}{\beta_{hj}S_h/p_j}
  =\frac{Y_h}{S_h}>1,
\]
其中严格不等式来自 \(p_M\bar m_h>0\)。把各兴趣品支出相加得到 \(E_h^I=\beta_h^IS_h\)，再对
\(s_h^I\) 求导即得正导数。严格凹的对数效用保证内点解唯一。证毕。
\end{proof}

\begin{corollary}[人工智能降低标准化必要品成本时的兴趣支出扩张]
\label{cor:ai-interest-expenditure-expansion}
若 \(p_M'(A)<0\)，偏好参数在声明区间内固定，且对每个相关消费者
\[
  S_h'(A)
  =Y_h'(A)-\bar m_hp_M'(A)\ge0,
\]
并至少对一个具有 \(\beta_h^I>0\) 的消费者严格为正，则聚合兴趣支出
\[
  E^I(A)=\sum_{h\in\mathcal H}\beta_h^IS_h(A)
\]
满足 \(dE^I/dA>0\)。若消费者还取得兴趣行业参与支付 \(T_h^I\)，且
\(Y_h=Y_h^0+T_h^I\)，则
\[
  \frac{\partial E^I}{\partial T_h^I}=\beta_h^I>0.
\]
因此，标准化必要品成本下降与参与收入增加都可以扩大兴趣消费预算；若自动化造成的可支配收入下降超过必要品
成本下降，则该结论不成立。
\end{corollary}

\begin{proof}[证明（聚合超额预算导数）]
由定理~\ref{thm:interest-luxury-engel}，
\[
  \frac{dE^I}{dA}
  =\sum_h\beta_h^I
  \left(Y_h'(A)-\bar m_hp_M'(A)\right).
\]
各项非负且至少一项严格为正，故总和严格为正。把 \(Y_h=Y_h^0+T_h^I\) 代入同一支出式，对
\(T_h^I\) 求偏导即得第二个结论。证毕。
\end{proof}

\begin{proposition}[兴趣品类进入与产业广度扩张命题]
\label{prop:interest-variety-entry}
令潜在兴趣服务供给者集合为有限集 \(\mathcal R_I\)，聚合超额预算为
\(S^I=\sum_hS_h\)。供给者 \(r\) 进入后的净利润为
\[
  \pi_r^I(S^I;A)
  =R_r^I(S^I)-C_r^I(R_r^I(S^I);A)-F_r,
\]
其中 \(F_r>0\) 为固定进入成本。若 \(\pi_r^I\) 关于 \(S^I\) 连续且严格递增，并且存在
\(\underline S_r<\overline S_r\) 使
\[
  \pi_r^I(\underline S_r;A)<0
  <\pi_r^I(\overline S_r;A),
\]
则存在唯一进入阈值 \(S_r^*(A)\)，使供给者在 \(S^I\ge S_r^*(A)\) 时进入。活跃品类数
\[
  N_I(S^I;A)
  =\sum_{r\in\mathcal R_I}
  \mathbf 1\{S^I\ge S_r^*(A)\}
\]
关于 \(S^I\) 弱单调增加，并在任何跨过至少一个进入阈值的区间严格扩张。若人工智能降低标准化组织、匹配、
传播或内容辅助成本，从而降低部分 \(C_r^I\) 或 \(F_r\)，相应进入阈值进一步下降。
\end{proposition}

\begin{proof}[证明（零利润进入阈值）]
连续性与两端异号由介值定理给出零点；严格递增保证零点唯一。零点以下净利润为负，零点以上为正，故自由进入
条件给出所述阈值规则。每个指标函数都关于 \(S^I\) 弱单调，求和后 \(N_I\) 仍弱单调；跨过至少一个进入
阈值时严格增加。成本或固定成本下降使给定 \(S^I\) 下的净利润弱增，零利润阈值因而弱降。证毕。
\end{proof}

\begin{corollary}[大众参与与稀缺体验溢价并存]
\label{cor:mass-access-authenticity-premium}
设兴趣服务 \(j\) 同时包含可低成本复制的标准化基础层与受容量
\(
  q_j^{\mathrm{auth}}\le\bar q_j^{\mathrm{auth}}
\)
约束的真人、现场、可信竞技、专属关系或历史性体验层。在竞争均衡、约束资格成立且不存在其他价格楔子时，
若真实性容量约束的 KKT 乘子 \(\mu_j^{\mathrm{auth}}>0\)，则高阶体验的边际稀缺溢价为正；与此同时，基础层
价格可以随复制边际成本下降。故“大众低门槛参与--专业持续消费--稀缺体验高溢价”的分层结构与人工智能
降低标准化成本并不矛盾。
\end{corollary}

\begin{proof}[证明（稀缺约束的影子价格）]
由命题~\ref{prop:scarcity-transfer} 的 KKT 与包络结论，
\[
  \mu_j^{\mathrm{auth}}>0
  \Longrightarrow
  \frac{\partial Y^*}{\partial\bar q_j^{\mathrm{auth}}}
  =\mu_j^{\mathrm{auth}}>0.
\]
因此真实性容量增加一单位具有正边际价值；在所列竞争均衡条件下，该影子价格进入稀缺层价格。标准化基础层
没有该绑定容量约束，二者价格可以分化。证毕。
\end{proof}

\begin{definition}[兴趣消费部门的净需求桥接规模]
令 \(D_{t+1}^{P,0}\) 为没有兴趣消费桥接时声明产品域的下一期有效需求，\(D_{t+1}^{\dagger}\) 为维持
既定产能利用、偿付与再生产所需的事前基准，定义需求缺口
\[
  \Omega_{t+1}
  =\max\{0,D_{t+1}^{\dagger}-D_{t+1}^{P,0}\}.
\]
令 \(\Delta D_{t+1}^{\mathrm{pref}}\) 为竞争性注意力引起的生产商品偏好与购买增量，
\(\Delta D_{t+1}^{\mathrm{inc}}\) 为兴趣行业参与支付引起的购买力与消费增量，\(\Lambda_{t+1}^{I}\ge0\) 为
可识别的既有消费替代、储蓄泄漏、进口泄漏与跨期挤出，定义
\[
  \Gamma_{t+1}^{I,\mathrm{net}}
  =\Delta D_{t+1}^{\mathrm{pref}}
  +\Delta D_{t+1}^{\mathrm{inc}}
  -\Lambda_{t+1}^{I}.
\]
三项必须使用同一货币单位、同一产品域和同一时间窗，生产者之间的份额转移不得重复计入总需求增量。
\end{definition}

\begin{theorem}[兴趣消费产业支持经济循环的充分条件定理]
\label{thm:interest-sector-macro-support}
若：
\begin{enumerate}[label=\textup{(L\arabic*)},leftmargin=4em]
  \item 兴趣消费部门通过定理~\ref{thm:interest-luxury-engel} 与命题~\ref{prop:interest-variety-entry}
  的需求和进入检验，并达到事前声明的广义产业阈值；
  \item 定理~\ref{thm:competitive-attention-injection} 与定理~\ref{thm:dual-channel-bridge} 的竞争性投入、
  偏好通道和收入通道条件成立；
  \item \(\Gamma_{t+1}^{I,\mathrm{net}}\ge\Omega_{t+1}\)；
  \item 兴趣部门的已到账收入覆盖组织成本、已货币化损害、基本与差异分配及到期债务；生产部门的产品销售
  回款覆盖产品成本，偏好增量销售利润覆盖赞助及配套营销成本并形成正净增利润，风险调整后追加投入回报为正；
  \item 动态更新在声明固定点邻域可微，且雅可比谱半径严格小于一；内部参与--分配和外部风险边界分别通过
  相应核验记录，
\end{enumerate}
则兴趣消费部门足以在声明产品域、时间窗和邻域内填补基准需求缺口，维持生产--消费、参与--分配与下一期
竞争性投入的循环。该结论是局部、条件性充分结论，不证明兴趣消费取代物质生产，也不证明其足以支撑所有
产品域、全部宏观经济或任意冲击状态。
\end{theorem}

\begin{Certificate}[兴趣消费产业宏观规模核验记录]
核验记录至少保存：最低物质支出和家庭超额预算；各兴趣品收入弹性、付费覆盖和真实收入；供给者进入退出与成本；
竞争性生产盈利投入；真实注意力、生产商品偏好与购买；参与者实际支付及其边际消费倾向；需求缺口
\(\Omega\)；替代、储蓄、进口和跨期泄漏；行业净余额；外部损害；动态雅可比与稳定邻域。观看量、交易总额、
赞助总额或单次盈利均不能单独确认“足以支持经济循环”的结论。
\end{Certificate}

\begin{proof}[证明（需求缺口覆盖与多期延续）]
由定义，兴趣消费桥接后的有效需求为
\[
  D_{t+1}^{P,I}
  =D_{t+1}^{P,0}+\Gamma_{t+1}^{I,\mathrm{net}}.
\]
由 L3，
\[
  D_{t+1}^{P,I}
  \ge D_{t+1}^{P,0}+\Omega_{t+1}
  \ge D_{t+1}^{\dagger}.
\]
故声明期的基准需求缺口被填补。L1 保证该结果来自达到规模和覆盖阈值的分散兴趣消费产业，而不是单笔高价
交易；L2 保证投入由生产商品盈利反馈内生维持，并经偏好和收入两条路径进入需求；L4 排除以未支付成本、
债务滚动或外部损害伪造循环；L5 使固定点附近的小偏离随期数收敛，并使内部支付与外部安全具备制度资格。
因此循环能在声明邻域内延续。任一条件缺失都只剩局部收入、短期刺激或经验相关，不能推出所述充分结论。
证毕。
\end{proof}

\begin{remark}[“新一代广阔奢侈品行业”的严格口径]
本节允许在 L1--L5 可通过的声明域内，把兴趣消费部门称为“新一代广阔的广义奢侈品行业”：其“奢侈品”
来自最低需要之上需求的收入弹性，其“新一代”来自价值载体由可复制物件部分转向真人时间、可信竞争、关系、
身份与不可重复经历，其“广阔”来自付费覆盖、收入规模、品类进入和分层供给，而不是人为封闭稀缺或单纯高价。
这是可检验的条件命题，不是关于所有兴趣活动、所有家庭或后 AI 历史路径的无条件预测。
\end{remark}

\begin{proposition}[兴趣消费奢侈品化的失败边界]
\label{prop:interest-luxury-failure-boundary}
若出现以下任一情形，就不能确认本节的相应结论：\(S_h\le0\) 使家庭没有超额预算；实证收入弹性
\(\eta_{hj}^{Y}\le1\)；兴趣收入主要来自行政摊派、一次性补贴、封闭积分或后来者资金；注意力不能改变生产商品
偏好与购买；参与收入没有形成购买力；收入增长只是挤走其他等额消费；高溢价来自欺诈、锁定、操纵或不可退出；
行业规模低于事前阈值；或者 \(\Gamma^{I,\mathrm{net}}<\Omega\)。前述失败只撤销“广义奢侈品产业”或
“宏观循环支撑”结论，不自动撤销某项兴趣服务本身可能具有的正经济价值。
\end{proposition}

\begin{proof}[证明（前提否定与结论分层撤销）]
若 \(S_h\le0\) 或 \(\eta_{hj}^{Y}\le1\)，定理~\ref{thm:interest-luxury-engel} 的超额预算或高收入弹性前提
失败，不能确认广义奢侈品需求结论。行政摊派、一次性补贴、封闭积分或后来者资金不满足竞争性盈利注入和真实
收入来源条件；注意力不改变购买或参与收入不形成购买力，分别使定理~\ref{thm:dual-channel-bridge} 的偏好或
收入通道为零。等额挤出、高风险操纵和不可退出分别进入 \(\Lambda^{I}\) 或使治理条件失败；规模低于阈值使 L1
失败；\(\Gamma^{I,\mathrm{net}}<\Omega\) 使 L3 失败。故相应充分定理的前件合取不成立，不能推出产业化或宏观
支撑。上述反例均可与某项服务对个体仍有正效用或正支付并存，因此不推出该服务经济价值为零。证毕。
\end{proof}

\section{全员底线承诺、充分救济与市场稳定条件}
\label{sec:sufficient-relief-floor}

\begin{definition}[最低生存--市场接入束与充分救济]
固定时期 \(t\) 和居民集合 \(\mathcal H_t\)。居民 \(h\) 的市场收入为 \(Y_{ht}^{m}\ge0\)，满足食物、住房、
能源、基本医疗和其他声明必要品所需的最低物质支出为 \(E_{ht}^{\min}>0\)。令
\(A_{ht}^{\mathrm{access}}\ge0\) 为取得基本身份、通信、交通、支付和求职或参与入口所需的最低市场接入成本，
并定义不含兴趣奢侈消费的底线束
\[
  F_{ht}
  =E_{ht}^{\min}+A_{ht}^{\mathrm{access}}.
\]
救济向量 \(\mathbf R_t=(R_{ht})_{h\in\mathcal H_t}\) 称为社会充分救济，当且仅当
\[
  R_{ht}\ge0,
  \qquad
  Y_{ht}^{+}=Y_{ht}^{m}+R_{ht}\ge F_{ht}
  \quad(\forall h\in\mathcal H_t),
\]
并且资金来源、储备、极端损失顺序和跨期预算通过本节核验记录。它只保证生存和真实进入市场所需的底线束；
兴趣品种类、消费数量、职业选择、投资方向、价格、利润和底线以上收入仍由分散主体选择。

在向量的逐坐标偏序下，若不存在另一个充分救济向量
\(\widetilde{\mathbf R}_t\le\mathbf R_t\) 且至少一个坐标严格更小，则称 \(\mathbf R_t\) 为坐标最小充分救济。
\end{definition}

\begin{definition}[救济资金、储备与激励相容融资域]
令生产部门的正盈利总额为
\[
  \Pi_t^{P,+}
  =\sum_{k\in\mathcal K}[\Pi_{kt}^{P}]_+,
\]
公共资产净收益和其他非重复收入为 \(Q_t^R\ge0\)，救济识别、支付和复核的真实管理成本为 \(C_t^R\ge0\)。
给定税费或共同缴款率 \(\tau_t^R\in[0,\bar\tau_t^R)\)、期初储备 \(K_t^R\ge0\) 与储备收益率
\(r_t^R>-1\)，救济账户按
\[
  K_{t+1}^{R}
  =(1+r_t^R)K_t^R
  +\tau_t^R\Pi_t^{P,+}+Q_t^R
  -\sum_hR_{ht}-C_t^R
\]
更新。称融资在时间窗 \(0,\ldots,T\) 内可持续，当且仅当
\[
  K_t^R\ge\underline K_t^R\ge0
  \quad(t=0,\ldots,T+1),
\]
不存在以未声明未来缴款滚动填补当期缺口，并且至少一个承担资本、组织或稀缺能力投入的市场主体在缴款后仍有
严格为正的预期剩余索取和边际投入回报。最后一项保证救济不以消灭市场投资激励为融资方式。
\end{definition}

\begin{theorem}[坐标最小充分救济的存在、唯一性与预算条件]
\label{thm:minimal-sufficient-relief}
对有限居民集合，给定 \(\mathbf Y_t^m\) 与 \(\mathbf F_t\)，满足底线束的坐标最小充分救济唯一，且为
\[
  R_{ht}^{*}
  =[F_{ht}-Y_{ht}^{m}]_+.
\]
记
\[
  \mathcal B_t^{R,*}
  =\sum_hR_{ht}^{*}+C_t^R.
\]
若
\[
  (1+r_t^R)K_t^R+\tau_t^R\Pi_t^{P,+}+Q_t^R
  -\underline K_{t+1}^R
  \ge\mathcal B_t^{R,*}
\]
且缴款后的正剩余激励条件成立，则该向量在当期既充分、可融资，又不要求收入平均化或剩余索取归零。
\end{theorem}

\begin{proof}[证明（逐坐标下界与账户恒等式）]
对任意居民 \(h\)，约束 \(Y_{ht}^{m}+R_{ht}\ge F_{ht}\) 与 \(R_{ht}\ge0\) 等价于
\[
  R_{ht}\ge\max\{0,F_{ht}-Y_{ht}^{m}\}=R_{ht}^{*}.
\]
因此任何充分向量都逐坐标不小于 \(\mathbf R_t^*\)，而 \(\mathbf R_t^*\) 本身使全部约束取等或保持原有
市场收入，故其存在、充分、坐标最小且唯一。把 \(\mathcal B_t^{R,*}\) 代入救济账户更新式，所列预算不等式
恰好给出 \(K_{t+1}^R\ge\underline K_{t+1}^R\)。缴款后的正剩余条件另行保证至少一条私人投入路径仍有
严格正回报。以上条件只确定救济下界，没有给出任何收入上界或相等约束。证毕。
\end{proof}

\begin{definition}[循环非退化阈值与救济需求增量]
令 \(D_{t+1}^{P,0}\) 为没有救济时的生产品有效需求，\(N_{t+1}^{0}\) 为能够真实进入兴趣与参与市场的有效
人数，\(D_{t+1}^{\dagger}\) 与 \(\bar N_{t+1}>0\) 分别为维持声明域再生产和双边匹配、注意力竞争所需的
事前阈值。居民 \(h\) 对生产品的边际消费倾向为 \(c_{ht}^{P}\in[0,1]\)，救济引起的直接需求增量为
\[
  \Delta D_{t+1}^{R}
  =\sum_hc_{ht}^{P}R_{ht}^{*}.
\]
若市场接入概率 \(n_h(Y;F)\in[0,1]\) 关于底线以上可支配资源弱增，则有效进入增量为
\[
  \Delta N_{t+1}^{R}
  =\sum_h
  \left[
  n_h(Y_{ht}^{m}+R_{ht}^{*};F_{ht})
  -n_h(Y_{ht}^{m};F_{ht})
  \right]\ge0.
\]
当 \(D_{t+1}^P\ge D_{t+1}^{\dagger}\)、\(N_{t+1}\ge\bar N_{t+1}\)，且生产盈利、竞争性投入、注意力、
参与收入与消费的局部更新满足稳定条件时，称该声明域的市场参与循环非退化。
\end{definition}

\begin{theorem}[全员底线承诺下的最小救济与稳定循环相容性]
\label{thm:relief-endogenous-stability}
假设：
\begin{enumerate}[label=\textup{(R\arabic*)},leftmargin=4em]
  \item 存在非空居民集合 \(\mathcal H_t^0\)，使 \(Y_{ht}^{m}<F_{ht}\)，且传统劳动收入传导不足以在声明期内
  自动补齐相应缺口；
  \item 没有救济时，
  \(D_{t+1}^{P,0}<D_{t+1}^{\dagger}\) 或 \(N_{t+1}^{0}<\bar N_{t+1}\)，而声明域的任一非退化稳定循环
  必须同时达到这两个阈值；
  \item 生产品需求关于底线收入弱增，市场接入概率关于可支配资源弱增；至少一个缺口居民的
  \(c_{ht}^{P}>0\)，且底线得到保障后存在真实市场进入路径；
  \item 坐标最小救济满足
  \[
    D_{t+1}^{P,0}+\Delta D_{t+1}^{R}
    +\Gamma_{t+1}^{I,\mathrm{net}}
    \ge D_{t+1}^{\dagger},
    \qquad
    N_{t+1}^{0}+\Delta N_{t+1}^{R}
    \ge\bar N_{t+1},
  \]
  其中 \(\Gamma_{t+1}^{I,\mathrm{net}}\) 通过定理~\ref{thm:interest-sector-macro-support} 的净额和稳定条件；
  \item 定理~\ref{thm:minimal-sufficient-relief} 的当期预算条件在整个声明时间窗逐期成立，且救济、税费和
  储备进入动态系统后，在声明固定点邻域可微且雅可比谱半径严格小于一。
\end{enumerate}
则任何同时承诺“全体成员不跌破底线束”和“市场参与循环非退化”的制度，都必须给每个居民至少
\(R_{ht}^{*}\)；在 R1--R5 下，\(\mathbf R_t^*\) 又与市场双通道共同足以维持声明域内的底线、有效需求、
参与密度和下一期循环。逐人的必要下界由全员底线承诺给出，需求、双边市场厚度与跨期条件则检验该底线方案能否与稳定循环相容。总量阈值本身不推出全员逐户补足，也不推出公共救济是所有经济环境中的唯一收入修复方式。
\end{theorem}

\begin{proof}[证明（必要下界、阈值矛盾与充分闭环）]
先证必要性。对任一 \(h\in\mathcal H_t^0\)，若 \(R_{ht}<R_{ht}^{*}\)，则
\(Y_{ht}^{m}+R_{ht}<F_{ht}\)，直接违反全体成员底线承诺。因此任何满足该承诺的制度都有
\(\mathbf R_t\ge\mathbf R_t^*\)。R2 只确认所声明的无救济基线未达到总量阈值；它不排除另一种转移、保险或参与收入方案补足该总量缺口。全员下界的必要性由逐人约束证明，不由 R2 代替。

再证充分性。由定理~\ref{thm:minimal-sufficient-relief}，\(\mathbf R_t^*\) 使全部居民达到底线且账户逐期
不低于储备下界。R3 给出 \(\Delta D^R,\Delta N^R\ge0\)，R4 使救济需求、兴趣行业的注意力偏好通道和
参与收入通道在扣除泄漏后共同越过需求与参与阈值；R5 排除以一次赤字、债务滚动或发散反馈伪造持续循环。
因此底线、需求、市场厚度、竞争性投入和下一期盈利在声明邻域内闭合。救济只修复阈值以下坐标，实际消费、
劳动、兴趣、价格、投资和赢家仍由市场路径决定。证毕。
\end{proof}

\begin{proposition}[充分救济与能者多得、贫富差距和市场配置兼容]
\label{prop:relief-not-equalization}
令救济后的底线以上净收入为
\[
  Z_{ht}^{+}=Y_{ht}^{m}+R_{ht}^{*}-F_{ht}
  =[Y_{ht}^{m}-F_{ht}]_+.
\]
若两名居民原本均高于各自底线，则救济前后的底线以上收入差完全不变；若一人低于底线，制度只把其净收入
提高到零而不设收入上限。因而 \(\mathbf R_t^*\) 与能力、贡献、资本、风险和市场声誉形成的差异化收入、
正利润、正剩余索取及贫富差距相容。它既不是按需分配全部资源，也不把兴趣消费变成公共配给；它只把生存和
真实进入从胜负对象中移出，再让下限以上的配置继续由供求、竞争、盈亏与进入退出形成。
\end{proposition}

\begin{proof}[证明（截断映射）]
结论直接来自截断映射 \(Z^+=[Y^m-F]_+\)。对所有 \(Y^m\ge F\) 的坐标，映射斜率为一，不改变任何两点
之间的差；只对 \(Y^m<F\) 的坐标填补负缺口。映射没有上界、等份或统一消费束约束，故不推出收入相等、
利润归零或计划配置。证毕。
\end{proof}

\begin{Certificate}[社会充分救济与循环底盘核验记录]
核验记录必须保存：最低物质束和市场接入束的项目、数量、价格与更新规则；居民市场收入和缺口；逐人救济与实际
到账；生产正盈利、缴款率、公共资产收益、管理成本、储备下界、压力情景和极端损失顺序；缴款后私人剩余
索取与边际投资回报；生产品边际消费倾向、有效进入、双边市场厚度、事前需求缺口、兴趣行业净桥接和动态
稳定邻域。只列预算拨款、平均补贴或总消费增长，均不能证明“充分”“最小”或“可持续”。
\end{Certificate}

\begin{remark}[自然推导的边界]
本节证明的是：给定全员底线目标，在 R1--R5 的融资、需求和参与约束下，最小充分救济与稳定循环相容。若仅要求总量稳定而不要求全员底线，逐人下界不再由本定理推出。若市场收入本身已覆盖全部底线、私人保险和家庭网络已可靠补齐缺口，或没有
需求与参与阈值失效，则最小公共救济可以退化为零；若生产剩余不足、冲击突破储备、缴款抹去全部正剩余回报，
则必须调整底线、税费、储备或其他收入来源，不能用无资金承诺冒充充分救济。
\end{remark}

\section{现有注意力经济反例与三项核心功能}
\label{sec:attention-counterexample-core-functions}

\begin{definition}[集中型注意力变现基线]
固定时期 \(t\) 和有限非空的有效参与者集合 \(\mathcal E_t^A\)。令 \(R_t^A>0\) 为由广告、订阅、打赏、
佣金或其他注意力渠道形成的已实现总收入，\(T_{it}^A\ge0\) 为参与者 \(i\) 实际取得的支付，并要求
\[
  \sum_{i\in\mathcal E_t^A}T_{it}^A\le R_t^A.
\]
令 \(Y_{it}^{0}\) 为该支付以前的市场收入，\(F_{it}\) 为最低物质--市场接入束的货币成本，定义正收益覆盖率与
剩余底线缺口
\[
  \kappa_t^A
  =\frac{\Abs{\{i\in\mathcal E_t^A:T_{it}^A>0\}}}
  {\Abs{\mathcal E_t^A}},
  \qquad
  D_t^{F,A}
  =\sum_{i\in\mathcal E_t^A}
  [F_{it}-(Y_{it}^{0}+T_{it}^A)]_+.
\]
若 \(R_t^A>0\) 而 \(\kappa_t^A<1\) 且 \(D_t^{F,A}>0\)，称该状态为集中型注意力变现基线。自媒体、
直播或平台创作只有在同口径数据满足这些条件时才是该定义的经验实例；行业名称、账户数量或总流量本身不构成
核验记录。
\end{definition}

\begin{proposition}[注意力货币化的不充分性反例]
\label{prop:attention-monetization-insufficient}
命题
\[
  R_t^A>0
  \Longrightarrow
  \bigl(\kappa_t^A=1\bigr)
  \wedge
  \bigl(D_t^{F,A}=0\bigr)
\]
为假。因此，“注意力可以变现”既不充分推出所有有效参与者取得正份额，也不充分推出生存--市场接入底线
得到覆盖；仅凭注意力收入为正，不能确认广泛吸收传统部门释放的就业与参与需求或修复收入--需求循环。
\end{proposition}

\begin{proof}[证明（有限反例）]
取 \(\mathcal E_t^A=\{1,\ldots,n\}\)、\(n\ge2\)，令 \(T_{1t}^A=R_t^A>0\)，并令
\(T_{it}^A=0\)（\(i\ge2\)）。于是预算约束成立而 \(\kappa_t^A=1/n<1\)。再取至少一个
\(i\ge2\) 满足 \(Y_{it}^{0}<F_{it}\)，则 \(D_t^{F,A}>0\)。故前件成立而后件为假。该反例只否定
注意力货币化的充分性，不否定注意力对偏好、收入或文化供给的局部作用。证毕。
\end{proof}

\begin{definition}[三项核心功能与实施坐标]
\label{def:three-core-functions}
对任一组织或机制载体 \(G\)，定义
\[
  \mathbf g^{\mathrm{core}}(G)
  =\bigl(g_{\mathrm{part}}(G),
  g_{\mathrm{diff}}(G),
  g_{\mathrm{relief}}(G)\bigr)\in\{0,1\}^3,
\]
其中：
\begin{enumerate}[label=\textnormal{(F\arabic*)},leftmargin=4em]
  \item \(g_{\mathrm{part}}(G)=1\) 当且仅当每名被接纳且完成约定的有效参与者均取得严格为正的基本份额；
  \item \(g_{\mathrm{diff}}(G)=1\) 当且仅当同一比较域内的已验证贡献严格有别时，支付严格同序；
  \item \(g_{\mathrm{relief}}(G)=1\) 当且仅当每名居民的市场收入已覆盖底线，或由坐标最小且资金可持续的充分
  救济补齐缺口。
\end{enumerate}
令 \(g_{\mathrm{impl}}(G)=1\) 表示资格、计量、验证、合同确权、实际支付、复算、申诉和治理成本核验记录全部
通过。前三项定义机制要实现的分配、激励与接入功能；\(g_{\mathrm{impl}}\) 定义这些功能是否已经获得可执行的
稳定形式。它们不得合并为一个“制度健全”指标。
\end{definition}

\begin{proposition}[三项核心功能的不可约性与载体非唯一性]
\label{prop:three-functions-irreducible}
在定义~\ref{def:three-core-functions} 的目标域内，任一核心坐标都不能由其余两个逻辑推出。若目标同时要求
广泛正份额、贡献差异激励与全员底线接入，则其必要条件为
\[
  \min\{g_{\mathrm{part}}(G),g_{\mathrm{diff}}(G),
  g_{\mathrm{relief}}(G)\}=1.
\]
该必要条件只规定功能等价类，不规定唯一载体：企业赞助、平台分成、合作组织、公共救济、私人互助或其组合，
都可以在分别通过相应核验记录时实现这些坐标。
\end{proposition}

\begin{proof}[证明（独立见证）]
等额正基本份额配合独立的充分救济，可以使
\((g_{\mathrm{part}},g_{\mathrm{relief}})=(1,1)\)，但在贡献严格有别时使 \(g_{\mathrm{diff}}=0\)。
只向贡献较高者支付并另行补齐全员底线，可以在两人严格排序例中使
\((g_{\mathrm{diff}},g_{\mathrm{relief}})=(1,1)\)，却使未获基本份额的有效参与者见证
\(g_{\mathrm{part}}=0\)。正基本份额与严格贡献差异可以同时成立，但若至少一名居民的市场收入仍低于底线且
缺口未补，则 \((g_{\mathrm{part}},g_{\mathrm{diff}})=(1,1)\) 而 \(g_{\mathrm{relief}}=0\)。
三个有限见证分别阻断任一坐标由其余两个坐标推出。目标三项并取时，每一坐标按定义都是必要的；但定义没有
限定实现主体或组织名称，故不推出单一制度形式或历史路径唯一。证毕。
\end{proof}

\begin{remark}[经济基础、演化方向与实施基础]
注意力影响产品偏好并形成真实回款时，偏好回流为参与收入提供市场内资金来源；传统劳动收入收缩、参与能力闲置
和集中型注意力分配无法覆盖多数参与者时，又分别形成扩大正收益覆盖、保留贡献激励和维持底线接入的选择压力。
因此，三项核心功能可以作为分散竞争中的条件性演化方向出现，不以中心主体一次设计为前提。正式规则的作用是
把已经出现的功能稳定为可执行、可复算和可申诉的安排。该方向不是历史必然：偏好不能形成真实回款、参与规模
不足、治理成本过高、平台锁定或进入封锁，都可能使经济长期停留在集中型注意力变现状态。
\end{remark}

\section{公共治理的外部边界：充分释放空间并抑制风险溢出}

\begin{definition}[安全合法域与空间释放极大性]
令 \(r_{\mathrm{spill}}(\gamma)\ge0\) 为市场路径 \(\gamma\) 在责任承担、保险、保证金、赔偿与其他可执行
内部化措施之后，仍可能留给非自愿第三方、公共系统或未来状态的可识别剩余风险外溢；它不同于已经货币化
并计入私人或社会成本的损害 \(H_t^B\)。令 \(\bar r_{\mathrm{spill}}\) 为经公开程序确定的可承受阈值。候选法律状态族
\(\mathfrak L\) 中的合法域 \(\Omega_\ell\)
称为安全，当且仅当
\[
  \sup_{\gamma\in\Omega_\ell}r_{\mathrm{spill}}(\gamma)
  \le\bar r_{\mathrm{spill}}.
\]
安全域 \(\Omega_{\ell^*}\) 称为空间释放极大，当且仅当不存在另一个安全域
\(\Omega_{\ell'}\) 使
\[
  \Omega_{\ell^*}\subsetneq\Omega_{\ell'}.
\]
这一定义以集合包含而非行政偏好衡量“充分释放”：所有能够在不突破风险阈值的情况下加入的合法竞争、进入、
商业模式与试验路径都不应被无故排除。
\end{definition}

\begin{theorem}[宽空间--低溢出公共治理定理]
\label{thm:maximal-safe-governance-space}
若候选法律状态族 \(\mathfrak L\) 有限且至少含有一个安全合法域，则至少存在一个空间释放极大的安全域。
若该域含有两条以上市场路径，公共治理仍只规定边界而不唯一决定市场结果。注意力桥接中的披露、反欺诈、
隐私、未成年人保护、公平竞争、金融杠杆、责任与风险外溢规则，应进入 \(h\) 或逐期合法域；生产者投入额、
赞助对象、消费者偏好和胜出行业不由该定理指定。
\end{theorem}

\begin{Certificate}[安全极大域核验记录]
核验记录列出有限候选域、每个域的风险上界、被排除路径的具体外溢理由以及所有安全严格扩张候选。验收条件为
\[
  \sup_{\gamma\in\Omega_{\ell^*}}r_{\mathrm{spill}}(\gamma)
  \le\bar r_{\mathrm{spill}},
  \qquad
  \nexists\ell':
  \Omega_{\ell^*}\subsetneq\Omega_{\ell'}
  \wedge
  \sup_{\gamma\in\Omega_{\ell'}}r_{\mathrm{spill}}(\gamma)
  \le\bar r_{\mathrm{spill}}.
\]
\end{Certificate}

\begin{proof}[证明（有限偏序极大元）]
按集合包含关系排列全部安全合法域。非空有限偏序集至少存在极大元；否则从任一安全域出发都能找到安全的
严格超集，并在有限集合中无限上升，矛盾。若极大元含有多条路径，则边界条件没有选定其中任何一条，实际
路径仍由市场主体分散选择。证毕。
\end{proof}

\section{规范治理的内部核心：参与--分配的规则内重放、复算一致与相对公正}

\begin{definition}[参与--分配判定域]
固定规则版本、时间窗与收入池，记声明判定域为
\[
  \mathcal D=(\Part_D,\Act_D,\mathcal Z_D,\succeq_D,\sim_D,B_D),
  \qquad B_D>0.
\]
其中 \(\mathcal Z_D\) 为已经声明的参与证据空间，\(i\sim_D j\) 表示两名参与者在全部合法补偿维度上
具有等价的已验证参与，\(i\succeq_D j\) 表示 \(i\) 在这些维度上不低于 \(j\)，\(B_D\) 为经市场交易
形成且已经扣除声明成本的可分配收入池。令
\[
  s_i(G)=\frac{T_i(G)}{B_D}
\]
为参与者的净收入份额。这里所谓“不以传统生产分配为核心”，仅指行业不以新增物质产出以及可直接相加的
工时或件数作为主要分配基准；赛事、游戏、文化、内容与社群仍然生产真实服务。
\end{definition}

\begin{definition}[规则内可重放、复算一致与相对公正核验记录]
给定误差容限 \(\epsilon_z,\epsilon_s\ge0\) 与失败概率 \(\delta\in[0,1)\)，机制 \(G\) 在
\(\mathcal D\) 上规则内可重放且复算一致，当且仅当：规则版本、合法输入和理由码事前可得；独立复算者 \(r,r'\) 对同一
证据得到不一致结果的概率满足
\[
  \Pr\!\left[d\!\left(\operatorname{Replay}_r(z;G),
  \operatorname{Replay}_{r'}(z;G)\right)>\epsilon_z\right]\le\delta;
\]
身份、行为、规则版本、估值、确权与支付之间均有可追踪证据锚点；错误可以经独立复核改变状态。

机制在该声明域内相对公正，当且仅当规则内可重放且复算一致，并且
\[
  i\sim_Dj\Longrightarrow |s_i(G)-s_j(G)|\le\epsilon_s,
\]
\[
  i\succeq_Dj\Longrightarrow s_i(G)+\epsilon_s\ge s_j(G),
\]
同时满足预算平衡、相关机会可进入、利益冲突披露、规则权力有界、可审计、可申诉以及市场可竞争性。
记上述合取核验记录为 \(\mathfrak C_{\mathrm{RPF}}(G;\mathcal D,\epsilon_z,\epsilon_s,\delta)\)。它确认的是相对于
声明目标、比较关系、规则版本与容限的程序性和横向公正，不是唯一伦理分配、收入相等或社会福利最优。
在核心竞技组织域内，核验记录还须合取 \(\mathfrak C_{\mathrm{part}}^+\)：有效参与的正基本份额与严格贡献差异
由定理~\ref{thm:base-differential-coverage} 给出。测量容限用于判定贡献是否可区分；一旦两项贡献被规则
确认严格有别，差异支付按该定理执行。
\end{definition}

\begin{theorem}[规范参与--分配判定定理]
\label{thm:relative-participation-fairness}
若 \(\mathfrak C_{\mathrm{RPF}}\) 的全部坐标均通过，则对任一参与者可以从原始证据重放
\[
  a_i\to z_i\to v_i\to w_i\to b_i\to T_i,
\]
等价参与者之间的收入份额差不超过 \(\epsilon_s\)，声明贡献次序不会在超过容限时被收入份额反转，且误判
具有状态修复路径。因此，只能在该声明域内确认
\(\operatorname{RelativelyFairParticipationDistribution}(G;\mathcal D)\)。任一坐标失败时，市场仍可形成
交易和收入池，但不能据此确认参与--分配相对公正结论。
\end{theorem}

\begin{Certificate}[参与--分配相对公正核验记录]
核验记录必须保存规则版本与生效时间、参与资格、原始行为记录、身份和真实性证据、独立复算样本、等价类与
贡献偏序、收入池来源和成本扣除、份额与支付、利益冲突、规则修改权限、竞争和替代入口、申诉结果及实际
修复。只公开最后余额、只证明预算平衡或只证明市场有人付费，均不足以确认本核验记录。
\end{Certificate}

\begin{proof}[证明（条件合取与反例）]
证据锚点和独立重放给出判定可复现性；两条容限不等式分别给出横向一致性与声明贡献单调性；预算、权限、
审计和申诉条件排除无来源支付、隐蔽改写与不可纠正误判；市场可竞争性保留分散选择而不把份额变为行政配置。
故全部条件合取时可以确认所述有限结论。反之，删去任一条件均有直接反例：相同贡献可被秘密规则区别支付，
真实行为可被伪造身份冒领，准确计算可执行偏置权重，或错误结果无法修复。因此其余条件不能替代失败坐标。
证毕。
\end{proof}

\begin{definition}[治理自动化状态与相对公正判定成本]
令 \(A_G\) 表示治理自动化状态，并与生产技术状态 \(A\) 严格区分。对参与规模 \(N\) 和声明判定域
\(\mathcal D_N\)，把确认并维持 \(\mathfrak C_{\mathrm{RPF}}\) 的总成本写为
\[
  C_{\mathrm{RPF}}(A_G,N)
  =F_G(A_G)
  +C_{\mathrm{std}}(A_G,N)
  +C_{\mathrm{res}}(A_G,N),
\]
其中固定成本 \(F_G\) 包括规则形式化、初始系统和责任结构建设；标准化成本
\[
  C_{\mathrm{std}}
  =C_{\mathrm{record}}+C_{\mathrm{measure}}
  +C_{\mathrm{verify}}+C_{\mathrm{execute}}
\]
包括行为留痕、结构化计量、计算复核和确定执行；残余成本
\[
  C_{\mathrm{res}}
  =C_{\mathrm{truth}}+C_{\mathrm{identity}}
  +C_{\mathrm{adversarial}}+C_{\mathrm{appeal}}
  +C_{\mathrm{legal}}+C_{\mathrm{legitimacy}}
\]
包括输入真实性、现实身份、对抗攻击、复杂申诉、法律责任与社会正当性。该分解按同一成本归属规则去重，
不把参与者支付本身重复记为治理资源耗费。
\end{definition}

\begin{proposition}[生产产能松弛不蕴含治理单位成本下降]
\label{prop:capacity-not-governance-cost}
对任意后 AI 生产适用域核验记录，均有
\[
  \PostAI_\epsilon(\Goods_0,A,r)
  \not\Longrightarrow
  \frac{C_{\mathrm{RPF}}(A_G,N)}{N}\to0.
\]
特别地，\(\CapacitySlack(j,A,r)\) 只约束产品 \(j\) 的生产能力与有限最优需求，不约束
\(A_G\)、事实输入质量、身份攻击、申诉复杂度、法律责任或制度认可。因而不能用总条件 C1 或生产成本下降
替代独立的治理自动化与成本核验记录。
\end{proposition}

\begin{proof}[证明（独立状态反例）]
固定任意满足 \(\PostAI_\epsilon\) 的生产状态 \(A\)，另取一个治理状态 \(A_G^0\)，使每名参与者都必须由
人工完成身份核验、事实调查和争议裁决，单位残余成本为常数 \(c_0>0\)。于是
\(C_{\mathrm{RPF}}(A_G^0,N)\ge Nc_0\)，尽管生产端全部产能约束已经松弛，单位治理成本仍不趋于零。
该构造证明两个状态之间不存在逻辑蕴含。证毕。
\end{proof}

\begin{theorem}[治理单位成本的上极限界与条件性收敛定理]
\label{thm:rpf-unit-cost-collapse}
考虑参与规模 \(N\to\infty\) 与治理自动化序列 \(A_G(N)\)。若
\[
  \frac{F_G(A_G(N))}{N}\to0,
  \qquad
  \frac{C_{\mathrm{std}}(A_G(N),N)}{N}\to0,
  \qquad
  \limsup_{N\to\infty}
  \frac{C_{\mathrm{res}}(A_G(N),N)}{N}
  \le\bar c_{\mathrm{res}}<\infty,
\]
则
\[
  \limsup_{N\to\infty}
  \frac{C_{\mathrm{RPF}}(A_G(N),N)}{N}
  \le\bar c_{\mathrm{res}}.
\]
若同一声明域的真实可分配收入池满足
\[
  \liminf_{N\to\infty}\frac{Y_N^*}{N}\ge\underline y>0,
\]
则还有
\[
  \limsup_{N\to\infty}
  \frac{C_{\mathrm{RPF}}(A_G(N),N)}{Y_N^*}
  \le\frac{\bar c_{\mathrm{res}}}{\underline y}.
\]
以上仅为上极限界。若另有
\[
 \frac{C_{\mathrm{res}}}{N}\to c_\infty,\qquad
 \frac{Y_N^*}{N}\to y_\infty>0,
\]
则 \(C_{\mathrm{RPF}}/N\to c_\infty\)，且
\(C_{\mathrm{RPF}}/Y_N^*\to c_\infty/y_\infty\)。只有当一个独立技术约束还保证
所有可行核验机制的单位成本不低于 \(c_\infty\) 时，才能将其称为可达到的成本下界。
当 \(c_\infty=0\) 时上述两个极限为零。更一般地，只要事前门槛
\[
  \frac{C_{\mathrm{RPF}}}{Y_N^*}\le\bar c_G<1,
  \qquad
  Y_N^*-C_{\mathrm{RPF}}\ge0
\]
持续通过，\(\mathfrak C_{\mathrm{RPF}}\) 就具有规模可负担性；总治理成本仍可随 \(N\) 上升，不能由单位
单位成本界反推出绝对成本下降。
\end{theorem}

\begin{proof}[证明（成本分解与规模归一化）]
将成本恒等式除以 \(N\)，前两项按假设趋于零，第三项的上极限不超过
\(\bar c_{\mathrm{res}}\)，即得第一式。再写成
\[
  \frac{C_{\mathrm{RPF}}}{Y_N^*}
  =
  \frac{C_{\mathrm{RPF}}/N}{Y_N^*/N},
\]
分子取上极限、分母取严格正下极限，得到第二式。另有两个极限时，由和与商的极限定理得到所列收敛。仅有原假设时，残余单位成本可取 $1+(-1)^N/2$，满足上极限界而不收敛，故不可省略极限条件。最后两条门槛
直接保证治理以后仍有非负可分配余额，故可确认规模可负担性，但不改变相对公正核验记录的其他条件。证毕。
\end{proof}

\begin{Certificate}[治理自动化与相对公正成本核验记录]
核验记录分别保存生产状态 \(A\) 与治理状态 \(A_G\)、参与规模、固定成本、标准化成本、残余成本、收入池、
成本归属与去重规则、自动化前后同口径单位成本、成本占收入池比例及压力情景。记录完整性与计算复核证明
所保存的输入及其规则内处理结果，输入真实性、规则相对公正以及法律与社会正当性仍须分别核验。
若只报告总成本下降而不控制规模，或只报告计算费用而遗漏身份、对抗、申诉、法律和正当性成本，不得确认
未经条件支持的成本收敛或规模可负担结论。
\end{Certificate}

\begin{proposition}[传统按劳分配是低维可判定特例]
\label{prop:labor-distribution-special-case}
若传统行业的可补偿劳动可以由可独立核验的工时、件数与质量向量 \(\ell_i\) 表示，工资表
\(w(\ell_i;p_L)\) 在行为前公布，劳动合同可执行，竞争市场给出可比较的劳动价格 \(p_L\)，且不存在未计入的
歧视、串谋、买方垄断、强制或实质退出锁定，则等价劳动得到等价工资、增加声明劳动不降低工资，因而可构造
\(\mathfrak C_{\mathrm{RPF}}\)。这说明传统按劳分配通常具有较低的判定维数，不说明现实劳动市场自动公正；
任一所列市场或证据条件失败，都撤销该特例结论。
\end{proposition}

\begin{proof}[证明（工资表实例化）]
取 \(z_i=\ell_i\)、\(T_i=w(\ell_i;p_L)\)，以合同和劳动价格定义 \(\sim_D\) 与 \(\succeq_D\)。独立工时、
件数和质量记录给出重放；工资表的同值性与单调性给出两条份额条件；合同执行、市场可竞争和纠错程序补齐
其余条件。若存在买方垄断或秘密歧视，等价劳动可以得到不同工资，故条件不可删除。证毕。
\end{proof}

\begin{proposition}[市场形成收入池不充分决定参与分配公正]
\label{prop:market-pool-not-fair-allocation}
对任意由市场形成的 \(B_D>0\) 和 \(n\ge2\)，预算平衡只要求份额
\(s\in\Delta^{n-1}\)。若非物质参与具有协同、网络或非加总贡献，而市场交易没有同时给出
\(\sim_D\)、\(\succeq_D\) 与组织估值规则，则同一收入池至少允许连续统多个分配，其中可以包含违反横向
一致或声明贡献单调性的点。因此，市场经济可以决定需求、收入池与组织竞争，却不自动决定参与--分配相对
公正；后者正是规范治理必须补齐的内部判定对象。
\end{proposition}

\begin{proof}[证明（单纯形反例）]
取两名声明等价参与者 \(1\sim_D2\)。份额 \((1/2,1/2,0,\ldots)\) 与
\((1,0,0,\ldots)\) 都非负且预算平衡，但后一份额在任意 \(\epsilon_s<1\) 时违反横向一致性。市场收入池
\(B_D\) 在两个分配中完全相同，故收入池形成不能单独选择公正点。证毕。
\end{proof}

\begin{proposition}[市场剩余调节与规范治理目标分离命题]
\label{prop:market-surplus-governance-separation}
\(\mathfrak C_{\mathrm{RPF}}=1\) 不推出任一归属基准下的剩余占有缺口 \(g_i^B=0\)，也不要求利润、租金或
收入差异消失。竞争、进入退出、声誉和价格可以改变议价地位与剩余份额；规范治理负责保证这些市场调节所
依赖的身份、合同、判定、披露、复核与竞争条件真实存在，而不预先规定唯一份额。若规则隐蔽、市场封锁或
退出不真实，则不是由治理直接把剩余改为零，而是相对公正或市场可竞争条件不能通过。
\end{proposition}

\begin{proof}[证明（非蕴含反例）]
构造一个规则事前公开、等价参与同份额、支付可重放且参与者可自由选择竞争平台的赛事；主办者在支付声明
份额并承担固定成本和风险后仍取得正利润。此时 \(\mathfrak C_{\mathrm{RPF}}=1\) 而某一归属基准下可以有
\(g_i^B>0\)，故相对公正不蕴含零剩余占有。反向地，一个平均分完收入池的机制仍可由虚假身份或秘密规则
操纵，故零利润或零差额也不蕴含相对公正。证毕。
\end{proof}

\begin{theorem}[规范治理双层条件定理]
\label{thm:dual-layer-governance}
参与--分配行业能够纳入本文意义下的规范治理，当且仅当所声明的内部结论通过
\(\mathfrak C_{\mathrm{RPF}}\)，所允许的外部路径属于空间释放极大的安全合法域，并且市场配置仍有至少两条
可行路径。内部核验记录回答“参与和分配能否在声明规则内重放、复算一致并相对公正地判定”，外部核验记录回答“在
风险约束内允许哪些合法试验路径，以及剩余风险是否受到约束”；二者相互独立且不可替代。
\end{theorem}

\begin{proof}[证明（双核验记录乘积）]
定义联合核验记录
\(\mathfrak C_{\mathrm{norm}}=\mathfrak C_{\mathrm{RPF}}\times\mathfrak C_{\mathrm{safe}}\times
\mathfrak C_{\mathrm{paths}}\)。三个坐标分别约束内部判定、外部风险和非计划市场选择。内部公正而外溢失控、
外部安全而内部任意分配、或只剩唯一行政路径，都构成单坐标通过而联合结论失败的反例。故只有乘积核验记录全部
通过时才能确认规范治理结论。证毕。
\end{proof}

\begin{remark}[治理型经济的构成性含义]
双层条件检验已经形成的制度是否合格；定义~\ref{def:constitutive-institution-operator} 还说明治理可以在市场
形成之前提供资格、证据、确权、支付、申诉与安全路径等基础条件。因此，本文的治理型经济同时包含构成性制度
供给和纠正性约束。前者使新型参与能够成为可执行的市场对象，后者限制内部任意分配与外部风险；两者都不把
分散选择改写为中心计划，也都不保证市场需求、均衡稳定或福利改进必然出现。
\end{remark}

\section{围棋行业作为参与型市场经济的存在模型}

\begin{definition}[围棋行业经济对象]
令围棋行业参与者集合为
\[
  \Part_G=P_G\cup O_G\cup L_G\cup S_G\cup K_G,
\]
其中分别表示棋手、观众、教学者、赛事与平台服务者、生产侧赞助者。行业可消费对象包括对局、训练、赛事、
观赏、解说、棋谱、软件和社群服务。围棋行业自身提供服务生产，但不要求新增物质产出；当
\(|K_G|\ge2\) 且赞助者来自相互竞争的生产企业时，它还可以成为注意力型参与--分配桥接行业。
令 \(n\) 为有效参与规模，\(\lambda(n)\) 为匹配率，定义净剩余
\[
  Y_G(n)
  =U_{\mathrm{activity}}(n)+U_{\mathrm{watch}}(n)
  +U_{\mathrm{learn}}(n)+U_{\mathrm{community}}(n)-C_G(n)-H_G(n).
\]
普通参与者 \(i\) 的行业边际价值为
\[
  \nu_i^G=Y_G(n)-Y_G(n-e_i).
\]
其中 \(e_i\) 是参与者 \(i\) 对应的单位参与向量；若参与不可分，\(n-e_i\) 表示删除该参与者后的反事实状态。
\end{definition}

\begin{proposition}[围棋行业的双边注意力桥接命题]
\label{prop:go-attention-bridge}
若至少两家生产企业 \(k\in K_G\) 竞争性投入 \(I_k^G\)，赛事与传播形成可验证注意力
\(a_k^{G,\mathrm{att}}\)，该注意力改变消费者对赞助企业及其生产商品的偏好状态 \(\vartheta_k^P\)，且赛事收入池
按基本与差异分配合同向完成约定的棋手、教师、观赛、赛事和传播参与者支付，则围棋行业同时存在
\[
  \Pi_k^P\to I_k^G\to a_k^{G,\mathrm{att}}\to\vartheta_k^P
  \to D_k^P\to\Pi_k^{P\prime}
\]
和
\[
  \sum_k I_k^G\to\mathcal B_G^{\mathrm{pool}}\to T_i^G\to C_i^P\to D^P
\]
两条桥接路径。第一条形成消费者关于生产侧产品的偏好指引，第二条形成参与者收入与购买力；二者都必须由
真实注意力、后续购买、实际支付及成本数据分别验收。
销售回款覆盖产品成本并形成销售利润，赞助企业的偏好增量销售利润扣除赞助及配套营销成本后形成净增利润；
每名有效参与者的正基本份额及贡献差异按定理~\ref{thm:base-differential-coverage} 验收。
\end{proposition}

\begin{Certificate}[围棋注意力桥接核验记录]
核验记录至少包含两家互为竞争者的赞助投入、赛事曝光口径、机器人与重复观看剔除规则、赞助披露、赞助前后的
产品考虑集或购买变化、产品销售利润与营销成本、赛事成本、收入池、基本与差异份额和参与者实际支付。
若赞助只提高围棋本身的报名或观看，却不能形成
对赞助企业产品的偏好指引，则它可以是围棋服务需求，却不确认生产侧注意力桥接结论。
\end{Certificate}

\begin{proof}[证明（定理实例化）]
把定理~\ref{thm:competitive-attention-injection} 中的桥接行业取为围棋赛事及其传播网络，把生产者集合取为
\(K_G\)。第一条路径由注意力、偏好与产品需求的链式映射得到，第二条路径由赛事收入池及参与支付得到。
缺少任一路径的数据只撤销相应通道，不撤销围棋作为服务行业的经济价值。证毕。
\end{proof}

\begin{proposition}[参与者正价值不要求取得收入]
\label{prop:go-participation-value}
若加入参与者 \(i\) 使匹配率、对局可得性、观赏需求或教学网络中的至少一项增加，且其增量效用超过新增
资源成本与外部损害，则 \(\nu_i^G>0\)。该结论与 \(T_i\) 的符号无关；参与者可以取得收入、无偿参与，
也可以支付会费或学费而仍对行业总剩余作正贡献。
无偿服务、购买服务与约定参与分配属于不同合同关系；进入核心参与分配合同并完成约定者取得正基本份额，
其经济价值与收入的概念区分不解除该支付义务。
\end{proposition}

\begin{Certificate}[围棋参与边际核验记录]
登记加入前后的
\[
  \Delta U_{\mathrm{activity}},\Delta U_{\mathrm{watch}},
  \Delta U_{\mathrm{learn}},\Delta U_{\mathrm{community}},
  \Delta C_G,\Delta H_G,
\]
并验证
\[
  \sum_k\Delta U_k-\Delta C_G-\Delta H_G>0.
\]
\end{Certificate}

\begin{proof}[证明（边际剩余展开）]
由定义
\[
  \nu_i^G
  =\Delta U_{\mathrm{activity}}+\Delta U_{\mathrm{watch}}
  +\Delta U_{\mathrm{learn}}+\Delta U_{\mathrm{community}}
  -\Delta C_G-\Delta H_G.
\]
代入核验记录严格不等式即得 \(\nu_i^G>0\)。式中没有 \(T_i\)，故支付不是正经济价值存在的必要条件。
证毕。
\end{proof}

\begin{proposition}[人工智能棋力与人类参与需求不等同命题]
\label{prop:ai-go-human-participation}
令 \(A_G\) 表示机器棋力与机器棋谱供给，\(h_G\) 表示真人身份、竞争悬念、成长经历和社群关系，围棋需求为
\[
  D_G=D_G(A_G,h_G,p,z).
\]
即使 \(A_G\) 达到超人水平，只要在相关状态仍有
\[
  \frac{\partial D_G}{\partial h_G}>0,
\]
就不能由机器棋力提高推出真人参与需求为零。AlphaGo 的结果证明机器高水平棋力可以存在
\cite{Silver2016AlphaGo}；国际围棋组织与持续赛事只能作为真人组织参与仍存在的边界化事实例
\cite{IGF2026}，不能单独证明全部需求函数、收入规模或未来趋势。这里的 \(D_G\) 是对围棋服务本身的需求，
不等同于消费者对赞助生产者商品的偏好 \(\vartheta^P\) 或需求 \(D^P\)；\(D_G>0\) 单独不能确认生产侧
注意力桥接结论。
\end{proposition}

\begin{proof}[证明（坐标独立性）]
机器棋力是需求函数的一项坐标，而真人参与属性是另一项坐标。若 \(\partial_{h_G}D_G>0\)，保持其他坐标
不变而增加 \(h_G\) 会提高需求。因此
\[
  A_G\uparrow
  \not\Longrightarrow
  D_G(A_G,h_G,p,z)=0.
\]
经验材料分别见证机器能力与组织参与的存在，偏导符号仍需具体市场数据识别。证毕。
\end{proof}

\begin{definition}[围棋规则层的固定规则可重放模型]
令标准棋盘为有限图 \(\mathcal B_G=(\mathcal V_{19},\mathcal E_{19})\)，对局状态、合法落子、历史记录与
终局结果分别为 \(x_t^G\)、\(A_G(x_t^G)\)、\(H_t^G\) 与 \(o_G(H_T^G)\)。在固定规则版本 \(R_G\) 下，
状态转移
\[
  x_{t+1}^G=\tau_G(x_t^G,a_t;R_G)
\]
以及提子、禁入、终局和计分均可由完整棋谱独立重放。令 \(\chi_G\) 表示由长期文化传承、公开竞赛和共同
使用形成的规则认同状态。\(\chi_G\) 可以降低参与者协调规则与接受比较基准的成本，但不替代棋谱、裁判、
合同和支付证据。
\end{definition}

\begin{proposition}[围棋参与--分配的治理就绪命题]
\label{prop:go-governance-readiness}
若围棋项目同时满足：规则版本事前固定；棋谱、用时、犯规、赛果与等级数据可独立重放；身份、反作弊与分组
条件可验证；基本份额、贡献差异池及奖金、课酬、版权、裁判和组织服务的合同映射事前声明且资金足额预留；
每名完成约定者的基本份额不受输赢影响；等价赛果或等价服务按同一规则
处理；争议可由独立裁判或机构复核并修正支付；市场存在替代赛事、教学或传播入口，则可以构造
\(\mathfrak C_{\mathrm{RPF}}^G\)，并在声明容限内确认围棋参与--分配相对公正结论。

棋局数学模型直接提高的是落子、赛果与竞技表现的可判定性；教师、解说、传播和组织者的经济贡献仍须分别
声明服务指标与合同，不能由棋盘规则自动推出唯一收入份额。历史认同提高规则合法性的共同预期，也不证明
任一具体赛事分账天然公正。
\end{proposition}

\begin{Certificate}[围棋规范治理核验记录]
核验记录保存棋盘与竞赛规则版本、完整棋谱和计时、身份与反作弊结果、分组和等级依据、裁判决定、收入池来源、
成本扣除、各角色合同、奖金或课酬映射、等价类测试、独立复算、申诉与更正、市场替代入口及文化规则依据。
棋谱只能确认规则层结果；若收入合同或服务贡献证据缺失，只撤销相应参与--分配公正结论。
\end{Certificate}

\begin{proof}[证明（规则模型实例化）]
有限棋盘、有限历史与固定 \(R_G\) 使合法状态转移和终局结果可重放；身份、反作弊和裁判证据把结果归于正确
主体；事前合同把赛果及其他已验证服务映射到收入份额；等价类、单调性、预算、复核和市场入口分别实例化
\(\mathfrak C_{\mathrm{RPF}}\) 的相应坐标。因此全部条件通过时结论成立。对局规则没有定义教师或平台的
边际收入，故不能越过服务合同直接推出全行业唯一分账。证毕。
\end{proof}

\section{社会文化背景、心理效用与内生文化生成}
\label{sec:endogenous-cultural-background}

\begin{definition}[可观测文化状态与参与沉积]
固定共同体与声明时间窗。令
\[
  C_t
  =
  \Tuple{
  C_t^{\mathrm{history}},
  C_t^{\mathrm{meaning}},
  C_t^{\mathrm{identity}},
  C_t^{\mathrm{institution}},
  C_t^{\mathrm{recognition}}
  }
  \in\RR_+^{d_C}
\]
为社会文化状态，五类坐标分别表示：可验证共同历史的存量；术语、隐喻和叙事的共享程度；角色、流派与群体
身份的可识别程度；礼仪、赛事、教育和传承制度的持续性；共同体对前述对象的实际承认范围。各坐标必须由
事前指标测量，例如可重放记录、持续参与者、跨期复用术语、制度存续、独立共同体采用和代际传播，不以平台
自行发布的世界观文字或注册账户数量替代。

令 \(Z_t^C\in\RR_+^{d_Z}\) 汇总本期已经验证且自愿发生的对局、创作、协作、探索、教学、复盘、公共叙事、
规则内争议与解决记录。文化状态按
\[
  C_{t+1}
  =
  (I-\Delta_C)C_t+\Phi_C(Z_t^C;R_t),
  \qquad
  \Delta_C=\operatorname{diag}(\delta_1,\ldots,\delta_{d_C}),
  \quad
  0<\delta_\ell\le1
\]
更新，其中 \(\Phi_C\) 是在规则版本 \(R_t\) 下由参与记录形成的非负文化沉积，\(\delta_\ell\) 表示遗忘、
失效或共同认同流失。若行为无法归属、历史无法重放、记录可被静默删除，或者所谓认同仅由平台单方宣布，则
相应 \(\Phi_C\) 坐标不得记为正。
\end{definition}

\begin{definition}[文化意义映射与心理--市场效用]
令 \(\mathcal X_h^L(B_h^L)\) 为居民 \(h\) 在生活性超额预算 \(B_h^L\) 下可负担且合法进入的选择集，
\(G_h\) 为其自愿关联的共同体。文化意义映射
\[
  \mu_{C_t}:
  \mathcal X_h^L(B_h^L)\times G_h
  \longrightarrow\RR^{d_\mu}
\]
把参与对象映射为共同体能够理解的成就、身份、归属、声誉、传统延续和共同记忆特征。定义文化中介的生活性
效用
\[
  W_h(x;C_t,G_h)
  =
  u_h^{\mathrm{intrinsic}}(x)
  +v_h\!\left(\mu_{C_t}(x,G_h)\right),
  \qquad x\in\mathcal X_h^L(B_h^L).
\]
若 \(v_h\) 关于至少一个文化意义坐标严格递增，并且该坐标由两个以上相互独立的真实参与主体承认，则称该
对象对居民 \(h\) 具有正文化中介效用。该定义不把高价格、平台稀缺、强制曝光或自报身份直接定义为文化价值。
\end{definition}

\begin{theorem}[重复参与形成非零文化存量的递推定理]
\label{thm:cultural-stock-accumulation}
若规则版本在每个结算期内固定，\(\Phi_C(Z_t^C;R_t)\) 可由参与日志独立重放，且
\(0<\delta_\ell\le1\)，则对任意 \(T\ge1\)，
\[
  C_T
  =
  (I-\Delta_C)^TC_0
  +
  \sum_{t=0}^{T-1}
  (I-\Delta_C)^{T-1-t}
  \Phi_C(Z_t^C;R_t).
\]
若存在至少一个坐标 \(\ell\) 和无限多个时期 \(t\)，使
\(\Phi_{C,\ell}(Z_t^C;R_t)>0\)，则该坐标不由技术复制能力提高自动归零。若规则与参与环境稳定且
\(\Phi_C(Z_t^C;R_t)\to\bar\Phi_C\)，则
\[
  C_t\longrightarrow C^*=\Delta_C^{-1}\bar\Phi_C.
\]
只要 \(\bar\Phi_C\ne0\)，长期文化状态非零。因而，稳定底层规则与可沉积的长期参与足以生成一个非零文化
存量；文化不需要由中心机构一次性预制，但也不会仅因场景上线而自动出现。
\end{theorem}

\begin{proof}[证明（线性非齐次递推）]
对文化更新式逐期代入，得到所列有限和。由于 \(I-\Delta_C\) 的全部特征值位于 \([0,1)\)，旧文化存量按
坐标衰减，而每一期经验证参与产生的非负沉积进入后续全部状态。稳定环境下，固定点满足
\[
  C^*=(I-\Delta_C)C^*+\bar\Phi_C,
\]
移项即得 \(C^*=\Delta_C^{-1}\bar\Phi_C\)。若 \(\bar\Phi_C\ne0\)，至少一个坐标严格为正。结论只证明
文化存量的存在与递推，不证明其内容善良、福利最优或不会发生文化冲突。证毕。
\end{proof}

\begin{theorem}[社会文化背景中介心理效用与市场需求定理]
\label{thm:culture-mediated-utility-demand}
固定居民 \(h\)、共同体 \(G_h\) 与生活性对象 \(j\)。令可验证参与强度 \(z_j\) 同时进入对象的内在体验与
文化状态，并明确写成
\[
  W_{hj}(z_j)
  =u_h^{\mathrm{intrinsic}}(j,z_j)
  +v_h\!\left(\mu_{C_t(z_j)}(j,G_h)\right).
\]
若
\[
  \frac{\partial u_h^{\mathrm{intrinsic}}}{\partial z_j}\ge0,
  \qquad
  \frac{\partial C_t}{\partial z_j}\ge0,\qquad
  D_C\mu_{C_t}(j,G_h)\frac{\partial C_t}{\partial z_j}\ge0,
  \qquad
  \nabla v_h\cdot
  D_C\mu_{C_t}(j,G_h)\frac{\partial C_t}{\partial z_j}>0,
\]
则
\[
  \frac{\partial W_{hj}}{\partial z_j}>0.
\]
进一步，若居民在对象 \(j\) 与外部选项之间的选择概率为
\[
  P_{hj}
  =
  \frac{\exp(\lambda_hW_{hj})}
  {1+\exp(\lambda_hW_{hj})},
  \qquad\lambda_h>0,
\]
则
\[
  \frac{\partial P_{hj}}{\partial z_j}
  =
  \lambda_hP_{hj}(1-P_{hj})
  \frac{\partial W_{hj}}{\partial z_j}>0.
\]
因此，在声明模型中存在严格传导
\[
  \text{兴趣参与}
  \longrightarrow
  \text{文化沉积}
  \longrightarrow
  \text{共同体意义与认同}
  \longrightarrow
  \text{心理效用}
  \longrightarrow
  \text{消费选择与市场需求}.
\]
文化不是需求之外的修辞变量，而是把参与经历转化为身份、归属、声誉和共同记忆效用的内生中介。
\end{theorem}

\begin{proof}[证明（链式法则与二项 Logit 导数）]
由效用定义，
\[
  \frac{\partial W_{hj}}{\partial z_j}
  =
  \frac{\partial u_h^{\mathrm{intrinsic}}}{\partial z_j}
  +
  \nabla v_h\cdot
  D_C\mu_{C_t}(j,G_h)
  \frac{\partial C_t}{\partial z_j},
\]
第一项非负，第二项严格为正，故总导数严格为正。对二项 Logit 概率求导得到
\[
  \frac{\partial P_{hj}}{\partial W_{hj}}
  =\lambda_hP_{hj}(1-P_{hj})>0.
\]
两式相乘即得需求选择概率对参与强度的严格正导数。该证明不要求所有居民共享同一文化意义，也不把选择概率
提升为跨主体可比较的幸福量。证毕。
\end{proof}

\begin{theorem}[文化背景下生活质量差异的选择集定理]
\label{thm:culture-life-quality-choice-set}
对同一居民类型、同一文化状态、同一价格和同一合法入口，令
\[
  V_h(B;C_t)
  =
  \max_{x\in\mathcal X_h^L(B)}W_h(x;C_t,G_h).
\]
若 \(B_2\ge B_1\)，则
\[
  \mathcal X_h^L(B_1)\subseteq\mathcal X_h^L(B_2),
  \qquad
  V_h(B_2;C_t)\ge V_h(B_1;C_t).
\]
若新增预算使至少一个具有严格正文化中介效用的选项首次可负担，并且其效用严格超过原最优值，则
\[
  V_h(B_2;C_t)>V_h(B_1;C_t).
\]
故在充分救济已经把生存底线移出竞争以后，收入和财富差异可以通过文化背景下不同规模的生活性可选择集合，
形成同一类型主体的生活质量与心理效用差异。该结论允许兴趣型广义奢侈品承载底线以上的差异化生活，但不
允许由消费金额直接推出不同人格主体之间的幸福排序。
\end{theorem}

\begin{proof}[证明（预算集合包含与间接效用）]
预算增加不删除任何原可行选择，故选择集包含关系成立。在较大集合上最大化同一效用函数，最优值不可能下降。
若新增集合中存在一个效用严格超过原最大值的文化有效选项，新最大值严格提高。证毕。
\end{proof}

\begin{proposition}[围棋文化存量的存在模型与新参与组织的文化形成]
\label{prop:go-cultural-stock-reference}
在固定围棋规则版本 \(R_G\) 下，完整棋谱、名局、术语、段位、流派、赛事、教学、礼仪、复盘文本和跨代共同体
构成 \(Z_t^{C,G}\) 的可验证输入；长期规则认同 \(\chi_G\) 是
\(C_t^{G,\mathrm{recognition}}\) 的一个观测坐标。若这些输入在声明窗口内持续形成正沉积，则定理
~\ref{thm:cultural-stock-accumulation} 给出非零围棋文化存量；若气、眼、势、地、厚薄、先手、弃子、劫争、
虚联通和回补等棋理经共同体使用形成稳定意义映射，则定理~\ref{thm:culture-mediated-utility-demand} 允许
它们进入身份、审美、归属和人生隐喻的文化中介效用。

围棋由此提供“简单稳定规则--可重放参与--共同历史--文化认同--持续需求”的存在模型。新参与组织可以参考这种
生成结构，却不能继承围棋已经积累的文化存量或合法性；每个新形成的参与共同体必须用自身可验证参与、共同记忆和
代际延续形成自己的 \(C_t\)。
\end{proposition}

\begin{proof}[证明（既有输入嵌入文化递推）]
把固定规则版本下的棋谱、赛事、术语、段位、流派、教学、礼仪与共同体采用记录组成
\(Z_t^{C,G}\)，并把规则认同保存为 \(C_t^G\) 的一个坐标。若声明窗口内相应沉积函数在至少一个坐标持续为正，
则定理~\ref{thm:cultural-stock-accumulation} 直接给出非零文化存量。若共同体对棋理形成稳定意义映射，且该映射
满足定理~\ref{thm:culture-mediated-utility-demand} 的导数条件，则文化存量进入心理效用和选择概率。反之，
新参与组织的初态不能把围棋的 \(C_t^G\) 当作自身状态输入；没有自身的可验证事件与正沉积，只能得到零或外生标签，
不能由围棋案例推出其自身文化已经形成。证毕。
\end{proof}

\begin{Certificate}[文化生成、心理效用与市场传导核验记录]
核验记录保存规则版本、参与主体、原始行动、历史记录、文化沉积函数及衰减率、共同体采用样本、文化意义映射、
身份与归属指标、生活性预算和选择集、心理结果、消费选择与实际支付。至少分别检验：平台宣布与共同体承认；
短期曝光与跨期沉积；内在体验与文化中介效用；消费金额与同类型间接效用；围棋既有文化存量与新参与组织的初生
文化存量。任一映射缺失，只撤销依赖该映射的文化或心理结论，不撤销已经证明的规则合法性或局部消费价值。
\end{Certificate}

\section{共同形成的剩余、价值归属及其与规范治理的分离}

\begin{definition}[合作剩余与贡献基准]
令 \(v:2^{\Part}\to\RR\) 为已经扣除共同资源成本与可识别外部损害的合作剩余函数，且 \(v(\varnothing)=0\)。
为避免把协同价值重复计算，选定并公开一种归属基准 \(B\)。例如 Shapley 基准为
\[
  \sigma_i(v)
  =\sum_{S\subseteq\Part\setminus\{i\}}
  \frac{|S|!(n-|S|-1)!}{n!}
  \bigl[v(S\cup\{i\})-v(S)\bigr],
\]
并满足
\[
  \sum_i\sigma_i(v)=v(\Part).
\]
Shapley 值提供一种具有明确公理的归属基线，不是唯一合法的伦理分配标准
\cite{Shapley1953Value}。
\end{definition}

\begin{definition}[合同剩余索取与差异化收入]
\label{def:neutral-surplus-exploitation}
给定事前公开的固定补偿 \(\omega_i\ge0\)，令扣除真实资源成本和已经货币化外部损害后的合作剩余为
\(v(\Part)\)，可供剩余索取者分配的正余额为
\[
  R^B
  =\left[v(\Part)-\sum_{i\in\Part}\omega_i\right]_+.
\]
令 \(\rho=(\rho_i)_{i\in\Part}\in\Delta^{n-1}\) 为合同、产权和市场议价形成的剩余索取份额，定义
\[
  X_i^B=\rho_iR^B,
  \qquad
  \sum_iX_i^B=R^B.
\]
当 \(X_i^B>0\) 时，主体 \(i\) 依合同取得正剩余收入。资本所有者、组织者、平台、承担风险者
或高能力参与者都可以成为剩余索取者。份额可以相同也可以不同；其来源、激励作用与分配公正分别检验。
利润、租金、能力溢价和组织者剩余须按具体合同与成本归属区分，不能只凭正收入推断强制占有或不公正。
\end{definition}

\begin{theorem}[正剩余索取在声明域内是非退化自由市场激励的必要条件]
\label{thm:residual-claim-incentive-necessity}
令承担资本、组织、风险或稀缺能力投入的主体集合为 \(\mathcal I^R\)。主体 \(i\) 选择
\(e_i\in[0,\bar e_i]\)，投入成本 \(c_i(e_i)\) 连续可微、严格递增且 \(c_i(0)=0\)，期望正余额
\(\EE[R^B(e)]\) 连续可微。固定补偿 \(\omega_i\) 不随 \(e_i\) 变化，主体按公开合同取得份额
\(\rho_i\ge0\)，其期望收益为
\[
  U_i(e)
  =\omega_i+\rho_i\EE[R^B(e)]-c_i(e_i).
\]
若分散、自愿的市场选择要实现一个非退化内点均衡 \(e_i^*>0\)，且
\[
  \frac{\partial\EE[R^B(e^*)]}{\partial e_i}>0,
  \qquad
  c_i'(e_i^*)>0,
\]
则必有
\[
  \rho_i
  \frac{\partial\EE[R^B(e^*)]}{\partial e_i}
  =c_i'(e_i^*)>0,
  \qquad
  \rho_i>0.
\]
只要成功状态以正概率满足 \(R^B(e^*)>0\)，就有
\(\Pr(X_i^B>0)>0\)。因此，在本文的分散投资--组织域内，禁止所有私人主体取得正剩余的制度不能依靠自由
市场内生实现该非退化投入；正剩余索取是资本投入、风险承担、组织建设和稀缺能力
在该定理域内发挥作用的必要激励条件。
\end{theorem}

\begin{proof}[证明（私人最优的一阶条件与零索取反例）]
内点最优的一阶条件为
\[
  \frac{\partial U_i}{\partial e_i}
  =
  \rho_i
  \frac{\partial\EE[R^B(e)]}{\partial e_i}
  -c_i'(e_i)
  =0.
\]
由于边际余额贡献和边际成本均严格为正，等式要求 \(\rho_i>0\)。若制度禁止全部私人正剩余索取，则
\(\rho_i=0\)，从而
\[
  U_i(e_i)-U_i(0)=-c_i(e_i)<0
  \quad(e_i>0),
\]
主体的唯一最优选择是零投入，与非退化内点均衡矛盾。固定工资、行政命令或外部捐赠可以另行诱导活动，但
它们不构成本定理所研究的剩余反馈型自由市场激励。成功状态中 \(R^B>0\) 与 \(\rho_i>0\) 直接给出
\(X_i^B>0\)。证毕。
\end{proof}

\begin{proposition}[剩余分享合同诱导正投入的充分条件]
\label{prop:residual-share-positive-effort}
设有限投入者的合同剩余为
\[
 R(e)=R_0+\sum_i\left(a_i e_i-\frac{b_i}{2}e_i^2\right),\qquad
 e_i\in[0,\bar e_i],\quad R_0\ge0,\ a_i>0,\ b_i\ge0,
\]
并在 \(b_i>0\) 时取 \(\bar e_i\le2a_i/b_i\)，故 \(R(e)\ge0\)。固定份额
\(\rho_i>0,\sum_i\rho_i\le1\)，其余剩余归属另行登记；投入者的效用为
\[
 U_i=\omega_i+\rho_iR(e)-f_ie_i-\frac{k_i}{2}e_i^2,\qquad f_i>0,\ k_i>0.
\]
这里投入成本表示未由固定补偿抵销的私人机会成本，货币支付仍受同一剩余总额约束。若
\[
 0<\rho_i a_i-f_i<(\rho_i b_i+k_i)\bar e_i\quad(\forall i),
\]
则存在唯一内点纳什均衡
\[
 e_i^*=\frac{\rho_i a_i-f_i}{\rho_i b_i+k_i}>0.
\]
自愿参与还要求 \(U_i(e^*)\ge u_i^{\mathrm{out}}\)，须连同固定补偿预算单独检查。
\end{proposition}

\begin{proof}[证明]
自身效用的二阶导数为 \(-(\rho_i b_i+k_i)<0\)，一阶导数为
\(\rho_i a_i-f_i-(\rho_i b_i+k_i)e_i\)。所列严格不等式使唯一零点位于行动区间内部；
其他人的投入只进入与 \(e_i\) 无关的常数项，故这些唯一最优回应共同构成唯一纳什均衡。
若仅有 \(\rho_i>0\) 而 \(\rho_i a_i\le f_i\)，最优投入仍可为零。例如
\(U(e)=\tfrac12(1+e)-2e\) 在 \([0,1]\) 上严格递减。因而正剩余份额是前一定理域内的必要条件，
正边际净回报与严格凹性才构成本命题的充分条件。证毕。
\end{proof}

\begin{proposition}[剩余份额的单人投入比较]
\label{prop:residual-share-monotone-response}
固定他人行动，令自身选择域为紧区间，\(R(e)\)、\(c(e)\) 连续，\(R\) 非减。
对效用 \(\omega+\rho R(e)-c(e)\)，最优投入集合的最大与最小选择均随 \(\rho\) 弱增。
若 \(R\) 严格递增，则任取 \(\rho_H>\rho_L\) 下的最优解，均有 \(e_H\ge e_L\)。
若 \(c(0)=0<c(e)\) 对所有 \(e>0\) 成立，零份额下唯一最优投入为零。
该比较固定成本、补偿和其他回报，不直接给出策略相互作用下的均衡比较。
\end{proposition}
\begin{proof}[证明]
连续性与紧性保证最优解存在。将两份额下的最优性不等式相加，得
\((\rho_H-\rho_L)[R(e_H)-R(e_L)]\ge0\)。严格递增时即排除 \(e_H<e_L\)。
非严格递增时，若 \(e_H<e_L\)，非减性迫使两点 \(R\) 相同，两个最优性不等式再迫使成本相同；
故两点同时为两个份额的最优解，不可能使最大或最小最优选择逆序。零份额结论直接来自严格正成本。
\end{proof}

\begin{corollary}[能者多得与差异化收入]
\label{cor:ability-differential-residual}
若能力参数 \(a_i>0\) 使
\[
  \frac{\partial^2\EE[R^B(e,a)]}{\partial e_i\partial a_i}>0,
\]
在投入和其余条件相同、边际贡献使用共同计量规则的比较域内，更高能力形成更高边际剩余产出。
若已验证贡献得分严格保留这一顺序，并采用定义~\ref{def:base-differential-distribution} 的正差异池，
则更高贡献者取得严格更多的组织分配；所有完成约定者的正基本份额同时保留。其他剩余合同在支付对可比
已验证贡献严格递增时具有同样结论。利润、能力溢价和收入差异由公开的贡献与合同映射形成，并受自愿交易、
竞争、价格、声誉、议价和进入退出调节。
\end{corollary}

\begin{proof}[证明（递增差异与严格贡献分配）]
在共同可比域内对能力参数积分，正交叉偏导给出更高能力对应的严格更高边际剩余产出；验证映射保留其严格
顺序。定理~\ref{thm:base-differential-coverage} 的差额公式给出 \(T_{it}^B>T_{jt}^B\)，正基本份额
由同一定理保证。其他合同的结论由严格递增支付直接得到。能力标签本身不替代实际贡献证据；缺少贡献差异、
正差异池或严格递增合同之一时，不确认严格多得结论。证毕。
\end{proof}

\begin{definition}[竞争性市场中的剩余索取核验记录与非市场强制占有]
记合同剩余索取的市场资格核验记录为
\[
  \mathfrak C_{\mathrm{NE}}
  =
  \left\langle
  \begin{gathered}
  \operatorname{VoluntaryContract},
  \operatorname{ExAnteRuleDisclosure},
  \operatorname{CompetitiveAlternatives},\\
  \operatorname{ComparablePrices},
  \operatorname{RealExit},
  \operatorname{CostRiskTraceability},
  \operatorname{ResidualPaymentReplay}
  \end{gathered}
  \right\rangle.
\]
全部坐标通过时，\(X_i^B>0\) 被识别为竞争性市场中的剩余索取。若正余额主要由欺诈、行政特许、秘密改写、
垄断封锁、不可携带身份或虚假退出维持，则记为
\(\operatorname{NonMarketCoerciveAppropriation}\)；它不是“剥削数额过大”的同义语，而是自由市场形成条件
失效的独立结论。
\end{definition}

\begin{proposition}[合同剩余索取与非市场强制占有不可混同命题]
\label{prop:appropriation-exploitation-separation}
平台利润、组织者剩余、高能力收入或其他 \(X_i^B>0\) 已经构成定义
~\ref{def:neutral-surplus-exploitation} 下的合同剩余索取，但均不单独推出非市场强制占有。若
\(\mathfrak C_{\mathrm{NE}}\) 通过，则正剩余索取与自由市场、能者多得和规范参与--分配相容；若任一关键
市场形成坐标失败，只能撤销相应市场资格并调查强制占有，不能把所有正利润或收入差异一并归零。
\end{proposition}

\begin{Certificate}[剩余索取、市场形成与强制占有分离核验记录]
核验记录保存合作剩余函数及基准版本、固定补偿与真实成本、资本和风险承担、能力与边际贡献、剩余份额和实际支付、
合同知情与自愿、规则修改权限、替代平台和可比价格、数据与身份可携带性、退出损失、行政特许、市场集中度和
可能的欺诈证据。它分别回答“谁取得正剩余”“该剩余是否形成投入激励”“市场条件是否真实”，不得用任何
一个问题的答案替代另外两个。
\end{Certificate}

\begin{proof}[证明（三个谓词的独立反例）]
构造平台承担固定成本和失败风险、规则事前公开、参与者可选择多个平台且可真实退出的赛事。平台取得
\(X_P^B>0\)，故存在合同剩余索取；同时 \(\mathfrak C_{\mathrm{NE}}=1\)，不存在非市场强制占有。反之，一个
平台可以账面零利润，却利用秘密关联交易、不可携带身份或行政排他把价值转移给外部主体，故零利润不推出
市场条件成立。再取具有正份额却不增加边际剩余的主体，可见正索取也不自动证明激励有效。三个谓词相互独立，
必须分别验收。证毕。
\end{proof}

\begin{remark}[治理不以本文定义的合同剩余索取归零为目标]
规范治理不把 \(X_i^B=0\)、收入相等或财富相等设为市场准入条件，也不直接规定“适当”的剩余比例。竞争、
议价、声誉、价格和进入退出负责调节剩余份额；社会充分救济承接生存与市场接入下限；监管只维护合同真实性、
公平竞争、可复核分配和风险不向无关主体溢出的边界。由此，在定理~\ref{thm:residual-claim-incentive-necessity}
的声明域与竞争性市场中的剩余索取核验记录同时通过时，非退化私人投入需要正剩余索取；竞争调节其强度，最低救济守住
声明的生存--市场接入底线，规范治理阻止非市场强制占有和风险外溢。
\end{remark}

\chapter{市场内制度选择、有效参与与初始收益权}
\label{chap:governance-participation}

\section{组织合同在固定收入池上的分配选择}

\begin{definition}[预算平衡分配单纯形]
在产能约束已松弛的产品域，令可分配真实剩余为 \(Y^*>0\)。定义
\[
  \DeltaN=\left\{s\in\RR_+^n:\sum_{i=1}^n s_i=1\right\},
\qquad
  y_i(G)=s_i(G)Y^*.
\]
其中 \(Y^*\) 已由市场交易和真实成本共同形成，\(s(G)\) 是声明组织依据合同或章程在其内部收入池上
选择的份额向量。该映射不决定市场总产量、市场价格或组织外主体的配置。
在核心参与合同中取 \(Y^*=N_t\beta_t^{\mathrm{base}}+D_t^{\mathrm{diff}}\)，允许的份额还满足
\(s_iY^*\ge\beta_t^{\mathrm{base}}\) 及可比贡献的严格差异约束；定义
~\ref{def:base-differential-distribution} 给出该受约束域中的明确实现。一般单纯形描述预算可行性，
核心合同同时落实正基本份额与能者多得。
\end{definition}

\section{固定收入池的预算可行支付}

\begin{proposition}[预算平衡支付向量的构造]
\label{thm:any-simplex-allocation-implementable}
对任意 \(s\in\DeltaN\)，令 \(T_i=s_iY^*\)，则 \(T_i\ge0\) 且
\(\sum_iT_i=Y^*\)。若 \(n\ge2\)，取 \(s(t)=(t,1-t,0,\ldots,0)\)、\(t\in[0,1]\)，
即得连续统多个预算可行支付。该结论不包含私人信息下的激励相容、参与或均衡实施。
\end{proposition}
\begin{proof}[证明]
由 \(s_i\ge0\) 与 \(\sum_i s_i=1\)，非负性及总额恒等式直接成立。不同 \(t\) 给出不同支付向量。
当某人的外部选择高于指定支付或报告能够改变其份额时，预算可行仍不保证其参与或真实报告。
\end{proof}
\begin{Certificate}[预算核对与行为实施分离]
\label{cert:any-simplex-allocation-implementable}
保存 \(Y^*,s,T\) 及总额核对；激励、参与和实际支付另提交对应行为模型与回执，不由本命题代替。
\end{Certificate}

\begin{theorem}[组织分配规则比较静态定理]
\label{thm:governance-comparative-statics}
保持市场形成的 \(Y^*\) 不变，只把组织分配规则从 \(G_1\) 改为 \(G_2\)，则
\[
  \Delta y_i
  =
  y_i(G_2)-y_i(G_1)
  =
  Y^*\bigl[s_i(G_2)-s_i(G_1)\bigr].
\]
因此，只有在 \(\RuleSensitive_i(G_1,G_2)\), 即
\(s_i(G_1)\ne s_i(G_2)\), 成立时，规则变化才对参与者 \(i\) 产生非零分配效应。
\end{theorem}

\begin{Certificate}[固定总量与份额差]
核验记录为
\[
  \mathfrak C_{\Delta G,i}
  =
  \Tuple{Y^*,G_1,G_2,s_i(G_1),s_i(G_2)}.
\]
非零效应验收谓词为
\[
  Y^*>0\wedge s_i(G_1)\ne s_i(G_2).
\]
\end{Certificate}

\begin{proof}[证明]
由定义
\[
  y_i(G_k)=s_i(G_k)Y^*,\qquad k\in\{1,2\}.
\]
两式相减得到结论。进一步，
\[
  Y^*>0\wedge s_i(G_1)\ne s_i(G_2)
  \Longrightarrow
  \Delta y_i\ne0.
\]
若份额规则不变，则 \(\Delta y_i=0\)；故不能仅由 \(G_1\ne G_2\) 推出每个个体收入必变，
必须附加规则敏感性见证。证毕。
\end{proof}

\section{有效参与到收益权的有限状态链}

\begin{definition}[参与确权机制]
参与者选择 \(a_i\in\Act_i\)，计量信号为
\[
  z_i=M_G(a_i,\varepsilon_i).
\]
真实性条件为 \(v_i=V_G(z_i,\omega_i)\in\{0,1\}\)，结构或贡献权重为
\[
  w_i=W_G(z_i)v_i\ge0.
\]
初始收益权和支付分别为
\[
  b_i=\phi_G(w_i),\qquad
  T_i=T_G(b_i).
\]
完整状态链定义为
\[
  a_i
  \xrightarrow{E_G}
  e_i
  \xrightarrow{M_G}
  z_i
  \xrightarrow{V_G}
  v_i
  \xrightarrow{W_G}
  w_i
  \xrightarrow{\phi_G}
  b_i
  \xrightarrow{T_G}
  y_i.
\]
\end{definition}

\begin{lemma}[逐项条件不可跳越引理]
\label{lem:atomic-gates-no-skip}
若
\[
  b_i=\phi_G(W_G(z_i)V_G(z_i,\omega_i))
\]
且 \(\phi_G(0)=0\)，则
\[
  \neg\Verified_i(z_i;G)
  \Longrightarrow
  b_i=0.
\]
一般收益映射中，\(\Verified_i\) 还须与正估值共同作用才能产生正权利。核心竞技组织对
\(i\in\mathcal E_t\) 使用
\[
  W_G(z_i)=\beta_t^{\mathrm{base}}+D_t^{\mathrm{diff}}\lambda_{it},
  \qquad \phi_G(w)=w,
\]
其中 \(z_i\) 包含个人记录及可复算的公共资金池、贡献汇总和规则引用。于是完成约定且通过验证的有效参与
具有 \(b_i\ge\beta_t^{\mathrm{base}}>0\)，与定理~\ref{thm:base-differential-coverage} 一致。
\end{lemma}

\begin{Certificate}[零吸收与正值双条件]
核验记录条件为
\[
  V_G\in\{0,1\},\qquad W_G\ge0,\qquad \phi_G(0)=0.
\]
正收益权核验记录另需
\[
  V_G=1\wedge W_G(z_i)>0\wedge\phi_G(W_G(z_i))>0.
\]
\end{Certificate}

\begin{proof}[证明]
若 \(\neg\Verified_i\)，则 \(V_G=0\)，所以
\[
  V_G=0
  \Longrightarrow
  W_G(z_i)V_G=0
  \Longrightarrow
  b_i=\phi_G(0)=0.
\]
在一般映射中，\(V_G=1\) 且 \(W_G(z_i)=0\) 仍会产生零权利；核心竞技组织由正基本份额保证每名有效
参与者的 \(W_G(z_i)>0\)，恒等确权映射随即给出正收益权。两类结论分别服从其已声明的机制域。证毕。
\end{proof}

\begin{proposition}[激励相容参与命题]
\label{prop:participation-incentive-compatibility}
令参与者效用为
\[
  u_i(a_i,a_{-i};G)
  =
  E_G(i,a_i)\,
  \EE\!\left[
  T_G\!\left(
  \phi_G\!\left(
  W_G(z_i(a_i))V_G(z_i(a_i),\omega_i)
  \right)\right)\right]
  -c_i(a_i)-\rho_i(a_i,a_{-i}),
\]
其中 \(z_i(a_i)=M_G(a_i,\varepsilon_i)\)，\(\rho_i\) 为操纵、处罚或拥堵风险。行为 \(a_i^*\) 是激励相容选择，当且仅当
\[
  u_i(a_i^*,a_{-i};G)
  \ge
  u_i(a_i,a_{-i};G)
  \quad
  (\forall a_i\in\Act_i).
\]
个体理性另要求
\[
  u_i(a_i^*,a_{-i};G)\ge u_i^{\mathrm{out}}.
\]
\end{proposition}

\begin{Certificate}[逐偏离不增益]
对有限行动集，核验记录为偏离差向量
\[
  \mathfrak C_{\mathrm{IC},i}
  =
  \bigl(
  u_i(a_i^*)-u_i(a_i)
  \bigr)_{a_i\in\Act_i}.
\]
若每一坐标非负，则 \(\IC_i\)；若再有
\(u_i(a_i^*)-u_i^{\mathrm{out}}\ge0\)，则 \(\IR_i\)。
\end{Certificate}

\begin{proof}[证明]
逐偏离核验记录满足
\[
  \forall a_i,\quad
  u_i(a_i^*)-u_i(a_i)\ge0
  \Longrightarrow
  a_i^*\in\argmax_{a_i\in\Act_i}u_i(a_i,a_{-i};G).
\]
这正是给定他人行为下的激励相容条件。再由
\[
  u_i(a_i^*)\ge u_i^{\mathrm{out}}
  \Longrightarrow
  \text{参与弱优于退出},
\]
得到个体理性。机制设计中的激励相容必须由这种逐偏离约束证明，不能由“规则公开”替代
\cite{Myerson1981OptimalAuction}。证毕。
\end{proof}

\begin{corollary}[二元参与阈值推论]
\label{cor:binary-participation-threshold}
若 \(a_i\in\{0,1\}\)、不参与效用归一为零，则持续参与的必要条件为
\[
  \EE[T_G\mid a_i=1]-\EE[T_G\mid a_i=0]
  \ge
  c_i(1)+\rho_i(1,a_{-i}).
\]
若该不等式严格成立且状态过程平稳，则参与是严格最优回应。
\end{corollary}

\begin{Certificate}[参与净增益]
定义
\[
  \Delta u_i^{\mathrm{part}}
  =
  \EE[T_G\mid a_i=1]-\EE[T_G\mid a_i=0]
  -c_i(1)-\rho_i(1,a_{-i}).
\]
核验记录状态为
\[
  \Delta u_i^{\mathrm{part}}\ge0
  \quad\text{或}\quad
  \Delta u_i^{\mathrm{part}}>0.
\]
\end{Certificate}

\begin{proof}[证明]
二元行动下只有一个偏离：
\[
  u_i(1)-u_i(0)=\Delta u_i^{\mathrm{part}}.
\]
故
\[
  \Delta u_i^{\mathrm{part}}\ge0
  \Longrightarrow
  u_i(1)\ge u_i(0),
\]
严格不等式则推出唯一严格最优回应。证毕。
\end{proof}

\begin{remark}[劳动分配机制与经济价值的关系]
传统按劳分配若主要依赖可加总工时、件数和质量，可以调用命题~\ref{prop:labor-distribution-special-case} 的
低维核验记录；劳动市场的价格比较和合同执行降低判定负担，但不自动排除买方垄断、歧视或退出锁定。劳动、创作、
审核、传播、数据提供、教学、游戏竞技或行业规则维护都可以成为 \(a_i\) 的元素；后几类参与往往具有协同、
网络和非加总性，不能直接复制“工时乘工资率”。只有通过资格、计量、验证、组织估值、确权与支付的完整链条，
行为才转化为可执行收入。市场内制度规则决定何种可验证参与能够依合同生成收入请求权；参与是否已经产生正
经济价值仍由净社会剩余的基线差判定，并受需求、成本、真实性和外部损害约束；规范治理是否成立，则另由
\(\mathfrak C_{\mathrm{RPF}}\) 判定。
\end{remark}
\clearpage
\chapter{收益权确立、货币支付与可选金融安排}
\label{chap:primary-finance}

组织依约定完成核验、确权和货币支付，即可使参与收入进入一般市场。下列收益权转让、流动性和金融工具
分析均以当事人实际选择相应业务为适用前提。没有二级转让时取 \(\Delta b_i=0\) 及该项转让成本
\(f_i=0\)，直接支付已确定的基本与差异份额；证券化与金融工具采用比例不构成该直接支付路径的条件。

\section{一级分配与二级转移}

\begin{definition}[初始收益权与二级净转移]
市场内组织制度只对满足 \(V_G(z_i,\omega_i)=1\) 的可验证参与指标生成初始收益权
\[
  b_i^0=\phi_G(W_G(z_i))\ge0,\qquad
  B^0=\sum_{i=1}^n b_i^0.
\]
在本节的封闭转移模型中，二级市场只允许参与者之间转移既有收益权，令
\[
  \delta_{ij}=-\delta_{ji},
\qquad
  \Delta b_i=\sum_{j=1}^n\delta_{ji},
\qquad
  b_i^1=b_i^0+\Delta b_i\ge0.
\]
每单位收益权支付 \(D\ge0\)，交易、清算和合规成本为 \(f_i\ge0\)，则
\[
  y_i=D(b_i^0+\Delta b_i)-f_i.
\]
若 \(b_i^0\) 不构成证券或可转让投资合同，本文将该层称为“初始确权层”或“一级分配层”。
\end{definition}

\section{参与收益权、货币支付与可选融资安排}
\label{sec:certified-participation-fiat-cycle}

\begin{definition}[治理认证的合规参与]
\label{def:governance-certified-participation}
固定司法辖区、结算期和规则版本 \((\mathcal L,t,R^\nu)\)。对参与者 \(i\) 的行为 \(a_{it}\)，定义参与治理核验记录
\[
  \mathfrak C_{it}^{\mathrm{part}}
  =
  \Tuple{
  \mathrm{id}_{it},a_{it},\ell_{it},R^\nu,
  E,M,V,W,\phi,
  \mathfrak C_{\mathrm{legal\text{-}part}},
  \mathfrak C_{\mathrm{RPF}},Q
  },
\]
其中 \(\ell_{it}\) 是不可静默覆盖的事件记录，\(E,M,V,W,\phi\) 依次给出资格、计量、真实性验证、贡献估值
和初始确权，\(\mathfrak C_{\mathrm{legal\text{-}part}}\) 证明行为、数据取得、合同和参与资格处于适用法律许可域，
\(\mathfrak C_{\mathrm{RPF}}\) 是相对公正核验记录，\(Q\) 给出审计、申诉和错误修复。令各逐项条件
\(c_{kit}^{\mathrm{part}}\in\{0,1\}\)，并定义
\[
  g_{it}^{\mathrm{part}}
  =\prod_{k\in\mathcal K_{\mathrm{part}}}c_{kit}^{\mathrm{part}}.
\]
只有 \(g_{it}^{\mathrm{part}}=1\) 时，行为才称为“经治理认证的合规参与”。这里的治理认证可以由合同授权的
组织、独立验证者、审计者或适用法律指定的主管主体分别完成；它不是政府对价格、收益或项目成功的担保。
\end{definition}

\begin{definition}[合同收益权、权利登记与可选证券化]
\label{def:rights-registration-finance}
通过参与条件的初始收益权为
\[
  b_{it}^{0}
  =\phi_t\!\left(g_{it}^{\mathrm{part}}W_t(z_{it})\right)\ge0.
\]
把参与事实依合同转化为可执行资源请求权称为\emph{权利化}。把权利主体、数量或份额、付款义务、期限、限制、
来源规则和申诉状态写入标准化记录称为\emph{凭证化}：
\[
  \beta_{it}=\tau_t\!\left(b_{it}^{0},\kappa_{it}\right),
\]
其中 \(\kappa_{it}\) 是上述法律和合同属性。登记可以由合同台账、结算清单或其他可核验凭证承担，必须
保持权利内容、责任主体和更正轨迹。登记行为本身不增加收益权金额、交易价格或付款来源。

令适用法律分类映射为
\[
  \chi_{\mathcal L,t}(\beta)
  \in
  \left\{
  \begin{gathered}
    \text{非转让参与权益},\ \text{合同支付请求权},\\
    \text{证券或集合投资权益},\\
    \text{其他受监管金融工具},\ \text{其他}
  \end{gathered}
  \right\}.
\]
只有当既有权利被汇集、分层、份额化，或者形成可转让且以未来收益为核心的投资请求权，并通过适用法律下的
发行、登记或许可、披露、适当性、托管、交易、市场完整性和风险管理条件时，本文才把
\[
  s_{kt}=\sigma_{\mathcal L,t}\bigl((\beta_{it})_{i\in I_k},\Omega_{kt}\bigr)
\]
称为\emph{条件证券化}。\(\Omega_{kt}\) 记录现金流、损失顺序和投资者权利。具体工具是否属于证券，始终由
适用法律依其合同权利、交易功能和风险确定。该分类作为声明制度域的外部约束，
不由本模型预设某项具体工具的法律地位。
\end{definition}

\begin{definition}[法币来源、支付能力与开放循环]
\label{def:fiat-source-open-cycle}
令经营性外部法币收入为
\[
  F_t^{\mathrm{op}}
  =C_t^{\mathrm{user}}+S_t^{\mathrm{producer}}+R_t^{\mathrm{contract}},
\]
分别表示消费者自愿支付、生产企业基于注意力与商品偏好回报作出的赞助或采购，以及其他可执行市场合同收入。
另令 \(K_t^{\mathrm{risk}}\) 为本期净新增且明确披露损失承担顺序的风险资本或信用，
\(H_t^{\mathrm{res}}\) 为期初未受限储备，\(U_t^{\mathrm{trans}}\) 为依法取得的公共转移或捐赠。四类来源采用
同一时点的非重复口径，必须与经营收入分账，不能冒充已经实现的市场需求。可支付法币池为
\[
  F_t^{\mathrm{pay}}
  =F_t^{\mathrm{op}}+K_t^{\mathrm{risk}}+H_t^{\mathrm{res}}+U_t^{\mathrm{trans}}.
\]
法币结算核验记录定义为
\[
  \mathfrak C_t^{\mathrm{fiat}}
  =\left\langle
  \begin{gathered}
  \mathfrak C^{\mathrm{part}},\mathfrak C^{\mathrm{claim}},
  \chi_{\mathcal L,t},\mathrm{payer},\\
  F_t^{\mathrm{op}},K_t^{\mathrm{risk}},H_t^{\mathrm{res}},U_t^{\mathrm{trans}},\\
  \mathrm{unitOfAccount},\mathrm{paymentRail},\mathrm{finality},\\
  \mathrm{custody},\mathrm{disclosure},\mathrm{KYC},
  \mathrm{AML/CFT},\mathrm{tax},\mathrm{riskIsolation}
  \end{gathered}
  \right\rangle.
\]
其中 \(\mathfrak C^{\mathrm{claim}}\) 证明付款义务、受益主体、期限、损失顺序、撤销和救济在适用法下可执行。
令 \(g_t^{\mathrm{fiat}}\) 为上述必要条件之积，实际法币毛支付为
\[
  p_{it}^{F}=g_t^{\mathrm{fiat}}\Psi_{it}(\mathbf b_t^1,F_t^{\mathrm{pay}}),
  \qquad
  \Psi_{it}\ge0,
  \qquad
  \sum_i\Psi_{it}\le F_t^{\mathrm{pay}}.
\]
令事前公开的税费和结算成本为 \(d_{it}^{F}\ge0\)，参与者实际收到
\[
  y_{it}^{F}=[p_{it}^{F}-d_{it}^{F}]_+.
\]
只有当该收入具有最终性回执并能够在项目外按法定货币单位
支付、储蓄或购买一般商品与服务时，才称其进入\emph{法币开放循环}。封闭积分、平台估值、未到账余额和只能
依靠后来者出价的内部凭证不满足该定义。若存在固定赎回承诺，还必须指定负债主体、兑付资产、期限与最终性；
若只对应可变收入池份额，则不得宣传为固定法币债权。如发生金融业务，核验记录按声明司法辖区逐项列出适用的身份识别、记录、税务、
托管、交易报告和风险隔离要求；不适用项须记录理由。
\end{definition}

\begin{theorem}[收益权登记与支付资金的双守恒定理]
\label{thm:rights-payment-conservation}
假定：
\begin{enumerate}[label=\textup{(\roman*)},leftmargin=3em]
  \item \(\phi_t(0)=0\)，且对每个支付状态 \(\omega\)，凭证化映射满足
  \(\operatorname{Payoff}(\beta_{it};\omega)=
  \operatorname{Payoff}(b_{it}^{0};\omega)\ge0\)；
  \item 若发生条件证券化，则对每个支付状态 \(\omega\)，扣除公开成本后的证券化总请求不超过基础权利总请求：
  \[
    \sum_k\operatorname{Payoff}(s_{kt};\omega)
    \le
    \sum_i\operatorname{Payoff}(\beta_{it};\omega);
  \]
  \item 法币支付满足 \(\sum_i\Psi_{it}\le F_t^{\mathrm{pay}}\)。
\end{enumerate}
则：
\begin{enumerate}[label=\textup{(\alph*)},leftmargin=3em]
  \item \(g_{it}^{\mathrm{part}}=0\Longrightarrow b_{it}^0=0\)，该行为不能生成可归属于自己的合规权利凭证；
  \item 权利登记和可选证券化都不能增加全部基础权利在任一状态下可索取的总资源；
  \item \(F_t^{\mathrm{op}}=K_t^{\mathrm{risk}}=H_t^{\mathrm{res}}=U_t^{\mathrm{trans}}=0\Longrightarrow
  \sum_i p_{it}^{F}=0\)；
  \item 若 \(F_t^{\mathrm{op}}=0\) 而其他来源至少一项为正，支付来自已披露资本、储备、公共转移或捐赠，
  不能记作市场经营收入或需求已经成立。
\end{enumerate}
因此，合同登记、权利转让及融资安排均受既有权利与实际付款来源约束；凭证数量的变化不增加可供
支付的资源，也不构成净社会剩余。
\end{theorem}

\begin{Certificate}[权利来源与法币来源双守恒]
\label{cert:right-fiat-dual-conservation}
双守恒核验记录为
\[
  \mathfrak C_{t}^{\mathrm{source}}
  =\left\langle
  \begin{gathered}
  (\mathfrak C_{it}^{\mathrm{part}})_i,
  (b_{it}^{0},\beta_{it})_i,
  (s_{kt},\Omega_{kt})_k,\\
  F_t^{\mathrm{op}},K_t^{\mathrm{risk}},H_t^{\mathrm{res}},U_t^{\mathrm{trans}},
  (p_{it}^{F})_i
  \end{gathered}
  \right\rangle.
\]
验收必须逐状态复算请求权守恒、逐账户复算法币来源与支付上界，并把经营收入、资本、信用、储备、公共转移和捐赠
分别列示。任一来源缺失、重复抵押、跨池重复计数或付款超过可支付池，核验记录即失败。
\end{Certificate}

\begin{proof}[证明（零吸收、请求权保持与预算约束复合）]
若 \(g_{it}^{\mathrm{part}}=0\)，则由 \(\phi_t(0)=0\) 有
\[
  b_{it}^{0}
  =\phi_t\!\left(g_{it}^{\mathrm{part}}W_t(z_{it})\right)
  =\phi_t(0)=0.
\]
\(\tau_t\) 仅记录该请求权，不能把零请求权变为正请求权。条件证券化的逐状态不等式进一步保证汇集、分层或
份额化不提高基础总偿付。最后，
\[
  \sum_i p_{it}^{F}
  =g_t^{\mathrm{fiat}}\sum_i\Psi_{it}
  \le F_t^{\mathrm{pay}}
  =F_t^{\mathrm{op}}+K_t^{\mathrm{risk}}+H_t^{\mathrm{res}}+U_t^{\mathrm{trans}}.
\]
四项来源均为零时总支付只能为零；只有资本、储备、公共转移或捐赠为正时，其会计身份由来源分账保持，不会因
进入支付函数而变成经营收入。证毕。
\end{proof}

\begin{theorem}[治理认证参与进入法币开放循环定理]
\label{thm:certified-participation-fiat-cycle}
对有限参与集合，若：
\begin{enumerate}[label=\textup{(\roman*)},leftmargin=3em]
  \item 每一笔正权利都具有 \(g_{it}^{\mathrm{part}}=1\) 的参与核验记录和可执行的
  \(\mathfrak C_{it}^{\mathrm{claim}}\)；
  \item 权利的法律分类、发行或付款主体、收入池、期限、损失顺序和争议救济已经确定；
  \item \(F_t^{\mathrm{pay}}>0\)，且经营收入、风险资本、信用、储备、公共转移和捐赠没有混同；
  \item 支付通道以法定货币为记账单位，具有适用法律承认的结算最终性，并完成该司法辖区对本项支付适用的身份、税务、披露
  及其他结算核验；
  \item \(\Psi_t\) 可重放、预算可行，支付回执与项目外法币账户一一对应；
  \item 实际到账集合
  \[
    I_t^+=\{i:b_{it}^{0}>0,\ p_{it}^{F}>d_{it}^{F}\}
  \]
  非空，且其中每笔支付均有结算最终性回执；
\end{enumerate}
则复合映射
\[
\begin{aligned}
  a_{it}
  &\xrightarrow{\mathfrak C^{\mathrm{part}}}
  b_{it}^{0}
  \xrightarrow{\tau_t}
  \beta_{it}
  \xrightarrow{\sigma_{\mathcal L,t}\ \mathrm{optional}}
  \widehat\beta_{it}\\
  &\xrightarrow{\Psi_t,F_t^{\mathrm{pay}}}
  p_{it}^{F}
  \xrightarrow{\mathrm{finality}}
  y_{it}^{F}
  \xrightarrow{\mathrm{general\ market}}
  C_{i,t+1}
\end{aligned}
\]
对 \(i\in I_t^+\) 存在且可审计，其中 \(\widehat\beta_{it}=\beta_{it}\) 表示不发生证券化。因而法律意义上的证券化不是参与收入
进入法币循环的必要条件；治理认证、可执行权利、真实付款来源和合规结算才是共同必要制度环节。
\end{theorem}

\begin{proof}[证明（有限函数复合与终端回执见证）]
条件 (i) 使参与事实唯一进入初始收益权，条件 (ii) 使 \(b_{it}^{0}\) 成为可执行而非仅技术可读的请求权，
条件 (iii) 给出非零且来源分离的支付域，条件 (iv) 使支付能够在法币体系内最终结算，条件 (v) 使每一笔终端
回执可以逆向追踪到原始参与、规则版本、权利和资金来源。条件 (vi) 与
\(y_{it}^{F}=[p_{it}^{F}-d_{it}^{F}]_+\) 给出 \(y_{it}^{F}>0\)，排除“只有正支付池而没有任何参与者实际到账”的
空循环。有限集合 \(I_t^+\) 上的逐项复合因此存在。若不进行
\(\sigma_{\mathcal L,t}\)，取恒等映射 \(\widehat\beta_{it}=\beta_{it}\)，其余箭头仍然成立，故证券化非必要。
实际到账 \(y_{it}^{F}\) 进入一般商品市场后，按个体预算约束成为下一期消费、储蓄或投资的可用法币资源，
开放循环成立。证毕。
\end{proof}

\begin{proposition}[治理认证非价值、价格、流动性或偿付能力充分条件命题]
\label{prop:certification-not-value-liquidity}
一般地，
\[
  g_{it}^{\mathrm{part}}=1
  \not\Longrightarrow
  \{b_{it}^{0}>0,\ p_{it}^{F}>0,\ P_t^{\mathrm{market}}>0,
  \ \mathrm{Liquid},\ \mathrm{Solvent}\}.
\]
一般合同中的正权利还需正估值与正确确权；核心参与合同由基本份额保障正收益权；正法币支付还需真实付款来源；正二级价格还需买方需求；流动性和偿付能力还需
期限、资产负债、储备和市场深度条件。反向亦不成立：投机价格为正或平台内部积分活跃，不能推出参与核验记录、
经济价值或法币偿付能力。权利转让、托管、支付及融资业务的适用约束分别列入合同与结算核验记录，
不由组织的名称或登记工具决定。
\end{proposition}

\begin{proof}[证明（相互独立的缺失后件反例）]
保持 \(g_{it}^{\mathrm{part}}=1\)：取 \(W_t(z_{it})=0\) 或令 \(\phi_t\) 在该点取零，得到
\(b_{it}^{0}=0\)；取 \(F_t^{\mathrm{pay}}=0\)，得到 \(p_{it}^{F}=0\)；取二级买方需求为零，得到正基础权利但
市场价格为零或无成交；取买卖盘深度趋零，流动性失败；取到期负债大于可用资产，偿付能力失败。每个状态均与
参与核验记录通过相容，故认证不充分推出集合中的任一后件。反向取没有参与来源的内部积分或仅由后来者出价维持的
交易价格，可以有正报价和活跃转手而 \(g_{it}^{\mathrm{part}}=0\)、基础权利或法币来源为零，故反向蕴含也失败。
证毕。
\end{proof}

\begin{Certificate}[合规参与到法币循环的分层确认]
参与核验、收益权登记与实际支付分别验收：
\[
  \mathfrak C^{\mathrm{part}},\quad
  \mathfrak C^{\mathrm{claim}},\quad
  \mathfrak C^{\mathrm{registration}},\quad
  \mathfrak C^{\mathrm{securities}}\text{（若适用）},\quad
  \mathfrak C^{\mathrm{fiat}},\quad
  \mathfrak C^{\mathrm{open\ cycle}}.
\]
前一核验记录通过不自动推出后一核验记录。\(\mathfrak C^{\mathrm{open\ cycle}}\) 另需实际到账、项目外可用性、一般消费
或投资去向和下一期市场回流见证；只有内部转账记录、内部报价或兑付宣传时不得确认。
\end{Certificate}

\begin{remark}[失败边界]
下列情形分别阻断不同箭头：参与不可验证，阻断初始确权；权利无法律义务主体，阻断可执行请求；平台把内部
积分称作资产但没有外部付款来源，阻断法币支付；合规参与存在但没有需求，阻断正市场价格；以风险资本持续
填补经营亏损并记作收入，阻断可持续性；可转让收益凭证绕过证券、市场完整性或反洗钱规则，阻断合法交易；
固定兑付缺少储备或最终性，阻断偿付核验记录。治理认证解决“何种参与有资格成为权利”，市场解决“这种权利以何种
价格被选择”，货币支付与清算安排解决“权利如何最终支付”，三者不可互相替代。
\end{remark}

\section{二级转移守恒}

\begin{theorem}[有限反对称转移守恒]
\label{thm:secondary-transfer-conservation}
若有限参与者之间的二级转移满足 \(\delta_{ij}=-\delta_{ji}\) 与 \(\delta_{ii}=0\)，则
\[
  \sum_{i=1}^n\Delta b_i=0,
\qquad
  \sum_{i=1}^n b_i^1=B^0.
\]
\end{theorem}

\begin{Certificate}[有限构造：参与者集合扩张下的净额守恒]
\label{cert:secondary-transfer-conservation}
令 \(P(k)\) 表示前 \(k\) 个参与者内部的全部成对转移净额为零：
\[
  P(k):\quad
  \sum_{i=1}^{k}\sum_{j=1}^{k}\delta_{ji}=0.
\]
基例 \(P(1)\) 由 \(\delta_{11}=0\) 给出。归纳步把第 \(k+1\) 个参与者加入：
\[
  \sum_{i=1}^{k}\bigl(\delta_{k+1,i}+\delta_{i,k+1}\bigr)=0.
\]
\end{Certificate}

\begin{proof}[证明（有限递推）]
基例
\[
  \delta_{11}=0\Longrightarrow P(1).
\]
若 \(P(k)\) 成立，则加入第 \(k+1\) 个参与者产生的新增总净额为
\[
  \sum_{i=1}^{k}
  \bigl(\delta_{k+1,i}+\delta_{i,k+1}\bigr)
  +\delta_{k+1,k+1}=0
\]
，因为每对转移反对称且对角项为零。因此
\[
  P(k)\Longrightarrow P(k+1).
\]
归纳得到 \(P(n)\)，即 \(\sum_i\Delta b_i=0\)。再由
\[
  \sum_i b_i^1
  =
  \sum_i b_i^0+\sum_i\Delta b_i
  =B^0
\]
得到收益权总量守恒。证毕。
\end{proof}

\begin{theorem}[受限流动性下的初始确权支配定理]
\label{thm:limited-liquidity-primary-dominance}
若流动性规则满足
\[
  \Abs{\Delta b_i}\le\lambda_i,
\]
则个人最终货币结果与初始确权支付之间满足
\[
  \Abs{y_i-Db_i^0}
  \le D\lambda_i+f_i.
\]
当 \(\lambda_i\to0\) 且 \(f_i\to0\) 时，
\[
  y_i\to Db_i^0=D\phi_G(W_G(z_i)).
\]
\end{theorem}

\begin{Certificate}[个体偏离上界]
核验记录元组为
\[
  \mathfrak C_{\mathrm{liq},i}
  =
  \Tuple{D,b_i^0,\Delta b_i,\lambda_i,f_i,y_i},
\]
并验证
\[
  D\ge0,\quad \Abs{\Delta b_i}\le\lambda_i,\quad
  y_i=Db_i^0+D\Delta b_i-f_i.
\]
\end{Certificate}

\begin{proof}[证明]
由收益恒等式，
\[
  y_i-Db_i^0=D\Delta b_i-f_i.
\]
三角不等式给出
\[
  \Abs{y_i-Db_i^0}
  \le D\Abs{\Delta b_i}+f_i
  \le D\lambda_i+f_i.
\]
右侧趋零即推出 \(y_i\to Db_i^0\)。该极限证明初始确权决定最终分配的程度上升，
不证明零流动性是福利最优。证毕。
\end{proof}

\begin{lemma}[二级市场工具性功能引理]
\label{lem:secondary-market-instrumental-functions}
二级转移不增加 \(B^0\)，但在风险厌恶、异质期限或异质信息条件下，可以通过风险分担、价格发现与
期限转换改变参与者期望效用。因此
\[
  \sum_i\Delta b_i=0
  \not\Longrightarrow
  \Delta\sum_i\EE[u_i]=0.
\]
\end{lemma}

\begin{Certificate}[两状态风险分担见证]
取两名风险厌恶参与者、两个等概率状态，初始收益权在状态间负相关。若二级合约使每人的对冲后消费是对冲前
消费的严格均值保持收缩，即
\[
  c_i^{\mathrm{hedged}}\preceq_{\mathrm{cx}}c_i^{\mathrm{unhedged}},
  \qquad
  \EE[c_i^{\mathrm{hedged}}]=\EE[c_i^{\mathrm{unhedged}}],
\]
且两个分布不相同，则严格凹效用 \(u_i\) 给出
\[
  \EE[u_i(c_i^{\mathrm{hedged}})]
  >
  \EE[u_i(c_i^{\mathrm{unhedged}})].
\]
\end{Certificate}

\begin{proof}[证明]
守恒只给出
\[
  \sum_i b_i^1=\sum_i b_i^0,
\]
并不固定各状态下的个体消费。若合约使消费形成严格均值保持收缩，则凹序定义与严格凹性给出
\[
  c_i^{\mathrm{hedged}}\preceq_{\mathrm{cx}}c_i^{\mathrm{unhedged}}
  \Longrightarrow
  \EE[u_i(c_i^{\mathrm{hedged}})]
  >\EE[u_i(c_i^{\mathrm{unhedged}})].
\]
因此二级市场可以产生工具性福利，却不创造初始收益权总量。证毕。
\end{proof}

\begin{proposition}[一级分配优先而非二级市场消灭命题]
\label{prop:primary-priority-not-elimination}
若合同目标是使分配主要反映可验证参与，则可选择有限正流动性上界 \(\lambda_i>0\)，并最小化
\[
  \mathcal L(\lambda)
  =
  \omega_1\sum_i\Abs{y_i-Db_i^0}
  +\omega_2\mathcal R_{\mathrm{illiquidity}}(\lambda)
  +\omega_3\mathcal R_{\mathrm{manipulation}}(\lambda).
\]
最优 \(\lambda^*\) 一般不由逻辑强制为零。
\end{proposition}

\begin{Certificate}[三目标权衡]
若
\[
  \mathcal R_{\mathrm{illiquidity}}'(\lambda)<0,\qquad
  \mathcal R_{\mathrm{manipulation}}'(\lambda)>0
\]
在相关区间成立，则内部解满足
\[
  0\in\partial\mathcal L(\lambda^*).
\]
这构成“强约束而非完全禁止”的一阶条件核验记录。
\end{Certificate}

\begin{proof}[证明]
\(\lambda\downarrow\) 由定理~\ref{thm:limited-liquidity-primary-dominance} 压缩最终分配对初始确权的偏离，
但也可能提高流动性风险；\(\lambda\uparrow\) 降低流动性风险，却可能增加操纵和资本优势对最终分配的影响。
因此
\[
  \text{贡献锚定}\uparrow
  \Longleftrightarrow
  \text{流动性保险}\downarrow
\]
不是单调福利关系。对三项损失求最小值才得到制度参数选择，零流动性不是无条件推论。证毕。
\end{proof}

\section{金融规则在资源请求权层面的趋近}

\begin{definition}[金融化分配比例]
令个人可支配资源为
\[
  c_i=(1-\theta)r_i+\theta y_i(G_F),
\qquad \theta\in[0,1],
\]
其中 \(r_i\) 是不受所比较金融安排影响的基线资源，可包括直接支付的参与收入；
\(y_i(G_F)\) 是由可选信用、投资、风险转移等金融合同规则 \(G_F\) 配置的资源请求权，
\(\theta\) 是本节比较模型中的金融化分配比例。一般记账和货币支付本身不使 \(\theta\) 必然上升。
\end{definition}

\begin{theorem}[金融规则趋近经济分配规则定理]
\label{thm:financial-rule-convergence}
保持 \(r_i\) 与 \(\theta\) 不变，比较 \(G_F^1,G_F^2\)，有
\[
  \Delta c_i
  =
  \theta\bigl[y_i(G_F^2)-y_i(G_F^1)\bigr].
\]
当 \(\theta\to1\) 时，
\[
  \Delta c_i\to
  y_i(G_F^2)-y_i(G_F^1).
\]
\end{theorem}

\begin{Certificate}[金融化权重极限]
核验记录为
\[
  \mathfrak C_{\theta,i}
  =
  \Tuple{r_i,\theta,G_F^1,G_F^2,y_i(G_F^1),y_i(G_F^2)}.
\]
误差精确为
\[
  \Abs{\Delta c_i-\Delta y_i}
  =(1-\theta)\Abs{\Delta y_i}.
\]
\end{Certificate}

\begin{proof}[证明]
两种规则下的 \(c_i\) 相减，非金融项消去：
\[
\begin{aligned}
  c_i(G_F^2)-c_i(G_F^1)
  &=
  (1-\theta)r_i+\theta y_i(G_F^2)\\
  &\quad-(1-\theta)r_i-\theta y_i(G_F^1)\\
  &=\theta\Delta y_i.
\end{aligned}
\]
再由 \(\theta\to1\) 得极限。该结论只涉及资源请求权分配，不覆盖生产、污染、数据产权、竞争法或
产品责任等全部经济规则。证毕。
\end{proof}

\begin{corollary}[金融规则非本体等同推论]
\label{cor:finance-not-total-economy}
若存在非金融资源 \(r_i\ne0\)、生产可行集 \(\Feas(A,r)\) 或外部性约束不由 \(G_F\) 决定，则即使
\(\theta\) 接近一，也不能推出
\[
  G_F\equiv\text{全部经济规则}.
\]
\end{corollary}

\begin{Certificate}[未被金融规则覆盖的坐标]
给定至少一个坐标
\[
  x\in\{\Feas,\text{环境约束},\text{数据产权},\text{竞争规则},\text{公共服务}\}
\]
满足 \(\partial x/\partial G_F=0\) 且 \(\partial W/\partial x\ne0\)，即构成非等同见证。
\end{Certificate}

\begin{proof}[证明]
\[
  \partial_{G_F}x=0
  \wedge
  \partial_xW\ne0
  \Longrightarrow
  G_F\text{ 不能决定 }W\text{ 的全部坐标}.
\]
故金融规则至多在高 \(\theta\) 时逼近分配规则，不能覆盖经济系统的全部状态转移。证毕。
\end{proof}

\section{交换价格、使用价值与稀缺租金}

\begin{proposition}[可复制品价格趋近的条件命题]
\label{prop:replicable-price-limit}
对可复制信息品，若市场可竞争、无持续排他壁垒、产品同质且复制边际成本
\(MC(A)\to0\)，则竞争均衡价格满足
\[
  p(A)\to MC(A)\to0.
\]
该结论针对交换价格，不推出使用效用 \(V(q)\) 趋零。
\end{proposition}

\begin{Certificate}[价格--边际成本与效用分离]
核验记录同时登记
\[
  \Abs{p(A)-MC(A)}\to0,\qquad
  MC(A)\to0,\qquad
  V(q_0)>0\text{ 对某 }q_0>0.
\]
\end{Certificate}

\begin{proof}[证明]
竞争条件给出
\[
  \Abs{p-MC}\to0,
\]
低复制成本给出 \(MC\to0\)，故三角不等式推出 \(p\to0\)。但 \(V(q_0)>0\) 是独立的使用效用见证，
所以
\[
  p\to0\not\Longrightarrow V(q_0)\to0.
\]
现实信息品还可能因固定成本、差异化、排他权与注意力稀缺而偏离该极限
\cite{Varian1998InformationGoods}。证毕。
\end{proof}

\begin{remark}[价值锚与金融工具性价值]
金融凭证不是最终消费效用本身，但可以通过支付结算、跨期配置、风险分担和激励设计产生工具性价值。
其长期可执行性仍需真实资源、服务、可持续收入、合同履行或其他价值锚。可复制品价格下降时，租金可能
迁移到可信身份、注意力、算力、能源、分发入口、品牌、责任承担和支付系统控制权，而不是“价值无限消失”。
\end{remark}
\clearpage
\chapter{可验证参与计价与贡献分配}
\label{chap:platform-contribution}

\section{参与收益映射的参数化表达}

\begin{definition}[参与支付函数]
对市场内制度机制中的参与者 \(i\)，定义
\[
  T_i(G)
  =
  \alpha_G R_G m_G(z_i)
  +B_G(z_i)-F_G(i,z_i),
\]
其中 \(R_G\) 为交易收入、服务收入、会员收入、赞助或其他合规市场来源形成的可分配池，\(\alpha_G\) 为规则分配率，
\(m_G\) 为有效参与认定函数，\(B_G\) 为合同奖励或其他收益，\(F_G\) 为费用、门槛、税费与扣减。
参与收益机制还包括资格谓词
\[
  \Eligible_i^{\mathrm{mechanism}}(G)\in\{0,1\}.
\]
其中 \(T_i(G)\in\RR\) 表示一般市场合同的净结算额，涵盖收入合同及另行成立的付费合同。核心竞技参与合同
取 \(T_i(G)=T_{it}^B\ge\beta_t^{\mathrm{base}}>0\)：令 \(\alpha_GR_G=N_t\beta_t^{\mathrm{base}}+D_t^{\mathrm{diff}}\)，
\(m_G(z_i)=T_{it}^B/(\alpha_GR_G)\)，且本项分配的 \(B_G=F_G=0\)，即可由上述参数化表达实现。
组织费用已经在资金池形成阶段归集；独立购买、借款或其他付费合同不得追溯改写已完成参与的基本份额。
一般合同的净结算、
实际向外支付与反向应付负债分别为
\[
\begin{aligned}
  \widetilde T_i^{\mathrm{net}}(G)
  &=\Eligible_i^{\mathrm{mechanism}}(G)\,T_i(G),\\
  \widetilde T_i^{\mathrm{pay}}(G)
  &=\Eligible_i^{\mathrm{mechanism}}(G)[T_i(G)]_+,\\
  \widetilde L_i(G)
  &=\Eligible_i^{\mathrm{mechanism}}(G)[-T_i(G)]_+.
\end{aligned}
\]
只有 \(\widetilde T_i^{\mathrm{pay}}\) 经结算最终性回执确认后才是实际向参与者支付；
\(\widetilde L_i\) 必须另有合法、公开且可申诉的扣费或负债依据。
\end{definition}

\section{参与收益链的有限构造}

\begin{theorem}[合格参与到合同支付]
\label{thm:platform-payment-chain}
若资格、计量、验证、分配、结算五层算子均已定义，则参与支付由以下有限复合唯一确定：
\[
  a_i
  \xrightarrow{E_G}
  e_i
  \xrightarrow{M_G}
  z_i
  \xrightarrow{V_G}
  \widehat z_i
  \xrightarrow{\alpha_GR_Gm_G+B_G-F_G}
  T_i
  \xrightarrow{\operatorname{Settle}_G}
  (\widetilde T_i^{\mathrm{net}},\widetilde T_i^{\mathrm{pay}},\widetilde L_i).
\]
\end{theorem}

\begin{Certificate}[有限构造：参与支付层级]
\label{cert:platform-payment-chain}
定义
\[
\begin{aligned}
P(0)&:a_i\text{ 已提交},&
P(1)&:e_i=E_G(a_i)\text{ 已判定},\\
P(2)&:z_i=M_G(a_i)\text{ 已计量},&
P(3)&:\widehat z_i=V_G(z_i)\text{ 已验证},\\
P(4)&:T_i\text{ 已计算},&
P(5)&:(\widetilde T_i^{\mathrm{net}},\widetilde T_i^{\mathrm{pay}},\widetilde L_i)\text{ 已结算}.
\end{aligned}
\]
核验记录转移为 \(P(k)\wedge\operatorname{Defined}(O_{k+1})\Rightarrow P(k+1)\)。
\end{Certificate}

\begin{proof}[证明（有限递推）]
\hfill\par\noindent
基例由已提交行为给出 \(P(0)\)。归纳假设为 \(P(k)\) 成立。若下一层算子在上一层输出上可定义，则
\[
  P(k)\Longrightarrow O_{k+1}(P(k))\Longrightarrow P(k+1).
\]
依次得到
\[
  a_i\Longrightarrow e_i\Longrightarrow z_i
  \Longrightarrow\widehat z_i\Longrightarrow T_i
  \Longrightarrow(\widetilde T_i^{\mathrm{net}},\widetilde T_i^{\mathrm{pay}},\widetilde L_i).
\]
任何一层未定义都使后续支付不具备核验记录。因此“发生行为”不直接蕴含“取得收入”，必须经过资格、计量、
真实性验证、规则分配与结算。证毕。
\end{proof}

\begin{theorem}[参与规则比较静态定理]
\label{thm:platform-rule-comparative-static}
在相同底层技术与同一参与信号 \(z_i\) 下，两个组织分配规则 \(G_1,G_2\) 的支付差为
\[
\begin{aligned}
\Delta\widetilde T_i^{\mathrm{net}}
&=
e_i^{(2)}\!
\bigl[\alpha_2R_2m_2(z_i)+B_2(z_i)-F_2(i,z_i)\bigr]\\
&\quad-
e_i^{(1)}\!
\bigl[\alpha_1R_1m_1(z_i)+B_1(z_i)-F_1(i,z_i)\bigr],
\end{aligned}
\]
因而规则参数足以在技术不变时改变参与者收入。
\end{theorem}

\begin{Certificate}[同信号双规则支付]
核验记录固定
\[
  a_i,\ z_i,\ \text{底层技术},
\]
并登记两组
\[
  (e_i^{(k)},\alpha_k,R_k,m_k,B_k,F_k),
  \qquad k\in\{1,2\}.
\]
若计算得到 \(\Delta\widetilde T_i^{\mathrm{net}}\ne0\)，则规则敏感性被见证。
\end{Certificate}

\begin{proof}[证明]
由支付定义分别代入 \(G_1,G_2\)，两式相减即得。由于 \(z_i\) 与技术固定，差异只来自资格、收入池、
有效性口径、分成、激励与扣减参数。因此
\[
  \text{技术相同}\wedge G_1\ne G_2
  \wedge\Delta\widetilde T_i^{\mathrm{net}}\ne0
  \Longrightarrow
  \text{组织分配规则改变收入}.
\]
但不同规则不必令每个个体收入都改变，仍需非零参数差见证。证毕。
\end{proof}

\begin{proposition}[双边平台价格结构命题]
\label{prop:two-sided-platform-price-structure}
固定两侧价格的开邻域。设连续参与需求映射
\((n_C,n_U)\mapsto(D_C(p_C,n_U),D_U(p_U,n_C))\) 把非空紧矩形映入自身，
在该矩形及价格域的开邻域上连续可微，并满足
\[
  \begin{gathered}
  u:=\partial_{p_C}D_C<0,\quad v:=\partial_{p_U}D_U<0,\\
  0<a:=\partial_{n_U}D_C\le\bar a,\quad
  0<b:=\partial_{n_C}D_U\le\bar b,\quad \bar a\bar b<1.
  \end{gathered}
\]
则参与均衡存在且唯一。在内点均衡处，令 \(\Delta=1-ab>0\)，有
\[
 \frac{\partial(n_C,n_U)}{\partial p_C}=\frac{(u,bu)}{\Delta},
 \qquad
 \frac{\partial(n_C,n_U)}{\partial p_U}=\frac{(av,v)}{\Delta}.
\]
因此，保持另一侧价格不变时，本侧降价增加两侧参与；保持总价格不变而调整价格结构时，
至少一侧的均衡参与量发生一阶变化。创作者分成、用户价格和广告主价格须联合进入平台分析。
\end{proposition}

\begin{Certificate}[两侧参与固定点]
令两侧参与量满足
\[
  n_C=D_C(p_C,n_U),\qquad
  n_U=D_U(p_U,n_C),
\]
记录价格偏导、跨侧偏导及反馈余量
\[
  -u>0,\quad -v>0,\quad 1-\bar a\bar b>0.
\]
需求的自映射区间、固定点、价格可变范围及上述导数共同构成平台结构核验记录；
边界参与、不可微需求或反馈余量非正时，不使用本命题的导数结论。
\end{Certificate}

\begin{proof}[证明（加权收缩与隐函数求导）]
以两侧权重比 \(\sqrt{\bar a/\bar b}\) 定义加权最大范数，需求映射的收缩系数不超过
\(\sqrt{\bar a\bar b}<1\)。紧矩形自映射与收缩性给出唯一固定点。对均衡方程求微分得
\[
 \begin{pmatrix}1&-a\\-b&1\end{pmatrix}
 \begin{pmatrix}dn_C\\dn_U\end{pmatrix}
 =
 \begin{pmatrix}u\,dp_C\\v\,dp_U\end{pmatrix}.
\]
行列式为 \(\Delta>0\)，逆矩阵直接给出所列导数。若 \(dp_U=-dp_C\)，则两侧响应分别为
\((u-av)/\Delta\) 与 \((bu-v)/\Delta\)。两者同时为零会推出 \(u=ab\,u\)，
与 \(u<0,\Delta>0\) 矛盾，故同一总价格的不同结构具有不同参与结果。这给出双边价格结构的一个
显式充分模型~\cite{RochetTirole2003TwoSided}，不把正跨侧效应本身视为稳定性条件。证毕。
\end{proof}

\section{登记--评价--分配的职责分离}

\begin{definition}[三层可验证贡献机制]
对一般贡献分配机制，依据交易成本、制度约束与可验证机制的分工逻辑
\cite{Coase1937Firm,North1990Institutions,Myerson1981OptimalAuction}，定义
\[
  U:\text{贡献对象}\to\text{登记与材料转交记录},
\]
\[
  V:\text{经验证的贡献信号}\to w^{(d)}(O)\in\RR_+,
\]
\[
  D:\bigl(v^{(d)},w^{(d)}\bigr)
  \to
  \bigl(T_{\mathrm{participant}}^{(d)},T_{\mathrm{operator}}^{(d)}\bigr).
\]
在 \(\sum_{O'}w^{(d)}(O')>0\) 的非退化域内，一个典型固定规则为
\[
\begin{aligned}
T_{\mathrm{participant}}^{(d)}(O)
&=
\alpha^{(d)}v^{(d)}
\cdot
\frac{w^{(d)}(O)}{\sum_{O'}w^{(d)}(O')},\\
T_{\mathrm{operator}}^{(d)}(O)
&=
(1-\alpha^{(d)})v^{(d)}
\cdot
\frac{w^{(d)}(O)}{\sum_{O'}w^{(d)}(O')}.
\end{aligned}
\]
若总权重为零，则该归一式未定义，机制必须停止分配或调用事前公开且另行通过相对公正核验记录的零权重处置规则。
在核心竞技组织中，上述对象权重机制只用于已预留基本份额以外的贡献差异部分；逐人基本份额按有效参与
名单独立确权。差异权重全零时执行定义~\ref{def:base-differential-distribution} 的等份处置，基本支付
仍按期执行。组织者收益与差异分配共同受已公开的总预算约束。
\end{definition}

\begin{theorem}[登记--评价--分配职责分离定理]
\label{thm:uvd-separation}
若登记或材料转交算子 \(P\) 只能调用 \(U\)，认证后的完整输入集合
\(\mathcal O_U\) 由不可静默覆盖的追加日志固定，\(P\) 不能删除、插入、替换或重排其中的合格对象，评价由
公开且身份无关的 \(V\) 决定，分配由不可被 \(P\) 单方改写的 \(D\) 决定，则对固定
\(\mathcal O_U\)，登记或材料转交主体不能仅凭其制度环节角色单方面改变贡献对象的权重或分配份额。
\end{theorem}

\begin{Certificate}[职责权限不相交]
核验记录登记权限集合
\[
  \operatorname{Control}(P)\subseteq\operatorname{Dom}(U),
\quad
\operatorname{Control}(P)\cap
  \bigl(\operatorname{Rules}(V)\cup\operatorname{Rules}(D)\bigr)=\varnothing,
\]
并登记 \(\mathcal O_U\) 的完整性承诺、逐对象来源引用以及 \(V,D\) 的公开版本承诺。
\end{Certificate}

\begin{proof}[证明]
输入完整性先给出
\[
  \operatorname{Commit}(\mathcal O_U)
  \Longrightarrow
  P\text{ 不能经省略、伪造或排序改变评价输入},
\]
权限分离再给出
\[
  P\text{ 提交或材料转交 }O
  \Longrightarrow U(O),
\]
但
\[
  P\not\Longrightarrow\operatorname{Modify}(V)
  \quad\text{且}\quad
  P\not\Longrightarrow\operatorname{Modify}(D).
\]
因此
\[
  U(O)\Longrightarrow V(O)=w(O)
  \Longrightarrow D(v,w)
\]
由固定输入和公共规则闭合，而不是由单一登记或材料转交主体自由定价。若 \(P\) 能省略、插入或改写
\(\mathcal O_U\)，即使不能修改 \(V,D\) 的规则内容，也可能改变输出分布，因此输入完整性是不可省略的前提。
该结论是设计条件下的权限定理，不是任何现实
组织已经满足权限分离的经验断言。证毕。
\end{proof}

\begin{corollary}[待认领承诺账户的对象先确权推论]
\label{cor:unclaimed-subaccount}
若贡献对象标识 \(c_O\) 唯一绑定承诺账户，且经验证事件按认证权重增加资源请求权余额
\[
  B(c_O)\mathrel{+}=
  \Delta R
  \frac{w_{c_O}}{\sum_{c}w_c},
\]
并且 \(\sum_cw_c>0\)，则在参与者尚未完成身份认领时，资源请求权可以先归集到贡献对象，之后再经权利证明
绑定主体。总权重为零时该归一式不得执行。
\end{corollary}

\begin{Certificate}[状态机与余额承诺]
子账户状态为
\[
  \operatorname{status}(c_O)\in
  \{\operatorname{Unclaimed},\operatorname{Claimed}\},
\]
核验记录记录
\[
  \Tuple{c_O,B(c_O),\operatorname{status}(c_O),
  \operatorname{claimant}(c_O),\operatorname{proof}}.
\]
\end{Certificate}

\begin{proof}[证明]
初始化
\[
  c_O\Longrightarrow
  \operatorname{status}(c_O)=\operatorname{Unclaimed}
  \wedge B(c_O)=0.
\]
认证调用触发
\[
  \Verified(w_{c_O})
  \Longrightarrow
  B(c_O)\mathrel{+}=\Delta B_{c_O},
\]
该转移不要求已有确定的权利人。之后
\[
  \operatorname{VerifyProof}(c_O,\omega)=1
  \Longrightarrow
  \operatorname{claimant}(c_O):=a
  \wedge
  \operatorname{status}(c_O):=\operatorname{Claimed}.
\]
故对象级初始确权与主体级后验认领可以分离。证毕。
\end{proof}

\begin{remark}[机制的经济学位置]
登记、独立评价、预算分配与待认领承诺账户把组织边界、制度承诺、可验证信号和预算平衡组合为一个
贡献分配机制。该机制可用于游戏、审核、传播、社区维护和行业规则维护，并在生产约束松弛与受限二级流动的
条件下分析初始收益权如何生成。
\end{remark}
\clearpage
\chapter{参与活动、合同履行与组织分配公正}
\label{chap:activity-governance}

\section{组织参与经济对象}

\begin{definition}[组织参与行为与验证信号]
令组织参与行为集合为
\[
  \Act_i^{\mathrm{activity}}
  =
  \{\text{竞技},\text{创作},\text{测试},\text{审核},
    \text{教学},\text{社区维护},\text{项目规则维护}\}.
\]
组织事前公布的合同与内部规则 \(G\) 对行为记录 \(r_i\)、身份凭据 \(\omega_i\) 与独立事实证明 \(o_i\) 作用，形成
\[
  z_i=M_G(r_i,o_i),\qquad
  v_i=V_G(z_i,\omega_i)\in\{0,1\},
\qquad
  w_i=v_iW_G(z_i)\ge0.
\]
核心组织参与合同以 \(\mathcal E_t\) 为已完成约定且验证通过的参与者集合，取
\[
  W_G(z_i)=\beta_t^{\mathrm{base}}+D_t^{\mathrm{diff}}\lambda_{it},
  \qquad \phi_G(w)=w,\qquad
  \Pi=N_t\beta_t^{\mathrm{base}}+D_t^{\mathrm{diff}}\le B_t^B.
\]
公共资金与贡献汇总随个人记录一同进入可复算输入，\(\lambda_{it}\) 包括全零差异得分的等份情形。
收益权与资金池分配为
\[
  b_i=\phi_G(w_i),\qquad
  T_i=
  \Pi\frac{b_i}{\sum_{j\in\mathcal E_t}b_j}
  =\beta_t^{\mathrm{base}}+D_t^{\mathrm{diff}}\lambda_{it}
  \ge\beta_t^{\mathrm{base}}>0
  \quad(i\in\mathcal E_t).
\]
这里的 \(T_i\) 是项目收入池上的应付分配，与实际到账收入分别核算。只有另行通过
\(\mathfrak C_t^{\mathrm{fiat}}\) 并取得具有最终性的项目外支付回执时，才有
\(p_{it}^{F}>0\) 或 \(y_{it}^{F}>0\)。
\end{definition}

\begin{definition}[账务与合同履行核验、组织分配公正谓词]
定义账目、合同履行与规则公开的联合核验条件
\[
  \mathcal C_{\mathrm{execution}}
  =
  \LedgerConsistent
  \wedge
  \ExecutionConsistent
  \wedge
  \RuleTransparent.
\]
账目一致要求资金来源、已确认收益权和支付记录相互对应；履约一致要求依事前公布的合同条款办理，
并保存授权裁量的依据与结果；规则公开要求参与者事前可知、事后可核查。人工办理与自动办理采用同一组
核验标准。合同、裁判记录、服务验收单、独立复核意见与支付回执均可作为相应事实的证据。
定义项目制度公正条件
\[
\begin{aligned}
  \mathcal C_{\mathrm{fair}}
  ={}&
  \IC\wedge\IR\wedge\BudgetBalanced
  \wedge\DuplicateClaimResistant\wedge\CollusionResistant\\
  &{}\wedge\OpportunityFair
  \wedge\Auditable\wedge\Appealable
  \wedge\operatorname{BoundedRulePower}.
\end{aligned}
\]
只有 \(\mathcal C_{\mathrm{execution}}\wedge\mathcal C_{\mathrm{fair}}\) 同时通过时，本文才确认
\(\GovernanceFair(G)\) 的条件核验记录。
\end{definition}

\begin{proposition}[经济结论对实施载体的条件不变性]
\label{prop:participation-carrier-independence}
设两种组织实施方式 \(h,h'\) 在同一期间和规则版本下，具有相同的真实参与事实、有效合同、
可用资金、资格与贡献计量、核验结果及可执行支付义务，且均具备本节所要求的独立审计与申诉条件。
若其原始记录经核验后投射为同一经济事实向量 \(z\)，则按同一分配函数得到的收益权与应付分配相同。
只有在两者另有相同的实际结算结果与支付最终性证据时，到账收入结论才相同。
\end{proposition}

\begin{proof}[证明]
令 \(\pi_h,\pi_{h'}\) 为核验后的经济事实提取映射。由
\(\pi_h(\ell_h)=\pi_{h'}(\ell_{h'})=z\)，同一资格函数、贡献函数和分配合同的复合给出同一
\(\lambda_{it}\) 与 \(T_{it}\)。实际支付函数还以结算结果为输入，故到账结论须另核验该输入，
不能由应付相同直接推出。证明仅使用经济事实、制度条件及函数关系，因而线下赛事、协会、企业或普通平台
均可在条件满足时适用；任何具体组织是否满足条件，仍须分别举证。证毕。
\end{proof}

\section{活动核验与跨活动收益权衔接}

\begin{definition}[组织活动与经济相关路径]
令组织实际开展的活动集合为 \(\mathcal M_U=\{m_1,\ldots,m_M\}\)。每个活动单元写成
\[
  m=(X_m,A_m,\tau_m,R_m,O_m,M_m,V_m,W_m,\phi_m,T_m,Q_m),
\]
依次表示状态、行动、状态转移、规则版本、结果、计量、真实性验证、贡献估值、初始确权、支付以及审计申诉。
若活动中的行动或跨活动变换会改变收入请求权、权利凭证、可交易资产、法币支付、排名准入、注意力分配或其他具有经济后果的状态，
则称其路径 \(\gamma_m\) 为经济相关路径。纯展示且不改变任何经济状态的界面变化不进入本定义。
\end{definition}

\begin{definition}[组织参与的权利生成与货币支付双阶段映射]
\label{def:activity-rights-map}
组织在真实参与、治理核验记录和可执行合同的联合域上确认权利：
\[
  \mathcal R_U^{\mathrm{claim}}:
  (a_{it},\ell_{it},\mathfrak C_{it}^{\mathrm{part}},
  \mathfrak C_{it}^{\mathrm{claim}})
  \longmapsto
  (b_{it}^{0},\beta_{it}).
\]
其定义域不包含“仅由组织声明稀缺”的空输入；参与核验记录或可执行权利任一缺失时，终端只能是
\(\operatorname{Rejected}\) 或 \(\operatorname{AppealPending}\)，不能标记为合规权利。法币支付是第二个独立条件：
\[
  \mathcal R_U^{\mathrm{fiat}}:
  (\beta_{it},\mathfrak C_t^{\mathrm{fiat}},F_t^{\mathrm{pay}})
  \longmapsto p_{it}^{F}.
\]
资金来源或法币核验记录缺失只使权利处于未支付状态，不抹去已经依法形成的 \(b_{it}^{0}\) 或 \(\beta_{it}\)；反之，
存在资金也不能替代第一阶段的参与与权利核验记录。两阶段分别回答“谁依何种参与取得请求权”与“谁以何种
已到账资金履行付款义务”；凭证登记、内部估值或单方承诺均不能补出缺失的参与事实或付款来源。
\end{definition}

\begin{definition}[组织文化状态、事实记录与规则生效]
\label{def:organization-cultural-extension}
对能够形成共同历史、身份或文化意义的活动 \(m\)，把原状态扩展为
\[
  \widehat X_m=X_m\times C_m,
\]
其中 \(C_m\) 是第~\ref{sec:endogenous-cultural-background}节中的文化状态。固定规则版本
\(R_m^\nu\) 时，扩展状态转移为
\[
  \widehat\tau_m:
  (x_{m,t},C_{m,t},a_{m,t},\xi_{m,t};R_m^\nu)
  \longmapsto
  \left(
  \tau_m(x_{m,t},a_{m,t},\xi_{m,t};R_m^\nu),
  \kappa_m(C_{m,t},\ell_{m,t};R_m^\nu)
  \right),
\]
其中 \(\xi_{m,t}\) 保存外部事实、授权裁量及随机实现，并作为输入纳入不可静默覆盖的参与事件记录
\[
  \ell_{m,t}
  =
  \Tuple{
  \mathrm{id},t,R_m^\nu,x_{m,t},a_{m,t},
  \xi_{m,t},o_{m,t},x_{m,t+1},
  \mathrm{entitlement},\mathrm{payment},
  \mathrm{power},\mathrm{appeal}
  }.
\]
组织必须保存主体、时间、适用规则、参与行为、事实依据、确认结果、收益权、支付及申诉更正的连续记录。
合同、签收单、裁判记录、审核意见与结算台账等证据须能够相互核对，并保留原件或可验证副本及更正轨迹。文化更新 \(\kappa_m\) 可以随参与
历史开放演化，但同一版本下不得暗改合法行动、结果、确权或结算函数。

规则可以升级而非永久冻结。版本变更
\[
  \Upsilon_m^{\nu\to\nu+1}:
  (X_m,H_m,B_m,C_m;R_m^\nu)
  \longmapsto
  (X_m',H_m',B_m',C_m';R_m^{\nu+1})
\]
必须事前公开授权主体、生效时间、影响范围、状态迁移、既有权利处理、风险测试、回退或补偿以及申诉路径；
旧版本历史和迁移映射必须保留。未经该程序的秘密、追溯或针对特定参与者的规则修改，记为
\(\mathfrak C_{\mathrm{rule}}=0\)。
\end{definition}

\begin{theorem}[稳定底层规则与开放文化演化兼容定理]
\label{thm:stable-rule-open-culture}
若在版本 \(R_m^\nu\) 的有效期内，\(\widehat\tau_m\) 与 \(\kappa_m\) 在已记录事实、授权裁量结果及
随机实现给定后具有确定输出，全部 \(\ell_{m,t}\) 按时间顺序留存，且任何版本变更通过
\(\Upsilon_m^{\nu\to\nu+1}\) 保存历史与权利来源，则：
\begin{enumerate}[label=\textup{(\roman*)},leftmargin=3em]
  \item 任一时点的约定可观测的活动状态、文化指标、确认结果、收益权和项目内应付状态均可从初始状态、规则版本与事件日志独立
  重放；实际法币支付只有在日志另含法币核验记录和外部结算最终性回执时才可重放为已到账；
  \item \(C_{m,t+1}\ne C_{m,t}\) 不构成规则变更，因为其差异由同一
  \(R_m^\nu\) 下的参与事件经 \(\kappa_m\) 生成；
  \item 规则版本可以经公开迁移更新，而既有参与不被事后规则重新解释；
  \item 若文化沉积满足定理~\ref{thm:cultural-stock-accumulation}，参与者可以在不破坏结算可审计性的条件下
  共同形成持续文化背景。
\end{enumerate}
因此，在上述状态扩展、留痕和版本迁移条件下，“底层规则严格稳定”和“上层社会文化开放生成”可以同时成立；
该兼容性结论不单独证明任何现实组织已经满足规范化条件。
这里可复算的是事前约定的经济状态及可观测文化指标，参与者的全部心理体验不属于记录重建的结论。
\end{theorem}

\begin{proof}[证明（状态扩展与版本分段归纳）]
固定版本内，对 \(t\) 作归纳。给定 \((x_{m,t},C_{m,t})\)、规则版本和完整
\(\ell_{m,t}\)，\(\widehat\tau_m\) 依据完整事实、已记录裁量及随机实现唯一确定下一状态，故从初始状态可以逐步重放整个
版本区间。文化变化来自 \(\kappa_m(C_{m,t},\ell_{m,t};R_m^\nu)\)，没有改变该映射本身，因而不等于规则
暗改。跨版本时，\(\Upsilon_m^{\nu\to\nu+1}\) 保存旧状态、历史、权利和迁移依据，把下一版本的初始状态
唯一连接到上一版本终态；对有限版本序列再次归纳即可得到全历史可追踪。若规则被秘密或追溯修改，则相同前态
和日志可以产生不同结果，独立重放失败，反证其不能通过规则核验记录。证毕。
\end{proof}

\begin{definition}[活动治理核验记录与跨活动一致性]
经济相关活动 \(m\) 的治理核验记录定义为
\[
  \mathfrak C_m^{\mathrm{activity}}
  =\left\langle
  \begin{gathered}
  \mathfrak C_{\mathrm{state}},\mathfrak C_{\mathrm{rule}},
  \mathfrak C_{\mathrm{id}},\mathfrak C_{\mathrm{anti\text{-}cheat}},
  \mathfrak C_{\mathrm{result}},\\
  \mathfrak C_{\mathrm{RPF}},\mathfrak C_{\mathrm{budget}},
  \mathfrak C_{\mathrm{claim}},\mathfrak C_{\mathrm{fiat}},
  \mathfrak C_{\mathrm{finance}},\\
  \mathfrak C_{\mathrm{power}},\mathfrak C_{\mathrm{appeal}},
  \mathfrak C_{\mathrm{culture}}
  \end{gathered}
  \right\rangle.
\]
它分别要求状态与行动可重放、规则事前固定、身份可归属、作弊可识别、结果可复算、参与--分配相对公正、
收入池闭合、权利法律义务可执行、法币来源与结算可追踪，并在实际发生证券化、托管、收益权转让或其他金融交易时通过
相应金融规则，且规则权受限、错误可修复。不产生法币请求的路径可以把 \(\mathfrak C_{\mathrm{fiat}}\) 标为
不适用，但须证明其状态不会直接或经跨活动映射改变任何法币权利；不发生金融化的路径可以把
\(\mathfrak C_{\mathrm{finance}}\) 标为不适用，但须证明没有投资、转让、托管、抵押或兑换功能。若活动声明形成文化状态，还要求参与事件完整留痕、文化更新可重放、
组织声明与共同体承认分离以及规则修改与文化演化分离。若资产或权利从活动 \(m\) 进入活动 \(n\)，跨活动核验记录还须验证
\[
  \Xi_{mn}:(x_m,b_m,R_m)\longmapsto(x_n,b_n,R_n)
\]
保持主体、权利数量、限制条件、来源和规则版本可追踪，且不把未通过的参与结果洗成已确权资产。
\end{definition}

\begin{theorem}[多活动组织的完整核验定理]
\label{thm:multi-activity-governance}
在至少存在一条经济相关路径且全部核验记录采用同一项目边界、规则版本与结算窗口时，参与组织在该窗口内的参与--分配
在本文声明口径下能够整体纳入规范治理，当且仅当每一条经济相关路径所属活动的
\(\mathfrak C_m^{\mathrm{activity}}\) 通过，使所有实际使用的跨活动映射 \(\Xi_{mn}\) 通过，并满足规范治理双层
条件。若文化状态会改变身份、声誉、注意力、需求、准入或收入，它的形成与跨活动传播同样属于经济相关路径。
若某一活动失败而各收入池、资产、文化影响与权利来源可完全隔离，只撤销该活动及其下游路径的核验记录；若失败活动
与其他活动共用不可追踪收入池或资产，则整个混合分配不能确认相对公正结论。
\end{theorem}

\begin{Certificate}[组织活动完整核验记录]
核验记录列出全部经济相关活动、每项活动的状态转移与规则版本、参与行为和身份、反作弊、结果认定依据、贡献与份额
规则、收入池、支付、权限、申诉、文化状态更新、事件留痕、版本迁移，以及所有跨活动资产、文化影响和权利映射。覆盖率必须按经济相关路径而非活动宣传
名称计算。令 \(\Gamma_{U,t}\) 为结算窗口内实际发生且具有经济后果的路径集合：
\[
  \operatorname{Coverage}_U
  =\frac{\#\{\gamma_m\in\Gamma_{U,t}:\mathfrak C_m^{\mathrm{activity}}=1
  \text{ 且全部下游 }\Xi=1\}}
  {\#\Gamma_{U,t}}.
\]
分母必须是有限正整数；若项目没有经济相关路径，则本核验记录标为不适用而不是把空集比率记为一。整体核验记录要求
\(\operatorname{Coverage}_U=1\)，并保存未纳入路径为零的清单见证。该结论只覆盖已核验的结算窗口，
不预先认证未来行为；新活动、规则变更或新结算窗口须更新核验记录。
\end{Certificate}

\begin{proof}[证明（逐路径覆盖与污染传播）]
任一经济支付或权利必须沿某条活动路径和跨活动映射产生。若每条源路径均通过活动核验记录、每个变换均保持来源
与规则限制，且双层治理核验记录通过，则每项最终分配都可回溯到规则内可重放的参与判定与安全合法域。反之，取一条未通过
路径：若可隔离，其失败不改变其他路径的证据；若与其他收入、文化影响或资产不可区分地混合，则最终份额无法判断来自
已验证还是未验证参与，故横向一致性和收入池来源均不可复核。围棋规则层是固定规则下可独立重放的实例，但不能替其他活动
自动补齐状态机、文化沉积、贡献规则和分配核验记录。证毕。
\end{proof}

\begin{corollary}[组织参与收入进入一般市场循环推论]
\label{cor:activity-payment-cycle}
若多活动组织的完整核验定理的全部路径核验记录通过，所有产生收益权的路径另通过
第~\ref{sec:certified-participation-fiat-cycle}节的权利来源与法币来源双守恒核验记录，并满足
定理~\ref{thm:certified-participation-fiat-cycle}，则
\[
\begin{aligned}
  \text{组织参与}
  &\longrightarrow
  \text{治理认证的合规参与}
  \longrightarrow
  \text{可执行权利凭证}\\
  &\longrightarrow
  \text{实际法币收入}
  \longrightarrow
  \text{一般市场消费或投资}
\end{aligned}
\]
是一条可重放、可审计且不依赖内部资产自我报价的开放经济路径。若权利只在项目内部循环、外部法币来源为零、
或只有二级价格而没有可执行基础请求，则该推论不成立。证券化只在需要汇集、份额化或二级投资交易时作为条件
分支出现，不是所有组织参与取得法币收入的必经步骤。
\end{corollary}

\begin{proof}[证明（定理嵌入）]
逐活动核验记录把每项终端权利回溯到真实参与和固定规则，双守恒核验记录排除权利与付款来源凭空增加，
定理~\ref{thm:certified-participation-fiat-cycle}再把可执行请求权复合到具有最终性的法币支付。实际到账可以进入
项目外一般预算约束，故开放路径成立。任一前提失败时，相应复合映射未定义或支付为零，不能由内部报价补足。
证毕。
\end{proof}

\section{有限参与集合的依约分配}

\begin{theorem}[有限参与集合的归一分配与基本份额保持]
\label{thm:participant-normalized-allocation}
若 \(\Pi\ge0\)、\(b_i\ge0\) 且 \(B=\sum_i b_i>0\)，则
\[
  T_i=\Pi\frac{b_i}{B}
\]
满足
\[
  T_i\ge0,\qquad
  \sum_iT_i=\Pi.
\]
在核心组织参与合同中，\(b_i=\beta_t^{\mathrm{base}}+D_t^{\mathrm{diff}}\lambda_{it}\)，
\(B=\Pi=N_t\beta_t^{\mathrm{base}}+D_t^{\mathrm{diff}}\)，故每名有效参与者的
\(T_i\ge\beta_t^{\mathrm{base}}>0\)，且严格贡献顺序由分配保持。
\end{theorem}

\begin{Certificate}[有限构造：参与者分配总额闭合]
\label{cert:participant-normalized-allocation}
令
\[
  P(k):\quad
  \sum_{i=1}^{k}T_i
  =
  \frac{\Pi}{B}\sum_{i=1}^{k}b_i.
\]
基例 \(P(1)\) 由支付定义给出；归纳步为
\[
  P(k)\wedge T_{k+1}=\frac{\Pi b_{k+1}}{B}
  \Longrightarrow P(k+1).
\]
\end{Certificate}

\begin{proof}[证明（有限递推）]
基例
\[
  T_1=\Pi b_1/B
  \Longrightarrow
  \sum_{i=1}^{1}T_i=(\Pi/B)\sum_{i=1}^{1}b_i.
\]
若 \(P(k)\) 成立，则
\[
  \sum_{i=1}^{k+1}T_i
  =
  \frac{\Pi}{B}\sum_{i=1}^{k}b_i
  +\frac{\Pi}{B}b_{k+1}
  =
  \frac{\Pi}{B}\sum_{i=1}^{k+1}b_i.
\]
因此 \(P(k)\Rightarrow P(k+1)\)。取 \(k=n\) 并用 \(\sum_i b_i=B\)，得到
\(\sum_iT_i=\Pi\)。非负性由 \(\Pi,b_i,B\) 的符号直接得到。在核心合同中 \(\Pi/B=1\)，
故 \(T_i=b_i\)，正基本份额与严格贡献顺序由定理~\ref{thm:base-differential-coverage} 保持。证毕。
\end{proof}

\begin{theorem}[参与--记录--确权--项目内分配复合定理]
\label{thm:activity-participation-composition}
对核心合同中 \(i\in\mathcal E_t\)，若 \(\operatorname{Recordable}(a_i)=1\)，事件日志追加完整且不可静默覆盖，
\(M_G,V_G,W_G,\phi_G\) 在前一步输出上依次定义，并采用已经获得资金覆盖的基本与差异分配映射，则
组织参与行为通过以下资格至支付的履约链生成正初始分配：
\[
\begin{gathered}
  a_i\in\Act_i^{\mathrm{activity}}
  \wedge\operatorname{Recordable}(a_i)=1
  \Longrightarrow
  r_i=\operatorname{Record}(a_i)
  \Longrightarrow
  z_i=M_G(r_i,o_i)\\
  \Longrightarrow
  V_G(z_i,\omega_i)=1
  \Longrightarrow
  w_i=W_G(z_i)
  \Longrightarrow
  b_i=\phi_G(w_i)
  \Longrightarrow
  T_i=\Pi b_i/\sum_jb_j\ge\beta_t^{\mathrm{base}}>0.
\end{gathered}
\]
\end{theorem}

\begin{Certificate}[逐条件事件日志]
核验记录为
\[
  \mathfrak C_{\mathrm{activity},i}
  =
  \Tuple{
  \operatorname{id}(i),a_i,r_i,z_i,
  \omega_i,v_i,w_i,b_i,T_i,
  \operatorname{ruleVersion}(G)}.
\]
每个字段包含前一字段的引用承诺。任一条件失败时，终端状态属于
\[
  \{\operatorname{Rejected},\operatorname{AppealPending}\},
\]
不得伪造为项目内已分配，更不得伪造为实际法币已支付。
\end{Certificate}

\begin{proof}[证明]
由记录可行与追加完整性，
\[
  \operatorname{Recordable}(a_i)
  \Longrightarrow r_i.
\]
由计量规则与独立事实证明，
\[
  r_i\wedge o_i\Longrightarrow z_i.
\]
由身份与行为验证，
\[
  z_i\wedge\omega_i\wedge V_G=1
  \Longrightarrow \Verified_i.
\]
再由核心合同的正基本份额与恒等确权映射，
\[
  i\in\mathcal E_t
  \Longrightarrow w_i\ge\beta_t^{\mathrm{base}}>0,
  \qquad
  \phi_G(w)=w
  \Longrightarrow b_i=w_i>0.
\]
再由 \(\Pi=\sum_jb_j\) 得 \(T_i=b_i\ge\beta_t^{\mathrm{base}}>0\)，闭合项目内应付分配链。
若任一箭头前提缺失，则相应后件未被证明；执行程序或人工结算不得自行补足缺失的身份、事实依据真实性
或项目制度合法性。即使 \(T_i>0\)，实际法币支付仍须另行满足
定理~\ref{thm:certified-participation-fiat-cycle}。证毕。
\end{proof}

\section{合同履行一致性与分配公正的区分}

\begin{lemma}[履约一致性不充分推出组织分配公正引理]
\label{lem:execution-not-fairness}
即使
\[
  \mathcal C_{\mathrm{execution}}=1,
\]
也不能仅由此推出 \(\GovernanceFair(G)\)。特别地，若投票权或规则修改权集中，执行者可以严格履行一个
对特定主体偏置的分配条款。
\end{lemma}

\begin{Certificate}[集中控制反例]
考虑一个按简单多数决定分配规则的组织。令某主体掌握的表决权满足
\[
  \nu_1>\frac12\sum_i\nu_i,
\]
且规则修改采用简单多数。则实体 \(1\) 可单独决定规则版本，尽管每次投票和执行均可公开验证。
取两名具有相同资格和已验证贡献的参与者，令控制者事前公布规则，把正资金池全部分给其中一人，
排除另一人并关闭独立申诉。若记录与支付均准确遵守该规则，则履约一致成立，而横向一致性与可申诉性失败。
\end{Certificate}

\begin{proof}[证明]
\[
  \nu_1>\frac12\sum_i\nu_i
  \Longrightarrow
  \operatorname{Decisive}(1).
\]
再有
\[
  \operatorname{Decisive}(1)
  \wedge
  \ExecutionConsistent
  \Longrightarrow
  \text{单方选择的规则被一致执行}.
\]
一致履行只证明账务与支付符合已公布的合同条款。它不证明机会公平、抗合谋、申诉有效或权力受限。因此
\[
  \mathcal C_{\mathrm{execution}}
  \not\Longrightarrow
  \mathcal C_{\mathrm{fair}}.
\]
账务记录准确、合同履行一致与分配规则公正分别成立，不能由其中一项推定其余项目。证毕。
\end{proof}

\begin{proposition}[重复申领防范与真实主体识别]
\label{prop:duplicate-claim-gate}
设固定资金池 \(P>0\) 按登记身份等份分配，除某一主体外还有 \(N\ge1\) 名各登记一次的真实参与者。
若该主体可以零成本登记 \(k\ge2\) 个互不关联且被视为合格的身份，则其所得份额严格超过单次登记时的份额。
因此，身份登记列表本身不足以证明真实参与覆盖与单人单次申领；还须核验真实主体、参与资格及事件的唯一性。
\end{proposition}

\begin{Certificate}[固定预算下的重复申领增益]
令主体以 \(k\) 个身份申领的总成本为 \(c(k)\ge0\)，其净收益为
\[
  U(k)=P\frac{k}{N+k}-c(k).
\]
重复申领具有严格净增益当且仅当
\[
  c(k)-c(1)
  <
  \frac{PN(k-1)}{(N+k)(N+1)}.
\]
核验记录须保存主体与申领的关联依据、有效参与证据及重复事件排除结果，不能仅保存身份数量。
\end{Certificate}

\begin{proof}[证明（固定预算份额比较）]
直接相减得
\[
  U(k)-U(1)
  =\frac{PN(k-1)}{(N+k)(N+1)}
  -\bigl(c(k)-c(1)\bigr).
\]
当 \(k\ge2\)、\(N\ge1\) 且登记成本为零时，右端严格为正；一般成本下，正收益条件恰为核验记录中的不等式。
主体所得总份额始终小于 \(P\)，增加部分来自挤占其他真实参与者的份额，预算并未增长。
故账务平衡和记录一致仍不足以排除重复申领，真实主体识别与参与事件去重必须在资格确认前完成。
基本份额合同亦须以核验后的真实接纳人数计算资金覆盖。证毕。
\end{proof}

\begin{corollary}[可申诉项目制度公正推论]
\label{cor:appealable-governance-fairness}
若机制同时满足
\[
  \begin{gathered}
    \mathcal C_{\mathrm{execution}}
    \wedge\IC\wedge\IR\wedge\BudgetBalanced
    \wedge\DuplicateClaimResistant\wedge\CollusionResistant\\
    {}\wedge\OpportunityFair
    \wedge\Auditable\wedge\Appealable
    \wedge\operatorname{BoundedRulePower},
  \end{gathered}
\]
则可以在这些明示维度上确认条件项目制度公正核验记录。
\end{corollary}

\begin{Certificate}[公正条件向量]
定义二元条件向量
\[
  \mathbf g_{\mathrm{fair}}
  =
  (g_{\mathrm{execution}},g_{\mathrm{IC}},g_{\mathrm{IR}},g_{\mathrm{BB}},
  g_{\mathrm{duplicate}},g_{\mathrm{coll}},g_{\mathrm{opp}},
  g_{\mathrm{audit}},g_{\mathrm{appeal}},g_{\mathrm{power}}).
\]
只有
\[
  \min_k g_k=1
\]
时，核验记录状态才是 \(\operatorname{ConditionallyFair}\)。
\end{Certificate}

\begin{proof}[证明]
每个 \(g_k\) 对应一个不可由其他条件替代的必要维度。合取语义给出
\[
  \min_k g_k=1
  \Longleftrightarrow
  \bigwedge_k(g_k=1).
\]
由定义，全部条件成立正是 \(\mathcal C_{\mathrm{execution}}\wedge\mathcal C_{\mathrm{fair}}\)。
故条件核验记录可确认；任一条件为零时只能登记具体缺口，不能把其余条件的成功冒充整体公正。
制度的监测、制裁、冲突解决和参与式修订边界与共同资源制度研究相容
\cite{Ostrom1990Commons}。证毕。
\end{proof}

\begin{remark}[参与组织的经济机制]
竞技、创作、教学、审核和服务可以在组织合同中共同形成参与收益。每项有经济后果的活动都须明确完成标准、
贡献比较、收益权、付款来源及权力边界；跨活动转移保留权利来源与限制，申诉程序能够修复误判及越权决定。
组织形态与记录工具可以不同，活动核验、预算覆盖、合同履行和独立复核仍是同一组经济要求。
\end{remark}
\clearpage
\chapter{动态稳定、福利边界、可证伪预测与总定理}
\label{chap:stability-synthesis}

\section{市场参与经济的动态系统}

\begin{definition}[状态、组织规则更新与参与更新]
令离散时间状态为
\[
  x_t=(N_t,B_t^B,\Pi_t^P,\mathbf I_t^B,
  \mathbf a_t^{\mathrm{att}},\boldsymbol\vartheta_t^P,D_t^P,
  \mathbf R_t,K_t^R,\tau_t,\mathbf b_t,\mathbf y_t,G_t,\ell_t,
  \beta_t^{\mathrm{base}},D_t^{\mathrm{diff}},\mathbf M_t,
  \boldsymbol\Gamma_t^{\mathrm{mkt}},\mathbf L_t^{\mathrm{cash}}),
\]
其中 \(N_t\) 为有效参与规模，\(B_t^B\) 为桥接行业可分配池，\(\Pi_t^P\) 为生产侧盈利，
\(\mathbf I_t^B\) 为生产者竞争性注意力投入，\(\mathbf a_t^{\mathrm{att}}\) 为相对注意力，
\(\boldsymbol\vartheta_t^P\) 为消费者对生产者及其商品的偏好状态，\(D_t^P\) 为生产品有效需求，
\(\mathbf R_t\) 为充分救济向量，\(K_t^R\) 为救济储备，\(\tau_t\in[0,1]\) 为信任与履约状态，
\(\mathbf b_t\) 为初始收益权向量，\(\mathbf y_t\) 为最终分配，
\(G_t\) 为组织规则版本，\(\ell_t\) 为公共法律状态；分配坐标记录基本份额和差异池，营销坐标记录
完整营销成本、偏好营销净增利润与已实现利润中的可用现金。\(\Pi_t^P\) 在再投入制度环节使用已实现净利润。
定义
\[
  x_{t+1}=F(x_t,\eta_t),
\]
其中 \(\eta_t\) 是需求、技术、攻击、法律或资源冲击。参与更新为
\[
  N_{t+1}
  =
  N_t+\operatorname{Entry}(\Delta u_t)
  -\operatorname{Exit}(-\Delta u_t),
\]
\[
  \Delta u_t
  =
  \EE[T_{G_t}]-c_t-\rho_t-u^{\mathrm{out}}.
\]
注意力桥接的状态更新必须保留顺序
\[
  \Pi_t^P\to\mathbf I_t^B\to\mathbf a_t^{\mathrm{att}}
  \to\boldsymbol\vartheta_{t+1}^P\to D_{t+1}^P
  \to\Pi_{t+1}^P\to\mathbf I_{t+1}^B,
\]
以及
\[
  \mathbf I_t^B\to B_t^B\to\mathbf y_t
  \to C_{t+1}^P\to D_{t+1}^P.
\]
\[
  (\Pi_t^{P,+},K_t^R,Q_t^R)
  \to\mathbf R_t^*
  \to(D_{t+1}^P,N_{t+1})
  \to(\Pi_{t+1}^P,\mathbf I_{t+1}^B),
\]
其中第三条路径只在定理~\ref{thm:relief-endogenous-stability} 的底线、预算、正剩余激励和阈值条件下确认。
第一条路径的利润由销售回款扣除产品成本及全部营销成本形成；偏好增量利润与营销成本之差为正时才确认
营销扩大利润结论。第二条路径覆盖每名有效参与者的基本与差异分配，并以实际支付回执连接消费。
下一期营销支出受 \(M_k(I_{k,t+1}^B)\le L_{k,t+1}^{\mathrm{cash}}\le\Pi_{k,t+1}^P\) 约束，
预期回款不得计作当期可用现金。
若经验系统只保存曝光而不保存生产商品偏好、购买和参与者收入用途，就不能识别完整更新映射。
\end{definition}

\begin{definition}[制度机制竞争绩效]
对市场内制度机制 \(G\)，定义可审计绩效
\[
\begin{aligned}
  \mathcal J(G)
  ={}&
  \omega_1 W(G)
  +\omega_2\operatorname{Participation}(G)\\
  &+\omega_3\operatorname{ContributionIncentive}(G)
  +\omega_4\operatorname{FloorAccess}(G)\\
  &+\omega_5\operatorname{CulturalChoice}(G)
  +\omega_6\operatorname{DemandStability}(G)\\
  &+\omega_7\operatorname{Innovation}(G)
  -\omega_8\operatorname{Manipulation}(G)\\
  &-\omega_9\operatorname{Concentration}(G)
  -\omega_{10}\operatorname{GovernanceCost}(G),
\end{aligned}
\]
其中 \(\omega_k\ge0\)，所有坐标必须事先规范到共同量纲并由声明的观测口径估计。参与覆盖、贡献激励、底线
接入和文化选择分别对应三项核心功能及其目标效应；治理成本作为实现代价单列。制度机制竞争是对
\(\mathcal J\) 的可比较竞争，不是把任何制度标签预设为优越。
\end{definition}

\begin{definition}[机制份额的条件性选择动态]
令有限候选载体族为 \(\mathfrak G\)，\(s_t(G)>0\) 为载体 \(G\) 在共同适用域内吸收的参与、交易或投入份额，
且 \(\sum_{G\in\mathfrak G}s_t(G)=1\)。取 \(\epsilon>0\)，定义正适应度
\[
  \Phi_t(G)
  =\mathcal J_t(G)-\min_{H\in\mathfrak G}\mathcal J_t(H)+\epsilon,
\]
以及离散选择更新
\[
  s_{t+1}(G)
  =\frac{s_t(G)\Phi_t(G)}
  {\sum_{H\in\mathfrak G}s_t(H)\Phi_t(H)}.
\]
该动态只表示参与者、赞助者、生产者与组织在可进入候选之间的条件性选择，不要求任何中心主体预先设计全部载体。
\end{definition}

\begin{proposition}[三项核心功能的条件性演化优势]
\label{prop:core-functions-evolutionary-advantage}
设 \(G^*,H\in\mathfrak G\) 使用共同观测口径。若 \(G^*\) 通过
\(\mathbf g^{\mathrm{core}}(G^*)=(1,1,1)\) 与必要的实施核验，使参与保持、贡献激励、底线接入、文化选择和
需求稳定的加权增益严格超过新增治理成本、操纵和集中损失，以致
\(\mathcal J_t(G^*)>\mathcal J_t(H)\)，则
\[
  \frac{s_{t+1}(G^*)}{s_{t+1}(H)}
  =\frac{s_t(G^*)}{s_t(H)}
  \frac{\Phi_t(G^*)}{\Phi_t(H)}
  >\frac{s_t(G^*)}{s_t(H)}.
\]
故三项核心功能在这些条件下具有相对扩张的演化动力。若真实回款不足、功能增益不显著、治理成本过高或候选
进入被封锁，使上述绩效严格序失败，则本命题不推出其份额上升，更不推出历史必然收敛。
\end{proposition}

\begin{proof}[证明（份额比）]
公共归一化分母在份额比中约去。由 \(\mathcal J_t(G^*)>\mathcal J_t(H)\) 与共同的平移量及
\(\epsilon>0\)，有 \(\Phi_t(G^*)>\Phi_t(H)>0\)，故适应度之比大于一，结论立即成立。证毕。
\end{proof}

\section{有限时间状态轨迹存在}

\begin{theorem}[有限制度轨迹递归存在]
\label{thm:finite-governance-trajectory}
若初态 \(x_0\in\States\) 已给定，且对每个 \(t<T\)，映射 \(F_t(\cdot,\eta_t)\) 在
\(\States\) 上为全函数，则唯一有限轨迹
\[
  (x_0,x_1,\ldots,x_T)
\]
由递归 \(x_{t+1}=F_t(x_t,\eta_t)\) 决定。
\end{theorem}

\begin{Certificate}[有限构造：时间步递归闭合]
\label{cert:finite-governance-trajectory}
令
\[
  P(k):\quad x_0,\ldots,x_k\text{ 已唯一构造}.
\]
基例 \(P(0)\) 由初态给出；归纳转移为
\[
  P(k)\wedge\operatorname{Function}(F_k)
  \Longrightarrow
  x_{k+1}=F_k(x_k,\eta_k)
  \Longrightarrow P(k+1).
\]
\end{Certificate}

\begin{proof}[证明（有限递推）]
\(x_0\) 已给定，故 \(P(0)\) 成立。若 \(P(k)\) 成立，则 \(x_k\) 唯一。又因 \(F_k\) 是全函数，
\[
  x_k\wedge\eta_k
  \Longrightarrow
  \exists!\,x_{k+1}=F_k(x_k,\eta_k).
\]
故 \(P(k)\Rightarrow P(k+1)\)。数学归纳法给出 \(P(T)\)。这只证明给定规则下轨迹可定义，
不证明轨迹稳定、公平或福利最优。证毕。
\end{proof}

\begin{theorem}[局部动态稳定充分条件定理]
\label{thm:local-dynamic-stability}
设 \(x^*=F(x^*,0)\) 是无冲击固定点，\(F\) 在 \(x^*\) 邻域连续可微，且雅可比矩阵满足
\[
  \rho\bigl(DF(x^*,0)\bigr)<1,
\]
其中 \(\rho\) 为谱半径。则 \(x^*\) 局部渐近稳定。
\end{theorem}

\begin{Certificate}[谱半径与收缩邻域]
给定范数 \(\Norm{\cdot}_*\) 与常数 \(\kappa<1\)，使某邻域 \(U\) 内
\[
  \Norm{F(x,0)-F(x^*,0)}_*
  \le
  \kappa\Norm{x-x^*}_*.
\]
元组 \(\Tuple{x^*,U,\kappa,\Norm{\cdot}_*}\) 构成局部收缩核验记录。
\end{Certificate}

\begin{proof}[证明]
谱半径小于一允许选择等价范数，使线性化算子范数小于一；由连续可微性可缩小邻域 \(U\)，使
\(F\) 在 \(U\) 为收缩。于是
\[
  x_t\in U
  \Longrightarrow
  \Norm{x_{t+1}-x^*}_*
  \le
  \kappa\Norm{x_t-x^*}_*.
\]
递推得到
\[
  \Norm{x_t-x^*}_*
  \le
  \kappa^t\Norm{x_0-x^*}_*\to0.
\]
因此固定点局部渐近稳定。稳定性必须由此类动态核验记录给出，不能由静态预算平衡或合同履行替代。证毕。
\end{proof}

\section{公平、效率与预算平衡的不可自动兼得边界}

\begin{lemma}[条件合取不推出一般福利最优引理]
\label{lem:gates-not-welfare-optimality}
即使某机制满足
\[
  \IC\wedge\IR\wedge\BudgetBalanced
  \wedge\DuplicateClaimResistant\wedge\Appealable,
\]
也不能在没有偏好、技术、信息和外部性假设时推出 Pareto 最优或社会福利最大。
\end{lemma}

\begin{Certificate}[遗漏外部性的反例]
取两个参与者，机制内部转移预算平衡且真实报告最优，但行为产生未计价损害 \(h(a)>0\)。
若机制目标未包含 \(h\)，则其选择 \(a^M\) 可以满足全部条件，却有
\[
  \sum_i u_i(a^M)-h(a^M)
  <
  \sum_i u_i(a')-h(a')
\]
对某个可行 \(a'\) 成立。
\end{Certificate}

\begin{proof}[证明]
IC、IR 与预算平衡只约束报告、参与和转移：
\[
  \IC\wedge\IR\wedge\BudgetBalanced
  \Longrightarrow
  \text{报告与资金条件闭合}.
\]
若社会目标还含未进入机制的 \(h(a)\)，则
\[
  h\notin\operatorname{Objective}(G)
  \Longrightarrow
  \argmax \text{机制目标}
  \ne
  \argmax \text{社会福利}
\]
可以发生。一般机制设计中效率、激励相容与预算条件存在真实权衡，不能把多项条件当作自动兼得
\cite{GreenLaffont1979Incentives}。证毕。
\end{proof}

\begin{proposition}[一般均衡与福利定理的条件缺口命题]
\label{prop:general-equilibrium-gap}
市场参与经济若要确认一般均衡存在核验记录，至少需要明确商品空间、偏好、禀赋、生产集合、收益权市场、
价格系统与机制可行域，并验证相应的非空性、凸性、连续性或替代条件。若要确认福利定理，还需处理
外部性、市场势力、信息不对称和制度运行成本。
\end{proposition}

\begin{Certificate}[Arrow--Debreu 对照清单]
定义一般均衡核验记录向量
\[
  \mathbf g_{\mathrm{GE}}
  =
  (g_{\mathrm{commodity}},g_{\mathrm{preference}},g_{\mathrm{endowment}},
  g_{\mathrm{production}},g_{\mathrm{price}},g_{\mathrm{clearing}},
  g_{\mathrm{governance}}).
\]
只有全部坐标具有形式定义并满足选定存在定理的假设时，才允许声称均衡存在。
经典竞争均衡存在背景见~\cite{ArrowDebreu1954Equilibrium}。
\end{Certificate}

\begin{proof}[证明]
市场出清需要
\[
  \text{需求对应}\wedge\text{供给对应}\wedge\text{价格空间}
  \Longrightarrow
  \text{超额需求对应}.
\]
固定点存在还需超额需求对应满足所用固定点定理的定义域、闭值、凸值和连续性等假设。
若任一对象未定义，则对应映射不存在；若假设未验证，则存在定理不能实例化。因此局部分配公式不能替代
一般均衡证明；福利比较、经验校准与动态学习也必须分别补充偏好、技术、识别和状态转移假设。证毕。
\end{proof}

\begin{corollary}[替代制度优越性不可由内部一致性推出]
\label{cor:no-superiority-without-comparison}
即使市场内参与确权机制逻辑一致，也不能据此推出它优于数据分红、平台合作社、工资合同或传统市场。
同理，本文既不否定按需分配、基本收入、公共供给、按劳分配及其他经济理论，也不由自身成立推出这些制度失效。
优越性需要共同结果指标、可比基线、识别策略和数据核验记录；组合可行性则由核验记录
~\ref{cert:minimal-compatible-path} 的预算、资源、权利、法律和制度环节一致性逐项判定。
\end{corollary}

\begin{Certificate}[可比较处理效应]
对候选机制 \(G\) 与基线 \(G_0\)，定义
\[
  \Delta\mathcal J
  =
  \mathcal J(G)-\mathcal J(G_0).
\]
只有在共同观测口径、可识别反事实与置信界
\[
  \underline{\Delta\mathcal J}>0
\]
成立时，才确认“在该指标与证据域内优于基线”的核验记录。
\end{Certificate}

\begin{proof}[证明]
\[
  \text{内部一致}
  \Longrightarrow
  \text{机制结果可计算}
\]
不蕴含
\[
  \mathcal J(G)>\mathcal J(G_0).
\]
后者需要同口径差值及其识别误差。若置信下界不为正，数据与模型不能排除无效或劣化，故不能确认
优越性结论。证毕。
\end{proof}

\section{可证伪预测与经验审计}

\begin{remark}[十八项预测的统一适用域前件]
\label{rem:prediction-post-ai-scope}
预注册以前，必须先固定有限非空产品域 \(\Goods_0\)、技术状态 \(A\)、互补资源状态 \(r\)、容限
\(\epsilon\) 与时间窗，并提交核验记录~\ref{cert:post-ai-applicability}。若
\(\kappa_{\mathrm{post}}^\epsilon=1\)，P1--P18 可以作为后 AI 理论的联合预测包接受检验；若
\(0<\kappa_{\mathrm{post}}^\epsilon<1\)，所有涉及产能松弛、稀缺迁移或总体循环的结论只能投影到
\(\Goods_{\mathrm{post}}^\epsilon\)，其余不依赖松弛的制度预测仍可作为局部命题检验；若
\(\kappa_{\mathrm{post}}^\epsilon=0\)，不得把整组结果表述为后 AI 理论的经验检验。

生产成本为正本身不构成域外见证；真正的域外见证是生产成本或其他生产约束仍在边际上持续截断有限最优点，且
既没有精确松弛，也没有通过渐近松弛核验记录。未先通过适用域前件时，观察结果只能说明前提尚未满足，不能被记作
后 AI 理论已经得到支持或遭到反驳。

若经验任务进一步评价“应否采用本文的制度扩展”，还必须固定现行制度基线、具体失灵点和候选路径族，并提交
\(\mathfrak C_{\mathrm{MCP}}\)。后 AI 适用域通过只证明技术经济前提成立，不单独证明本文路径已经必要；失灵点
通过只证明需要修复，也不单独证明某个改造是改动最小或优于其他制度。
\end{remark}

\begin{proposition}[市场参与与注意力桥接的可证伪预测命题]
\label{prop:falsifiable-predictions}
在适用域状态已经按备注~\ref{rem:prediction-post-ai-scope} 固定，并控制技术、需求和参与者构成后，该框架至少
产生以下可证伪预测：
\begin{enumerate}[label=\textnormal{P\arabic*:},leftmargin=3.5em]
  \item 规则变化在支付映射有非零响应时改变收入；只有跨越资格或参与阈值才预测离散变化，内点分成变化可以产生连续响应，参与供给还须由成本与外部选择模型检验；
  \item 产能约束松弛后，单纯扩产的边际效应下降；只有相应制度约束具有正影子价格且制度投入确实放松约束时，其相对边际作用才可能超过扩产，不能推出其绝对效应必然上升；
  \item 在实际采用二级转让的分支中，流动性上界下降时，最终持有量与初始收益权的偏离上界下降；
  \item 在按登记身份等份分配且缺少主体关联核验的固定资金池中，当重复申领额外成本低于
  命题~\ref{prop:duplicate-claim-gate} 的份额增益门槛时，重复申领具有正净收益；实际发生率还须控制
  被查获概率、处罚及其他行为成本；
  \item 在生产状态与治理状态分开记录、判定质量不下降且成本归属可比时，治理自动化提高应使标准化记录、
  计量、复算和执行的单位成本变化按同口径检验；总单位成本与收入占比只在相应极限条件下预测收敛。复核效果分别估计，真实参与者退出率的方向另依赖其效用与外部选择，不能由审计一致性直接推出；
  \item 在商品、价格和需求可比时，竞争者投入对本方注意力份额与最优投入的效应分别检验；内点严格凹模型中，最优回应方向由收益交叉偏导决定。定理~\ref{thm:constructive-equilibrium-cycle} 的均衡邻域预测最优回应斜率为 $-1/4$，不预测普遍加码；
  \item 在正必要品需要、至少两个可行载体和分散选择同时成立时，已实现购买份额应始终形成非零选择份额投射；
  但若注意力变化长期不能改变消费者对赞助生产者商品的考虑集、载体份额、偏好或购买，且不存在其他可货币化
  回报，则生产者注意力投入应收缩而不是持续增长；
  \item 竞争性注资经桥接行业形成的参与者净收入增加时，具有正边际消费倾向的参与者应增加生产品消费；
  \item 在风险阈值不变且没有新增外溢证据时，安全扩张弱增可选择的合法模式集合；实际采用模式的增加还要求新增模式有参与需求、融资与可行利润，不能由集合扩张直接推出；若风险外溢
  同步越过阈值，则空间扩张不满足安全极大域条件。
  \item 在收入池与参与构成可比时，提高独立重放一致率、等价参与横向一致性和申诉实际修复率，应降低无法
  解释的规则内份额差异；误判修复速度与退出效应另行识别。未减少退出本身不反驳规则内公正，只有对应行为前提独立通过时才检验其方向；
  \item 组织共用收入池或资产而丢失必要来源记录时，不能完成逐笔权利来源验证；分配争议发生率的变化另须行为与识别条件，不能由记录缺失直接推出。若混合池仍能逐笔复算全部来源，混合本身不使认证失败；
  \item 在最低物质支出、价格和偏好参数可比的收入区间内，超额预算型兴趣品需求的收入弹性应大于一，且兴趣
  支出占收入份额应随收入上升而提高；若稳定估计不支持该方向，广义奢侈品属性在该区间被反驳；
  \item 当家庭聚合超额预算增加或标准化进入成本下降时，达到零利润门槛的兴趣品类数应弱增；若成本和需求
  条件可比而活跃品类持续下降，品类进入命题需要修订；
  \item 只有兴趣行业带来的生产商品偏好购买增量与参与者消费增量，在扣除替代和泄漏后不小于事前需求缺口，
  并且净余额与动态稳定条件通过时，才能观察到声明域循环被持续支撑；低于该规模却仍长期支撑循环的稳定反例
  将反驳相应规模下界或缺口测量；
  \item 对市场收入低于底线束的居民，少于
  \([F_{ht}-Y_{ht}^{m}]_+\) 的到账救济不能实现百分之百底线覆盖；达到该值后，具有正边际消费倾向者的
  生产品消费和具有真实入口者的市场参与概率应弱增，否则相应消费或进入单调性前提被反驳；
  \item 在生产盈利、缴款率、公共收益和管理成本可比时，救济支出持续超过账户收入与可用储备的方案不能维持
  声明的跨期底线；若无外部资金仍能长期维持，则预算口径遗漏了收入或损失承担者；
  \item 固定其他主体行动、外部选择及成本，且固定补偿与努力无关、期望剩余随自身努力非减时，按命题~\ref{prop:residual-share-monotone-response} 比较唯一最优解或同一极端选择，最优投入随剩余份额弱增。份额降为零时严格正努力成本给出零投入；原投入已经为零时，零响应不构成反例。策略反馈下的均衡响应另证，不能把单人比较推广为所有均衡的比较；
  \item 在共同可比域内，能力提高边际剩余产出且验证得分保留贡献顺序时，采用正基本份额、正差异池与严格
  递增贡献权重的合同，应使每名有效参与者取得正分配、贡献更高者取得严格更多分配；任何已完成约定者被分配
  零基本份额，或严格贡献差异未产生分配差异，均反驳相应合同执行结论。
\end{enumerate}
\end{proposition}

\begin{Certificate}[预测依赖图与整体、部分、域外标签]
定义
\[
  \mathfrak C_{\mathrm{pred}}^\epsilon
  =
  \Tuple{
  \mathfrak C_{\mathrm{post}}^\epsilon,
  \mathcal D_{P1},\ldots,\mathcal D_{P18},
  \mathsf{Scope},
  \mathsf{Window}
  },
\]
其中 \(\mathcal D_{Pk}\) 保存第 \(k\) 项预测实际调用的局部前提与数据字段，且
\[
  \mathsf{Scope}
  =
  \begin{cases}
    \mathsf{FullPostAI},&\kappa_{\mathrm{post}}^\epsilon=1,\\
    \mathsf{PartialPostAI},&0<\kappa_{\mathrm{post}}^\epsilon<1,\\
    \mathsf{OutsidePostAI},&\kappa_{\mathrm{post}}^\epsilon=0.
  \end{cases}
\]
P2 只在 \(\Goods_{\mathrm{post}}^\epsilon\) 上检验：精确松弛允许检验边际效应为零，渐近松弛只允许检验其绝对
值不超过预注册容限并继续趋近于零。其他局部制度预测即使能在
\(\mathsf{PartialPostAI}\) 或 \(\mathsf{OutsidePostAI}\) 状态下检验，也不得反向补出缺失的后 AI 适用域核验记录。
\end{Certificate}

\begin{Certificate}[预注册估计式]
对规则变更时点 \(t_0\)，可预注册
\[
  y_{it}
  =
  \alpha_i+\gamma_t
  +\beta\,
  \mathbf 1\{t\ge t_0\}\operatorname{Treated}_i
  +\delta^\top X_{it}+\epsilon_{it}.
\]
核验记录记录处理定义、对照组、结果变量、协变量、窗口、标准误口径与排除规则。
\end{Certificate}

\begin{proof}[证明]
前五条预测把组织制度参数变化映射到可观测结果：
\[
  \Delta G
  \Longrightarrow
  \Delta(E,M,V,W,\phi,T,L,Q)
  \Longrightarrow
  \Delta(y,N,\operatorname{Manipulation},\operatorname{Exit}).
\]
P5 还要求把生产状态 \(A\) 与治理自动化状态 \(A_G\) 分开，并将固定、标准化和残余成本同口径拆分；生产
产能松弛不能替代治理成本核验记录。P6--P9 分别检验竞争反应、偏好传导、收入回流与安全空间。P7 先检验正刚需
与多选下非零选择份额是否存在，再独立检验注意力是否改变该份额；前者成立不补出后者。它们要求观察
\(\mathbf I^B,\mathbf a^{\mathrm{att}},\boldsymbol\vartheta^P,D^P,\mathbf T^B,C^P,
r_{\mathrm{spill}}(\gamma)\)，不能用总曝光量
替代全部坐标。P10--P11 检验参与--分配判定质量和组织跨活动来源隔离，另需保存独立重放差异、等价类
份额差、申诉修复、活动核验记录覆盖率及跨活动资产来源。P12--P14 分别检验兴趣需求的 Engel 弹性、超额预算
与品类进入阈值，以及净需求桥接是否跨过宏观规模门槛；它们需要最低物质支出、可支配收入、价格、进入退出、
替代泄漏和需求缺口数据。P15--P16 检验坐标最小救济、消费与进入响应以及跨期资金闭合，须保存逐户底线、
市场收入、实际到账、边际消费倾向、进入状态、生产盈利、缴款、公共收益、管理成本和储备。P17--P18 检验
合同剩余索取的剩余激励与能者多得机制，须区分固定补偿、剩余份额、投入成本、边际余额贡献、能力和已验证贡献。
若控制条件与识别假设成立而估计结果稳定地与预测方向相反，则相应局部命题被反驳或需要修订。
只有适用域核验记录与相应依赖图同时通过时，反向结果才构成对该后 AI 命题的有效反证；适用域失败只产生
\(\mathsf{PartialPostAI}\) 或 \(\mathsf{OutsidePostAI}\) 标签。因此该框架的局部规则效应由预注册比较和反事实
结果接受检验，而后 AI 总体解释另受适用域前件约束。证毕。
\end{proof}

\section{竞技参与经济的一个可解联合模型}
\label{sec:constructive-participation-economy}

本节固定两家生产者、一家赛事组织和两名合格参与者，联合给出产品需求、赞助决策、真实参与、
贡献核验、预算分配和跨期现金更新。计价单位固定，价格为一；背景消费者有来自域外部门的既有收入，
本节为开放市场中的部分均衡模型。背景购买力逐期作为明示禀赋登记，不归因于赞助创造货币。
外部损害在本实例中取零，产品产能在下述状态矩形内充足；
公共救济、文化形成和一般均衡结论仍分别服从其自身条件。

\begin{definition}[家庭购买、产品成本与赛事收入]
\label{def:constructive-market}
两家生产者 \(k=1,2\) 的单位产品成本为 \(1-m\)，其中 \(0<m<1\)。
背景消费者各期对每家产品有固定购买量 \(q_0\)，另把数量为 \(Q>0\) 的购买预算在两家产品与外部商品之间配置：
\[
 a_k(I)=\frac{I_k}{a+I_1+I_2},\qquad
 a_0(I)=\frac{a}{a+I_1+I_2},\qquad a>0.
\]
这些份额非负且和为一。消费者的归一化特征效用
\(\sum_{j=0}^2 a_j(I)\log x_j\)，在 \(\sum_jx_j\le Q\) 下给出 \(x_j=Qa_j(I)\)；
零权重对象取零购买量。注意力与商品联系通过这一明示偏好映射进入需求，偏好参数须另行识别。
品牌间及品牌与外部商品间的份额变化均不作为全市场新增购买力。

组织先收到自愿门票与服务收入 \(F=C+2\beta\)，其中 \(\beta>0\) 为每名完成基础任务者的基本份额，
\(C>0\) 为组织与全部核验成本。赞助承诺在接受参与合同前到账。于是
\[
 R=F+I_1+I_2,\qquad P=R-C=2\beta+I_1+I_2,\qquad D=I_1+I_2.
\]
最大接纳人数为二，基本份额始终由既有服务收入覆盖；正差异分配域取 \(D>0\)。
若两家均不赞助，已承诺的基本份额仍可支付，但该状态不属于正差异池的核心循环域。
门票购买者的自愿性要求其服务效用至少覆盖票款，票款同时列入其预算支出。
具体可取一名观众，其准线性效用为 \(v s+c\)，入场选择 \(s\in\{0,1\}\)，
预算为 \(Fs+c\le w\)。只要 \(w\ge F\) 且 \(v>F\)，购买门票便是严格最优选择。

两名参与者把到账收入的比例 \(\chi\in(0,1)\) 用于这两家商品，并在两家之间等额消费，
其余用于外部商品或储蓄。这可由三类数量的齐次 Cobb--Douglas 效用与支出权重
\((\chi/2,\chi/2,1-\chi)\) 得到；选取正的效用归一化系数，使价格为一时的间接消费效用等于货币收入，
从而与下文参与效用中的 \(T_i\) 一致。两家实际需求、利润与家庭收支为
\[
 \begin{aligned}
 q_k&=q_0+Qa_k+\frac{\chi P}{2},\\
 \pi_k&=mq_k-I_k=\pi_0+B a_k+\frac{\eta P}{2}-I_k,\\
 \pi_0&=mq_0,\qquad B=mQ,\qquad \eta=m\chi\in(0,1).
 \end{aligned}
\]
家庭消费不超过其收入；企业销售回款 \(q_k\) 分解为产品成本 \((1-m)q_k\)、营销支出 \(I_k\)
与净利润 \(\pi_k\)。组织账分解为成本 \(C\) 与参与支付 \(P\)，合并时赞助和分配两类转移抵销。
本实例无额外配套营销支出；如增加此项，必须重新求解投入目标与预算。
\end{definition}

\begin{definition}[可核验任务与两阶段纠错合同]
\label{def:audited-effort-contract}
基础任务为按公开规则完成对局与赛后确认，完成成本为 \(b_0>0\)。
参与者可以另选可观察的复盘服务作为贡献行动 \(e_i\in\{0,1\}\)，其私人努力成本为 \(k_i e_i\)。
共同服务质量给每名参与者带来效用 \(h(e_1+e_2)\)，其中 \(h\ge0\)；该效用不是货币收入。
两名主体按相同的公开资格条件接纳，合同期内规则固定，参与者可以真实退出。
原始资料分别记录真实主体、基础任务完成及贡献事件，同一主体与事件只确认一次，全部支付可独立核对。

核验采用两阶段制度：初核可有错误率 \(\varepsilon\in[0,1]\)；支付前对每份记录作独立复核，
依据已认证的原始资料逐项比对，修正初核差异，并允许参与者免费提交遗漏记录的更正申请。
本可核验任务域要求原始资料完整、可取得且不可由参与者伪造，复核按有限记录的确定规则给出
真实完成状态与 \(e_i\)。本实例的组织成本归于核验与结算服务，每人核验费 \(\kappa\) 覆盖身份核对、初核、全部复核及更正；
即使全部初核出错，总成本仍不超过 \(C=2\kappa\)。因此在该资料技术域内，终审误判率为零，
而非将初核错误隐去；若资料缺失或终审仍有残余误判，则不适用本实例的精确排序结论。

个人自行填写的贡献数字只用于核对，不进入支付函数。对两名真实完成基础任务者，
\[
 T_i=\beta+D\lambda_i,\qquad
 \lambda_i=
 \begin{cases}
 e_i/(e_1+e_2),&e_1+e_2>0,\\
 1/2,&e_1+e_2=0,
 \end{cases}
 \qquad
 U_i=T_i+h(e_1+e_2)-k_ie_i-b_0.
\]
虚假身份与未完成基础任务的申请不生成支付义务。若仅一人完成基础任务，该人取得基本份额与差异池，
未支付的另一份基本份额留在组织账户；若无人完成，全部参与预算留存。由此合同也定义了退出状态的支付。
对已真实完成任务的参与者，申报错误的修复不取消其基本份额。
原始资料与复核机构的独立性属于核验技术前提；反合谋结论的参与者范围限于上述两人，不包括能控制核验机构的联盟。
\end{definition}

\begin{theorem}[预算、真实努力、自愿参与与两人合谋约束的联合实现]
\label{thm:audited-effort-implementation}
若外部选择效用为零，且
\[
 \beta>b_0,\qquad
 k_1-h<D/2<k_2-h,\qquad
 k_1<2h<k_2,
\]
则上述合同以唯一严格努力均衡 \((e_1,e_2)=(1,0)\) 实现
\[
 T_1=\beta+D>T_2=\beta>0,\qquad T_1+T_2=P,
\]
且两人均愿意参与。虚报和重复身份不提高支付；两名参与者即使可以相互补偿，也不能通过联合改变努力
提高两人的总效用。真实努力、核验及支付都在 \(R=C+P\) 内完成，罚款收入不构成资金来源。
\end{theorem}

\begin{proof}[证明]
无论另一人的努力为零还是一，把自身努力从零改为一都会增加差异收入 \(D/2\)，
故私人效用增量为 \(D/2+h-k_i\)。所列严格不等式使第一人努力、第二人不努力均为严格占优行动。
第一人效用为 \(\beta+D+h-k_1-b_0>\beta+D/2-b_0>0\)，
第二人为 \(\beta+h-b_0>0\)，故外部选择不会取代该参与均衡。
更一般地，不论他人是否完成，完成基础任务而不选附加努力至少取得 \(\beta-b_0>0\)；
故所有赞助投入下的延续博弈都排除退出，这也保证赞助者偏离时的均衡总支付仍为 \(P\)。
核验后支付只依赖真实记录，虚报不改变支付，复制身份不能产生新权利。
联盟内部转移抵销后，总效用为
\[
 U_1+U_2=P-2b_0+(2h-k_1)e_1+(2h-k_2)e_2.
\]
由 \(2h-k_1>0>2h-k_2\)，联盟最优努力同样是 \((1,0)\)；单人偏离已由占优行动排除。
对只有一人完成的联盟方案，让另一人完成基础任务且不作附加努力，会增加联盟总收入 \(\beta\)，
只增加成本 \(b_0\)，并使共同服务效用非减，故退出也不能改善联盟收益；两人均退出的效用为零。
终审纠错保证使用的贡献等于真实贡献，基本与差异份额总额恒为 \(P\)，核验费已包含在 \(C\) 中。证毕。
\end{proof}

\begin{definition}[赞助最优回应与现金更新]
赞助者在已实现且未另作承诺的利润现金内选择 \(I_k\in[0,L_k]\)，其目标为
\(\pi_k(I_k,I_{-k})\)。合同依次完成资金承诺、参与选择、核验及结算，参与者行为不改变
总分配 \(P\)，因而上述需求函数适用于赞助者的偏离比较。令 \(d=1-\eta/2>0\)，则
\[
 \partial_{I_k}\pi_k
   =\frac{B(a+I_{-k})}{(a+I_1+I_2)^2}-d,\qquad
 \partial_{I_k}^2\pi_k
   =-\frac{2B(a+I_{-k})}{(a+I_1+I_2)^3}<0,
\]
\[
 BR_k(I_{-k})=
 \left[\sqrt{\frac{B(a+I_{-k})}{d}}-a-I_{-k}\right]_{[0,L_k]}.
\]
符号 \([x]_{[l,u]}=\min\{u,\max\{l,x\}\}\) 表示区间截断。偏好营销净收益保持消费者收入于同一
实现值 \(P\)，只比较本品牌有无赞助，故
\[
 \Gamma_k^{\mathrm{mkt}}
 =B a_k-I_k
 =I_k\left(\frac{B}{a+I_1+I_2}-1\right).
\]
赞助者总利润还包含收入通道 \(\eta P/2\)；这一项不计入上述固定收入的偏好增量。
因此最优回应中的总利润提高与偏好回报为正分别核对。
\end{definition}

\begin{theorem}[非空参数域中的联合均衡与局部稳定循环]
\label{thm:constructive-equilibrium-cycle}
取下列参数点：
\[
 \begin{gathered}
 a=1,\quad m=\chi=\tfrac12,\quad \eta=\tfrac14,\quad
 q_0=8,\quad Q=\tfrac{63}{8},\quad \pi_0=4,\quad B=\tfrac{63}{16},\\
 \beta=\tfrac12,\quad \kappa=\tfrac14,\quad C=\tfrac12,\quad F=\tfrac32,\quad
 b_0=\tfrac1{10},\quad h=\tfrac3{10},\quad
 k_1=\tfrac12,\quad k_2=2.
 \end{gathered}
\]
观众取 \(v=w=2\)，故 \(v>F\) 且 \(w>F\)，门票购买与预算约束同时成立。
初始两家可用现金均位于 \([4,5]\)，且在声明窗口内全部净利润留作下一期可用现金。
则阶段博弈存在唯一赞助均衡 \(I_1^*=I_2^*=1\)，对应唯一参与努力均衡 \((1,0)\)，并有
\[
 D^*=2,\quad P^*=3,\quad (T_1^*,T_2^*)=(\tfrac52,\tfrac12),\qquad
 \Gamma_k^{\mathrm{mkt},*}=\tfrac5{16},\quad \pi_k^*=\tfrac{75}{16}.
\]
令投入按阻尼最优回应调整：
\[
 I_{k,t+1}=(1-\delta)I_{k,t}+\delta BR_k(I_{-k,t})+\xi_{k,t},
 \qquad L_{k,t+1}=\pi_k(I_t),\quad 0<\delta\le1.
\]
对状态矩形
\[
 \mathcal V=[\tfrac9{10},\tfrac{11}{10}]^2\times[4,5]^2
\]
及投入调整扰动 \(\|\xi_t\|_\infty\le7\delta/100\)，该更新前向不变。
其中每期核验、基本与差异分配、正偏好营销净收益及下一期现金覆盖同时成立。
无冲击时唯一固定点局部渐近稳定；以上严格经济不等式也在满足会计定义关系的自由参数空间内、该点的某个开邻域内成立。
\end{theorem}

\begin{proof}[证明（阶段均衡、统一界与显式雅可比）]
这里 \(d=7/8\)，未截断最优回应的最大值为 \(B/(4d)=9/8<4\)，故现金上限不约束最优回应。
若一家零投入，另一家的最优投入为 \(\sqrt{9/2}-1<9/8\)，零投入者的边际收益严格为正，
因此均衡不在坐标轴。两家内点一阶条件相减给出 \(I_1=I_2\)，再由
\[
 B(1+I)=d(1+2I)^2
\]
得到唯一非负根 \(I=1\)。在该投入矩形上 \(D/2\in[9/10,11/10]\)，而
\(k_1-h=1/5<9/10\)、\(k_2-h=17/10>11/10\) 及
\(k_1=1/2<2h=3/5<k_2=2\)，故努力、参与和合谋约束在整个矩形严格成立。
代入组织账和利润式即得所列均衡支付与收益。

对 \(z\in[9/10,11/10]\)，最优回应导数为
\[
 b'(z)=\frac12\sqrt{\frac{9/2}{1+z}}-1,\qquad |b'(z)|<3/10,\qquad b(1)=1.
\]
故阻尼映射的最大范数收缩系数不超过 \(q_\delta=1-7\delta/10\)。
由 \(\|I-\mathbf1\|_\infty\le1/10\)，有
\[
 \|I_{t+1}-\mathbf1\|_\infty
 \le q_\delta/10+7\delta/100=1/10.
\]
矩形内 \(P\in[14/5,16/5]\)，并且
\[
 \Gamma_k^{\mathrm{mkt}}
 \ge\frac9{10}\left(\frac{63/16}{16/5}-1\right)
 =\frac{531}{2560}>0.
\]
利用 \(\pi_k=4+\eta P/2+\Gamma_k^{\mathrm{mkt}}\)，下界严格大于四；上界满足
\[
 \pi_k\le4+\frac25
       +\frac{11}{10}\left(\frac{63/16}{14/5}-1\right)<5.
\]
因而下一期可用现金仍在 \([4,5]\)，超过最大投入 \(11/10\)，组织接受参与前的基本与差异预算亦已到账。

固定点处 \(b'(1)=-1/4\)。以 \((I_1,I_2,L_1,L_2)\) 排列状态，现金上限不约束最优回应，
完整雅可比为分块下三角矩阵，其特征值为
\[
 1-\tfrac34\delta,\qquad 1-\tfrac54\delta,\qquad 0,\qquad0.
\]
当 \(0<\delta\le1\) 时均严格位于单位圆内，故无冲击固定点局部渐近稳定。
有持续有界扰动时只保证留在可行矩形及相应误差界，不声称收敛到同一固定点。
预算、收益、激励、现金和收缩条件在所列紧矩形上具有严格余量；连续性保证足够小的参数变动
仍存在邻近均衡和不变邻域，从而适用参数域非空且不限于孤立点。证毕。
\end{proof}

\begin{remark}[模型的经济边界]
该模型同时构造了赞助最优回应、参与供给、有限成本核验、实际支付、消费回流及稳定现金周转。
消费者的背景收入、门票需求、注意力偏好函数与资料可核验性均为明示原始条件；
它证明这些条件可以相容，不证明现实市场已经具有所用数值或所有主体都能进入相同任务。
终审仍有误差、核验机构被控制、背景需求下降、利润另作分配或冲击超出允许范围时，须重算预算、行为偏离与稳定界。
文化、社会救济和一般均衡不由这个局部实例自动推出。
本例两家商品的偏好需求总响应为 \(Qa/(a+I_1+I_2)^2>0\)，对应外部商品的减少须一并记录；
全商品域的这项份额重分配净额为零。收入回流另产生本产品域的 \(\chi\,dP\) 消费响应。
由于固定点附近现金上限不约束最优回应，本例证明正额的持续周转，不证明利润每增加一单位都增加赞助；
如需后一结论，仍须验证定理~\ref{thm:dual-channel-bridge} 的严格投入响应条件。
\end{remark}



\section{核验误差、组织合同与有限接纳规模}
\label{sec:contract-scale-extension}

本节在一个公开合同菜单内求解组织者、赞助者与参与者的选择。所有金额均为模型计价单位。
组织者可不开办，利润为零；开办时选择整数 \(n\ge1\) 个双人组、基本份额 \(\beta\)
和复核强度 \(r\ge0\)。每组有两类可观察资格角色，各有至少两名申请者，努力成本分别为
\(k_1=2/5,k_2=2\)。成本不同不改变同一合同的计价规则。各角色组内按公开随机顺序接纳，
每个真实主体只进入一个名额；申请人数构成额外上限。配对在贡献发生前公开固定，各双人任务组构成独立的贡献比较域。未接纳者没有已订立支付合同，
不能把其零支付记为已完成合同者的基本份额违约，也不能把接纳者收入外推为全员收入。

\begin{definition}[参与约束、残余误差与统一资金用途]
\label{def:noisy-contract-menu}
基础任务完成与真实身份可以准确核验。参与者没有可用于垫付本期必要支出的流动资产，
其接受合同的现金保留条件是每个结算状态均有 \(T_i\ge\ell=1/2\)；这是一类有限流动性参与者的偏好与约束，
不假定所有社会成员都具有同一保留条件。另要求期望效用不低于零，
\[
 U_i=\EE[T_i]+h(e_1+e_2)-k_i e_i-b_0,
 \qquad h=3/10,\quad b_0=1/10,\quad e_i\in\{0,1\}.
\]
接纳承诺不以赞助包售出为生效条件：三个候选合同均须在包括零赞助的全部投入组合下满足现金保留条件，不能以售包失败取消已承诺名额。
因此本节比较的是具有事前资金保障的合同菜单；若允许附条件接纳，固定报酬可改由赞助承担，组织者选择必须重新求解。
没有外部信用或保证保险替代该现金条件；若引入它们，须把保费、赔付和融资列入合同预算再比较。

附加贡献的最终信号 \(s_i\in\{0,1\}\) 独立地以概率
\[
 \varepsilon(r)=\tfrac14\exp(-r)
\]
偏离真实努力。这里的误差已经包含初核和申诉后的残余误差；不能继续称为零误判。
信号技术由独立复核机构提供，参与者不能控制记录或复核人员；个人自报不进入支付函数。
公开质量合同要求 \(\varepsilon(r)\le1/10\)。核验与组织的完整成本为
\[
 C(n,r)=n\left(\tfrac25+\tfrac{r^2}{100}\right)+\tfrac{n^2}{10}.
\]
它包含基础身份核对、全部复核、已承诺的纠错和组织服务；若出现域外操纵或额外成本，必须重算。

两家赞助者分别选择 \(I_k\in\{0,1\}\)，表示拒绝或购买一个标准赞助包，令 \(S=I_1+I_2\)。
赞助合同约定 \(S\) 全部进入参与分配池，组织者不能将其改作利润；门票收入覆盖基本份额、组织成本和剩余。
比较合同均将总额 \(P=2n\beta+S\) 用于参与支付。基本加差异合同每组的差异池为 \(d_n=S/n\)，且
\[
 T_i=\beta+\frac{d_n}{2}(1+s_i-s_j).
\]
固定报酬合同在同一预算下给每人 \(P/(2n)\)；纯奖金合同把同一 \(P\) 全部按上述信号权重分配，基本份额为零。
这三个公开合同构成本节比较域，努力按个人最优回应选择；额外可强制执行的组内再分配或努力奖励协议不在菜单内。
若允许这类协议，固定报酬也可能诱导合作，须把协议作为新的合同重算。本节不声称穷尽全部可行制度。
\end{definition}

\begin{proposition}[有误差时的预算、激励与贡献排序]
\label{prop:noisy-effort-contract}
在基本加差异合同中，每个信号状态都有
\(T_i\ge\beta\)、\(T_1+T_2=2\beta+d_n\)，且
\[
 \begin{aligned}
 \EE[T_i\mid e]&=\beta+\frac{d_n}{2}
       \bigl[1+(1-2\varepsilon)(e_i-e_j)\bigr],\\
 \Delta_{e_i}U_i&=\frac{d_n}{2}(1-2\varepsilon)+h-k_i.
 \end{aligned}
\]
若 \(d_n(1-2\varepsilon)/2+h-k_1>0>d_n(1-2\varepsilon)/2+h-k_2\)，
唯一严格努力均衡为 \((1,0)\)。两人完成时，真实努力较高者的期望支付严格较高；
同组实际支付随最终证据评分严格排序，但同组真实努力被反向支付的概率为 \(\varepsilon^2\)，
平额概率为 \(2\varepsilon(1-\varepsilon)\)。基本份额与资金总额不受这些误差影响。
\end{proposition}

\begin{proof}[证明]
二元信号下的归一权重恰为 \((1+s_i-s_j)/2\)，且
\(\EE[s_i\mid e_i]=\varepsilon+(1-2\varepsilon)e_i\)，故所列期望与边际效用直接成立。
努力条件分别使一人努力、另一人不努力成为严格占优选择。真实努力为 \((1,0)\) 时，
只有 \((s_1,s_2)=(0,1)\) 产生反向支付，独立性给出概率 \(\varepsilon^2\)。
若 \(k_1<2h<k_2\)，两名参与者的期望效用总和仍为
\(2\beta+d_n-2b_0+(2h-k_1)e_1+(2h-k_2)e_2\)，
联合改变努力也不能提高该总和。该结论不排除合法保险的风险分担收益，不覆盖操纵复核机构的联盟。
\end{proof}

\begin{remark}[精确贡献与带误差贡献的结论类型]
正基本份额、预算恒等式和同组最终可区分证据评分的严格排序逐状态成立。
不同组的支付还受各自同伴信号影响；本例不确认跨组绝对贡献排序。需要全组织排序时，须另设统一比较与支付函数并重证激励。
真实贡献的排序在本节为期望及错误概率结论，不能登记为逐状态绝对正确。
一般相对公正核验须分别保存规则内复算一致性、信号的校准误差、真实贡献排序误差和纠错成本；
一项通过不替代另一项。若合同承诺绝对正确的真实贡献排序，正的残余误差即使很小也不满足该承诺。
\end{remark}

\begin{definition}[贡献质量、付费需求与赞助收益]
\label{def:quality-feedback-demand}
令 \(z\in[0,1]\) 为各组低成本角色的平均努力。预期质量通过实际消费选择影响赛事收入和商品联系：
\[
 F(n,z)=\frac{(10+2z)n}{1+n},\qquad
 B(n,z)=\frac{13}{4}+\frac{nz}{1+n}.
\]
第一式可由观众对延续均衡预期质量的总保留价格 \(F(n,z)\) 与无差异时接受的购票规则实现；
基础票款 \(F(n,0)\) 先到账，赞助组合确定后，观众根据预期质量支付附加票款，再开始努力选择。
买者承担已披露的服务质量风险，票款不按有噪贡献信号重新结算。在下述严格努力均衡中，预期质量等于真实质量。
观众预算足以覆盖最高票款，其价值和资金来自明确的既有收入。附加票款未到账以前不得预先承诺给参与者。
因此准入融资必须满足
\[
 F(n,0)\ge C(n,r)+2n\beta.
\]
定义产品价格为一、单位成本 \(1/2\)、背景销量为八，参与收入的商品消费比例为 \(1/2\)。
背景消费者总注意力预算上限为 \(17/2\)：其中 \(2B(n,z)\) 按
\(a_k=I_k/(1+S)\)、\(a_0=1/(1+S)\) 在两家商品与外部商品间配置，其余购买外部商品。
这给出真实、有限的购买预算；质量提高带来的本产品需求增加对应外部商品变化，不能直接称作全市场新增需求。
两家生产者利润为
\[
 \pi_k=4+B(n,z)\frac{I_k}{1+S}+\frac{P}{8}-I_k.
\]
组织者利润为 \(\Omega=F(n,z)-C(n,r)-2n\beta\)，
满足 \(F+S=C+P+\Omega\)。\(\Omega\) 是正当登记的组织剩余，并非再次从已承诺份额中扣费。
时序为：组织者提出合同及接纳方案；基础票款与资金承诺到位；赞助者选择；观众确认附加票款；参与者接受并努力；
核验、纠错、支付及票款结算。赞助者预见各个投入组合下的延续行为，不能把质量固定在有利值后忽略偏离。
\end{definition}

\begin{theorem}[组织选择、误差精度、质量反馈与有限规模的联合实现]
\label{thm:endogenous-organization-contract}
若申请人数允许每类接纳两人，生产者可用现金至少为一，则上述合同菜单中有严格最优组织方案
\[
 n^*=2,\qquad \beta^*=\ell=\tfrac12,\qquad
 r^*=\log(5/2),\qquad I_1^*=I_2^*=1.
\]
这里 \(r^*\) 是在最优合同内的最小成本精度选择。每组努力为 \((1,0)\)，总接纳人数为四，
最终贡献信号错误率为 \(1/10\)。令 \(c=2/5+(\log(5/2))^2/100\)，则
\[
 \begin{gathered}
 P^*=4,\qquad F^*=8,\qquad C^*=2c+2/5,\qquad
 \Omega^*=8-2c-2/5-2>0,\\
 \EE(T_1,T_2)=(7/5,3/5)\quad\text{（每组）},\qquad
 \Gamma_k^{\mathrm{mkt}}=11/36>0,\quad \pi_k=173/36.
 \end{gathered}
\]
同预算固定报酬合同产生零附加努力，组织者最优接纳一个双人组；其最优利润严格小于上述方案。
纯奖金合同不满足本申请者群体的逐状态现金保留条件，不能取得同一参与供给。
因此该合同是在所列市场、保留条件和菜单内由组织者选择的结果，而非全体市场必然采用的唯一合同。
\end{theorem}

\begin{proof}[证明（逆向选择与全体候选规模比较）]
现金保留条件要求 \(\beta\ge1/2\)，质量条件要求 \(r\ge\log(5/2)\)。
前一必要条件适用于混合与固定报酬合同，因为零赞助仍不得取消现金承诺；纯奖金则没有可行现金底线。其余候选合同的支出均不低于 \((1+c)n+n^2/10\)。由于
\(F(n,0)/n=10/(1+n)\) 随 \(n\) 下降而单位最低支出 \(1+c+n/10\) 上升，
\(n\ge5\) 不满足事前融资；可行接纳规模有限。

对 \(1\le n\le4\)，两名赞助者都有严格占优选择 \(I_k=1\)：若对手已投入，
本方投入相对不投入的利润增加至少 \((13/4)/3-7/8=5/24>0\)；若对手未投入，
增量至少 \((13/4)/2-7/8=3/4>0\)。本方不投入时其注意力份额为零，
故即使本方投入改变延续质量，上述下界仍成立。每家投入额为一，现金条件覆盖它。
在 \(S=2\)、\(\varepsilon\le1/10\) 下，低成本者努力效用增量至少 \(4/(5n)-1/10>0\)，
高成本者的增量至多 \(1/n+3/10-2<0\)。各组故有唯一严格均衡 \((1,0)\)，\(z=1\)。
基本任务不作附加努力也有正效用；完整参与均衡给两类正效用，身份虚报不能增加名额或改变信号。

在该行为域内，增大 \(\beta\) 或 \(r\) 不提高均衡质量而严格增加组织成本，故二者取最小可接受值。
组织者剩余成为
\[
 f(n)=\frac{12n}{1+n}-(1+c)n-\frac{n^2}{10}.
\]
\(f(2)-f(1)=7/10-c>0\)，\(f(3)-f(2)=-1/2-c<0\)，
并且 \(f''(n)=-24/(1+n)^3-1/5<0\)，故整数最优规模唯一为二。
无论更大规模是否改变努力或采用更高精度，其利润都不超过 \(f(n)\)，不能改善该选择。
固定报酬下努力增量为 \(h-k_i<0\)，因而 \(z=0\)，利润上界为
\(f_0(n)=10n/(1+n)-(1+c)n-n^2/10\)。
\(f_0(2)-f_0(1)=11/30-c<0\) 且严格凹，故固定报酬最优为一组；
\(f(2)-f_0(1)=17/10-c>0\)。纯奖金具有正概率零支付，违反现金保留条件。
\(\Omega^*>0\) 又排除不开办。代入 \(n=2,z=1\) 得 \(B=47/12\)，
\(\Gamma_k=B/3-1=11/36\)，\(\pi_k=4+P/8+\Gamma_k=173/36\)，其余数值由预算和期望支付式直接得到。
\end{proof}

\begin{remark}[同预算比较与准入边界]
在最优规模下，每组总支付为二。固定报酬为每人一，基本加差异的最低支付为二分之一，
纯奖金的最低支付为零。基本加差异同时保留现金底线与努力激励；固定报酬满足现金条件但不诱导本例附加努力；
纯奖金若允许外部保证保险，可能重新获得参与，此时须加入保证成本后再比较。
本例不证明基本加差异优于所有风险偏好、劳动合同或保险制度。
贡献提高 \(F\) 与 \(B\)，组织规模却仍因有限需求、基本份额与凸成本而有上限；
提高接纳数不能被当作无成本的社会收入扩张。未接纳者的保障须另由市场收入、公共服务或底线机制处理。
\end{remark}

\begin{proposition}[有误差合同下的逐状态现金延续与局部稳健性]
\label{prop:noisy-cycle-robustness}
在上述最优方案内，两家生产者现金状态取 \([4,5]^2\)，下一期可支配额度取本期净利润。
允许每家实际及同口径无营销反事实的净销售利润冲击绝对值各不超过 \(1/10\)，组织结算净收入冲击绝对值不超过 \(1/10\)，
且不改变事前已覆盖预算、需求响应和核验技术。则基本分配、组织偿付和下一期赞助现金逐状态可持续。
无冲击时现金一步达到 \((173/36,173/36)\)；在自由参数的充分小邻域内，最优合同类别、规模和努力仍保持。
\end{proposition}
\begin{proof}[证明]
信号误差只重新分配固定的 \(P=4\)，不改变总商品回流、组织总支付或赞助收益。
偏好营销净收益至少为 \(11/36-2/10=19/180>0\)，故正营销回报也在该冲击域内保持。
\(4<173/36-1/10<173/36+1/10<5\)，组织剩余亦大于 \(1/10\)，
故所列现金域前向不变且全部承诺由当期既有资金覆盖。无冲击现金更新为常数映射，局部导数为零。
赞助占优差、努力效用差、组织规模利润差和资金覆盖都有严格正余量，有限候选规模与连续性保持这些符号。
质量约束取等时 \(r^*=\log(\varepsilon_0/\bar\varepsilon)\) 随参数连续变化，
不能把它误写成所有不等式都严格；稳健性在该可行约束面及其邻近经济参数中成立。
\end{proof}

\section{自动化、资金机会成本与收入循环的总量闭合}
\label{sec:aggregate-income-closure}

\begin{definition}[两个收入群体与同口径资金用途]
\label{def:aggregate-income-transfer}
固定已扣除真实中间投入成本的总收入与产出计价域。令 \(Y_t\) 为当期总收入，
劳动及参与群体的初始收入份额为 \(w(A_t)\in(0,1)\)，资本及其他剩余索取者占 \(1-w(A_t)\)。
\(w'(A)<0\) 表示声明域内自动化对初始劳动收入份额的净替代效应；它是独立的任务与收入假设，
不由产能上限松弛定理推出。若新任务或其他收入抵销该效应，须使用实际 \(w'(A)\)。

令 \(R_t=r_tY_t\)、\(T_t=\tau_tY_t\) 分别为底线转移和参与制度带来的跨群体净收入转移，
满足 \(r_t,\tau_t\ge0\)、\(r_t+\tau_t\le1-w(A_t)\)。\(T_t\) 必须合并企业、组织、
核验机构和家庭账后计算；赞助、门票、组织利润及同一群体内部付款不能重复计入。
\(T_t\) 不是总赞助额或参与者总到账额，票款支出、资金提供者损失和规则改变的既有收入都须入账。
两群体的可支配收入为
\[
 Y_t^L=[w(A_t)+r_t+\tau_t]Y_t,\qquad
 Y_t^K=[1-w(A_t)-r_t-\tau_t]Y_t,
 \qquad Y_t^L+Y_t^K=Y_t.
\]
其消费倾向为 \(0<c_K<c_L<1\)。令 \(J>0\) 为有融资来源的计划投资或其他自主最终支出，
采用有库存调整的需求决定产出过程；产能上限 \(\bar Y\) 在声明邻域内不约束需求。
\(\Lambda_t\ge0\) 为相对于比较基线的净支出挤出或泄漏，真实成本已从收入域扣除，不在此再次扣除。
\[
 \begin{aligned}
 C_t&=c_LY_t^L+c_KY_t^K,\\
 Y_{t+1}&=\min\{\bar Y,\ J+C_t-\Lambda_t\}.
 \end{aligned}
\]
这是一个明确的总量需求模型，\(J\)、消费倾向及任务收入函数是其原始条件，不构成一般均衡或福利最优证明。
同一资金若改用于分红、其他营销或投资，应分别改变群体收入、\(\Lambda_t\) 或 \(J\)，不能只记录参与者消费增加。
\end{definition}

\begin{theorem}[含机会成本的需求增量、自动化补偿阈值与稳定固定点]
\label{thm:aggregate-demand-closure}
令 \(d_c=c_L-c_K>0\)，取常数 \(A,r,\tau\)，并令 \(\Lambda_t=\lambda Y_t\)、\(\lambda\ge0\)。定义
\[
 k=c_K+d_c[w(A)+r+\tau]-\lambda.
\]
若 \(0\le k<1\)、\(Y^*=J/(1-k)<\bar Y\)，则存在唯一内点固定点 \(Y^*\)，
无冲击收入轨迹局部渐近稳定。零泄漏时收入、消费与含库存的投资账直接闭合；非零泄漏另须有对应收款及融资账，不由收入固定点确认金融存量稳定。与同一 \(A,r,J\) 下
\(\tau=0,\lambda=0\) 的资金用途基线相比，固定收入上的净需求增量为
\[
 \Delta D=(d_c\tau-\lambda)Y.
\]
因此只有 \(d_c\tau>\lambda\) 时，该渠道提高净需求和稳定产出；若投资用途改变，则总差额另加 \(\Delta J\)。
对 \(A_1>A_0\)、\(w(A_1)<w(A_0)\)，在 \(r,J\) 不变且原基线无新增泄漏时，恢复原消费系数至少需要
\[
 \tau\ge w(A_0)-w(A_1)+\lambda/d_c,
\]
并同时满足出资群体预算、微观合同接纳能力及真实支付约束。
\end{theorem}
\begin{proof}[证明]
代入两群体收入，得 \(C_t=[c_K+d_c(w+r+\tau)]Y_t\)。转移同时减少出资群体收入，
故净增量为消费倾向之差乘转移额，而非 \(c_L T_t\)。与基线相减并扣除 \(\Lambda_t\)，即得需求差。
产能未约束时，更新为 \(Y_{t+1}=J+kY_t\)，故固定点为所列值，偏离满足
\(Y_{t+1}-Y^*=k(Y_t-Y^*)\)。收入加总恒等于 \(Y_t\)，固定点满足最终支出等于产出；
当 \(\lambda=0\) 时，家庭储蓄 \(Y^*-C^*\) 恰等于 \(J\)。非零泄漏的收款主体及融资余额另入账，
不能把未解释残差当作闭合。例如泄漏为净进口时，固定点储蓄为 \(J-\lambda Y^*\)，
国外净储蓄及本经济对外金融负债增加额均为 \(\lambda Y^*\)，于是国内外储蓄合计等于投资；
持续正净进口下该负债存量增长，收入稳定不等于全部存量稳定。未证明持续外部融资时，只确认已有融资覆盖期间。
比较自动化前后系数，条件为
\(d_c[w(A_1)-w(A_0)+\tau]-\lambda\ge0\)，移项得到阈值。
该阈值是规模要求；微观组织无法产生足额净转移时，本定理不确认补偿目标可达。
\end{proof}

\begin{remark}[非空总量参数域与超出能力的边界]
取 \(c_L=4/5,c_K=2/5,J=40,\lambda=0\)，自动化前 \(w_0=12/25\)，之后 \(w_1=2/5\)，
且 \(r=1/50\)。原消费系数为 \(3/5\)，固定点为一百；不增加参与净转移时，新系数为 \(71/125\)，
固定点降为 \(2500/27\)。若经完整组织与家庭账核验可实现 \(\tau=2/25\)，系数恢复 \(3/5\)，
两群体收入在固定点各为五十，消费合计六十、储蓄四十，恰好融资 \(J\)。
可取 \(\bar Y=120\)，状态区间 \([90,110]\) 在冲击 \(|\xi_t|\le4\) 下仍满足
\(Y_{t+1}=40+(3/5)Y_t+\xi_t\in[90,110]\)。
这给出收入更新自身的非空稳定参数域；\(\tau=2/25\) 的供给能力必须另外由净转移账与组织容量实现，
不能用两人或四人的局部实例自动填入。若容量上界 \(\bar T(Y)<(2/25)Y\)，宏观恢复结论即失败，
但可行组织的局部支付和盈利结论仍然保留。
\end{remark}

\begin{proposition}[微观合同与总量模型的连接条件]
\label{prop:micro-macro-compatibility}
在同一计价期，将有限个已通过预算、参与和核验检验的组织及出资人逐笔合并，
得到可实现的净转移集合 \(\mathcal T(Y,A)\)。若
\[
 \tau Y\in\mathcal T(Y,A),\qquad
 r+\tau\le1-w(A),\qquad
 \text{全部先期付款均有未重复承诺的现金覆盖},
\]
并且相应总量消费、产能和动态条件成立，则微观合同与总量收入循环可以联合确认。
否则只能分别确认已证明的局部结论，不把总量阈值作为微观已实现的事实。
\end{proposition}
\begin{proof}[证明]
合并时，每笔内部付款在出资者与收款者账上抵销；跨群体净额定义为 \(T=\tau Y\)，
因此恰好进入两群体收入式且不创造额外总收入。逐组织预算保持基本与差异支付，
出资预算保持资本群体非负收入，现金覆盖排除同一资金的同时重复使用。
将该净额代入总量更新即可应用前一定理。若集合成员资格或任何一项预算失败，
则不存在所声称的共同会计状态，不能通过联合结论。
\end{proof}

\begin{theorem}[微观合同与总量收入的一个共同可行实例]
\label{thm:joint-micro-macro-witness}
令 \(g=4+2c+2/5\)，其中 \(c\) 如定理~\ref{thm:endogenous-organization-contract}，
故 \(5<g<6\)。取两个观众与赞助品牌互不交叉的赛事组织，第一个使用该定理的计价，第二个将全部货币量、
努力成本、效用与现金保留额按同一正比例 \(\theta=8/g-1\in(0,1)\) 缩放。
赞助包、门票、生产者产品销量及核验成本同时缩放，产品单位价格与单位成本不变。
第二组织按赞助包数量形成注意力：写 \(I'_k=\theta I_k\)，则
\[
 a'_k=\frac{I'_k}{\theta+I'_1+I'_2}=a_k,\qquad
 a'_0=\frac{\theta}{\theta+I'_1+I'_2}=a_0.
\]
因此 \(\pi'_k=\theta\pi_k\)、\(\Omega'=\theta\Omega\)、\(U'_i=\theta U_i\)，
全部选择不等式、最优规模及信号误差概率保持。不能只缩放货币投入而保留错误的注意力尺度。

在以下共同核算条件下，两组织合计实现每期净转移 \(T=8\)：组织成本全部是劳动群体的真实服务报酬，
参与者属于劳动群体，组织剩余与生产者利润属于资本群体；观众票款来自资本群体，
其被替代支出原本全部形成资本群体净收入。品牌之间及外部自动化产品之间的份额转移都在资本群体内部核算。
真实资源成本与部门增加值须先统一，不能把消费付款直接当作收入转移。

总量取 \(w_1=2/5,r=1/50,c_L=4/5,c_K=2/5,J=40\)，收入在 \([90,110]\) 内变化，
无额外泄漏。则联合更新为
\[
 Y_{t+1}=\frac{216}{5}+\frac{71}{125}Y_t,
\]
唯一收入固定点为一百，恢复 \(w_0=12/25\)、无参与净转移时的原固定点。
两组织的随机支付、正营销收益、赞助现金、观众预算和总量资源账在同一参数点同时可行。
\end{theorem}
\begin{proof}[证明（缩放、机会成本抵销与资源核对）]
单个原单位组织向劳动群体支付 \(P+C=g\)。资本群体取得组织剩余，但同时承担赞助及被替代的消费收入损失，
合并净变动为 \(\Omega-S-F=-g\)。这一步使用所列消费替代条件；若替代部门也支付劳动收入，
须扣除该劳动收入损失，不能仍使用 \(g\)。两组织净转移为 \((1+\theta)g=8\)，不存在重复计算。
缩放后劳动现金底线和努力效用全部同比例变化，故两个微观均衡保持；两组织参与人数合计八人，
核验服务人员另列，其收入计入成本收款账而不重复算作参与支付。

总量可支配收入为 \(Y^L=(21/50)Y+8\)、\(Y^K=(29/50)Y-8\)，
均为正且和为 \(Y\)。消费为 \((71/125)Y+16/5\)，加 \(J\) 得所列更新。
\(43.2+(71/125)100=100\)，斜率 \(71/125<1\)。
两组织票款总额为 \(64/g<64/5\)，低于整个收入区间资本群体消费的最小值 \(442/25\)。
每个原单位组织的两品牌消费合计为 \(209/9\)，外部商品的注意力预算余额为 \(59/18\)，
两者合计 \(53/2\)。将外部商品也完整计入，两组织对应商品及赛事消费总额为
\[
 (1+\theta)\left(\frac{53}{2}+8\right)=\frac{276}{g}<\frac{1358}{25}.
\]
最后一个严格不等式由 \(g>5.2\) 得到，右边正是总消费的最小值。其中参与者按收入一半购买指定产品，
指定商品可以包含必要消费；群体消费倾向不要求每个人有相同边际消费倾向。
先从资本消费中扣票款、从劳动消费中扣参与者指定购买，两个余额均为正，且余额之和足以覆盖其余固定背景需求；
故可以在两群体间分配背景购买，余款购买其他商品。两群体储蓄均为正；生产者须留存的现金合计为 \((8/g)(173/18)<15\)，
低于资本群体储蓄的区间最小值 \(663/25\)，因此留存与家庭消费不重复占用收入。背景购买者允许同时来自两个群体，
不得把全部品牌购买强加给单一群体。生产产能与先期现金留有余量，实物商品采用订单预付款覆盖当期可变成本，
赞助只使用另行预留的已实现利润，故没有把同一现金同时抵押给生产与赞助。
可将每家原单位生产者的银行余额取为 \(K_t=5\)，与次期赞助额度 \(L_t\) 区分。
令分红 \(d_t=\pi_t\)，则 \(K_{t+1}=K_t+\pi_t-d_t=5\)，始终覆盖 \(L_t\le5\)；
将本期利润预留为下一期额度时，相同金额从以前期间未承诺的资本中释放。
缩放组织的各银行余额与分红同比例缩放。额度更新不是额外现金收入，不与家庭分红重复登记。

为核对过渡期而非只核对固定点，令库存为 \(V_t\)，
\(V_{t+1}-V_t=Y_t-J-C_t=Y_t-Y_{t+1}\)，取 \(V_0\ge20\)。
区间不变使库存始终非负；家庭储蓄 \(Y_t-C_t\) 等于计划投资 \(J\) 加库存投资，
故每一期储蓄与投资都一致。库存与收入之和守恒，收入收敛后库存也在该守恒面上收敛，
不宣称不同初始库存都收敛到同一个库存水平。允许自主支出冲击 \(|\xi_t|\le4\) 时，
同一收入区间仍不变，库存式相应使用 \(J_t=40+\xi_t\)；持续冲击下只确认有界延续。
\end{proof}

\section{模型结论的分层与可识别试验}
\label{sec:identified-validation}

\begin{definition}[四类结论与独立适用域]
会计结论检验金额加总与资金用途；履约结论检验公开合同、资格、证据与实际支付；
行为结论检验在明确偏好、成本与信息条件下的选择；总量与福利结论另检验净转移、替代用途、规模、外部性与动态。
规则内复算不能替代真实信号准确性，私人盈利不能替代净福利，编译与有限数值样本不能替代全域证明。
适用域指标、成本归属、退出与缺失样本处理、检验窗口和失败阈值须在观察结果前固定，
不得在出现反向结果后增加排除条件。事后探索可以生成新假设，但须在独立样本重新检验。
\end{definition}

\begin{proposition}[收入与偏好通道的可识别有序分解]
\label{prop:factorial-channel-identification}
以无跨组干扰的组织或地区为随机化单位，对额外参与收入 \(Z_I\) 与品牌接触 \(Z_P\)
作独立的二乘二随机分配，四组均具有正概率且保持既有合同权利。潜在结果
\(Y(a,b)\) 可取同口径、已扣产品成本的销售利润。若处理执行一致，失访可忽略或已按预注册机制处理，
且组间不存在未计入的价格、曝光或消费溢出，则令 \(\mu_{ab}=\EE[Y(a,b)]\)，有
\[
 \Delta_{\mathrm{inc}}=\mu_{10}-\mu_{00},\qquad
 \Delta_{\mathrm{pref}\mid\mathrm{inc}}=\mu_{11}-\mu_{10},\qquad
 \mu_{11}-\mu_{00}=\Delta_{\mathrm{inc}}+\Delta_{\mathrm{pref}\mid\mathrm{inc}}.
\]
随机分组的均值差无偏识别相应平均处理效应，交互项按已声明的顺序归入后一项。
偏好营销净收益还须从第二项扣除对应窗口全部增量营销成本；试验收入处理及复核经费另列完整经营账，
不能因其用于识别而隐去。
\end{proposition}
\begin{proof}[证明]
随机分配独立于四组潜在结果，一致性给出
\(\EE[Y\mid Z_I=a,Z_P=b]=\EE[Y(a,b)]\)。取相应均值差即得识别；两差相加的中间项抵销。
没有随机分配时，该等式一般不成立；同期销量增加或品牌选择相关不能代替它。
有限样本只给估计与不确定性区间，正点估计不直接确认正净营销收益。
\end{proof}

\begin{remark}[观察性研究与停止条件]
采用规则变更的双重差分时，必须另声明未处理潜在结果的共同趋势、无提前反应、样本构成稳定、
无未处理溢出及处理时点可比；多期分批变更应按处理批次与事件时间估计，不能仅凭双向固定效应式确认因果。
参与收入、基本份额兑现、错误排序、申诉时间、净营销收益、普通参与者多期净收益与退出分别报告。
停止条件包括资金覆盖失败、既有权利受损、误判超过预定上界或完整成本后的收益下界不支持继续。
这些是待实施的识别设计与验收标准，文稿没有据此登记已完成的市场试验、估计值或外部有效性。
\end{remark}

\section{市场参与经济总定理}

\begin{remark}[联合结论的量词与模型兼容]
以下总定理将已经列明前提的局部结论合取；它们不把不同主体、不同反事实或互不兼容参数的结果拼接为一个均衡。
核心微观与总量联合实例见定理~\ref{thm:joint-micro-macro-witness}。
文化生成、金融安排、规范全员底线及附录拓扑比较各有额外条件，未启用的分支不确认其结论。
带误差信号只确认其误差域内的排序与激励，不以独立复算一致冒充真实贡献绝对正确。
\end{remark}

\begin{theorem}[市场参与循环扩展总定理]
\label{thm:market-participation-synthesis}
固定一个声明经济域。若分别满足：
\begin{enumerate}[label=\textnormal{M\arabic*:},leftmargin=3.5em]
  \item 存在持续正物质需要，并把宏观生产--消费内核与单项活动的经济价值分开；同一必要需要至少有两个
  可行载体且由消费者分散选择，因而已实现配置形成选择份额投射；经济价值按净社会剩余的基线差定义，物质产出、
  价格与支付作为不同坐标；
  \item 构成性制度算子通过 \(\mathfrak C_{\mathrm{inst}}\)，市场主体在合法路径空间 \(\Omega_\ell\) 内分散选择，
  \(\Omega_\ell\) 含有多条路径，监管算子不指定唯一产量、价格或个体配置；
  \item 生产盈利对消费、投资、参与--分配和其他需求通道的边际传导具有可识别口径，停滞风险只在联合证据
  通过时确认；
  \item 竞技及其支持网络构成注意力组织核心，至少两个生产者以已实现利润支持赞助，扣除完整营销成本；
  采用定理~\ref{thm:competitive-attention-injection} 的连续固定收入模型、
  定理~\ref{thm:constructive-equilibrium-cycle} 的连续联合模型，或
  定理~\ref{thm:endogenous-organization-contract} 的有限赞助包模型，并满足所选分支的正均衡与正偏好回报条件；
  \item 真实注意力能够形成消费者对投入方生产者商品的偏好指引，参与收入具有正边际消费倾向，两条通道
  均有时间顺序和反事实证据；严格盈利边际传导结论另须满足定理~\ref{thm:dual-channel-bridge} 的可微性与严格投入响应前提，离散投入或局部现金约束不绑定时只确认已证明的水平循环；
  \item 围棋行业参与的效用、成本、外部损害、生产侧赞助、真实注意力、产品偏好与收入分配具有可比较口径，
  且规则层和赛事参与--分配分别通过固定规则重放和独立复算核验记录；
  \item 若另外比较生产市场循环与围棋行业循环的拓扑结构，则提交有向同伦核验记录
  \(\mathfrak C_{dH}\)；不作该附加比较时，本项标为不适用且不确认同伦结论；
  \item 非传统按劳的参与--分配机制通过规则内重放、独立复算一致、横向一致、贡献单调、预算、权限、竞争与申诉组成的
  \(\mathfrak C_{\mathrm{RPF}}\)，并以 \(\mathfrak C_{\mathrm{part}}^+\) 保证每名有效参与者的正基本份额及
  可比贡献的严格差异分配；若声称该机制可大规模运行，还须把生产状态与治理自动化状态分开，并通过
  单位治理成本、收入占比的界及所需极限条件的成本核验记录；
  \item 合同剩余索取给出固定补偿、真实成本、能力和风险、正余额、剩余份额与支付，并通过
  \(\mathfrak C_{\mathrm{NE}}\) 验证自愿合同、公开规则、竞争、可比价格和真实退出；它与非市场强制占有及
  规范治理的相对公正核验记录分开；正投入另满足命题~\ref{prop:residual-share-positive-effort}
  的充分条件，或提交其他具有参与约束的正投入均衡；
  \item 组织每一条经济相关活动路径及其跨活动权利映射通过来源可追踪的逐活动核验记录；
  \item 候选公共法律域有限非空，并以风险外溢阈值筛选空间释放极大的安全合法域；
  \item 在市场收入使生存、需求或参与跌破非退化阈值的声明域内，坐标最小充分救济通过底线、资金、储备、
  正剩余激励、需求覆盖、市场厚度与动态稳定核验记录。
\end{enumerate}
则可以同时确认以下有限结论：
\[
  \operatorname{SustainedMaterialNeed}
  \Longrightarrow
  \operatorname{ProductionConsumptionCore},
\]
\[
  \bar q_{hg}>0,
  \quad |\mathcal C_{hg}^{F}|\ge2
  \Longrightarrow
  \boldsymbol\vartheta_{hg}^{M}\in
  \Delta(\mathcal C_{hg}^{F}),
  \quad \max_c\vartheta_{hgc}^{M}>0,
\]
\[
  q_M>0\not\Longrightarrow\nu>0,
  \qquad
  \nu>0\not\Longrightarrow q_M>0,
\]
\[
  \nu_i>0\not\Longrightarrow T_i>0,
  \qquad
  T_i>0\not\Longrightarrow\nu_i>0,
\]
\[
  X_P\simeq_dX_G
  \Longrightarrow
  U(X_P)\simeq U(X_G).
\]
并且在所列局部导数条件下，
\[
  \Pi^P\to\mathbf I^{B*}\to\mathbf a^{\mathrm{att}}
  \to\boldsymbol\vartheta^P\to D^P\to\Pi^{P\prime},
  \qquad
  \mathbf I^{B*}\to\mathbf T^B\to C^P\to D^P
\]
形成注意力与收入双通道桥接；产品回款形成销售利润，模型中的偏好增量销售利润覆盖完整营销成本并形成
\(\Gamma_k^{\mathrm{mkt}}>0\)。有效参与者有 \(T_{it}^B\ge\beta_t^{\mathrm{base}}>0\)，同一可比域内
\(c_{it}>c_{jt}\) 推出 \(T_{it}^B>T_{jt}^B\)。实际营销回报与下一期现金预算经闭环核验记录验收后，
已实现净利润继续支持下一期生产与赞助。若偏好传导乘积为零且无其他可货币化回报，均衡投入退化为零，不能声称存在
市场内生注资。规范治理先要求参与--分配在声明规则内可重放、复算一致且在声明容限内相对公正，再由公共治理选择安全极大合法域，
充分释放可行竞争与试验空间并把已识别风险外溢限制在阈值内，但仍不唯一决定市场路径。围棋参与只有在边际
效用增量超过成本与外部损害时才确认正经济价值，其赛事分配另须通过围棋治理核验记录；组织整体结论要求所有
经济相关活动路径覆盖。正余额的差异化索取构成合同剩余索取，并在本定理域内承担非退化投入激励；治理既不要求
它归零，也不把非市场强制占有伪装成自由市场结果。充分救济只把全体成员抬升到生存--接入底线，并在资金与
稳定条件下恢复需求和市场厚度，不取消下限以上的能者多得、利润和贫富差距。
若定理~\ref{thm:rpf-unit-cost-collapse} 的独立治理自动化前件同时成立，则相对公正判定的单位成本与收入占比
满足相应上极限界；只有另具残余单位成本与人均收入的极限条件时才得到收敛。该结论不由生产产能松弛自动推出，也不要求总治理成本下降。
\end{theorem}

\begin{Certificate}[十二项扩展条件与分层结论]
定义
\[
  \mathfrak C_{\mathrm{market}}
  =
  \left\langle
  \begin{gathered}
  \mathfrak C_{\mathrm{value}},
  \mathfrak C_{\mathrm{projection}},
  \mathfrak C_{\mathrm{inst}},
  \mathfrak C_{\mathrm{stagnation}},
  \mathfrak C_{\mathrm{attention}},
  \mathfrak C_{\mathrm{part}}^+,
  \mathfrak C_{\mathrm{dual}},
  \mathfrak C_{\mathrm{go}},\\
  \mathfrak C_{dH},
  \mathfrak C_{\mathrm{RPF}},
  \mathfrak C_{\mathrm{RPFcost}},
  \mathfrak C_{\mathrm{NE}},
  \mathfrak C_{\mathrm{activities}},
  \mathfrak C_{\mathrm{safe}},
  \mathfrak C_{\mathrm{relief}}
  \end{gathered}
  \right\rangle.
\]
核验记录允许逐项确认：任一坐标失败，不撤销其他已经证明的局部结论，但不得确认依赖该坐标的联合结论。
\end{Certificate}

\begin{proof}[证明（既有定理的条件合取）]
M1 与定理~\ref{thm:material-core-compatibility}、定理~\ref{thm:material-production-neither}、定理
~\ref{thm:essential-multichoice-projection}、推论~\ref{cor:projection-versus-attention-shift}、定理
~\ref{thm:value-price-payment-separation} 分别给出宏观物质内核和微观非蕴含关系。M2 与命题
~\ref{prop:market-regulation-nonsubstitution}、命题~\ref{prop:constitutive-governance-market-compatible} 给出制度构造、
监管约束与实际配置的非唯一性。M3 与命题
~\ref{prop:core-bridge-stagnation-risk} 限制停滞风险。M4--M5 与定理
~\ref{thm:competitive-attention-injection} 或定理~\ref{thm:constructive-equilibrium-cycle} 的相应均衡分支、
命题~\ref{prop:no-guidance-no-market-injection} 和定理
~\ref{thm:dual-channel-bridge} 给出竞争性注资、营销净增利润及其退化边界；命题
~\ref{prop:sales-marketing-accounting} 分离产品回款与营销回报。M6 与命题
~\ref{prop:go-participation-value}、命题~\ref{prop:go-attention-bridge}、命题
~\ref{prop:go-governance-readiness} 给出围棋参与价值、桥接和治理就绪条件。在 M7 启用并通过时，由定理
~\ref{thm:production-go-directed-homotopy} 给出严格有向同伦，推论
~\ref{cor:ordinary-homotopy-projection} 给出普通同伦投影。M8 与定理
~\ref{thm:relative-participation-fairness} 及定理~\ref{thm:base-differential-coverage} 给出相对公正、正基本
份额与严格贡献差异；命题
~\ref{prop:capacity-not-governance-cost} 与定理~\ref{thm:rpf-unit-cost-collapse} 分别排除从生产状态直接推出治理
成本，并给出可大规模运行时的成本条件。M9 与命题
~\ref{prop:appropriation-exploitation-separation}、定理~\ref{thm:residual-claim-incentive-necessity} 和命题
~\ref{prop:market-surplus-governance-separation} 给出合同剩余索取的激励必要性，并把它与非市场强制占有和治理
目标分离；正投入的充分性另由命题~\ref{prop:residual-share-positive-effort} 或相应完整均衡证明给出。
M10 与定理~\ref{thm:multi-activity-governance} 给出逐活动覆盖。M11 与定理
~\ref{thm:maximal-safe-governance-space}、定理~\ref{thm:dual-layer-governance} 给出安全极大合法域与双层治理。
M12 与定理~\ref{thm:minimal-sufficient-relief}、定理~\ref{thm:relief-endogenous-stability} 和命题
~\ref{prop:relief-not-equalization} 给出全员底线承诺下的救济下界、融资和非平均化边界。各结论依赖的条件已经
逐项列明，其中 M7 只在附加结构比较启用且其核验记录通过时确认同伦结论，其余经济结论不依赖该分支。
故合取后得到声明中的有限结论，而不产生市场自动最优、所有参与有价值、所有利润都来自强制占有
或救济等于计划分配等额外结论。
证毕。
\end{proof}

\begin{theorem}[后 AI 市场参与经济条件总定理]
\label{thm:governance-economy-synthesis}
固定产品域 \(\Goods_0\subseteq\Goods\)。若同时满足：
\begin{enumerate}[label=\textnormal{C\arabic*:},leftmargin=3.5em]
  \item 后 AI 精确适用域核验记录通过，即
  \(\PostAI^{\mathrm{exact}}(\Goods_0,A,r)\) 成立，等价地，对每个
  \(j\in\Goods_0\)，\(\CapacitySlack(j,A,r)\) 成立；
  \item 有效需求或互补稀缺要素使 \(q_j^*(r)\) 有限；
  \item 持续正物质需要使生产--消费构成宏观内核；同一必要需要具有至少两个预算可行且由消费者分散选择的
  载体，已实现配置通过选择份额投射核验记录；并对盈利进入消费、生产投资、参与--分配行业和其他需求通道的边际
  传导分别估计；
  \item 分散市场回应形成可分配剩余 \(Y^*>0\)，竞技组织依基本与差异分配合同配置其内部收入池，满足
  定理~\ref{thm:base-differential-coverage} 的接纳、资金覆盖与严格贡献比较条件；
  \item 至少两个生产者的竞争性广告赞助满足定理~\ref{thm:competitive-attention-injection}
  的固定收入正均衡条件，或定理~\ref{thm:constructive-equilibrium-cycle}、
  定理~\ref{thm:endogenous-organization-contract} 各自的收入内生联合均衡条件，
  完整营销成本由偏好增量销售利润覆盖并形成净增利润；
  \item 注意力能够形成消费者对投入方生产品的偏好指引，参与收入能够转化为生产品消费，并通过双通道核验记录；
  若确认严格盈利边际传导，还须满足定理~\ref{thm:dual-channel-bridge} 的全部微分与严格投入响应条件，否则该导数结论不适用；
  \item 每名完成约定的有效参与者均按期取得基本与差异分配，提交并通过定理~\ref{thm:typed-atomic-execution-chain}
  的同一对象、同一规则版本、连续来源引用逐项履约核验记录及实际支付最终性回执；
  \item 参与--分配通过规则内重放、独立复算一致、横向一致、贡献单调、预算、权限、竞争、审计与申诉的相对公正核验记录；生产
  状态与治理自动化状态分别记录，规模化实施另通过固定、标准化与残余成本分解及单位成本门槛；
  \item 若当事人选择二级转让，则转移守恒、受合同限制及流动性上界约束；未选择时直接结算初始收益权；
  \item 若采用金融工具配置部分资源请求权，则比例 \(\theta\in[0,1]\) 及其比较基线事前声明；
  比例上升或趋近一仅是可选金融分支的附加比较条件；
  \item 合同履行与项目制度公正条件被分别审计；
  \item 组织所有经济相关活动及跨活动权利映射通过逐路径规范化核验记录；
  \item 构成性制度算子先提供可执行且非单点的市场路径；该路径相对于声明的现行制度基线和已证失灵点通过
  最小兼容适用路径核验记录，公共治理再在风险外溢阈值内选择空间释放极大的非单点安全合法域；
  \item 动态稳定、一般均衡、福利与经验结论只在各自核验记录通过时确认；
  \item 在无救济状态跌破生存、需求或参与阈值时，坐标最小充分救济具有可持续资金、储备和极端损失顺序，
  且缴款后仍保留正剩余投入激励；
  \item 至少一个资本、组织、风险或稀缺能力投入者通过竞争性市场中的剩余索取核验记录，并满足
  命题~\ref{prop:residual-share-positive-effort} 的充分条件或其他完整的正投入均衡条件，实际取得正剩余收入。
\end{enumerate}
则可以推出：
\[
  \partial_A Y_j=0
  \quad(j\in\Goods_0)
\]
在精确松弛域成立；
\[
  \bar q_{hg}>0,
  \quad |\mathcal C_{hg}^{F}|\ge2
  \Longrightarrow
  \boldsymbol\vartheta_{hg}^{M}\in
  \Delta(\mathcal C_{hg}^{F}),
  \quad \max_c\vartheta_{hgc}^{M}>0;
\]
\[
  \Delta y_i
  =
  Y^*\bigl[s_i(G_2)-s_i(G_1)\bigr]
\]
在固定总剩余比较静态下成立。
核心参与合同还满足
\[
  T_{it}^B\ge\beta_t^{\mathrm{base}}>0,\qquad
  c_{it}>c_{jt}\Longrightarrow T_{it}^B>T_{jt}^B.
\]
结合 C15 中“市场收入已经覆盖底线或由坐标最小充分救济补齐”的分支，三项核心功能写为
\[
  \boxed{\mathbf g^{\mathrm{core}}
  =\bigl(g_{\mathrm{part}},g_{\mathrm{diff}},g_{\mathrm{relief}}\bigr)
  =(1,1,1)}.
\]
C5--C6 的注意力偏好回流提供竞争性资金与演化动力；C7--C8 的原子执行链、相对公正与成本核验提供实施
基础。两组条件分别回答“为何可能扩张”和“如何稳定执行”，不取代三项核心功能回答的分配、激励与接入问题。
\[
  \Abs{y_i-Db_i^0}\le D\lambda_i+f_i
\]
仅在所选择的受限二级转让分支中成立。若另采用金融工具，在固定基线与规则比较对象、令其配置比例趋近一的比较域内，有
\[
  \Delta c_i=\theta\Delta y_i(G_F)
  \longrightarrow
  \Delta y_i(G_F)
  \qquad(\theta\to1).
\]
若另通过 C6 所列可微性、非负通道及至少一项严格投入响应条件，由链式法则还成立
\[
  \frac{\partial D_{t+1}^P}{\partial\Pi_t^P}>0,
\]
从而生产盈利对参与--分配与下一期消费恢复正边际贡献。命题~\ref{prop:sales-marketing-accounting} 给出
\(\Pi_k^{P,\mathrm{net}}=\Pi_k^{P,0}+\Delta\Pi_k^{P,\mathrm{inc}}+\Gamma_k^{\mathrm{mkt}}\)，
产品回款与营销回报分账闭合。若另满足定理~\ref{thm:interest-luxury-engel} 与
定理~\ref{thm:interest-sector-macro-support} 的附加条件，还可确认
\[
  \eta_{hj}^{Y}>1,
  \qquad
  D_{t+1}^{P,I}\ge D_{t+1}^{\dagger},
\]
即兴趣消费在声明收入域内具有广义奢侈品属性，并在规模、净额与稳定性条件通过时填补声明需求缺口。因此，
在 C15 的适用域内还成立
\[
  R_{ht}^{*}=[F_{ht}-Y_{ht}^{m}]_+,
  \qquad
  Y_{ht}^{m}+R_{ht}^{*}\ge F_{ht},
\]
并在需求和市场厚度阈值上有
\[
  D_{t+1}^{P,0}+\Delta D_{t+1}^{R}
  +\Gamma_{t+1}^{I,\mathrm{net}}
  \ge D_{t+1}^{\dagger},
  \qquad
  N_{t+1}^{0}+\Delta N_{t+1}^{R}\ge\bar N_{t+1}.
\]
在 C16 的非退化内点投入域内还成立
\[
  \rho_i
  \frac{\partial\EE[R^B]}{\partial e_i}
  =c_i'(e_i)>0,
\]
故合同剩余索取承担自由市场投入激励，充分救济只承担生存、接入和循环底盘，二者不相互替代。因此，
在指定产品域内，复制产能的边际约束减弱时，
生产--消费仍是物质内核；市场通过分散回应形成价格、交易和剩余；多个生产者通过竞争性注意力投入形成对
生产商品的消费偏好指引，参与收入形成购买力回流；组织合同对内部资源请求权分配的边际作用上升。规范治理
先以构成性制度供给使新型参与取得非单点的可执行市场路径，但不指定实际市场结果；其内部关键是参与--分配
在声明规则内可重放、复算一致并在声明标准内相对公正，组织须逐一覆盖所有经济相关活动；公共治理的
外部边界是充分释放不突破风险外溢阈值的合法竞争与试验空间，而不指定具体市场结果、消灭合同剩余索取或抹平
贫富差距；充分救济以最小截断映射保障社会底线，却不规定底线以上配置。相对于已经通过核验记录的现行制度基线与
具体失灵点，这些新增制度环节构成候选族内的最小兼容扩展；它不否定按需分配、基本收入、公共供给、按劳分配或
其他制度，也不声称自身唯一或全局福利最优。
在 C8 的成本条件下，单位治理成本与收入占比具有相应上极限界；另有极限条件时才确认收敛，实际规模可负担性还须检查收入扣除治理成本后的支付预算，
但不推出总治理成本下降，也不把生产产能核验记录当作治理成本核验记录。
金融规则仅在实际采用金融工具且 \(\theta\to1\) 的声明域内趋近资源请求权分配规则；
直接支付基本与差异份额不以该可选分支成立为前提。
\end{theorem}

\begin{Certificate}[十六条件合取与分层结论投影]
定义总核验记录
\[
  \mathfrak C_{\mathrm{GEcon}}
  =
  \left\langle
  \begin{gathered}
  \mathbf g_C,
  \mathfrak C_{\mathrm{post}}^{\mathrm{exact}},
  \mathfrak C_{\mathrm{MCP}},
  \mathfrak C_{\mathrm{projection}},
  \mathfrak C_{\Delta G},
  \mathfrak C_{\mathrm{part}}^+,
  \mathfrak C_{\mathrm{RPF}},
  \mathfrak C_{\mathrm{RPFcost}},
  \mathfrak C_{\mathrm{liq}},
  \mathfrak C_{\theta},\\
  \mathbf g_{\mathrm{fair}},
  \mathfrak C_{\mathrm{activities}},
  \mathfrak C_{\mathrm{inst}},
  \mathfrak C_{\mathrm{safe}},
  \mathfrak C_{\mathrm{dyn}},
  \mathfrak C_{\mathrm{emp}},
  \mathfrak C_{\mathrm{relief}},
  \mathfrak C_{\mathrm{NE}}
  \end{gathered}
  \right\rangle.
\]
其中 \(\mathbf g_C=(g_{C1},\ldots,g_{C16})\)。只有
\[
  \min_k g_{Ck}=1
\]
时，总定理状态为 \(\operatorname{ConditionallyEstablished}\)；否则只保留已经通过的局部定理。
对于 C9、C10，另保存业务适用标记 \(a_j\in\{0,1\}\) 及其事实依据，令 \(v_j=1\) 当且仅当该适用性判定通过核验。
若业务未启用，须由合同与实际交易记录
共同证明，并令该分支审计值 \(c_j=0\) 作为占位；若业务已启用，则 \(c_j=1\) 当且仅当相应业务核验记录通过。
合取条件取
\[
  g_{Cj}=v_j(1-a_j+a_jc_j),\qquad j\in\{9,10\}.
\]
不适用分支只证明其条件义务未被触发，不确认二级转让或金融趋近结论，也不阻断基本份额与直接支付路径。
总核验记录中的 \(\mathfrak C_{\mathrm{liq}}\)、\(\mathfrak C_{\theta}\) 在未启用时保存经核验的不适用记录；
金融极限结论还须另外证明所声明的 \(\theta\to1\) 比较条件。
\end{Certificate}

\begin{remark}[渐近适用、部分适用与总定理降级]
定理~\ref{thm:governance-economy-synthesis} 的精确等式
\(\partial_A Y_j=0\) 使用 C1 的精确适用前件。若仅有
\(\PostAI_\epsilon(\Goods_0,A,r)\) 而非精确适用，则该结论降级为各产品的预注册容限界，并须把误差沿后续
比较静态和动态方程传播；不得保留精确零导数表述。若
\(0<\kappa_{\mathrm{post}}^\epsilon<1\)，总定理只能投影到
\(\Goods_{\mathrm{post}}^\epsilon\)，不得对整个 \(\Goods_0\) 确认十六条件合取结论。若
\(\kappa_{\mathrm{post}}^\epsilon=0\)，后 AI 总定理不适用，但不依赖产能松弛的局部制度定理仍可分别保留。
\end{remark}

\begin{proof}[证明]
由 C1--C2 与定理~\ref{thm:capacity-marginal-relaxation}，
\[
  \CapacitySlack
  \Longrightarrow
  Y_j(A,r)=S_j(q_j^*;r)
  \Longrightarrow
  \partial_A Y_j=0.
\]
由 C3 与定理~\ref{thm:material-core-compatibility}、定理~\ref{thm:essential-multichoice-projection}、推论
~\ref{cor:projection-versus-attention-shift}、命题~\ref{prop:core-bridge-stagnation-risk}，生产--消费的宏观
必要性、刚需多选下的选择份额投射与盈利边际传导分别成立。由 C4 与定理~\ref{thm:governance-comparative-statics}，
\[
  Y^*\text{ 固定}
  \wedge G_1\to G_2
  \Longrightarrow
  \Delta y_i=Y^*\Delta s_i.
\]
同一 C4 的资金与合同条件经定理~\ref{thm:base-differential-coverage} 给出全部有效参与者的正基本份额，
以及可比贡献严格有别时的严格差异分配。
在 C6 的严格响应分支内，由 C5 的相应均衡定理与定理~\ref{thm:dual-channel-bridge}，
\[
  \Pi^P\to\mathbf I^{B*}\to
  (\boldsymbol\vartheta^P,\mathbf T^B)
  \to D^P,
  \qquad
  \frac{\partial D_{t+1}^P}{\partial\Pi_t^P}>0.
\]
其中，固定消费者收入分支的营销净增利润由定理~\ref{thm:competitive-attention-injection} 的最优回应与零投入收益比较给出；
收入随分配变化的显式分支，则由定理~\ref{thm:constructive-equilibrium-cycle} 或定理~\ref{thm:endogenous-organization-contract} 各自证明偏好营销净收益的正下界。
产品成本、营销成本及净利润的闭合由命题~\ref{prop:sales-marketing-accounting} 给出。
若再通过兴趣消费的最低需要、Engel 弹性、进入阈值、宏观规模与稳定性附加条件，则由定理
~\ref{thm:interest-luxury-engel} 和定理~\ref{thm:interest-sector-macro-support} 直接得到
\(\eta_{hj}^{Y}>1\) 与 \(D_{t+1}^{P,I}\ge D_{t+1}^{\dagger}\)；它们不是 C1--C14 自动推出的无条件结论。
由 C7，每个声称完成参与支付的主体都必须提交定理~\ref{thm:typed-atomic-execution-chain} 的同一见证核验记录；因而
\[
  \mathfrak C_i^{\mathrm{atom}}=1
  \Longrightarrow
  \Eligible_i\wedge\Measured_i\wedge\Verified_i
  \wedge\Valued_i\wedge\Entitled_i\wedge\Paid_i.
\]
由 C8 与定理~\ref{thm:relative-participation-fairness}，内部参与--分配只能在规则内重放、独立复算一致、横向一致、贡献单调和
纠错条件全部通过时确认相对公正；命题~\ref{prop:capacity-not-governance-cost} 阻止从 C1 推出治理成本下降，
定理~\ref{thm:rpf-unit-cost-collapse} 再给出单位成本、收入占比的界及收敛所需的独立条件。
在二级转让实际启用时，由 C9 与定理
~\ref{thm:limited-liquidity-primary-dominance}，
\[
  \Abs{\Delta b_i}\le\lambda_i
  \Longrightarrow
  \Abs{y_i-Db_i^0}\le D\lambda_i+f_i.
\]
在金融工具实际启用、固定基线与规则比较对象且 \(\theta\to1\) 时，由 C10 与定理~\ref{thm:financial-rule-convergence}，
\[
  c_i=(1-\theta)r_i+\theta y_i(G_F)
  \Longrightarrow
  \Delta c_i=\theta\Delta y_i(G_F)
  \xrightarrow{\theta\to1}
  \Delta y_i(G_F).
\]
未启用上述业务时，C9、C10 以经核验的不适用记录通过条件义务检查，初始确权和直接支付结论由 C4、C7 独立给出。
C11 阻止把合同履行冒充公正；C12 与定理~\ref{thm:multi-activity-governance} 要求经济相关活动逐路径覆盖；
C13 与定理~\ref{thm:minimal-compatible-extension}、命题~\ref{prop:constitutive-governance-market-compatible}、
定理~\ref{thm:maximal-safe-governance-space}、定理~\ref{thm:dual-layer-governance} 给出相对于现行制度基线的
最小兼容扩展、市场形成条件、宽空间、低溢出的公共规则边界及内部公正的联合条件；
C14 阻止把局部代数恒等式冒充动态稳定、一般均衡、福利优越或经验
事实。由 C15 与定理~\ref{thm:minimal-sufficient-relief}，任一底线承诺都要求
\[
  \mathbf R_t\ge\mathbf R_t^*,
\]
而定理~\ref{thm:relief-endogenous-stability} 进一步把需求、市场厚度、资金和局部稳定闭合；命题
~\ref{prop:relief-not-equalization} 排除由底线推出收入平均。由 C16 与定理
~\ref{thm:residual-claim-incentive-necessity}，正成本内点投入要求正剩余份额；命题
~\ref{prop:appropriation-exploitation-separation} 再把合同剩余索取与非市场强制占有分开。合取十六条件即得到
分层投影结论及其边界。证毕。
\end{proof}

\section{条件性循环联合效应：需求传导、底线覆盖与生活性选择扩张}
\label{sec:strategic-value}

本节不把“支持文明发展”作为不可检验的价值判断直接写入结论，而把它分解为三个可观察结果：生产盈利到
有效需求的正边际传导得到恢复；全体成员不因传统劳动收入衰减而跌破生存--市场接入底线；底线以上由主体
自由选择的生活性活动集合扩张。只有三者与私人投入激励、相对公正和风险边界同时闭合时，才确认本节的
条件性循环联合效应结论。

\begin{definition}[生活性超额预算与可选择容量]
固定结算期 \(t\)。令 \(Y_{ht}^{b}\) 为居民 \(h\) 在本期新型参与收入计入前的市场收入，底线束仍为
\(F_{ht}\)，坐标最小充分救济为
\[
  R_{ht}^{*}=[F_{ht}-Y_{ht}^{b}]_+.
\]
令 \(T_{ht}^{I}\ge0\) 为经过资格、验证、初始确权和实际支付后取得的兴趣或其他合格参与收入，
\(\omega_{ht}\in[0,1]\) 为扣除税费、下一期救济调整和其他已声明负担后可由居民保留的边际份额。定义
底线以上生活性超额预算
\[
  B_{ht}^{L}
  =Y_{ht}^{b}+R_{ht}^{*}-F_{ht}+\omega_{ht}T_{ht}^{I}
  =[Y_{ht}^{b}-F_{ht}]_++\omega_{ht}T_{ht}^{I}\ge0.
\]
该式要求当期已经确认的底线支付不对同一期新取得的合格参与收入作超出声明比例的逆向追回；下一期可以按
新收入基线重算。若 \(\omega_{ht}=0\)，参与收入对当期生活性预算没有边际激励，本节不确认选择扩张结论。

令 \(\mathcal J_t^L\) 为在安全合法域内可供自由选择的竞技、文化、艺术、学习、社群、探索和其他兴趣服务
集合。对每个 \(j\in\mathcal J_t^L\)，令 \(\phi_{hjt}>0\) 为居民 \(h\) 取得一个事前声明的最小有效体验
单位所需的全成本，\(a_{hjt}\in\{0,1\}\) 表示该服务对该居民真实可进入且具有已显露正偏好，
\(w_{hjt}>0\) 为预先固定的比较权重。定义生活性可选择容量
\[
  \Lambda_t^L
  =\sum_{h\in\mathcal H_t}\sum_{j\in\mathcal J_t^L}
  w_{hjt}a_{hjt}
  \mathbf 1\{B_{ht}^{L}\ge\phi_{hjt}\}.
\]
它只计算居民负担得起且能够真实进入的声明选项，不把内容数量、注册账户或被迫参加记作生活性选择，也不
声称不同选项具有相同福利。
\end{definition}

\begin{proposition}[充分救济、参与收入与生活性选择的分工]
\label{prop:floor-income-life-choice}
坐标最小充分救济保证
\[
  Y_{ht}^{b}+R_{ht}^{*}\ge F_{ht},
\]
但本身不保证 \(B_{ht}^{L}>0\)。在 \(\omega_{ht}>0\) 时，
\[
  \frac{\partial B_{ht}^{L}}{\partial T_{ht}^{I}}
  =\omega_{ht}>0.
\]
因此，合格参与收入弱扩张居民的可负担选择集；若至少存在一对 \((h,j)\) 满足
\[
  B_{ht}^{L}(0)<\phi_{hjt}
  \le B_{ht}^{L}(T_{ht}^{I}),
  \qquad a_{hjt}=1,
\]
则 \(\Lambda_t^L\) 严格增加。充分救济负责把生存和市场接入移出胜负对象，参与收入负责把至少一部分居民从
“只能维持底线”推进到“能够在底线以上作出市场选择”。
\end{proposition}

\begin{proof}[证明（预算恒等式与集合包含）]
由正部函数恒等式，
\[
  Y_{ht}^{b}+[F_{ht}-Y_{ht}^{b}]_+-F_{ht}
  =[Y_{ht}^{b}-F_{ht}]_+,
\]
故生活性超额预算表达式成立。对 \(T_{ht}^{I}\) 求偏导得到 \(\omega_{ht}\)。当预算弱增时，任一原来满足
\(B_{ht}^{L}\ge\phi_{hjt}\) 的选项仍然可负担，故可负担选择集按包含关系弱扩张；若跨过至少一个正权重门槛，
相应指标函数从零变为一，\(\Lambda_t^L\) 严格增加。证毕。
\end{proof}

\begin{definition}[联合效应闭环条件]
定义联合效应闭环条件向量
\[
  \mathbf g_t^{S}
  =\Tuple{
  g_{\mathrm{floor}},
  g_{\mathrm{attention}},
  g_{\mathrm{income}},
  g_{\mathrm{choice}},
  g_{\mathrm{culture}},
  g_{\mathrm{incentive}},
  g_{\mathrm{fair}},
  g_{\mathrm{safe}},
  g_{\mathrm{scale}},
  g_{\mathrm{dyn}}
  }\in\{0,1\}^{10},
\]
其坐标依次表示：充分救济及融资可持续；竞争性注意力投入能够改变生产品偏好；参与收入实际到账并具有正
保留率；至少一个生活性选择门槛被跨越；可验证参与形成非零文化沉积并通过共同体意义进入心理效用与需求；
合同剩余索取保留正投入激励；参与--分配相对公正；风险外溢受控且
合法路径非单点；兴趣产业净规模足以覆盖声明需求缺口；局部动态稳定。只有
\[
  \underline g_t^{S}:=\min_k g_{k,t}^{S}=1
\]
时，称声明域通过联合效应闭环条件。
\end{definition}

\begin{theorem}[后 AI 市场参与经济的条件性循环联合效应定理]
\label{thm:post-ai-strategic-value}
在定理~\ref{thm:governance-economy-synthesis} 的 C1--C16、定理~\ref{thm:interest-sector-macro-support}
的 L1--L5 以及定理~\ref{thm:relief-endogenous-stability} 的 R1--R5 成立的共同声明域内，进一步假设：
\begin{enumerate}[label=\textup{(S\arabic*)},leftmargin=4em]
  \item 至少一个居民取得 \(T_{ht}^{I}>0\) 且 \(\omega_{ht}>0\)，并至少跨过一个真实生活性选择门槛；
  \item 生产企业的注意力投入、消费者偏好、参与收入、实际购买和下一期生产盈利均由不重复计算的支付与销售
  数据连接，不以行政注资、封闭积分或后来者资金替代；本定理的严格盈利导数结论另要求
  定理~\ref{thm:dual-channel-bridge} 的局部可微性、通道非负及至少一项严格投入响应条件；
  \item 兴趣产业收入扣除资源成本、挤出和已货币化外部损害后仍满足净规模门槛；
  \item 兴趣参与经可重放历史形成正文化沉积，共同体承认而非平台单方宣布其意义，并满足定理
  ~\ref{thm:culture-mediated-utility-demand} 的正传导条件；
  \item 公共治理只在非单点候选路径族内选择风险约束下的空间释放极大域，不指定具体产品、价格、投入量、
  分配比例或赢家。
\end{enumerate}
若 \(\mathbf g_t^S=\mathbf 1\)，则同时成立：
\[
  \frac{\partial D_{t+1}^{P}}{\partial\Pi_t^{P}}>0,
  \qquad
  D_{t+1}^{P}\ge D_{t+1}^{\dagger},
  \qquad
  N_{t+1}\ge\bar N_{t+1},
\]
即生产盈利经竞争性注意力和参与收入两条路径恢复对下一期有效需求的正边际传导，并越过声明域的需求和市场
厚度阈值；
\[
  Y_{ht}^{b}+R_{ht}^{*}\ge F_{ht}
  \quad(\forall h),
  \qquad
  \Lambda_t^{L}>\Lambda_t^{L,0},
\]
即全体成员的生存--接入底线得到保障，且相对于没有合格参与收入的同口径基线，生活性可选择容量严格扩张；
对至少一个文化有效生活性对象还成立
\[
  \frac{\partial W_{hj}}{\partial z_j}>0,
  \qquad
  \frac{\partial P_{hj}}{\partial z_j}>0,
\]
即参与通过文化意义、共同体认同和心理效用形成持续需求，而不只是一次曝光；
并且至少一个私人投入者满足
\[
  \rho_i\frac{\partial\EE[R^B]}{\partial e_i}
  =c_i'(e_i)>0,
\]
故资本、组织、风险承担和稀缺能力仍有正市场投入激励。若动态固定点的谱半径小于一，上述循环在声明邻域内
局部渐近稳定。

因此，本定理不替代生产、平均分配收入或由政府计划消费，也不排他替代按需分配、基本收入、公共供给、按劳
分配及其他经济理论；它只在所列前提与条件通过时推出：人工建立的资格、验证、确权、支付、纠错和风险边界
允许分散选择形成
\[
  \begin{aligned}
    \text{生产盈利}
    &\longrightarrow \text{竞争性兴趣投入}\\
    &\longrightarrow \text{注意力、文化意义与生产商品偏好}\\
    &\longrightarrow \text{共同体认同、参与收入和生活性选择}\\
    &\longrightarrow \text{有效消费与下一期生产盈利}.
  \end{aligned}
\]
的非退化市场循环。这里能够证明的生活性结论仅为：声明域内没有成员跌破生存--市场接入底线，且底线以上由
本人决定的可负担选择集合相对同口径基线严格扩张；它不等价于效用最大、绝对公平、无限增长或文明永久存续。
\end{theorem}

\begin{proof}[证明（既有定理合成与选择集严格扩张）]
由定理~\ref{thm:minimal-sufficient-relief} 和 R1--R5，\(\mathbf R_t^*\) 逐坐标保障底线，救济账户与储备
逐期闭合，并使需求和市场进入人数分别达到 \(D_{t+1}^{\dagger}\) 与 \(\bar N_{t+1}\)。由 C5--C6、S2
及定理~\ref{thm:dual-channel-bridge}，
\[
  \Pi_t^P
  \to\mathbf I_t^{B*}
  \to(\boldsymbol\vartheta_t^P,\mathbf T_t^I)
  \to D_{t+1}^P,
  \qquad
  \frac{\partial D_{t+1}^{P}}{\partial\Pi_t^{P}}>0.
\]
L1--L5 与 S3 排除把毛收入、注意力总量、外部损害或挤出消费记作净需求增量，并给出
\(D_{t+1}^{P}\ge D_{t+1}^{\dagger}\)。由 S1 和命题~\ref{prop:floor-income-life-choice}，至少一个正权重
可负担门槛被跨越，故 \(\Lambda_t^L>\Lambda_t^{L,0}\)。
由 S4 与定理~\ref{thm:culture-mediated-utility-demand}，至少一条
\[
  z_j\to C_t\to\mu_{C_t}\to W_{hj}\to P_{hj}
\]
链具有严格正导数，故生活性选项不只被技术列出，还通过共同历史和文化认同进入心理效用与市场选择。

再由 C16 和命题~\ref{prop:residual-share-positive-effort}，存在正成本内点投入及其最优回应；
定理~\ref{thm:residual-claim-incentive-necessity} 再给出该域内的正剩余份额必要性。坐标最小救济不设置收入上限，故不消灭该份额。C8、C12--C14 与 S5 分别保证内部
分配可复核、逐活动经济路径可追踪、市场路径非单点且风险受限。最后，定理
~\ref{thm:local-dynamic-stability} 的谱半径条件使上述状态更新在固定点邻域内收敛。十个条件的合取同时给出循环、底线、选择、文化、激励、公正和
风险结论；任一单项都不能代替其余单项。证毕。
\end{proof}

\begin{proposition}[系统性方案不能由平台内部盈利替代]
\label{prop:platform-profit-not-strategic-closure}
单个平台满足收入池为正、账目平衡或注册量增长，既不充分也不必要于联合效应闭环。具体地：
\begin{enumerate}[label=\textup{(\alph*)},leftmargin=3em]
  \item 令平台收入池为正而 \(\kappa^{\mathrm{pref}}=\kappa^{\mathrm{inc}}=0\)，则生产盈利到生产品需求的桥接
  导数为零；
  \item 令注意力投入为正而实际参与支付 \(T_h^I=0\)，则购买力通道和生活性预算增量同时为零；
  \item 令平台盈利为正而部分居民低于 \(F_h\)，且需求或市场厚度低于阈值，则平台盈利不能推出宏观循环
  非退化；
  \item 令剩余索取份额 \(\rho_i=0\) 且投入成本严格递增，则正成本内点投入不能持续；
  \item 令内部份额不可重放或风险损害超过净增量，则私人收入池不能确认相对公正或净社会可行。
\end{enumerate}
因此，本节方案同时跨越生产部门、兴趣部门、居民预算、组织分配和公共风险边界；它是系统传导结构，不是
对单个平台商业模式的局部修补。
\end{proposition}

\begin{proof}[证明（坐标独立反例）]
(a) 直接代入定理~\ref{thm:dual-channel-bridge} 的桥接导数；(b) 代入
\(B_h^L=[Y_h^b-F_h]_++\omega_hT_h^I\)；(c) 违反定理~\ref{thm:relief-endogenous-stability} 的必要阈值；
(d) 使 \(\rho_i\partial\EE[R^B]/\partial e_i=0<c_i'(e_i)\)；(e) 分别使相对公正条件或净收益条件为零。
每个见证都允许平台内部收入为正而联合效应闭环失败，故平台盈利不充分。反之，联合效应闭环可以由多个互操作行业和
组织共同形成而无单一支配平台，故单个平台盈利也非必要。证毕。
\end{proof}

\begin{proposition}[与计划经济、按需分配和传统劳动中心传导的严格区分及非排他兼容]
\label{prop:strategic-system-classification}
设构成性治理给出的候选市场路径族为 \(\mathcal P_t^{\mathrm{legal}}\)，实际路径为分散主体在价格、合同、竞争、
盈亏和进入退出约束下形成的 \(p_t^*\)。若
\[
  |\mathcal P_t^{\mathrm{legal}}|\ge2,
  \qquad
  p_t^*\in\mathcal P_t^{\mathrm{legal}},
\]
且治理算子不唯一指定 \(p_t^*\)，则该制度不是以行政算子预先指定产量、价格和赢家的计划配置。若救济只取
\(R_{ht}^*=[F_{ht}-Y_{ht}^{b}]_+\)，底线以上收入、兴趣品数量和投资方向不由需要函数统一分配，则它也不是
按需分配。由于参与收入 \(T_h^I\) 可以在传统劳动收入不增加时独立为正，所恢复的传导又不同于仅依赖
“生产--就业--工资--消费”的传统劳动中心路径。该制度应归类为：人为建造可验证且受风险约束的市场盘面，
再由自由市场决定盘面上的实际行动与差异化结果。

上述分类差异不构成逻辑排斥。对任一按需分配、基本收入、公共供给、按劳分配或其他制度 \(H\)，只要
\(\Sigma_D^*\oplus H\) 的共同预算、资源、权利、法律与制度环节一致性核验记录通过，二者即可在不同坐标或层级上
并存；是否优于 \(H\) 仍须依推论~\ref{cor:no-superiority-without-comparison} 作反事实比较。
\end{proposition}

\begin{proof}[证明（三个定义反例的合取）]
计划配置要求中心主体在分散选择之前唯一指定关键产量、价格或个体配置。这里
\(|\mathcal P_t^{\mathrm{legal}}|\ge2\)，且治理只给出合法路径族，实际 \(p_t^*\) 由价格、盈亏和进入退出形成，
所以不满足唯一行政指定的定义。按需分配要求由需要函数统一决定底线以上资源；坐标最小救济只补足
\([F_{ht}-Y_{ht}^{b}]_+\)，而底线以上收入、兴趣选择和投资仍可不同，故不满足按需分配定义。最后，构造
\(\Delta T_h^I>0\) 而传统工资增量为零的允许状态，双通道桥接仍可经参与收入增加购买力，故该传导不等同于
只含生产、就业、工资和消费的劳动中心链。三个分类差异分别由定义中的必要特征被否定，结论成立。
非排他兼容部分直接由定理~\ref{thm:minimal-compatible-extension} 的 M4 得到：制度类别不同不构成冲突见证，
只有实际预算、资源、权利、法律或制度环节矛盾才能使特定组合不可行；故“不是按需分配”不蕴含“否定按需分配”。
证毕。
\end{proof}

\begin{Certificate}[后 AI 条件性循环联合效应核验记录]
核验记录必须同时保存：必要品产能与成本、传统劳动收入传导基线、现行制度坐标与具体失灵点、最小兼容候选族及
其字典序极小见证、底线束及逐户救济缺口、救济资金与储备、至少
两个生产者的竞争性投入、真实注意力、投入企业商品偏好与实际销量、参与者实际到账收入及保留率、兴趣产业
净收入和需求缺口覆盖、生活性选择门槛及 \(\Lambda_t^L\)、剩余索取与私人投入响应、内部相对公正、外部损害、
合法路径数量和动态稳定估计。它须分别确认“循环恢复”“底线保障”“选择扩张”“投入激励”和“风险受控”；
不得以其中一项代替总核验记录。
\end{Certificate}

\begin{remark}[可证伪预测与停止边界]
本节至少给出以下可证伪预测：在其他条件可比时，合格参与收入及其净保留率提高应弱扩张可负担生活性选择集；
生产企业对兴趣行业的投入若不能改变其生产品偏好、销量或其他可货币化回报，应在竞争均衡中衰减；参与收入
若不能进入居民购买，购买力通道应接近零；剩余索取被压至零时，正成本私人投入应下降；底线覆盖失败并使需求
或匹配厚度跌破阈值时，循环应失去非退化性。若这些方向在预注册窗口内持续不成立，或者净需求增量不足以
覆盖缺口、救济账户不可持续、分配不可复核、市场路径退化为单点、风险损害超过收益、动态固定点不稳定，则
撤销条件性循环联合效应核验记录。被撤销的只是本节联合结论，不否认局部活动可能仍有消费价值或文化价值。
\end{remark}

\section{基础经济环节不可约性与外部理论位置假设}
\label{sec:theoretical-position}

条件性循环联合效应显著或声明语料库存在联合覆盖缺口，均不足以单独确定理论位置。本节不按声望、传播或经验
成熟度判定外部学术层级，只检验本系统是否同时覆盖经济对象、人的经济主体位置、收入--需求循环、分配规则、
市场构成方式与风险边界，以及这些制度环节能否被任一声明的部门投影无损替代。所得结果是模型内部的制度环节完备性与
不可约性证据，不直接确认外部学术定位。

\begin{definition}[基础制度环节域与理论结果向量]
定义后 AI 经济转型的基础制度环节域
\[
  \mathcal U^*
  =\{
  \mathsf{Prod},
  \mathsf{Demand},
  \mathsf{Human},
  \mathsf{PartDist},
  \mathsf{RightFiat},
  \mathsf{Relief},
  \mathsf{Incentive},
  \mathsf{Culture},
  \mathsf{MarketConst},
  \mathsf{Risk}
  \}.
\]
这些坐标依次表示：生产能力与盈利；有效需求；居民不再仅由劳动投入定义的经济主体位置；参与事实到初始
收益权；治理认证、可执行权利、条件证券化与法币开放循环；生存--市场接入底线；私人正剩余与差异化激励；共同历史到心理效用与持续需求；治理构造非单点
合法市场空间；以及外部损害和系统性风险隔离。

对声明状态 \(x\)，定义理论结果向量
\[
  \Omega^*(x)
  =\Tuple{
  \frac{\partial D_{t+1}^{P}}{\partial\Pi_t^{P}},
  \mathbf 1\{Y_h^b+R_h^*\ge F_h\}_{h\in\mathcal H},
  \mathbf T^I,
  \mathbf b,
  \Tuple{\mathfrak C_t^{\mathrm{source}},\mathbf p_t^F},
  \Lambda^L,
  \mathbf e,
  C_t,
  \mathcal P_t^{\mathrm{legal}},
  r_{\mathrm{spill}}
  }.
\]
结果向量不把各坐标等同为同一种福利量；它只保存本系统必须同时判定的传导、底线、收入、初始份额、权利与
法币来源、选择、投入、文化、市场路径和风险结果。
\end{definition}

\begin{definition}[部门可约性与基础制度不可约性]
对 \(J\subsetneq\mathcal U^*\)，令 \(\pi_J(x)\) 为只保留制度环节集合 \(J\) 的部门投影。若存在全函数
\(\Psi_J\)，使对声明域内每个允许状态 \(x\) 都有
\[
  \Omega^*(x)=\Psi_J(\pi_J(x)),
\]
并且 \(\Psi_J\) 同时保留联合效应条件、相对公正、投入激励、市场非单点和风险核验记录的通过或失败，则称本系统可约
为部门模型 \(J\)。若任何遗漏至少一个必要制度环节的真子集都不存在这种恢复映射，则称本系统在
\(\mathcal U^*\) 上具有基础制度不可约性。

该定义不禁止为研究目的使用部门子模型；它只禁止由某个部门子模型的成功，推出整个经济转型架构已经成立。
\end{definition}

\begin{definition}[基础经济范式候选的内部必要判据]
本文把下列条件定义为提出“后 AI 基础经济范式候选”外部评议假设的内部必要判据：
\begin{enumerate}[label=\textup{(P\arabic*)},leftmargin=4em]
  \item 独立研究对象：生产能力与传统劳动收入传导分离后的经济循环；
  \item 主体位置：人不仅是劳动投入，还作为消费者、可验证参与者、文化形成者和收益权主体进入模型；
  \item 核心机制：生产盈利、注意力和文化偏好、治理认证参与、可执行权利、法币收入、购买力与再生产闭合为有向循环；
  \item 分配宪制：生存--市场接入下限退出竞争，底线以上保留价格、盈亏、合同剩余索取、能力差异和自由选择；
  \item 市场宪制：治理在后 AI 适用域和具体失灵核验记录触发后，以相对于现行制度基线的最小兼容方式建立资格、
  证据、确权、权利与法币来源、支付、申诉和风险边界，但不唯一指定产品、价格、投入和赢家，也不排他否定其他
  分配制度；
  \item 研究程序：为每个关键制度环节给出状态、假设、失败见证、可审计核验记录、可逆试点与停止条件。
\end{enumerate}
P1--P6 与基础制度不可约性同时成立，只证明该系统具备跨制度环节且不可约的内部结构；它们不充分决定其外部学术
类别。“基础经济范式候选”在本文中因此只是一项可提交外部比较、同行评议和竞争性理论检验的定位假设，不是
由本文内部定义自行确认的结论。
\end{definition}

\begin{definition}[必要制度环节独立见证核验记录]
对每个 \(u\in\mathcal U^*\)，核验记录保存一对允许状态
\[
  W_u=(x_u^+,x_u^-),
  \qquad
  \pi_{\mathcal U^*\setminus\{u\}}(x_u^+)
  =
  \pi_{\mathcal U^*\setminus\{u\}}(x_u^-),
\]
但 \(u\) 坐标不同，并满足
\[
  \Omega^*(x_u^+)\ne\Omega^*(x_u^-)
  \quad\text{或至少一个必要核验记录的真值不同}.
\]
这称为制度环节 \(u\) 的独立见证。见证必须来自明确反例或可执行压力状态，不能只以术语不同代替结果差异。
\end{definition}

\begin{theorem}[后 AI 市场参与经济的基础经济环节不可约性定理]
\label{thm:nonsectoral-theoretical-position}
若定理~\ref{thm:governance-economy-synthesis} 的 C1--C16、联合效应闭环条件和必要制度环节独立见证核验记录均已定义，
并且每个 \(u\in\mathcal U^*\) 至少具有一个独立见证，则：
\begin{enumerate}[label=\textup{(T\arabic*)},leftmargin=4em]
  \item 不存在遗漏任一必要制度环节而仍能恢复 \(\Omega^*\) 及全部必要核验记录的部门投影；
  \item 本系统具有基础制度不可约性；
  \item 在 P1--P6 同时成立时，本系统具备提交为“后 AI 经济转型基础制度架构与基础经济范式候选”的内部
  结构条件；该条件不判定外部学术共同体应将其归入何种理论层级。
\end{enumerate}
该结论只判定声明制度环节相对于所列部门投影的不可约性，不推出现实参数已经满足，也不等于学术共同体已经接受
相应外部定位假设。
\end{theorem}

\begin{proof}[证明（投影不可识别与独立失败见证）]
反设存在真子集 \(J\subsetneq\mathcal U^*\) 及恢复映射 \(\Psi_J\)。取任一
\(u\in\mathcal U^*\setminus J\) 及其独立见证 \(W_u=(x_u^+,x_u^-)\)。由于这两个状态在除 \(u\) 外的全部
坐标上相同，必有
\[
  \pi_J(x_u^+)=\pi_J(x_u^-).
\]
全函数性要求
\[
  \Psi_J(\pi_J(x_u^+))
  =
  \Psi_J(\pi_J(x_u^-)),
\]
但独立见证定义给出两状态的 \(\Omega^*\) 或必要核验记录真值不同，矛盾。因此任何遗漏必要制度环节的部门投影都不能
无损恢复本系统，T1--T2 成立。

独立见证的候选情景包括：令偏好或收入桥接系数为零，生产盈利不能恢复需求传导；令参与支付为零，
居民购买力和生活性选择不扩张；保持同一合规参与和初始收益权而令法币来源或可执行义务为零，权利不能进入
开放法币循环；令救济低于缺口，部分居民跌破生存--接入底线；令同一收入池采用不同初始份额，
市场收入不能唯一确定相对公正；令剩余份额为零而投入具有正边际成本，私人内点投入消失；令文化沉积或共同体
意义映射为零，不能确认文化中介的心理效用与需求增量；令合法路径退化为单点，构成性治理变成预先指定；令外部损害
超过净增量，内部盈利不能确认社会可行。这些情景须逐一满足定义中的同投影约束，并验证预算与行为关系相容，才成为独立见证；仅列出可能失效的坐标不证明这些见证已经存在。本定理因此是对已提交独立见证的条件性判据。

最后，P1--P6 分别给出独立对象、主体位置、核心循环、分配规则、市场构成和证伪程序，结合 T2 即满足提交
外部定位假设所需的内部结构条件，所以 T3 成立；外部理论分类不在本证明的结论域内。证毕。
\end{proof}

\begin{proposition}[制度环节不可约性与经验成熟度正交]
\label{prop:position-maturity-orthogonal}
令 \(L\) 表示由对象范围和恢复映射决定的内部结构状态，其取值为“部门制度环节可约”或“基础经济环节不可约”。令
\[
  M
  =\Tuple{
  \mathfrak C_{\mathrm{parameter}},
  \mathfrak C_{\mathrm{causal}},
  \mathfrak C_{\mathrm{scale}},
  \mathfrak C_{\mathrm{external}}
  }
\]
表示经验成熟度核验记录。则 \(L\) 不能由 \(M\) 单独推出，\(M\) 也不能由 \(L\) 单独推出。制度环节不可约的系统可以
仍处于低经验验证阶段，一个高度成熟的部门模型也不会仅因证据充分自动取得外部基础范式定位。因此，经验未决
不改变已经证明的制度环节不可约性；反之，制度环节不可约也不能免除任何经验检验或直接决定外部学术分类。
\end{proposition}

\begin{proof}[证明（二维反例）]
取一个只解释成熟商品拍卖的部门模型，可以拥有充分参数和因果证据而不覆盖救济、文化、参与收入与市场构成，
故 \(M\) 高不推出 \(L\) 为基础经济环节不可约。再取本节满足 P1--P6 和独立见证、但尚未完成规模试点的系统，
其 \(L\) 已由结构判据确定，而 \(\mathfrak C_{\mathrm{scale}}=0\)，故制度环节不可约不推出经验成熟。两个方向均存在
反例，正交性成立。证毕。
\end{proof}

\begin{Certificate}[内部制度环节状态与外部定位假设核验记录]
核验记录必须保存：基础制度环节域、理论结果向量、P1--P6 的逐项定义来源、每个必要制度环节的成对独立见证、部门投影与
恢复映射挑战、不能恢复的结果或核验记录坐标、最小兼容适用路径及其非排他组合状态，以及参数、因果、规模、
外部有效性和同行审查状态。核验记录分别确认“基础经济环节不可约”“基础经济范式外部定位假设待评议”和“经验
验证状态”；不得用经验成熟度替代制度环节不可约性证明，也不得用内部制度环节状态替代现实验证、外部理论比较或学术
共同体确认。
\end{Certificate}

\section{前瞻试验条件、声明语料库联合覆盖缺口与关系结构差异}
\label{sec:urgency-scarcity-innovation}

本节不用分数表示理论地位。是否具备前瞻试验条件由行动时滞、潜在损失和可逆试验的信息价值判定；联合覆盖
缺口只相对于带日期的文献证据域判定；关系结构差异只记录既有研究尚未在同一机制中闭合的路径。三者都不直接
判定理论的普遍迫切性、历史稀缺性或外部创新地位。

\begin{definition}[带日期的外部证据域]
固定证据截止日 \(t_0=\text{2026--08--31}\)。令
\[
  \mathcal E_{t_0}
  =\Tuple{
  \mathcal E^{Q},
  \mathcal E^{L},
  \mathcal E^{F},
  \mathcal E^{A},
  \mathcal E^{R}
  }
\]
分别保存公众预期、已实现劳动市场信号、企业采用与预期、前沿经济学机制以及文化与参与经济研究。当前证据
域包含以下相互制约的事实：
\begin{enumerate}[label=\textup{(E\arabic*)},leftmargin=4em]
  \item 2025 年 Pew 调查中，52\% 的美国劳动者担忧人工智能未来在工作场所的应用，只有 6\%
  认为它会增加自己的长期就业机会；2026 年 Bentley--Gallup 调查中，79\% 的美国成年人预计人工智能
  会在未来十年减少美国岗位，只有 27\% 对企业负责任使用人工智能具有较多或一定信任
  \cite{Pew2025WorkersAI,Gallup2026AI}；
  \item 截至 2026 年 6 月的 ADP 工资记录没有显示全经济广泛就业替代，但 22--25 岁劳动者在高人工智能
  暴露职业中的相对就业缺口达到 19\%，主要经招聘减少形成；作者把它限定为描述性前导指标而非因果估计
  \cite{BrynjolfssonChandarChen2026}；
  \item ILO 的任务级测量估计全球四分之一劳动者所在职业具有某种生成式人工智能暴露，多数岗位更可能发生
  任务转型而非整体消失\cite{ILO2025GenAIJobs}；
  \item 对近六千名企业高管的跨国调查显示，约九成企业报告过去三年人工智能尚未影响本企业就业或生产率，
  但其未来三年平均预期为生产率提高 1.4\%、产出提高 0.8\%、就业下降 0.7\%
  \cite{YotzovEtAl2026}；
  \item 前沿宏观讨论已经把资本偏向的人工智能收益、收入不确定性和消费迟滞联系起来：若收益更多流向资本
  所有者，不平等可能限制需求在部门间扩张\cite{ECB2026AIMonetaryPolicy}。
\end{enumerate}
不同调查的样本、问题、国家和时间窗不得合并为单一“焦虑指数”；职业暴露、企业预期和描述性就业差异也不得
直接写成宏观失业的因果证明。
\end{definition}

\begin{definition}[准备时滞、未准备损失与前瞻试验条件]
令 \(T_G>0\) 为建立资格、证据、确权、支付、申诉、规则版本和风险隔离所需的准备时滞，
\(T_A\) 为声明冲击达到非退化阈值的随机到达时点。若冲击先于准备完成，未准备状态损失为 \(L_0\)，通过
可逆试验后损失为 \(L_1\)，其中 \(L_0>L_1\ge0\)。令 \(K_P>0\) 为当前进行有限试验的净成本，
\(V_P\ge0\) 为试验产生的参数识别、停止权和设计修订所形成的信息期权价值，
\(\beta\in(0,1]\) 为折现因子。定义前瞻试验条件
\[
  K_P
  \le
  \EE\!\left[
    \beta^{T_A}(L_0-L_1)
    \mathbf 1\{T_A<T_G\}
  \right]
  +V_P.
\]
只有该条件在声明压力场景、试验边界和损失口径下通过时，才称制度试验满足本文定义的前瞻试验条件。它要求的是可逆、
分阶段和可停止的市场及治理试验，不是因舆情而全面配置产业。
\end{definition}

\begin{theorem}[不等待宏观危机的条件性迫切定理]
\label{thm:conditional-anticipatory-urgency}
若 \(T_G>0\)、\(L_0>L_1\)、试验不预先指定市场赢家且前瞻试验条件成立，则相对于“保持其他条件不变并等待
到冲击显现后才开始准备”，当前进行有限试验的期望净损失不更高；若不等式严格，则当前试验严格占优。
该结论与“宏观危机已经发生”逻辑独立。
\end{theorem}

\begin{proof}[证明（期望损失差与信息期权）]
等待策略在事件 \(\{T_A<T_G\}\) 上承担额外未准备损失
\[
  \EE\!\left[
    \beta^{T_A}(L_0-L_1)\mathbf 1\{T_A<T_G\}
  \right].
\]
当前试验支付 \(K_P\)，同时取得可停止、可校准和可改写后续设计的信息价值 \(V_P\)。故“当前试验减等待”的
期望净损失不超过
\[
  K_P
  -\EE\!\left[
    \beta^{T_A}(L_0-L_1)\mathbf 1\{T_A<T_G\}
  \right]
  -V_P,
\]
由前瞻试验条件知其非正，严格条件下严格为负。该证明不需要断言 \(T_A\) 必然有限，也不把公众担忧当作
冲击到达；公众预期、青年招聘差异和企业预期只用于校准分布与损失，宏观危机是否已发生仍由独立数据判定。
证毕。
\end{proof}

\begin{proposition}[当前证据只确认前导风险而非宏观崩溃]
\label{prop:leading-risk-not-collapse}
在 \(\mathcal E_{t_0}\) 内，一方面，青年招聘差异、广泛职业暴露、企业未来减员预期和公众信任下降使
“风险概率必为零且无需准备”没有证据支持；另一方面，ADP 未发现全经济广泛替代、企业过去影响接近零且
ILO 以岗位转型为主要判断，使“人工智能已经造成宏观就业崩溃”同样不能确认。因此，当前最强结论是：需要
尽早估计 \(T_A,T_G,L_0-L_1,K_P,V_P\) 并开展受限试验，而不是把危机当成既成事实。
\end{proposition}

\begin{proof}[证明（双向非蕴含的证据域判定）]
在 \(\mathcal E_{t_0}\) 内，青年高暴露职业招聘缺口、职业任务暴露、企业未来减员预期和公众担忧至少使
“冲击概率已被证成零”不能成立；这为估计前瞻试验条件中的概率与损失提供正信息。另一方面，同一证据域还含有
全经济未见广泛替代、约九成企业报告过去就业和生产率未受影响、以及多数岗位更可能转型的相反约束。因此，
\(\mathcal E_{t_0}\) 既不蕴含零风险，也不蕴含宏观崩溃已经发生。两者之间可共同确认的最强交集是“存在需要
估计的前导风险”；是否立即试验再由定理~\ref{thm:conditional-anticipatory-urgency} 的成本--损失不等式决定。
证毕。
\end{proof}

\begin{definition}[问题条件、研究家族与覆盖核验记录]
\label{def:research-problem-gates}
把本系统必须同时回答的九项问题条件记为
\[
  \mathcal Q^*
  =\{q_D,q_F,q_E,q_R,q_I,q_C,q_K,q_M,q_S\},
\]
分别表示需求回传、充分救济、参与经济身份、相对公正分配、私人投入激励、生活性选择、文化生成、分散市场
配置与风险外溢控制。令带日期文献语料库为
\[
  \mathcal L_{t_0}
  =\{
  \mathcal L_{\mathrm{attention}},
  \mathcal L_{\mathrm{AI\text{-}distribution}},
  \mathcal L_{\mathrm{culture}},
  \mathcal L_{\mathrm{creator}}
  \}.
\]
若研究家族 \(r\) 明确定义某条件的对象、状态、因果或均衡映射、可观察量及失败边界，则记
\(A_{rq}=1\)；仅提及概念、风险或政策愿景，记 \(A_{rq}=0\)。若同一研究家族还在一个模型或可执行核验记录中
闭合从生产盈利、注意力和文化偏好、参与收入、居民购买力到下一期需求的完整有向路径，则记
\(B_r=1\)，否则 \(B_r=0\)。

当前覆盖证据显示：注意力经济前沿模型已经联合处理平台投资、商品发现、产品需求和广告收入
\cite{Chen2026Attention}；人工智能分配模型处理任务替代、生产率、劳动收入和资本回报
\cite{RockallTavaresPizzinelli2025}；文化消费研究识别文化资本、社会资本与体验参与的关系
\cite{GladstoneBellezza2025}；创作者经济研究观察到参与和生产扩张而多数创作者仍无显著回报
\cite{ZhangNegus2026}。最后一项结果与第~\ref{sec:attention-counterexample-core-functions}节的集中型注意力
变现基线相容：它支持“注意力货币化不充分”的现实动机，但不单凭一项研究证明全部平台均集中、三项核心功能
已经充分或某一载体历史必然胜出。这些研究分别提供本系统的前件或局部映射，但在当前语料库中
没有单一研究家族满足全部 \(A_{rq}=1\) 且 \(B_r=1\)。
\end{definition}

\begin{theorem}[声明语料库内的联合覆盖缺口定理]
\label{thm:declared-corpus-integrative-scarcity}
若覆盖核验记录证明
\[
  \forall r\in\mathcal L_{t_0},\qquad
  \sum_{q\in\mathcal Q^*}A_{rq}<|\mathcal Q^*|
  \quad\text{或}\quad B_r=0,
\]
则在语料库 \(\mathcal L_{t_0}\) 内，不存在一个单一研究家族同时闭合全部九项问题与完整市场循环；本系统
登记了一组相对于该语料库的联合覆盖差异。该结论不评价世界范围的历史首创、理论稀缺程度或语料库之外是否
存在等价理论。
\end{theorem}

\begin{proof}[证明（覆盖矩阵反证）]
假设语料库内存在一个单一研究家族 \(r^*\) 已闭合全部问题和完整循环。由覆盖定义，必有
\[
  \sum_{q\in\mathcal Q^*}A_{r^*q}=|\mathcal Q^*|,
  \qquad B_{r^*}=1.
\]
这与核验记录对每个 \(r\) 给出的析取条件矛盾。因此语料库内没有这样的单一家族。由于结论量词只遍历
\(\mathcal L_{t_0}\)，扩大检索域、发现遗漏文献或出现新研究都可以撤销该结论，故没有把有限检索误写为
绝对原创性证明。证毕。
\end{proof}

\begin{definition}[机制图、关键路径与关系结构差异]
把研究家族 \(r\) 的机制写成有向图 \(H_r=(V_r,E_r)\)，把本系统写成
\(H^*=(V^*,E^*)\)。节点可以来自既有经济学；关系结构比较不要求每个节点或每条局部边首次出现。令
\(\mathcal P^{\mathrm{crit}}(H^*)\) 为经过本系统全部必要制度环节的关键路径族，其中至少包括
\[
\begin{aligned}
  \Pi^P
  &\longrightarrow I^B
  \longrightarrow(\vartheta^P,T^I)
  \longrightarrow(D^P,B^L)
  \longrightarrow\Pi^{P\prime},\\
  Z^C
  &\longrightarrow C_t
  \longrightarrow\mu_{C_t}
  \longrightarrow W_h
  \longrightarrow P_h
  \longrightarrow D^I,\\
  (a_i,\varepsilon_i,\omega_i,h_i)
  &\xrightarrow{E_G=1,\,M_G}z_i
  \xrightarrow{V_G=1,\,W_G}w_i
  \xrightarrow{\phi_G}b_i\\
  &\xrightarrow{T_G}y_i
  \xrightarrow{Q_G=(1,1)}
  \Auditable_i\wedge\Appealable_i.
\end{aligned}
\]
定义相对于声明语料库的关系结构差异集合
\[
  \mathcal N_{\mathrm{rel}}
  =
  \mathcal P^{\mathrm{crit}}(H^*)
  \setminus
  \bigcup_{r\in\mathcal L_{t_0}}\mathcal P(H_r),
\]
其中一条路径只有在同一研究家族内保留节点语义、方向、失败支路和结算约束时才视为已存在；不得把若干互不
连接论文中的局部边拼接后冒充一个已经闭合的既有机制。
\end{definition}

\begin{theorem}[关系结构差异判据]
\label{thm:relational-architecture-innovation}
若覆盖核验记录可重放且 \(\mathcal N_{\mathrm{rel}}\ne\varnothing\)，则本系统相对于
\(\mathcal L_{t_0}\) 至少登记一条当前单一研究家族未闭合的关键有向路径。该结论只证明声明语料库中的路径
集合差，不自行确认外部创新性评价，也不把生产、消费、注意力、文化、利润、救济、市场或治理等基础概念据为
原创。
\end{theorem}

\begin{proof}[证明（有向路径集合差）]
取 \(p\in\mathcal N_{\mathrm{rel}}\)。由定义，\(p\in\mathcal P^{\mathrm{crit}}(H^*)\)，故本系统定义了该路径
的全部节点、方向、制度环节与失败条件；同时对任意 \(r\in\mathcal L_{t_0}\)，均有
\(p\notin\mathcal P(H_r)\)。所以该关系结构在声明语料库的任一单一家族中均未闭合，而在本系统中闭合，
关系结构差异成立。若将来发现某个 \(H_{r'}\) 已包含等价路径，则从集合差中删除 \(p\)，相应差异结论自动缩减。
证毕。
\end{proof}

\begin{proposition}[本系统登记的关系差异与既有概念来源]
\label{prop:innovation-location-boundary}
在当前覆盖核验记录下，本系统登记的关系结构差异集中于：把生产者竞争性盈利注入同时连接商品偏好与参与收入；
把参与收入接入
生活性预算和有效需求；把坐标最小充分救济与正剩余索取同时纳入局部稳定条件；把文化沉积作为心理效用和持续
需求的内生中介；把逐活动留痕、规则版本、相对公正和风险隔离接入同一组织核验记录；并用围棋提供
“稳定规则--可重放参与--文化沉积--持续市场”的存在模型。自动化任务、双边平台、广告发现、最低需要、剩余
索取、文化资本、机制设计、合同执行和组织风险本身均有既有来源，本命题不对这些概念提出原创性主张。
\end{proposition}

\begin{proof}[证明（关键路径集合差与节点来源分离）]
由定理~\ref{thm:relational-architecture-innovation}，只有属于
\(\mathcal N_{\mathrm{rel}}\) 的完整关键路径才能在声明语料库内确认关系结构差异。命题前半列出的各项都对应
\(H^*\) 中同时保留方向、制度环节、失败支路和结算约束的关键路径；覆盖核验记录把它们登记为当前单一研究家族未闭合
的集合差元素。后半列出的概念则分别已经是某个 \(H_r\) 的节点或局部边，故不能仅因进入 \(H^*\) 就取得节点
原创性。于是本证明只把差异定位于路径复合与条件关系，而不对既有基础概念作创新归属判断。若覆盖矩阵更新，结论随
\(\mathcal N_{\mathrm{rel}}\) 同步缩减。证毕。
\end{proof}

\begin{definition}[逻辑核验记录与经验验证核验记录分离]
定义
\[
  \mathfrak C_{\mathrm{status}}
  =\Tuple{
  \mathfrak C_{\mathrm{logic}},
  \mathfrak C_{\mathrm{parameter}},
  \mathfrak C_{\mathrm{causal}},
  \mathfrak C_{\mathrm{scale}},
  \mathfrak C_{\mathrm{external}}
  }.
\]
五项分别表示形式推演闭合；关键弹性、转化率、保留率、成本和时滞已估计；反事实因果效应已识别；净需求
缺口覆盖与动态稳定已在目标规模通过；跨时间、地区和制度环境具有外部有效性。当前文档可以对已列假设的
条件蕴含确认 \(\mathfrak C_{\mathrm{logic}}\)，但其余坐标必须由预注册试点、实际支付与销售、独立复算、
长期跟踪和跨域复制分别取得，不得由公式数量或声明语料库覆盖缺口代替。
\end{definition}

\begin{theorem}[形式证明不能替代经验成立]
\label{thm:formal-proof-not-empirical-validation}
对本系统任一形式结论 \(A\Rightarrow B\)，若现实证据尚未证明前件 \(A\) 在目标域成立，则
\(\mathfrak C_{\mathrm{logic}}=1\) 不能推出
\(\mathfrak C_{\mathrm{parameter}}=\mathfrak C_{\mathrm{causal}}=
  \mathfrak C_{\mathrm{scale}}=\mathfrak C_{\mathrm{external}}=1\)。
\end{theorem}

\begin{proof}[证明（实质蕴含的前件未定）]
命题 \(A\Rightarrow B\) 为真允许 \(A\) 为假。形式推演只证明一旦前件成立，结论按所列算子推出；它不从
逻辑内部生成现实参数、处理组反事实、规模阈值或外部环境不变性。因此，取一个满足形式语言但现实中
\(\kappa^{\mathrm{pref}}=0\)、参与收入未到账或风险损害超过净收益的状态，逻辑定理仍保持条件真，而相应经验
坐标失败。故逻辑核验记录不能替代经验验证。证毕。
\end{proof}

\begin{Certificate}[前瞻试验条件、联合覆盖缺口与关系结构差异核验记录]
核验记录保存证据截止日、数据库与检索式、纳入排除标准、全部命中文献和版本、公众调查样本与题目、劳动市场数据
口径、企业实现值与预期值、研究家族覆盖矩阵、机制图、关键路径集合差、准备时滞、损失场景、试验成本、信息
期权及五项经验状态。核验记录禁止数字评分，分别确认“前瞻试验条件通过”、“声明语料库内联合覆盖缺口已登记”、
“声明语料库内关系结构差异已登记”、“外部创新性评价未决”和“经验状态未决或已通过”。新增文献、相反数据、
参数变化或试点失败必须触发版本更新；任何一个局部结论撤销，不自动撤销其余独立结论。
\end{Certificate}

\section{经济脆弱性与市场排除暴露及多条件逻辑覆盖}
\label{sec:economic-anxiety-gap}

主观焦虑包含心理、医疗、家庭、文化和政治等多重原因，不能由经济模型整体定义或治愈。本节只定义可由收入、
市场接入、支付可靠性与生活性选择直接观测的经济脆弱性和市场排除暴露，并证明前述制度组合能够在声明域内
降低这些暴露。由此得到的是经济机制结论，而不是临床诊断或社会心理学的充分解释。

\begin{definition}[经济脆弱性与市场排除暴露指数]
令 \(\mathcal H_t^w\subseteq\mathcal H_t\) 为在时期 \(t\) 自愿寻求市场参与机会的居民集合。对每个
\(h\in\mathcal H_t^w\)，定义：
\[
  s_{ht}
  =
  \Pr\!\left(
  Y_{h,t+1}^{b}+R_{h,t+1}<F_{h,t+1}
  \,\middle|\,\mathcal I_t
  \right)
\]
为下一期跌破生存--市场接入底线的条件概率；
\[
  x_{ht}=1-n_h(Y_{ht};F_{ht})
\]
为无法真实进入至少一个合格市场入口的概率；令
\[
  z_{ht}
  =
  \Pr\!\left(
  \neg\Paid_{h,t+1}
  \,\middle|\,
  \Eligible_{h,t}=1,\,
  \Verified_{h,t}=1,\,
  \Valued_{h,t}=1
  \right)
\]
为已合格、已验证且已取得正估值的参与未按声明合同支付的条件概率。

若居民 \(h\) 至少显露一个正偏好的合法生活性选项，定义其最小有效体验门槛
\[
  \phi_{ht}^{\min}
  =
  \min\{\phi_{hjt}:j\in\mathcal J_t^L,\ a_{hjt}=1\}>0
\]
及无选择暴露
\[
  c_{ht}^{L}
  =
  \min\left\{
  1,\,
  \frac{[\phi_{ht}^{\min}-B_{ht}^{L}]_+}{\phi_{ht}^{\min}}
  \right\}.
\]
给定预先声明且非负的权重
\(\alpha_h,\beta_h,\gamma_h,\delta_h\)，其中每名居民至少一个权重严格为正，定义个体与聚合经济脆弱性及市场排除暴露
\[
  A_{ht}^{E}
  =
  \alpha_hs_{ht}
  +\beta_hx_{ht}
  +\gamma_hz_{ht}
  +\delta_hc_{ht}^{L},
  \qquad
  \mathcal A_t^{E}
  =
  \sum_{h\in\mathcal H_t^w}A_{ht}^{E}.
\]
该指数度量生存不确定、市场排除、有效参与不获支付和无法负担任何已显露偏好选项四类经济暴露。它不度量
主观焦虑，不把收入差异、竞争失败或未选择某项兴趣本身定义成病理。
\end{definition}

\begin{assumption}[制度传导对经济暴露的弱单调性]
\label{ass:economic-anxiety-monotonicity}
在共同观察窗口与同口径基线下假设：
\begin{enumerate}[label=\textup{(A\arabic*)},leftmargin=4em]
  \item 可持续且可信承诺的坐标最小充分救济使
  \(s_{ht}^{1}\le s_{ht}^{0}\)，在原基线存在正底线缺口而支付可信时严格下降；
  \item 市场接入概率 \(n_h(Y;F)\) 关于可支配资源和真实入口数量弱增，故
  \(x_{ht}^{1}\le x_{ht}^{0}\)；
  \item 完整参与支付链、独立复核和有效申诉不提高合格正估值参与的未支付概率，故
  \(z_{ht}^{1}\le z_{ht}^{0}\)；
  \item 参与收入及其净保留率不降低 \(B_{ht}^{L}\)，故
  \(c_{ht}^{L,1}\le c_{ht}^{L,0}\)。
\end{enumerate}
上标 \(0,1\) 分别表示无制度闭环基线与通过联合效应闭环条件的处理状态。上述方向必须由实际支付、入口、申诉和
消费数据识别；不得用制度文本存在替代。
\end{assumption}

\begin{theorem}[经济脆弱性与市场排除暴露的弱降低与严格降低条件]
\label{thm:economic-anxiety-exposure-reduction}
若假设~\ref{ass:economic-anxiety-monotonicity} 成立，且定理~\ref{thm:post-ai-strategic-value} 的联合效应闭环
条件通过，则
\[
  \mathcal A_t^{E,1}\le\mathcal A_t^{E,0}.
\]
若至少存在一个居民和一个正权重坐标，使 A1--A4 中相应不等式严格成立，则
\[
  \mathcal A_t^{E,1}<\mathcal A_t^{E,0}.
\]
特别地，若可信充分救济使原有正底线风险归零、至少一个自愿参与者取得正保留参与收入并跨过首个生活性选择
门槛，且其相应权重为正，则即使不假设所有居民获得相同收入，聚合经济脆弱性与市场排除暴露也严格下降。
\end{theorem}

\begin{proof}[证明（非负加权的逐坐标单调性）]
由 A1--A4，对每个 \(h\in\mathcal H_t^w\) 有
\[
  s_{ht}^{1}-s_{ht}^{0}\le0,\quad
  x_{ht}^{1}-x_{ht}^{0}\le0,\quad
  z_{ht}^{1}-z_{ht}^{0}\le0,\quad
  c_{ht}^{L,1}-c_{ht}^{L,0}\le0.
\]
分别乘以非负权重并求和，得到
\[
  \mathcal A_t^{E,1}-\mathcal A_t^{E,0}
  =
  \sum_h
  \left[
  \alpha_h\Delta s_{ht}
  +\beta_h\Delta x_{ht}
  +\gamma_h\Delta z_{ht}
  +\delta_h\Delta c_{ht}^{L}
  \right]
  \le0.
\]
若至少一项在正权重下严格为负，其余项均非正，则总和严格为负。充分救济与选择跨越只是保证至少两个候选
严格项的充分条件，不要求所有坐标或所有居民同时改善。证毕。
\end{proof}

\begin{remark}[从经济暴露到主观焦虑的识别边界]
若另有经验证的心理或行为关系
\[
  \EE[\mathrm{Anxiety}_{h,t+1}\mid\mathcal I_t]
  =\psi_h(A_{ht}^{E},\mathbf Z_{ht}),
  \qquad
  \frac{\partial\psi_h}{\partial A_{ht}^{E}}>0,
\]
且处理没有通过 \(\mathbf Z_{ht}\) 产生抵消性伤害，才可进一步预测主观焦虑下降。没有这一独立识别，形式系统
只能证明经济性暴露降低，不能声称治疗焦虑、消除意义危机或保证社会稳定。
\end{remark}

\begin{definition}[九项联合经济条件与候选机制族]
令关键问题条件为
\[
  \mathcal G^{*}
  =
  \{
  g_D,g_F,g_E,g_R,g_I,g_C,g_K,g_M,g_S
  \},
\]
分别表示：有效需求传导、充分救济底盘、非传统参与取得可执行经济身份、参与--分配相对公正、私人正投入
激励、生活性选择扩张、共同文化内生形成、非计划的分散市场配置、风险外溢受控。令候选机制族为
\[
  \mathcal M^{*}
  =
  \{
  M_A,M_F,M_E,M_R,M_I,M_C,M_K,M_M,M_S
  \},
\]
其中 \(M_A\) 为竞争性注意力双通道，\(M_F\) 为坐标最小充分救济，\(M_E\) 为资格--验证--确权--支付链，
\(M_R\) 为相对公正与申诉机制，\(M_I\) 为竞争性市场中的剩余索取，\(M_C\) 为兴趣产业与正保留参与收入，
\(M_K\) 为可验证参与沉积、共同体认同与文化中介需求，\(M_M\) 为非单点合法路径和市场选择，\(M_S\) 为安全极大域和风险外溢约束。
\end{definition}

\begin{theorem}[联合经济条件的充分覆盖与局部失效边界]
\label{thm:joint-economic-conditions}
在本形式系统声明的机制类别内，若各机制均满足其实际调用定理的全部前提，
包括正投入的充分条件、真实支付、选择门槛、文化沉积和风险约束，则九项经济条件同时成立。
任何局部支付、价格或收益结论都不能自动补足未验证的其他条件。
这些结论不单独证明九类机制相互独立，亦不证明该机制族具有全局或集合包含意义上的最小性。
\end{theorem}

\begin{proof}[证明（逐项充分条件与适用边界）]
定理~\ref{thm:dual-channel-bridge} 在其需求与投入响应条件下给出 \(g_D=1\)；
定理~\ref{thm:relief-endogenous-stability} 在底线、融资及稳定条件下给出 \(g_F=1\)；
定理~\ref{thm:typed-atomic-execution-chain} 的真实履约与支付条件给出 \(g_E=1\)；
定理~\ref{thm:relative-participation-fairness} 的完整条件给出 \(g_R=1\)。
私人正投入须由命题~\ref{prop:residual-share-positive-effort} 的严格边际回报、凹性与参与条件
或另一完整均衡证明给出 \(g_I=1\)，仅有正剩余份额不足以替代这些前提。
命题~\ref{prop:floor-income-life-choice} 在净收入跨过实际选择门槛时给出 \(g_C=1\)；
定理~\ref{thm:cultural-stock-accumulation} 与定理~\ref{thm:culture-mediated-utility-demand}
在正文化沉积、承认及严格传导条件下给出 \(g_K=1\)；
命题~\ref{prop:constitutive-governance-market-compatible} 与
定理~\ref{thm:maximal-safe-governance-space} 分别给出 \(g_M=1,g_S=1\)。
上述结论逐项合取即得充分覆盖。

局部失效边界分别为：需求传导系数为零，不能确认正需求响应；收入不足而未补足底线，不能确认充分救济；
参与没有确权或未支付，不能确认实际收入；预算平衡但等价参与被区别对待，不能确认相对公正；
正剩余份额的边际收入低于努力成本，不能确认正投入；净参与收入不足以跨越体验门槛，不能确认选择扩张。
文化中介消失时，
\[
 D_C\mu_{C_t}\frac{\partial C_t}{\partial z_j}=0
 \quad\Longrightarrow\quad
 \frac{\partial W_{hj}}{\partial z_j}
 =\frac{\partial u_h^{\mathrm{intrinsic}}}{\partial z_j},
\]
右端仍可为正，故只能撤销文化中介传导结论，不能撤销全部体验价值或需求。
最后，合法选择退化为单点，不能确认分散市场配置；风险超过声明上限，不能确认风险可控。
这些失效情景只说明各项结论的前提不可省略，并未逐一固定其余全部经济条件。

若另欲确认某机制族的不可删减性，须对每个机制 \(M_j\) 构造一个允许状态，使
\[
 \bigwedge_{\ell\ne j}\mathfrak C_{M_\ell}=1,\qquad g_j=0,
\]
并检查共同预算、行为均衡、核验和跨期状态没有随之失效。若条件相互依赖，则应把必要依赖视为一个组合
再比较，不能用单一坐标归零的示意代替联合可行见证。证毕。
\end{proof}

\begin{proposition}[五类单维应对不能推出系统闭合]
\label{prop:one-dimensional-responses-insufficient}
下列单维应对均不足以推出 \(\min_{g\in\mathcal G^*}g=1\)：
\begin{enumerate}[label=\textup{(\roman*)},leftmargin=3em]
  \item 只提高人工智能产能：可使供给扩张，却不推出有效需求、参与收入、文化认同或分配公正；
  \item 只冻结旧岗位或虚构任务：可暂时维持工资，却不证明活动具有市场需求，也不能长期抵消技术替代；
  \item 只提供充分救济：可以保障 \(g_F\)，却不自动生成 \(g_D,g_E,g_I,g_C,g_K\)；
  \item 只扩大平台流量和收入池：可以形成私人收入，却不推出参与者实际支付、生活性选择、宏观需求或风险
  受控；
  \item 只由行政计划指定兴趣项目和资金：可以令部分账户非零，却使 \(g_M=0\)，并失去分散偏好对项目存亡
  的检验。
\end{enumerate}
因此，本系统所登记的多条件联合覆盖问题不是“过去没有生产理论、救济理论、平台理论或监管理论”，而是这些
单坐标工具不能单独推出本文定义的九项条件同时闭合。
\end{proposition}

\begin{proof}[证明（五个允许反模型）]
(i) 令复制产能趋于无穷，同时取 \(\kappa^{\mathrm{pref}}=\kappa^{\mathrm{inc}}=0\)，则需求条件失败；
(ii) 用外部转移维持旧工资但令相应商品无自愿购买且净成本持续为正，则市场需求与跨期可持续条件失败；
(iii) 令 \(\mathbf R=\mathbf R^*\) 而不存在参与确权、正剩余索取和文化沉积，只有 \(g_F=1\)；
(iv) 令平台收入池为正但 \(\widetilde T_i^{\mathrm{pay}}=0\)、份额不可复算且风险外部化，则经济身份、公正、生活选择与
安全条件均可失败；(v) 令行政算子唯一指定项目、价格和赢家，则无论账户是否为正，合法路径退化为单点并有
\(g_M=0\)。每个反模型都满足相应单维措施而至少一个条件为零，故没有任何单维措施充分推出九条件合取。证毕。
\end{proof}

\begin{Certificate}[经济脆弱性与市场排除暴露及多条件闭合核验记录]
核验记录保存：自愿寻求参与者集合、底线缺口概率、真实市场入口及拒绝原因、合格正估值参与的未支付率、首个
正偏好生活性选项的全成本、生活性超额预算、四项权重与基线、各条件处理前后值、主观焦虑指标及其独立识别
设计。多条件部分逐项保存九个机制核验记录和九个失败见证。若只观察到问卷焦虑下降而没有经济坐标变化，不能据此
证明本机制；若经济暴露下降而主观焦虑不变，只保留经济机制结论；若任一条件失败，只撤销依赖该条件的联合
结论。
\end{Certificate}

\begin{corollary}[制度机制竞争推论]
\label{cor:governance-competition}
在生产约束松弛且分配规则敏感的共同观测域 \(\Omega\) 内，若技术可行集和外部冲击保持可比，机制差异只经
参与验证、操纵损失、预算平衡、申诉纠错、参与、创新和长期稳定进入同一事前声明绩效函数 \(\mathcal J\)，且
核验记录中的稳健正下界成立，则 \(G_1\succ_{\mathcal J}G_2\)。因此，在这些条件下，上述治理维度可以成为制度机制
竞争的可识别排序来源；本文不由生产约束松弛本身推出历史上必然“越来越”如此。
\end{corollary}

\begin{Certificate}[同口径机制排序]
给定共同观测域 \(\Omega\) 与同一绩效函数 \(\mathcal J\)，若
\[
  \inf_{\omega\in\Omega}
  \bigl[\mathcal J(G_1;\omega)-\mathcal J(G_2;\omega)\bigr]>0,
\]
则在 \(\Omega\) 内确认 \(G_1\succ_{\mathcal J}G_2\) 的稳健排序核验记录。
\end{Certificate}

\begin{proof}[证明]
严格正下界给出
\[
  \forall\omega\in\Omega,\quad
  \mathcal J(G_1;\omega)>\mathcal J(G_2;\omega).
\]
因此排序在声明证据域内稳健。核验记录不外推到 \(\Omega\) 之外，也不把单一绩效函数的排序提升为所有价值
维度上的绝对优越。证毕。
\end{proof}

\section{竞技参与、营销回报与再投入的闭环循环}
\label{sec:competition-closed-cycle}

\begin{remark}[闭环验收的模型分支]
连续投入基准适用定理~\ref{thm:constructive-equilibrium-cycle}；组织选择与残余误差分支适用
定理~\ref{thm:endogenous-organization-contract} 及命题~\ref{prop:noisy-cycle-robustness}。
后一分支逐状态保持基本支付和预算，真实贡献排序使用期望与误判概率结论，不能替换成绝对正确排序。
涉及总量恢复时还须通过命题~\ref{prop:micro-macro-compatibility}；定理~\ref{thm:joint-micro-macro-witness}
给出了这一连接的一个共同可行实例。一般福利优越性与经验成立仍需对应证据，不由联合会计状态直接推出。
\end{remark}

\begin{definition}[组织账、产品账与营销账的跨期闭合]
\label{def:closed-cycle-ledgers}
固定竞技组织、至少两家竞争生产者、产品域、合同与归属窗口。组织收入、组织成本、参与分配与留存余额
\(Q_t^B\ge0\) 满足
\[
\begin{aligned}
  R_t^B&=F_t^B+\sum_k I_{k,t}^B,\\
  R_t^B&=C_t^B+H_t^B
       +\sum_{i\in\mathcal E_t}T_{it}^B+Q_t^B,\\
  T_{it}^B&=\beta_t^{\mathrm{base}}+D_t^{\mathrm{diff}}\lambda_{it}.
\end{aligned}
\]
留存余额依公开合同用于储备、组织者剩余或再投入，参与分配不重复计入组织成本。产品账与营销账满足
\[
\begin{aligned}
  S_{k,t+1}^P
    &=C_{k,t+1}^{P,\mathrm{prod}}+M_{k,t}+\Pi_{k,t+1}^{P,\mathrm{net}},\\
  M_{k,t}&=I_{k,t}^B+J_k(I_{k,t}^B),\\
  \Gamma_{k,t}^{\mathrm{mkt}}
    &=\Delta\Pi_{k,t+1}^{P,\mathrm{pref}}-M_{k,t}>0.
\end{aligned}
\]
销售按本轮营销的归属窗口结算；跨期折现沿用定义~\ref{def:preference-marketing-return}。
偏好与收入增益按命题~\ref{prop:sales-marketing-accounting} 的同一分解顺序记录。
已实现净利润扣除已承诺用途后的可用现金记为 \(L_{k,t+1}^{\mathrm{cash}}\)，下一期营销满足
\[
  0<M_k(I_{k,t+1}^B)\le L_{k,t+1}^{\mathrm{cash}}
    \le\Pi_{k,t+1}^{P,\mathrm{net}},
  \qquad I_{k,t+1}^B>0.
\]
赞助在生产者账上为营销支出，在组织账上为资金来源；参与分配在组织账上为支付，在家庭账上为收入。
各转移在主体账间逐笔对应，合并核算时抵销；产品回款与产品成本分别登记。同一资金在不同时间支持下一轮
交易构成周转，不能在同一时点同时支持两项排他的付款承诺。
\end{definition}

\begin{theorem}[竞技组织参与经济的闭环循环定理]
\label{thm:competition-closed-cycle}
固定声明状态域 \(\mathcal V\)、允许冲击集合 \(\mathcal H_{\mathrm{shock}}\) 与按期结算的状态更新
\(x_{t+1}=F(x_t,\eta_t)\)。若初态 \(x_0\in\mathcal V\)，且对该域内每个状态及允许冲击均满足：
\begin{enumerate}[label=\textnormal{L\arabic*:},leftmargin=3.5em]
  \item 赛事与支持网络通过定义~\ref{def:competition-led-domain} 的竞技主导识别，至少两家生产者的
  注意力竞争符合定理~\ref{thm:competitive-attention-injection} 的固定收入模型，或
  定理~\ref{thm:constructive-equilibrium-cycle} 的连续联合模型，或定理~\ref{thm:endogenous-organization-contract}
  与命题~\ref{prop:noisy-cycle-robustness} 的有限赞助包模型，投入由既有已实现利润现金支持；
  \item 接纳、完成与验证规则事前固定，资金覆盖人数上限、正基本份额及正差异池；每名有效参与者的
  应付额与实际支付均通过 \(\mathfrak C_{\mathrm{part}}^+\) 及支付最终性核验记录。依法发生的代扣单独登记，
  对每名参与者另有可验证下界 \(y_{it}^F\ge\underline y_t>0\)，下界与代扣已在资金预算中覆盖；
  \item 真实竞技注意力改变赞助产品的购买偏好，参与到账收入具有正消费传导；产品价格、质量、收入与
  竞争者投入的反事实口径可识别，生产者份额转移与总需求增量分别记录；
  \item 产品销售回款覆盖产品成本，\(P_k(\vartheta^0,Y^1)\ge0\)，实际偏好增量销售利润覆盖全部营销
  成本并形成 \(\Gamma_{k,t}^{\mathrm{mkt}}>0\)。产品、组织与营销账满足定义~\ref{def:closed-cycle-ledgers}，
  已实现净利润中的可用现金足以支持下一期正营销投入及其分配承诺；
  \item 制度公正、申诉、真实退出与适用风险边界通过对应核验记录，且状态域具有前向不变性：
  \(F(\mathcal V,\mathcal H_{\mathrm{shock}})\subseteq\mathcal V\)。若确认生存底线、宏观规模或动态稳定
  结论，还分别满足相应救济、规模与稳定性定理的前提。
\end{enumerate}
定理~\ref{thm:constructive-equilibrium-cycle} 给出了这些核心资金、激励和状态条件同时成立的显式局部模型；更广的适用域仍须分别证明其前提。在上述前提下，对所有有限结算期 \(t\ge0\)，核心循环保持
\[
\begin{gathered}
  \forall i\in\mathcal E_t:\quad
  T_{it}^B\ge\beta_t^{\mathrm{base}}>0,\quad y_{it}^F>0,\\
  c_{it}>c_{jt}\Longrightarrow T_{it}^B>T_{jt}^B,\qquad
  \sum_{i\in\mathcal E_t}T_{it}^B\le B_t^B,\\
  \Gamma_{k,t}^{\mathrm{mkt}}>0,\quad
  \Pi_{k,t+1}^{P,\mathrm{net}}>0,\quad I_{k,t+1}^B>0.
\end{gathered}
\]
此处 \(c_{it}\) 是合同声明比较域内的最终可比贡献评分；真实贡献排序仅在核验准确时逐状态成立，残余误差分支只确认已证明的期望及概率界。
因此竞技汇聚注意力，广告竞争注入赞助，组织以资金池覆盖基本与差异分配，参与收入转化为商品购买；
商品回款覆盖产品成本并形成销售利润，偏好增量利润覆盖营销投入并形成新增净利润，已实现净利润支持
下一期生产与赞助。新增净利润相对于同一窗口的无本次营销反事实计量，不要求每一期利润均高于上一期。
若无冲击固定点另满足定理~\ref{thm:local-dynamic-stability}，则该闭环在相应邻域局部渐近稳定。
\end{theorem}

\begin{Certificate}[闭环循环的逐期验收核验记录]
\label{cert:competition-closed-cycle}
记 \(\mathfrak C_{\mathrm{cycle}}\) 为上述五项条件与以下记录的合取：赛事及支持事件关系、赞助披露、
真实注意力、资金到账时序、接纳人数与完成名单、基本及差异分配、逐人回执与依法代扣、产品购买与成本、
偏好反事实、完整营销成本、营销净增利润、可用现金和下一期合同。前向不变性另提交状态域、允许冲击范围、
资金与参与下界及状态转移的统一界证明；只验证历史上若干期收支，不构成未来全部状态的证明。
经验验收逐期核对实际值，模型预期与实际结算分别保存。存在未支付的有效参与者、无法归属的营销回报或
下一期资金缺口时，联合核验记录不通过，已成立的支付义务仍须履行。
\end{Certificate}

\begin{proof}[证明（分账恒等式与时间归纳）]
初态属于 \(\mathcal V\)，故 L1--L5 在首期成立。由 L2 和定理~\ref{thm:base-differential-coverage}，
有效参与者取得正基本份额，合同声明比较域内的最终贡献评分严格排序产生严格差异分配，总额不超过资金池；
带误差分支对真实努力另使用命题~\ref{prop:noisy-effort-contract} 的期望排序与误判概率，不以评分复算代替真实排序。最终性与代扣下界保证实际
正到账。L3 使注意力与收入分别连接产品选择和购买力，因果分解按同一销售口径计量。由 L4 和命题
~\ref{prop:sales-marketing-accounting}，
\[
  \Pi_{k,t+1}^{P,\mathrm{net}}
    =P_k(\vartheta^0,Y^1)+\Gamma_{k,t}^{\mathrm{mkt}}>0,
\]
已实现回款扣除产品、营销成本与已承诺用途后，为下一期提供已覆盖的正营销现金预算；下一期基本与差异
分配随其合同再次预留。假设第 \(t\) 期已成立，L5 给出 \(x_{t+1}\in\mathcal V\)，故同样论证适用于
第 \(t+1\) 期。数学归纳法使各期分配、成本覆盖、营销净增利润和再投入依次闭合。局部稳定结论由独立的
固定点与谱半径条件给出；上述收支与时间归纳本身不替代稳定性证明。证毕。
\end{proof}

\begin{remark}[闭环的失效边界]
产品销售获利与营销扩大利润分别检验。例如无本次营销的销售利润为 1000，营销后为 1300，而本次营销成本
为 400，则产品仍有正利润，营销净增利润却为 \(300-400=-100\)，不满足 L4。若资金只能覆盖少数获胜者，
或实际完成约定的普通参与者被分配零基本份额，则 L2 失败。若收入仅形成应收账款而没有可支配回款，或
下一期投入必须挪用已承诺的参与分配，跨期现金条件失败。若允许冲击使状态越出 \(\mathcal V\)，则持续
闭环结论仅截至该域内的已验证期间；下一期新承诺须重新验收。
\end{remark}

\begin{remark}[市场参与经济的循环结构]
生产--消费提供物质与利润基础，竞技组织把真实参与汇聚为注意力与可验证贡献，广告营销竞争把生产利润
转为赞助资金。组织依约保障参与即分配，并以贡献差异实现能者多得；参与收入增加购买力，产品偏好引导
购买方向；当市场收入不足以维持生存和真实进入时，坐标最小救济补齐底线。三者构成核心功能：
\[
  \boxed{\text{参与式分配}\;+\;\text{能者多得}\;+\;\text{坐标最小救济}}.
\]
产品销售回款覆盖产品成本并形成销售利润，偏好增量销售利润覆盖营销成本并扩大利润。
在闭环核验记录的声明域内，这些关系收束为
\[
\boxed{\begin{aligned}
  \text{已实现生产净利润}
  &\longrightarrow\text{竞争性广告赞助}\\
  &\longrightarrow\text{竞技组织与注意力汇聚}\\
  &\longrightarrow\text{资金池覆盖基本、差异分配}\\
  &\longrightarrow\text{参与收入与产品偏好}\\
  &\longrightarrow\text{商品购买与销售回款}\\
  &\longrightarrow\text{产品成本覆盖与销售利润}\\
  &\longrightarrow\text{营销成本覆盖与新增净利润}\\
  &\longrightarrow\text{下一期生产与竞争性赞助}.
\end{aligned}}
\]
长期参与形成的文化、规则信任和组织经验支持后续参与及偏好传导。注意力偏好回流提供市场资金和演化选择
压力；治理保证资格、计量、确权、支付和申诉可执行。资金覆盖、真实回款与可持续再投入把这些作用接入同一个
竞技主导、组织分配、消费回流和利润再投入的闭环循环。
\end{remark}

\appendix
\chapter{市场循环结构的条件性数学扩展}
本节研究经济过程的一种附加结构表示。收入、预算、参与激励与均衡由各自经济模型决定；
阶段名称对应不推出这些经济量相等。下面的有向映射和收缩条件须单独构造或核验，
不作为基本分配、营销回报或现金循环成立的必要前提。


\begin{definition}[有向空间、普通投影与最大有向嵌入]
有向空间 \(X=(|X|,dX)\) 由拓扑空间 \(|X|\) 与允许路径集 \(dX\) 组成；允许路径包含常值路径，并对
保向重参数化与首尾拼接封闭。有向映射把允许路径映为允许路径。有向同伦使用有向柱对象保存过程方向，
相关基础见~\cite{Grandis2003}。

令
\[
  U:\mathbf{dTop}\to\mathbf{Top}
\]
为遗忘方向函子，令
\[
  J:\mathbf{Top}\hookrightarrow\mathbf{dTop}
\]
给每个拓扑空间赋予全部连续路径可行的最大有向结构，则 \(U\circ J=\operatorname{id}_{\mathbf{Top}}\)。
因此普通同伦既可在最大有向结构下作为特例，也可由遗忘方向得到底层投影。有向框架更一般、保存的信息
更多；在固定有向结构内，有向等价的条件通常比遗忘方向后的普通同伦等价更强。

由于单次有向同伦一般不可逆，本文以 \(f\simeq_d g\) 表示由端点相容的单次有向同伦及其有限折线
组合生成的等价闭包；等价地，核验记录可以同时给出两个方向的保向变形。这样，\(\simeq_d\) 是本文用于
“有向同伦等价”的对称关系，而单次未来向变形仍保留其不可逆性。
\end{definition}

\begin{definition}[市场循环的公共有向骨架]
令 \(K\) 为具有七个顶点的有向胞腔复形：
\[
  k_I\to k_O\to k_R\to k_X\to k_S\to k_A\to k_{I'}\to k_I,
\]
依次表示投入或参与、形成可回应对象、分散回应、交换与价格发现、剩余形成、剩余分配、再投入。
失败、退出与监管阻断可以作为从相应顶点伸出的终止支路加入 \(K\)，但不能把时间方向倒置。

令 \(\mathcal C_P\) 为物质生产活动进入市场循环的有向状态范畴，\(\mathcal C_G\) 为围棋行业及其声明的
生产侧注意力桥接路径进入市场循环的有向状态范畴，定义其有向神经几何实现
\[
  X_P=|N\mathcal C_P|_d,
  \qquad
  X_G=|N\mathcal C_G|_d.
\]
其中神经中的一个有向单形由一串首尾可复合的合法状态转移生成，几何实现沿该转移顺序赋予允许方向。
二者的阶段对应为
\begin{center}
\begin{tabular}{p{0.10\textwidth}p{0.38\textwidth}p{0.42\textwidth}}
\toprule
公共阶段 & 生产市场循环 & 围棋行业及注意力桥接循环 \\
\midrule
\(k_I\) & 生产资料、劳动与资本投入 & 棋手时间、组织投入与多个生产者的竞争性盈利注入 \\
\(k_O\) & 商品或服务形成 & 对局、教学、赛事与观赏服务形成 \\
\(k_R\) & 消费者和交易对手回应 & 对手、观众与学员参与，注意力在赞助生产者之间分配 \\
\(k_X\) & 成交、价格与声誉更新 & 学费、门票、版权、赞助以及对生产商品的偏好和购买更新 \\
\(k_S\) & 净剩余形成 & 围棋服务剩余与生产侧回流收益形成 \\
\(k_A\) & 工资、利润、税费与其他分配 & 棋手、教师、平台、赛事方和传播者取得合同收入 \\
\(k_{I'}\) & 维护、消费与再生产 & 参与者消费、训练办赛与生产者下一轮竞争性投入 \\
\bottomrule
\end{tabular}
\end{center}
\end{definition}

\begin{assumption}[可收缩纤维与保向截面]
对 \(D\in\{P,G\}\)，存在满射有向映射 \(r_D:X_D\to K\) 与有向截面
\(s_D:K\to X_D\)，满足
\[
  r_Ds_D=\operatorname{id}_K,
  \qquad
  s_Dr_D\simeq_d\operatorname{id}_{X_D}.
\]
经济含义是：同一循环阶段内被商去的产品形态、棋局内容、赞助品牌和身份差异形成有向可收缩纤维；阶段之间的
代表路径与实际路径之间存在保持因果次序的变形。该条件必须逐纤维和逐过渡验证，不能由名称相似代替。
\end{assumption}

\begin{theorem}[生产市场循环--围棋行业循环有向同伦定理]
\label{thm:production-go-directed-homotopy}
若可收缩纤维与保向截面假设成立，则
\[
  X_P\simeq_d K\simeq_d X_G,
  \qquad
  X_P\simeq_d X_G.
\]
即物质生产活动与满足所列阶段和桥接条件的围棋行业路径不是在经验内容上相同，而是它们进入市场循环后的
声明子空间在有向同伦意义下等价。
\end{theorem}

\begin{Certificate}[有向同伦核验记录]
核验记录为
\[
  \mathfrak C_{dH}
  =\Tuple{X_P,X_G,K,r_P,s_P,r_G,s_G,H_P,H_G},
\]
其中 \(H_D\) 逐点见证
\[
  s_Dr_D\simeq_d\operatorname{id}_{X_D},
\]
并验证映射连续、保向、端点固定、合法域一致以及失败支路不被错误压缩进成功循环。
\end{Certificate}

\begin{proof}[证明（显式互逆映射）]
定义
\[
  f=s_Gr_P:X_P\to X_G,
  \qquad
  g=s_Pr_G:X_G\to X_P.
\]
由截面恒等式，
\[
\begin{aligned}
  gf
  &=s_Pr_Gs_Gr_P
  =s_P\operatorname{id}_Kr_P
  =s_Pr_P
  \simeq_d\operatorname{id}_{X_P},\\
  fg
  &=s_Gr_Ps_Pr_G
  =s_G\operatorname{id}_Kr_G
  =s_Gr_G
  \simeq_d\operatorname{id}_{X_G}.
\end{aligned}
\]
故 \(f,g\) 构成有向同伦互逆。该证明不是把两个原始领域按名称宣告相同，而是以
\(r_D,s_D,H_D\) 的可检查核验记录证明两个市场循环子空间具有共同的有向形变收缩核。证毕。
\end{proof}

\begin{corollary}[普通同伦特例与遗忘投影]
\label{cor:ordinary-homotopy-projection}
在最大有向结构 \(J\) 上，普通同伦作为有向同伦的特例出现；对上述有向同伦等价应用遗忘函子，则
\[
  U(X_P)\simeq U(X_G).
\]
反向一般不成立，因为普通同伦可以遗忘先参与后结算、先交换后分配等不可逆顺序。
\end{corollary}

\begin{proof}[证明（函子投影）]
有向映射和有向同伦在遗忘方向后仍分别是连续映射和普通连续同伦，所以有向同伦互逆映为普通同伦互逆。
而 \(U\circ J=\operatorname{id}\) 表明普通空间可嵌入最大有向结构。若仅知道底层普通同伦，则无法恢复被
遗忘的允许路径集，故不能一般地反推出有向同伦。证毕。
\end{proof}

\begin{remark}[严格同伦结论的失效条件]
若某一阶段为空、阶段投影不连续、纤维含有不可收缩锁定结构、债务或责任形成不可逆支路、退出状态被强行
并入再投入，或者映射改变交换与分配的先后顺序，则 \(\mathfrak C_{dH}\) 失败，不能确认严格同伦结论。
围棋行业只为非物质参与进入市场循环提供存在模型，不证明所有娱乐、平台或生产行业都满足同一核验记录。
\end{remark}

\begin{remark}[综合结论]
本文的核心贡献对象是参与式分配、能者多得与坐标最小救济组成的三功能结构。它依次扩大有效参与的正份额
覆盖，保留能力与贡献的差异激励，并把生存和真实市场接入移出胜负对象；在共同目标固定为广泛正份额、差异
激励和全员底线接入时，任何可行载体都必须实现与三者等价的功能，但平台、企业、合作组织、公共机制及其组合
均可能成为载体。注意力偏好回流是市场资金与条件性演化动力，核验、确权、支付、复算和申诉是实施基础。
现有集中型注意力经济证明注意力可以变现，却由命题~\ref{prop:attention-monetization-insufficient} 表明其本身不足以
广泛吸收就业与参与或修复收入--需求循环；命题~\ref{prop:core-functions-evolutionary-advantage} 只在共同口径下的
功能增益超过治理与集中成本时推出相对扩张，不宣称历史必然。

生产--消费是宏观物质再生产的必要内核，但经济不被狭义物质生产穷尽；狭义物质生产与单项活动的正经济价值
既非充分也非必要关系。生产盈利对参与、收入和需求的边际传导趋近于零时出现停滞风险。多个生产者只有在
注意力能够形成消费者对其产品的偏好指引时，才具有持续竞争性注资的内生市场动力；注资再经注意力和收入
两条通道把生产盈利接回消费。在最低物质需要之上的 Stone--Geary 声明域内，兴趣消费具有大于一的收入
弹性；超额预算扩大还会提高兴趣支出并使更多可盈利品类进入，使传统生产部门释放的劳动与能力有机会转入
兴趣导向的文化生产、竞技、服务和共同体活动。只有当付费覆盖、行业净余额、净需求桥接和
动态稳定同时越过事前门槛，才能进一步确认广义奢侈品产业足以支持声明域经济循环的结论。参与只要对净社会
剩余作正边际贡献，就具有经济价值，是否被定价、支付或由贡献者取得属于后续识别与分配问题。规范治理的
内部关键是参与事实、贡献次序和收入份额能否按声明规则重放、由独立复算者一致复核并在
声明容限内相对公正；外部边界是在风险外溢阈值内释放尽可能多的合法路径，不替市场指定结果。围棋行业既是
服务生产和参与分配的现实存在模型，也因有限棋盘、固定规则、可重放棋谱与历史认同而提供规则内可独立复算的
具体样本，但
具体赛事分账仍须合同和支付核验记录。只有满足阶段、纤维和保向核验记录的声明子空间，才与生产市场循环共享严格的
有向形变收缩核。共同剩余由谁取得仍取决于归属基准、规则控制权、议价和退出条件；治理不把全部差额归零，
市场形成交易与剩余也不等于市场自动保证相对公正。
\end{remark}

\clearpage

\clearpage
\begin{thebibliography}{99}
\phantomsection
\addcontentsline{toc}{chapter}{\bibname}

\bibitem{AcemogluRestrepo2019}
Daron Acemoglu and Pascual Restrepo.
\newblock Automation and New Tasks: How Technology Displaces and Reinstates Labor.
\newblock \emph{Journal of Economic Perspectives}, 33(2):3--30, 2019.
\newblock \url{https://doi.org/10.1257/jep.33.2.3}.

\bibitem{Varian1998InformationGoods}
Hal R. Varian.
\newblock Markets for Information Goods.
\newblock University of California, Berkeley, revised October 1998.
\newblock \url{https://people.ischool.berkeley.edu/~hal/Papers/japan/}.

\bibitem{Stone1954LES}
Richard Stone.
\newblock Linear Expenditure Systems and Demand Analysis: An Application to the Pattern of British Demand.
\newblock \emph{The Economic Journal}, 64(255):511--527, 1954.
\newblock \url{https://doi.org/10.2307/2227743}.

\bibitem{Lancaster1966Characteristics}
Kelvin J. Lancaster.
\newblock A New Approach to Consumer Theory.
\newblock \emph{Journal of Political Economy}, 74(2):132--157, 1966.
\newblock \url{https://doi.org/10.1086/259131}.

\bibitem{Myerson1981OptimalAuction}
Roger B. Myerson.
\newblock Optimal Auction Design.
\newblock \emph{Mathematics of Operations Research}, 6(1):58--73, 1981.
\newblock \url{https://doi.org/10.1287/moor.6.1.58}.

\bibitem{RochetTirole2003TwoSided}
Jean-Charles Rochet and Jean Tirole.
\newblock Platform Competition in Two-Sided Markets.
\newblock \emph{Journal of the European Economic Association}, 1(4):990--1029, 2003.
\newblock \url{https://doi.org/10.1162/154247603322493212}.

\bibitem{Ostrom1990Commons}
Elinor Ostrom.
\newblock \emph{Governing the Commons: The Evolution of Institutions for Collective Action}.
\newblock Cambridge University Press, 1990.
\newblock \url{https://doi.org/10.1017/CBO9780511807763}.

\bibitem{Coase1937Firm}
Ronald H. Coase.
\newblock The Nature of the Firm.
\newblock \emph{Economica}, 4(16):386--405, 1937.

\bibitem{North1990Institutions}
Douglass C. North.
\newblock \emph{Institutions, Institutional Change and Economic Performance}.
\newblock Cambridge University Press, 1990.

\bibitem{ArrowDebreu1954Equilibrium}
Kenneth J. Arrow and G\'erard Debreu.
\newblock Existence of an Equilibrium for a Competitive Economy.
\newblock \emph{Econometrica}, 22(3):265--290, 1954.
\newblock \url{https://doi.org/10.2307/1907353}.

\bibitem{Hayek1945Knowledge}
Friedrich A. Hayek.
\newblock The Use of Knowledge in Society.
\newblock \emph{American Economic Review}, 35(4):519--530, 1945.
\newblock \url{https://www.jstor.org/stable/1809376}.

\bibitem{Shapley1953Value}
Lloyd S. Shapley.
\newblock A Value for n-Person Games.
\newblock In H. W. Kuhn and A. W. Tucker, editors,
\emph{Contributions to the Theory of Games II}, pages 307--317.
Princeton University Press, 1953.
\newblock \url{https://doi.org/10.1515/9781400881970-018}.

\bibitem{Grandis2003}
Marco Grandis.
\newblock Directed Homotopy Theory, I: The Fundamental Category.
\newblock \emph{Cahiers de Topologie et G\'eom\'etrie Diff\'erentielle Cat\'egoriques},
44(4):281--316, 2003.
\newblock \url{https://eudml.org/doc/91674}.

\bibitem{Silver2016AlphaGo}
David Silver, Aja Huang, Chris J. Maddison, et al.
\newblock Mastering the Game of Go with Deep Neural Networks and Tree Search.
\newblock \emph{Nature}, 529:484--489, 2016.
\newblock \url{https://doi.org/10.1038/nature16961}.

\bibitem{IGF2026}
International Go Federation.
\newblock Official Website and President's Message.
\newblock Accessed 2026-08-30.
\newblock \url{https://intergofed.org/}.

\bibitem{GreenLaffont1979Incentives}
Jerry Green and Jean-Jacques Laffont.
\newblock \emph{Incentives in Public Decision-Making}.
\newblock North-Holland, Amsterdam, 1979.

\bibitem{Pew2025WorkersAI}
Luona Lin and Kim Parker.
\newblock U.S. Workers Are More Worried Than Hopeful About Future AI Use in the Workplace.
\newblock Pew Research Center, 2025.
\newblock \url{https://www.pewresearch.org/social-trends/2025/02/25/workers-views-of-ai-use-in-the-workplace/}.

\bibitem{Gallup2026AI}
Julie Ray.
\newblock Americans Cool Toward AI.
\newblock Bentley University--Gallup Business in Society research, 2026.
\newblock \url{https://news.gallup.com/poll/712751/americans-cool-toward.aspx}.

\bibitem{BrynjolfssonChandarChen2026}
Erik Brynjolfsson, Bharat Chandar, and Ruyu Chen.
\newblock Canaries in the Coal Mine? Six Facts about the Recent Employment Effects of Artificial Intelligence.
\newblock Stanford Digital Economy Lab working paper, revised 2026.
\newblock \url{https://digitaleconomy.stanford.edu/publication/canaries-in-the-coal-mine-six-facts-about-the-recent-employment-effects-of-artificial-intelligence/}.

\bibitem{ILO2025GenAIJobs}
Pawe\l{} Gmyrek, Janine Berg, Karol Kami\'nski, et al.
\newblock Generative AI and Jobs: A Refined Global Index of Occupational Exposure.
\newblock ILO Working Paper 140, 2025.
\newblock \url{https://www.ilo.org/publications/generative-ai-and-jobs-refined-global-index-occupational-exposure}.

\bibitem{YotzovEtAl2026}
Ivan Yotzov, Jose Maria Barrero, Nicholas Bloom, et al.
\newblock Firm Data on AI.
\newblock NBER Working Paper 34836, 2026.
\newblock \url{https://doi.org/10.3386/w34836}.

\bibitem{ECB2026AIMonetaryPolicy}
European Central Bank.
\newblock AI and Monetary Policy.
\newblock 2026.
\newblock \url{https://www.ecb.europa.eu/press/key/date/2026/html/ecb.sp260706~b81aa4e329.en.html}.

\bibitem{Chen2026Attention}
Daniel Chen.
\newblock A Model of the Attention Economy: Data, Platforms, and Discovery.
\newblock \emph{American Economic Review}, forthcoming, 2026.
\newblock \url{https://doi.org/10.1257/aer.20240651}.

\bibitem{RockallTavaresPizzinelli2025}
Emma J. Rockall, Marina Mendes Tavares, and Carlo Pizzinelli.
\newblock AI Adoption and Inequality.
\newblock IMF Working Paper 2025/068, 2025.
\newblock \url{https://doi.org/10.5089/9798229006828.001}.

\bibitem{GladstoneBellezza2025}
Joe J. Gladstone and Silvia Bellezza.
\newblock More Than Money: The Relative Importance of Cultural, Social, and Economic Capital for Highbrow Cultural Experiences.
\newblock \emph{Journal of the Association for Consumer Research}, 10(4), 2025.
\newblock \url{https://doi.org/10.1086/737202}.

\bibitem{ZhangNegus2026}
Qian Zhang and Keith Negus.
\newblock Introduction to Artificial Intelligence in the Creator Economy: Creativities, Adaptation, Agency, and Trust.
\newblock 2026.
\newblock \url{https://doi.org/10.1177/20594364261466625}.

\bibitem{Gao2025PureMathLocal}
Gao, Z.（高政）.
\newblock \emph{Meta-Mathematical Theory based on Pan-Logic Analysis and Pan-Iterative Analysis
(GaoZheng G-Framework) and the Principal-Bundle-Based Generalized Noncommutative Lie Algebra
(GaoZheng G-Algebra): An Integrated Construction}.
\newblock Zenodo, v1.0, 2025.
\newblock \url{https://doi.org/10.5281/zenodo.17651584}.
\newblock 本地源文件：\code{docs/project_docs/Pub_GFramework_PureMath.tex}。
本地文件的版本标注为 Version 1.1, 2025；上述引用版本为 v1.0。

\bibitem{Gao2025AppMathVol1Local}
Gao, Z.（高政）.
\newblock \emph{GaoZheng G-Framework and GaoZheng G-Algebra: Applied Mathematics Volume I
-- Law-Space Geometry, GRL, Quantum Computing, and Superconductivity}.
\newblock Zenodo, v1.2, 2025.
\newblock \url{https://doi.org/10.5281/zenodo.17686133}.
\newblock 本地源文件：\code{docs/project_docs/Pub_GFramework_AppMath_Vol1.tex}。
本地文件的版本标注为 Version 1.3, 2025；上述引用版本为 v1.2。

\bibitem{Gao2025AppMathVol2Local}
Gao, Z.（高政）.
\newblock \emph{GaoZheng G-Framework and GaoZheng G-Algebra: Applied Mathematics Volume II
-- LBOPB, Life-Science Monoids, and Generative Precision Medicine}.
\newblock Zenodo, v1.1, 2025.
\newblock \url{https://doi.org/10.5281/zenodo.17744672}.
\newblock 本地源文件：\code{docs/project_docs/Pub_GFramework_AppMath_Vol2.tex}。

\bibitem{Gao2025AppMathVol3Local}
Gao, Z.（高政）.
\newblock \emph{GaoZheng G-Framework and GaoZheng G-Algebra: Applied Mathematics Volume III
-- HACA, PACER, and Certificate-Based AI}.
\newblock Zenodo, v1.0, 2025.
\newblock \url{https://doi.org/10.5281/zenodo.17762295}.
\newblock 本地源文件：\code{docs/project_docs/Pub_GFramework_AppMath_Vol3.tex}。

\bibitem{Gao2026NotebookLocal}
Gao, Z.（高政）.
\newblock \emph{Meta-Mathematical Theory based on Pan-Logic Analysis and Pan-Iterative Analysis
(GaoZheng G-Framework) and the Principal-Bundle-Based Generalized Noncommutative Lie Algebra
(GaoZheng G-Algebra) 中文工作笔记}.
\newblock Zenodo, Version v5.2, 2026.
\newblock \url{https://doi.org/10.5281/zenodo.22739417}.
\newblock 本地源文件目录：\code{docs/project_docs/notebook_tex_v5.2/}。

\end{thebibliography}

\end{CJK*}
\end{document}
